% Ensamblaje causal de la Parte III del sucesor de 102 capítulos.
% Los propietarios rectores se consumen completos.  Los títulos de capítulo
% se sustituyen únicamente por la jerarquía aprobada; las ampliaciones se
% añaden dentro de la residencia científica correspondiente.

% Punto de montaje redefinible por el compilador único.
% La ruta por defecto es válida cuando el módulo se consume desde
% ``manuscrito/``; el editor único puede redefinirla antes del \input.
\providecommand{\HMTParteIIIRoot}{../colaboracion/parte_iii}
\providecommand{\HMTParteIIISource}{\HMTParteIIIRoot/source/public_final}
\providecommand{\HMTParteIIICapitulos}{\HMTParteIIIRoot/tex/capitulos_finales}
\providecommand{\APP}{\mathrm{APP}}
\providecommand{\TRIT}{\mathrm{TRIT}}
\providecommand{\TPK}{\mathrm{TPK}}
\providecommand{\etaRet}{\eta_{\mathrm{ret}}^{(\mathrm{deg})}}
\providecommand{\etaRetrad}{\eta_{\mathrm{ret}}^{(\mathrm{rad})}}
\providecommand{\alphaAngle}{\alpha_{\angle}}


El orden de los capítulos que siguen es expositivo, no causal. La conexión
nonádica y la prolongación completa

\[
w_6\longrightarrow w_{12}\longrightarrow w_{18}\longrightarrow w_{24}
\longrightarrow w_{30}\longrightarrow R_{36}\longrightarrow G_9
\]

estabilizan un único estado coinductivo de APP--TRIT--TPK. De ese estado se
publican correlacionadamente \(\pi\), \(e\), \(\varphi\) y \(\alpha\): los
capítulos sucesivos separan sus lectores para poder estudiarlos con precisión,
pero no convierten esa separación tipográfica en una generación serial. La
clausura de \(\pi\), la propagación de \(e\), la autoescala de \(\varphi\) y
la publicación firmada y dodecafásica de \(\alpha\) conservan estatutos
propios dentro de una misma generación coinductiva correlacionada.

\newcommand{\HMTSuccessorChapterInput}[2]{%
  \chapter{#1}%
  \begingroup
  \renewcommand{\chapter}[2][]{}%
  \input{#2}%
  \endgroup}

\HMTFlushCurrentChapter
\HMTSuccessorChapterInput
  {Álgebra y clasificación genealógica de las constantes fundamentales}
  {../colaboracion/parte_iii/tex/capitulos_finales/27_algebra_clasificacion.tex}
% Rectificación quirúrgica: generación conjunta, profundidad arbitraria y
% estatuto tipado de los transportes ternarios.  Este bloque pertenece al
% capítulo 27 y precede a la separación expositiva de los cuatro lectores.

\section{Teorema conjunto de generación coinductiva y profundidad arbitraria}
\label{sec:c27-generacion-coinductiva-profundidad-arbitraria}

La estructura de partida de este teorema no es el conjunto desnudo de nueve
símbolos. Es el sistema tipado
\[
 \mathfrak H_9=
 \bigl(\mathrm{APP},\mathrm{TRIT},\mathrm{TPK};
       \widehat{\mathcal X}^{\mathrm{enr}},\mathcal U,
       \Gamma_9,\operatorname{Ext},\operatorname{tr}\bigr),
\]
construido sobre la doble realización posicional y aritmética de
\(\mathcal D_9=\{1,\ldots,9\}\). APP acumula y pliega mediante sus hojas
\(\Sigma\) y \(\Pi\), conservando residuo y cociente; TRIT orienta cada
transición en uno de sus tres regímenes; y TPK selecciona, transporta,
memoriza y retorna sobre la fibra enriquecida. En particular,
\[
 M_{\rm ph}:\mathcal D_9\longrightarrow\{+1,-1,0\},
 \qquad
 m_{\rm ph}:=\bigl(M_{\rm ph}(1),\ldots,M_{\rm ph}(9)\bigr)^{\mathsf T},
\]
\[
 \operatorname{Tra}_t(y,d):=\operatorname{Lift}(T_d)(y),
 \qquad
 \operatorname{Upd}_t:=\operatorname{Rec}_t\circ\operatorname{Mem}_t,
\]
y, para todo \(x\in\widehat{\mathcal X}^{\mathrm{enr}}\),
\[
 \mathcal U_t(x)
 =\operatorname{Upd}_t\!\left(
   \operatorname{Tra}_t\!\left(
    \operatorname{Sel}_t\!\left(x,M_{\rm ph}(g(t))\right),d(t)
   \right)\right).
\]
Esta composición es el dato dinámico efectivo: la carta de nueve símbolos
es su soporte finito, no un sustituto de la dinámica. En particular,
\(M_{\rm ph}(g(t))\) es el valor escalar que recibe el selector, mientras
que \(m_{\rm ph}\) es el vector de los nueve valores de la función; ninguno
se compone directamente con un operador de estado.

Para \(r\in6\mathbb N_{>0}\), sean
\(\widehat{\mathcal X}^{\mathrm{enr}}_r\) los estados truncados con hoja,
orientación, residuo, cociente, acarreo, ruta, frontera y memoria, y sean
\[
 \operatorname{Ext}_{r+6,r}\subseteq
 \widehat{\mathcal X}^{\mathrm{enr}}_r\times
 \widehat{\mathcal X}^{\mathrm{enr}}_{r+6},
 \qquad
 \operatorname{tr}_{r+6,r}:
 \widehat{\mathcal X}^{\mathrm{enr}}_{r+6}
 \longrightarrow
 \widehat{\mathcal X}^{\mathrm{enr}}_r
\]
la relación de prolongación admisible y el truncamiento. Para
\(x_r\in\widehat{\mathcal X}^{\mathrm{enr}}_r\), la fibra local
\[
 \operatorname{Ext}_{r+6,r}(x_r)
 :=\{x_{r+6}:(x_r,x_{r+6})\in\operatorname{Ext}_{r+6,r}\}
\]
puede contener varios hijos. La compatibilidad correcta es
\[
 (x_r,x_{r+6})\in\operatorname{Ext}_{r+6,r}
 \Longrightarrow
 \operatorname{tr}_{r+6,r}(x_{r+6})=x_r.
\]
En \(R_{36}\), los tres renglones recuperados
\(222220,021222,102011\), transportados al calibre activo por la permutación
\(p_0=(0,1,2,4,5,3)\) en índices desde cero ---equivalentemente,
\(p_1=(1,2,3,5,6,4)\) en índices desde uno---, portan antes de toda
realización real los caracteres
internos de clausura, propagación y autoescala. Su triple, junto con el
registro firmado \(U_{12}\) y el sello \(K\), porta la coordenada interna de
\(\alpha\). Para cada
\(\chi\in\{\pi,\varphi,e,\alpha\}\), escribimos el estado así enriquecido
como
\[
 x_r^\chi=(\widetilde x_r,\chi_{\rm HMT},
            \varepsilon_r^\chi,N_r^\chi).
\]
Aquí \(\varepsilon_r^\chi\) es el residuo interno producido por el cociente
APP en el estado, no \(\operatorname{frac}(P_{\rm ar}(\chi))\). La división
euclídea interna fija
\[
 d_r^\chi=\operatorname{Dig}_{729}(\varepsilon_r^\chi),
 \quad 0\le d_r^\chi<729,
 \quad
 \varepsilon_{r+6}^\chi=729\varepsilon_r^\chi-d_r^\chi,
 \quad
 N_{r+6}^\chi=729N_r^\chi+d_r^\chi.
\]
La selección, el transporte y la actualización sobre la sección
\(\chi\) se tipan por
\[
 \operatorname{Sel}_r^\chi(x)
 :=\operatorname{Sel}\!\left(
 x,M_{\rm ph}(\rho_9(r/6+1))\right),
 \qquad
 \operatorname{Tra}_r^\chi(y,d):=\operatorname{Lift}(T_d)(y),
 \qquad
 \operatorname{Upd}_r^\chi:=\operatorname{Rec}_r^\chi\circ
 \operatorname{Mem}_r^\chi.
\]
La transición completa es
\[
 \mathcal U_r^\chi(x)
 :=\operatorname{Upd}_r^\chi\!\left(
  \operatorname{Tra}_r^\chi\!\left(
   \operatorname{Sel}_r^\chi(x),d_r^\chi
  \right)\right),
\]
y sólo lee el estado enriquecido anterior: fase, memoria, acarreo, hoja,
orientación, ruta, frontera, firma, registro, carácter y residuo internos.
No lee una cifra futura, \(P_{\rm ar}\), ni una horquilla de Machin,
factorial o Perron. Definimos la subrelación forward
\[
 \operatorname{Ext}^{\chi,\rightarrow}_{r+6,r}
 :=\{(x_r^\chi,\mathcal U_r^\chi(x_r^\chi))\}
 \subseteq\operatorname{Ext}_{r+6,r}.
\]
No es la selección retrospectiva
\(\operatorname{Ext}\cap\{P_{\rm ar}=\chi\}\): el carácter ya es una
coordenada causal de \(x_{36}^\chi\). Las ecuaciones anteriores dan
\(\lfloor N_{r+6}^\chi/729\rfloor=N_r^\chi\) y conservan hoja, orientación,
acarreo, frontera, ruta, firma, registro y genealogía.

La relación completa \(\operatorname{Ext}\) sigue siendo multivaluada. La
unicidad requerida es la de cada sección forward sobre su semilla R36
producida,
\[
 \#\Bigl\{(x_r^\chi)_{r\ge36}:x_{36}^\chi\text{ fijado},\ 
 (x_r^\chi,x_{r+6}^\chi)
 \in\operatorname{Ext}^{\chi,\rightarrow}_{r+6,r}\Bigr\}=1,
\]
no la igualdad falsa \(\#\operatorname{Ext}_{r+6,r}(x_r)=1\) en cada puerta
local. Esta compatibilidad se prolonga a todo nivel. La notación
\[
 w_6\longrightarrow w_{12}\longrightarrow w_{18}\longrightarrow w_{24}
 \longrightarrow w_{30}\longrightarrow R_{36}\longrightarrow G_9
\]
designa estas aplicaciones y sus cambios de tipo, no una lista de rótulos.
La holonomía \(\Gamma_9\) satisface
\[
 g(k+9)=g(k),
 \qquad q^{(9)}(k+9)=q^{(9)}(k)+1,
\]
de modo que retorna la fase mientras avanzan el cociente, la ruta y la
memoria del estado.

Una historia canónica genérica basada en \(a_{(1,1)}\) y
\(t_j=M_{\rm ph}(\rho_9(j+1))\) acredita la no vacuidad de
\(\operatorname{Ext}\); no individualiza por sí sola las cuatro secciones
anteriores. Éstas quedan empalmadas por sus coordenadas internas R36 y por la
regla forward que acaba de definirse.

\begin{theorem}[Construcción inaugural desde las semillas]
\label{thm:c27-construccion-inaugural-semillas}
Sea
\[
 I_9=\{0,\ldots,8\},\qquad D=\{N,E,S,O\},\qquad
 \mathcal S=I_9^2\times D,
\]
y sea \(\Omega_{\mathrm{seed}}=\mathcal S_+\times\mathcal S_\times\) el
dominio de los dos cursores que realizan las hojas aditiva y multiplicativa
de APP. Entonces
\[
 |\Omega_{\mathrm{seed}}|=(9^2\cdot4)^2=104\,976,
\]
y la recurrencia de \(54\) pasos y la reducción ternaria definen
\[
 \Omega_{\mathrm{seed}}
 \xrightarrow{\ T_{54}^{\rm seed}\ }
 \mathcal U_6^{\mathrm{cat}}
 \xrightarrow{\ \rho_3\ }
 \mathcal W_6
 \hookrightarrow\mathbb F_3^6,
\]
con
\[
 |\mathcal U_6^{\mathrm{cat}}|=468,
 \qquad |\mathcal W_6|=243,
 \qquad |\mathbb F_3^6|=729.
\]
Las fibras satisfacen exactamente
\[
 432\cdot144+18\cdot432+18\cdot1944=104\,976.
\]
La relación de extensión del TPK prolonga las palabras supervivientes a la
familia de secciones compatibles de profundidad arbitraria. Las subrelaciones
forward ancladas en R36 individualizan la generación coinductiva
correlacionada de \(\pi,\varphi,e\) y \(\alpha\); la relación completa puede
seguir teniendo varias prolongaciones locales. Sobre esa familia se
construyen, conjuntamente, la sombra arquimediana y la realización
excepcional del \cref{thm:c27-tesis-rectora-doble-proyeccion}.
\end{theorem}

\begin{proof}
El producto cartesiano anterior enumera exhaustivamente las posiciones y
orientaciones de ambos cursores. La aplicación \(T_{54}^{\rm seed}\) se obtiene al
componer en cada paso la selección de hoja, el transporte cardinal, la
actualización del cociente y la memoria de acarreo. La enumeración de su
imagen contiene \(468\) sextetos distintos; la suma de las cardinalidades de
sus fibras es la identidad expuesta. La aplicación componente a componente
\(\rho_3\) contiene \(243\) palabras distintas. Finalmente, la no vacuidad en
todo horizonte y la compatibilidad de \(\operatorname{Ext}\) con los
truncamientos producen el límite inverso. La funcionalidad de cada
\(\operatorname{Ext}^{\chi,\rightarrow}\), a partir del estado R36 que ya
porta \(\chi_{\rm HMT}\), demuestra por inducción la unicidad de esa sección
global sin exigir que cada fibra local de \(\operatorname{Ext}\) sea
unitaria. Cada afirmación es una identidad interna o una consecuencia de la
compatibilidad proyectiva; ninguna utiliza como entrada un número real, una
constante reconocida o una tabla decimal.
\end{proof}

\begin{theorem}[Tesis rectora: un continuo y dos proyecciones]
\label{thm:c27-tesis-rectora-doble-proyeccion}
Sea \(\Omega_\Gamma^{\mathrm{cal}}\subset
X_\infty^{\mathrm{cal}}\) el espacio de secciones supervivientes de las
vueltas enriquecidas del TPK. El acarreo de las tres vueltas
\(\gamma_-,\gamma_0,\gamma_+\) define, sin evaluación arquimediana previa,
el homeomorfismo
\[
 \beta:\Omega_\Gamma^{\mathrm{cal}}\xrightarrow{\ \sim\ }
 (\mathbb F_3^6)^{\mathbb N}=: \mathcal W_\infty.
\]
Las formas equilibradas normalizadas
\[
 a_k+c_k=\rho_k+3c_{k+1},\qquad
 \rho_k\in\{-1,0,+1\},
\]
construyen el anillo ordenado \(\mathbb Z_{\rm HMT}\), su cuerpo de
fracciones \(\mathbb Q_{\rm HMT}\) y la completación
\[
 \mathbb R_{\rm HMT}:=
 \operatorname{Cau}(\mathbb Q_{\rm HMT})/{\asymp}.
\]
La estructura discreta conjunta es el producto fibrado
\[
\mathcal C_{\mathrm{cont}}^{\mathrm{disc}}
=\bigl(\mathbb Z_{\rm HMT}\times
\Omega_\Gamma^{\mathrm{cal}}\bigr)
\times_{\mathcal W_\infty}X_\infty^{\mathrm{cal}}.
\]
Sus dos aplicaciones naturales son
\[
 P_{\mathrm{ar}}:\mathbb Z_{\rm HMT}\times
 \Omega_\Gamma^{\mathrm{cal}}\longrightarrow\mathbb R_{\rm HMT},
 \qquad
 Q_\infty:X_\infty^{\mathrm{cal}}\longrightarrow
 (\mathscr C_W^{\mathrm{cal}})^{\mathbb N}.
\]
Si \(\beta(\xi)=(b_q)_{q\geq0}\),
\(\nu(b)=\sum_{i=1}^{6}b_i3^{6-i}\) y
\[
 q_n(m,\xi)=m+\sum_{q=0}^{n-1}
 \frac{\nu(b_q)}{729^{q+1}},
\]
entonces \((q_n)\) es de Cauchy y \(P_{\mathrm{ar}}(m,\xi)\) es su clase
en \(\mathbb R_{\rm HMT}\). Los cilindros racionales
\[
 I_n(m;b_0,\ldots,b_{n-1})=
 \left[q_n,q_n+729^{-n}\right]
\]
son compatibles, tienen diámetro tendente a cero y la partición recurrente
en \(729\) subcilindros construye, para toda clase de Cauchy, una cinta
\((b_q)\) que la representa. En consecuencia,
\[
 P_{\mathrm{ar}}\ \text{es sobreyectiva},\qquad
 (\mathbb Z_{\rm HMT}\times\Omega_\Gamma^{\mathrm{cal}})
 /{\sim_{\rm ar}}\ \cong\ \mathbb R_{\rm HMT}.
\]
Sólo después de esta construcción interna, la unicidad del cuerpo ordenado
completo arquimediano reconoce \(\mathbb R_{\rm HMT}\cong\mathbb R\).

La segunda proyección conserva la incidencia calibrada de cada nivel; en las
coordenadas \(x=((w_j,f_j))_{j\geq1}\) actúa por
\[
 Q_\infty(x)
 =\bigl(f_j,\Gamma_{f_jA_Wf_j^{-1}}(w_j)\bigr)_{j\geq1}.
\]
Así, \(\mathbb R_{\rm HMT}\), reconocido posteriormente como \(\mathbb R\),
es la sombra arquimediana obtenida al contraer la genealogía completa de
cilindros. La proyección excepcional conserva esa genealogía como incidencia
de frontera y admite la realización sucesiva
Paley--Witt--Golay--Leech--\(V^\natural\)--Monstruo. Las dos ramas proceden
del mismo estado y ninguna se reconstruye retrospectivamente desde la otra.
\end{theorem}

\begin{proof}
Las tres vueltas enriquecidas existen después de cada prefijo y su acarreo
las distingue por \(-1,0,+1\). La concatenación y el truncamiento son
inversos, lo que prueba la biyectividad y continuidad de \(\beta\). Para una
clase de Cauchy de \(\mathbb Q_{\rm HMT}\), se elige primero su cilindro
unitario ordenado y después, recursivamente, el único subcilindro semiabierto
de la partición interna en \(729\) piezas que contiene la clase; en una
frontera común, las dos escrituras se identifican por \(\sim_{\rm ar}\).
La sucesión de extremos izquierdos es precisamente \((q_n)\), converge a la
clase dada y prueba la sobreyectividad sin presuponer una expansión de
\([0,1]\) ni un punto de \(\mathbb R\). La identificación con el cuerpo real
es el teorema de unicidad aplicado después a \(\mathbb R_{\rm HMT}\).

Los subespacios que satisfacen el cierre de un carácter concreto son filtros
de supervivencia dentro de la capacidad global. Su función es individualizar
las secciones de \(\pi,e,\varphi,\alpha\); la sobreyectividad del continuo
procede de la completación de \(\mathbb Q_{\rm HMT}\). Esta separación de
dominios impide
convertir una condición de selección de constantes en una premisa de la
cobertura de \(\mathbb R\).

En la segunda rama, la compatibilidad
\(r_{n,m}Q_n=Q_m p_{n,m}\) permite pasar al límite y define \(Q_\infty\).
La conjugación por \(f_j\) transporta el calibre sin borrar la incidencia.
Los funtores posteriores de código, retículo y álgebra de vértices realizan
la cadena excepcional. La igualdad de las cintas en el producto fibrado
construye la fuente conjunta y conserva informaciones complementarias: valor
arquimediano en la primera; incidencia, frontera y genealogía en la segunda.
\end{proof}

\begin{corollary}[Expansión infinita como objeto matemático]
\label{cor:c27-expansion-infinita-objeto}
Para cada carácter numérico \(\chi\) publicado por el estado coinductivo, la
lectura decimal es una secuencia completa
\[
 \operatorname{Dec}_\chi(x_\infty)
 =(d_{\chi,1},d_{\chi,2},\ldots)
 \in\{0,\ldots,9\}^{\mathbb N}.
\]
Su prefijo de longitud \(N\) es la proyección
\(\rho_N\operatorname{Dec}_\chi(x_\infty)\). Por ello, mil, diez mil o un
millón de cifras son cortes del mismo objeto infinito; la generación no
consiste en elevar sucesivamente una cota experimental.
\end{corollary}

\begin{theorem}[Generación coinductiva correlacionada de \(\pi\), \(\varphi\), \(e\) y \(\alpha\)]
\label{thm:c27-cuadrirrelacion-generada}
El perfil completo del sistema \(\mathfrak H_9\) determina la familia
correlacionada de secciones forward compatibles
\[
 \boldsymbol x_\infty=
 (x_\infty^\pi,x_\infty^\varphi,x_\infty^e,x_\infty^\alpha)
 \in
 \widehat\Omega_\pi\times_{R_{36}}
 \widehat\Omega_\varphi\times_{R_{36}}
 \widehat\Omega_e\times_{R_{36}}\widehat\Omega_\alpha,
 \qquad
 \widehat\Omega_\chi:=
 \varprojlim_{r\ge36}
 (\widehat{\mathcal X}^{\mathrm{enr},\chi}_r,
  \operatorname{Ext}^{\chi,\rightarrow}_{r+6,r}),
\]
y cuatro caracteres correlacionados
\[
 \bigl(\Pi_{\mathrm{cl}},\Pi_{\mathrm{aut}},
       \Pi_{\mathrm{prop}},\Pi_{\mathrm{dod}}\bigr)
       (\boldsymbol x_\infty)
 =\bigl(\pi_{\mathrm{HMT}},\varphi_{\mathrm{HMT}},
        e_{\mathrm{HMT}},\alpha_{\mathrm{HMT}}\bigr).
\]
Los tres primeros caracteres proceden, respectivamente, de la única
órbita de clausura, la única órbita de autoescala y la única
órbita de propagación del catálogo enriquecido. El cuarto procede de la
coordenada primitiva firmada de la familia terminal correlacionada y de su
cierre dodecafásico. Las cuatro coordenadas son, por tanto, salidas de la
generación coinductiva correlacionada de profundidad arbitraria gobernada por
la misma dinámica APP--TRIT--TPK, con retorno de fase, avance de memoria y
conservación genealógica, aunque sus aplicaciones de publicación sean
distintas.
\end{theorem}

\begin{proof}
El censo exhaustivo
\(104,976\to468\to243\subset\mathbb F_3^6\) descompone la imagen realizada
en \(43\) órbitas diedrales. Los invariantes de fibra dejan una sola órbita
en cada una de las clases de clausura, propagación y autoescala; el calibre
interno orienta en ellas, sin consulta decimal, los representantes
\[
 w_6^{\mathrm{cl}}=010211,
 \qquad w_6^{\mathrm{prop}}=201101,
 \qquad w_6^{\mathrm{aut}}=121200.
\]
Las subrelaciones forward conservan el carácter interno y los predicados que
definen las tres clases; Machin, la recurrencia factorial y Perron no eligen
ningún hijo de \(\operatorname{Ext}\). En el
canal de clausura, la composición pentarregional con retorno unitario produce
las razones \(1/5\), \(1/239\) y la relación \(120/119\); la identidad de
Machin reconoce exactamente la media vuelta. En el canal de propagación, la
ley semigrupal normalizada fuerza
\((n+1)a_{n+1}=a_n\), luego \(a_n=1/n!\) y \(P_1=e\). En el canal de
autoescala, la incidencia
\[
 F_{\mathrm{av}}=\begin{pmatrix}0&1\\1&1\end{pmatrix}
\]
posee un único rayo positivo y su ecuación propia fuerza
\(x^2-x-1=0\), \(x>0\), esto es, \(x=\varphi\).

En el perfil TPK completo, la lectura terminal
\(x_{\mathrm{term}}\mapsto(B_{\mathrm{term}},U_{12}^{\mathrm{sgn}})\)
publica simultáneamente la carta de los tres canales y la coordenada
primitiva firmada. La
segunda se identifica reversiblemente con el sello \(K\) y sus diferencias
dodecafásicas. La recurrencia euclídea con acarreo
\[
 s_m=P_m+E_m-\Phi_m-K_m+c_{m+1},\qquad
 A_m=s_m\bmod1000,\qquad
 c_m=(s_m-A_m)/1000
\]
produce \(\alpha_{\mathrm{HMT}}\) sin recibirla como entrada. La
caracterización analítica independiente comienza con el núcleo finito
\(E_{21/22}^{(9)}\), que determina una raíz simple \(\xi_9\) en
\((0,10^{-2})\) sin consultar la salida dodecafásica. La iteración del mismo
lector graduado produce el sistema inverso
\(E_{21/22}^{(6)}\leftarrow E_{21/22}^{(9)}\leftarrow
E_{21/22}^{(12)}\leftarrow\cdots\). Sus truncamientos conmutan con los de la
vía dodecafásica y ambas familias determinan en cada profundidad el mismo
cilindro superviviente. Como sus diámetros tienden a cero, la intersección es
unitaria y se obtiene exactamente
\(\alpha_A=\alpha_B=\alpha_{\mathrm{HMT}}\). La igualdad
\(\operatorname{Tr}_9(\xi_9^{-1})=
\operatorname{Tr}_9(\alpha_{12}^{-1})\) es únicamente el certificado finito
de orden nueve de esa identidad de secciones completas. Así, la separación
de lectores conserva el origen común de las cuatro coordenadas sin introducir
una prolongación pendiente ni un valor objetivo.
\end{proof}

\begin{proposition}[Registro dodecafásico del estado completo anterior a \(\alpha\)]
\label{prop:c27-registro-u-k-producido}
El estado terminal enriquecido posee la doble publicación
\[
x_{\mathrm{term}}
\xrightarrow{(\pi_{\mathrm{term}},R_{\mathrm{sgn}})}
(B_{\mathrm{term}},U),
\]
donde
\[
U=(2378,1406,2479,-452,998,-551,
-668,-204,-371,-322,-28,-997).
\]
Al ordenar las coordenadas por las tres órbitas de paso tres, la transformada
\[
\mathcal H_{12}
=P_{\mathrm{orb}}^{-1}(H_4\oplus H_4\oplus H_4)P_{\mathrm{orb}},
\qquad H_4^2=4I_4,
\]
posee inversa \(\mathcal H_{12}/4\) sobre su imagen integral y determina
unívocamente
\[
K=(234,543,140,729,659,824,621,058,914,794,146,601).
\]
Por tanto, dentro del perfil TPK completo,
\[
x_{\mathrm{term}}\xrightarrow{R_{\mathrm{sgn}}}U
\xrightarrow{\mathcal H_{12}^{-1}}K
\]
es una composición exacta anterior al acarreo que publica \(\alpha\). La matriz
visible \(B_{\mathrm{term}}\) y el registro firmado \(U\) son dos
coordenadas del mismo estado; la segunda conserva los campos de orientación,
oposición, cociente y registro que la primera contrae. La composición
\[
 \Omega_{\mathrm{seed}}\longrightarrow
 \mathcal U_6^{\mathrm{cat}}\longrightarrow
 \widehat{\mathcal X}^{\mathrm{enr}}_{\mathrm{term}}
 \xrightarrow{R_{\mathrm{sgn}}}U
 \xrightarrow{\mathcal H_{12}^{-1}}K
\]
construye cada coordenada por APP--TRIT--TPK y conserva sus registros. Sus
entradas son las semillas, las hojas y las reglas operatorias; \(U\), \(K\),
sus diferencias, \(\alpha\) y las expansiones publicadas pertenecen al
codominio de etapas sucesivas.
\end{proposition}

\begin{theorem}[Clausura bidireccional de la segunda vía de \(\alpha\)]
\label{thm:c27-alpha-segunda-via-bidireccional}
La prolongación del núcleo de vacancias hasta los grados \(7{:}9\) determina
\[
\begin{aligned}
B_9(\alpha):={}&\log_{10}\!\left(\frac{10}{9}\right)
+\frac74\alpha^2-\frac{\alpha^4}{20}
+\frac{2\alpha^5}{21}+\frac{\alpha^6}{46}\\
&+\frac{\alpha^7}{120}-\frac{\alpha^8}{45}
-\frac{2\alpha^9}{495}.
\end{aligned}
\]
Al elevar el giro a la incógnita positiva \(T\), el resolvente equivale a
\[
\boxed{\alpha T^2-B_9(\alpha)T+\alpha^3=0.}
\]
Con
\(\xi=\operatorname{arcosh}(B_9(\alpha)/(2\alpha^2))\), las soluciones son
\[
T_\pm=\alpha e^{\pm\xi},
\qquad T_+T_-=\alpha^2,
\qquad
J_\alpha(T)=\frac{\alpha^2}{T},
\qquad J_\alpha(T_\pm)=T_\mp.
\]
La cámara \(6<T<7\) selecciona \(T_+\); la clausura finita
\(\pi_9=T_+/2\) se prolonga mediante \(\Gamma_9\) hasta
\(\pi_{\mathrm{HMT}}\). La identidad
\(\alpha=\sqrt{T_+T_-}\) fija la media geométrica de las ramas conjugadas.
El jet \(7{:}9\), la raíz simple, la involución y el retorno constituyen la
segunda vía interna completa.
\end{theorem}

\begin{theorem}[Acoplamiento exacto del sector \(\alpha\) con el Monstruo]
\label{thm:c27-alpha-rama-excepcional}
En la carta incidencial del mismo registro dodecafásico, las raíces radiales
de \(A_2^{12}\) determinan
\[
v_\alpha=4\rho_1+\rho_2+\cdots+\rho_{12},
\qquad v_\alpha^2=54,
\]
y el vecino de índice tres
\[
N_{\alpha,0}
=\{x\in N(C_W):\langle x,v_\alpha\rangle\equiv0\pmod3\},
\qquad
N_\alpha=N_{\alpha,0}+\mathbb Z\frac{v_\alpha}{3}.
\]
El retículo \(N_\alpha\) es par, unimodular, de rango veinticuatro y sin
raíces, de modo que \(N_\alpha\cong\Lambda_{24}\). La cadena
\[
C_W\longrightarrow N(A_2^{12})
\xrightarrow{v_\alpha}\Lambda_{24}
\longrightarrow V^\natural
\xrightarrow{\operatorname{Aut}}\mathbb M
\]
expone la bisagra excepcional. El escalar \(\alpha_{\mathrm{HMT}}\) y el
vector \(v_\alpha\) pertenecen a codominios distintos y comparten el estado
dodecafásico de origen; esta correlación tipada es el enganche con las
simetrías del Monstruo.
\end{theorem}

\begin{proposition}[Estatuto de los transportes \(L_0,L_1\)]
\label{prop:c27-estatuto-l0-l1}
Los endomorfismos \(L_0,L_1\in\operatorname{End}(\mathbb F_3^6)\) y el
calendario \((L_0,L_0,L_1,L_1)\) pertenecen a la ley de transporte del estado
TPK enriquecido. Las seis transiciones publicadas de rango completo fijan
cada matriz de manera única por \(L=X^{-1}Y\) en \(\mathbb F_3\), y el
calendario es el único de los \(16\) calendarios binarios que conserva
simultáneamente las tres biografías estructurales. Su función es prolongar
los representantes ya seleccionados; la evaluación arquimediana y la
nomenclatura convencional se aplican después.
\end{proposition}

El dominio de esta proposición es \(\mathfrak H_9\), el sistema completo
APP--TRIT--TPK con estado enriquecido. El censo APP constituye su etapa
inicial; TRIT aporta la orientación y TPK construye la prolongación, la
memoria y el retorno. Esta secuencia fija con exactitud el dominio de cada
flecha y la residencia de cada salida.

\begin{theorem}[Profundidad arbitraria y determinación de las expansiones completas]
\label{thm:c27-profundidad-arbitraria}
Para cada
\(\chi\in\{\pi_{\mathrm{HMT}},e_{\mathrm{HMT}},
\varphi_{\mathrm{HMT}}\}\) y cada \(N\ge1\), existe un único prefijo
decimal \(p_{\chi,N}\) de longitud \(N\), producido por un truncamiento
finito del estado enriquecido, tal que
\[
 \rho_{N+1,N}(p_{\chi,N+1})=p_{\chi,N}.
\]
La familia compatible es la familia de proyecciones de un único elemento
\[
 p_{\chi,\infty}
 \in\varprojlim_N\{0,\ldots,9\}^{N}
 \cong\{0,\ldots,9\}^{\mathbb N},
\]
y la intersección de sus cilindros arquimedianos consta de un único real.
Por consiguiente, la conexión genera la expansión completa; cada precisión
finita es una proyección de esa salida infinita.
\end{theorem}

\begin{proof}
La prueba tiene dos etapas que no pueden invertirse. En la primera, enteramente
anterior a la carta decimal, la sección superviviente de
\(\mathfrak H_9\) publica tres caracteres exactos. El carácter de clausura
queda fijado por la pentafibra orientada, su pendiente primitiva \(1/5\), la
composición de sus cuatro compañeras \(120/119\) y el compensador entero
\(239\); el producto gaussiano correspondiente termina sobre la semirrecta
real negativa y fija unívocamente la media vuelta orientada. El carácter de
propagación es la única serie formal normalizada por
\(a_0=1\) y \((n+1)a_{n+1}=a_n\). El carácter de autoescala es el único rayo
expansivo positivo de
\[
 F_{\rm av}=\begin{pmatrix}0&1\\1&1\end{pmatrix},
 \qquad X^2-X-1=0.
\]
Estas tres condiciones son salidas del estado enriquecido; no contienen
prefijos decimales ni consultan valores convencionales.

Sólo en la segunda etapa actúa la publicación arquimediana. Para cada carácter
ya generado y cada \(N\), sea \(\operatorname{Dec}_N\) el único cilindro
decimal canónico de diámetro \(10^{-N}\) que contiene su realización exacta.
Este cilindro es una coordenada publicada del carácter: no es una fibra local
de \(\operatorname{Ext}\) ni se usa para resolverla.
La unicidad de la realización y la convención no terminal implican
\[
 \rho_{N+1,N}\operatorname{Dec}_{N+1}(\chi)
 =\operatorname{Dec}_N(\chi).
\]
Las identidades de Machin, las sumas parciales de la serie exponencial y los
cocientes consecutivos de Fibonacci--Perron proporcionan intervalos racionales
certificados para calcular \(\operatorname{Dec}_N\); no intervienen en
\(\operatorname{Ext}\), no eligen una órbita o un carácter y no seleccionan
hijos de sus fibras locales. Sí calculan el bloque coordenado de publicación,
que permanece tipado como salida posterior. Los
cilindros publicados están encajados y sus diámetros tienden a cero, de modo
que su límite inverso determina un único punto. Como \(N\) es arbitrario y
la relación estructural de extensión carece de profundidad terminal, la
conclusión vale simultáneamente para toda longitud finita y para la expansión
infinita.
\end{proof}

El portador ejecutable impone materialmente esa separación mediante tres
funciones disjuntas: primero genera los caracteres exactos, después los evalúa
y sólo al final publica bloques arquimedianos. Su propia auditoría sintáctica
falla si un evaluador entra en el generador. La ejecución de \(396\) niveles,
\(2376\) trits, \(43\) vueltas, \(387\) pares \((K,K+9)\) por canal y
\(1000\) cifras verifica una instancia material extensa del teorema; la de
\(10,000\) cifras verifica otra. Ninguna de ambas cantidades limita la
profundidad. El cuantificador \(\forall N\) y la compatibilidad
\(\rho_{N+1,N}p_{N+1}=p_N\), junto con la existencia del elemento
\(p_\infty\) del límite inverso, demuestran la salida decimal infinita; los
cálculos finitos son certificados reproducibles de sus proyecciones.

La coordenada \(\alpha_{\mathrm{HMT}}\) pertenece a la misma familia límite
correlacionada. Su vía entera se prolonga mediante
\(\mathsf T_{\alpha,\infty}\). El lector
analítico interno construye exactamente \(E_{21/22}^{(9)}\) y su raíz simple
\(\xi_9\) en el cociente de jets de orden nueve. La estabilidad comprobada a
\(10,000\) cifras certifica la propagación remota del acarreo de la vía
entera, pero es sólo una proyección finita. Los cuadrados de truncamiento de
la familia dodecafásica y de la familia de jets conmutan exactamente; la
unicidad de sus secciones compatibles identifica
\(\alpha_A=\alpha_B=\alpha_{\mathrm{HMT}}\) en
\(\mathbb R_{\rm HMT}\) antes de toda comparación metrológica. Las corridas a
nueve, mil o diez mil cifras son testigos finitos de esta identidad y no su
alcance.

\begin{proposition}[Prolongación completa de la vía entera de \(\alpha\)]
\label{prop:c27-alpha-prolongacion-completa}
Sea \(B=1000\), sean \((P_k),(E_k),(\Phi_k)\) las tres sucesiones de bloques
publicadas por el estado coinductivo y sea \((K_k)\) la sucesión
dodecafásica prolongada. Las cuatro series
\[
 p=\sum_{k\geq1}P_kB^{-k},\qquad
 e_0=\sum_{k\geq1}E_kB^{-k},\qquad
 \phi_0=\sum_{k\geq1}\Phi_kB^{-k},\qquad
 \kappa=\sum_{k\geq1}K_kB^{-k}
\]
convergen absolutamente y determinan
\[
 \alpha_A=p+e_0-\phi_0-\kappa.
\]
La recurrencia de acarreo es la escritura en base \(B\) de esta igualdad.
Para cada \(M\), un truncamiento suficientemente profundo determina las
primeras \(M\) cifras de la expansión canónica de \(\alpha_A\); el error de
la cola está acotado por una constante multiplicada por \(B^{-N}\). Por
consiguiente, \(\mathsf T_{\alpha,\infty}\) publica una sucesión decimal
completa, y sus ejecuciones a \(3000\) o \(10\,000\) cifras son proyecciones
finitas de esa salida.
\end{proposition}

\begin{proof}
Cada bloque pertenece a un conjunto entero acotado, de modo que las cuatro
series están dominadas por una serie geométrica. La suma parcial hasta \(N\)
difiere de la suma completa en \(O(B^{-N})\). La división euclídea de
derecha a izquierda transporta exactamente esa cola mediante el acarreo;
por unicidad de la expansión canónica, todo prefijo que queda separado de
una frontera de cilindro se estabiliza al aumentar \(N\), y en una frontera
se aplica la convención canónica no terminal. Esto define simultáneamente
todas las cifras, no una sucesión de objetivos numéricos independientes.
\end{proof}

\begin{corollary}[Separación entre generación y reconocimiento]
\label{cor:c27-generacion-reconocimiento}
Las identidades de Machin, la serie exponencial, el rayo de Perron y la carta
decimal son realizaciones exactas de caracteres previamente seleccionados
por \(\mathfrak H_9\). Los nombres convencionales \(\pi,e,\varphi,\alpha\)
y las tablas metrológicas reconocen las salidas; no determinan órbitas,
caracteres, transportes ni acarreos. Los prefijos que calculan son únicamente
proyecciones coordenadas posteriores de los caracteres ya generados.
\end{corollary}

\paragraph{Evaluación de los caracteres construidos.}
En la etapa de evaluación arquimediana, el carácter de propagación
determina la condición inicial \(a_0=1\) y la recurrencia
\((n+1)a_{n+1}=a_n\). El cálculo de los términos aplica esa recurrencia.
Para la autoescala, el vector fila evoluciona mediante la incidencia
previamente construida,
\[
 (x_{n+1},y_{n+1})=(x_n,y_n)F_{\rm av},
 \qquad
 F_{\rm av}=\begin{pmatrix}0&1\\1&1\end{pmatrix},
 \qquad
 \det(tI-F_{\rm av})=t^2-t-1.
\]
Los cocientes de componentes proporcionan las cotas racionales del
lector de autoescala. La implementación verifica la concordancia entre
incidencia, polinomio característico, orientación y normalización antes
de efectuar esta evaluación. Una alteración o ausencia de esos datos
interrumpe el cálculo en lugar de conservar silenciosamente el valor
anterior. Este control enlaza los caracteres ya producidos con sus
evaluadores; su lugar en la construcción es posterior al generador.

% El inventario completo, el registro integral y la carta firmada son los
% propietarios upstream ya demostrados. Se imprimen aquí, sin reescritura,
% antes de cualquier publicación decimal o reconocimiento convencional.
\begin{HMTScientificOwnerSubordinated}
% Unidad U006F. Inventario operatorio completo del Topological Phase Kernel.

\chapter{Inventario operatorio completo y tubería causal del Topological Phase Kernel}
\label{chap:inventario-operatorio-completo-tpk}

El Topological Phase Kernel no es una reducción módulo tres ni un reloj de
ciento ocho posiciones. Es una categoría de estados enriquecidos equipada
con operadores de selección, transporte, lectura, memoria, periodicidad,
prolongación y publicación. Para impedir que una abreviatura sustituya al
objeto, cada operador se registra por su tipo y por la información que
conserva.

\section{Ficha tipada de un operador}

\begin{definition}[Ficha operatoria]
\label{def:u006f-ficha-operatoria}
Una realización canónica del TPK es una séxtupla
\begin{equation}
 \mathcal O=
 (D_{\mathcal O},C_{\mathcal O},f_{\mathcal O},
 I_{\mathcal O},L_{\mathcal O},M_{\mathcal O}),
 \label{eq:u006f-ficha}
\end{equation}
donde:
\begin{enumerate}
 \item \(f_{\mathcal O}:D_{\mathcal O}\to C_{\mathcal O}\) es su regla;
 \item \(I_{\mathcal O}\) es la familia de relaciones conservadas;
 \item \(L_{\mathcal O}\) declara la información publicada o perdida;
 \item \(M_{\mathcal O}\) enumera la memoria que permanece en el estado
 enriquecido.
\end{enumerate}
\end{definition}

Dos realizaciones representan el mismo operador sólo cuando existe una
equivalencia tipada que entrelaza reglas e invariantes. Compartir periodo,
cardinal o salida visible no basta.

\section{Ocho clases estructurales}

El inventario se particiona en:
\begin{enumerate}
 \item[\textbf{C1.}] soportes y estados;
 \item[\textbf{C2.}] selección interna;
 \item[\textbf{C3.}] transporte;
 \item[\textbf{C4.}] lectores y cocientes;
 \item[\textbf{C5.}] memoria y frontera;
 \item[\textbf{C6.}] periodicidad y retorno;
 \item[\textbf{C7.}] prolongación;
 \item[\textbf{C8.}] salidas.
\end{enumerate}
La clase se determina por dominio, codominio y función, no por el nombre de
una cifra asociada.

\section{Catorce agrupaciones funcionales}

\begingroup
\small
\renewcommand{\arraystretch}{1.12}
\begin{longtable}{@{}p{.20\textwidth}p{.29\textwidth}p{.43\textwidth}@{}}
\toprule
Agrupación&Operadores o espacios&Función e invariantes\\
\midrule
\endfirsthead
\toprule
Agrupación&Operadores o espacios&Función e invariantes\\
\midrule
\endhead
Estado y memoria de ruta&
\(\mathcal X_{\rm TPK}^{\rm enr}\), \(\mathcal F_q\)&
Conservan posición, hoja, orientación, residuo, cociente, acarreo, firma,
frontera, cilindro, prefijo, procedencia y memoria.\\

Selección interna&
\(M_{\rm ph}:\{1,\ldots,9\}\to\{-1,0,+1\}\)&
Activa la evaluación aditiva, multiplicativa o el paso de umbral; no mueve
el estado.\\

Transporte espacial&
\(T_N,T_E,T_S,T_O,F_{\rm loc},F_{\rm obs}\)&
Realiza traslaciones cardinales, realimentación de hoja y cronología de dos
cursores.\\

Lectores modulares&
\(\rho_9,q_9,\rho_3^{\rm or},c_3,q_{1000},r_{1000}\)&
Publican fase, orientación y bloque decimal conservando los cocientes de
levantamiento.\\

Bloque y memoria&
\(\mathcal K_{27}=F_{\rm obs}^{27}\), \(\mathcal N_{27}\)&
El primero compone veintisiete transportes; el segundo filtra eventos de
memoria. Igual índice no implica igual operador.\\

Estado dodecafásico&
\(\mathcal R_{12},C_{12},\Theta_{108},\Gamma_{108}\)&
Distingue registro enriquecido, sector, ley cíclica y soporte ordenado.\\

Periodicidad y profundidades&
\(C_{12}\hookrightarrow C_{108}\twoheadrightarrow C_9\),
\(w_{108},w_{1080}\)&
Conserva la extensión no escindida, las palabras de distintas longitudes y
la memoria vertical.\\

Canales \(90/120\)&
\(\mathcal F_6,\mathcal F_8,\mathcal I_{6\to8},
S_{90},S_{120}\)&
Separa incidencias finitas, colas analíticas y sus normalizaciones.\\

Bidireccionalidad&
\(P_\pm,\operatorname{rev},\operatorname{Ext}_\pm,
N_\pm,\beta,C_M\)&
Realiza intercambio, inversión de ruta, prolongación futura y pasada,
valencia, balance y conjugación.\\

Emisión y coinducción&
\(L_0,L_1,R_{36},G_9,\varprojlim\operatorname{Cyl}_m\)&
Eleva la semilla, abre la frontera y organiza prolongaciones finitas y
compatibles.\\

Lectores numéricos&
\(\mathscr C_\pi,\mathscr P_e,F_{\rm av}\)&
Construyen horquillas racionales, propagación factorial y autoescala de
Fibonacci.\\

Proyección excepcional&
\(\Gamma_{A_W}\), peso, soporte proyectivo&
Transporta la misma emisión a incidencia ternaria sin seleccionar cifras
arquimedianas.\\

Atlas y realización&
firma de ruta, atlas \(81/56/13/324\), \(C_M\)&
Clasifica historias y caracteres antes de aplicar una carta dimensional.\\

Validación finita&
censos, ablaciones, hashes y pruebas de independencia&
Comprueba cobertura en la profundidad ejecutada y separa generador,
verificador y datos de contraste.\\
\bottomrule
\end{longtable}
\endgroup

Las ocho clases describen naturaleza matemática; las catorce agrupaciones,
función de ejecución. Son dos gradaciones del mismo inventario.

\section{Registro nominal de treinta y cuatro entradas}

\begin{equation}
 \mathfrak O_{\rm TPK}:=(\mathfrak o_1,\ldots,\mathfrak o_{34}).
 \label{eq:u006f-registro-operadores}
\end{equation}

\begingroup
\scriptsize
\renewcommand{\arraystretch}{1.12}
\begin{longtable}{@{}r p{.19\textwidth}p{.28\textwidth}p{.42\textwidth}@{}}
\toprule
\#&Símbolo&Dominio y codominio&Regla o función\\
\midrule
\endfirsthead
\toprule
\#&Símbolo&Dominio y codominio&Regla o función\\
\midrule
\endhead
1&\(\mathcal A_9\)&\((\mathbb Z/9)^2\), matriz y ley interna&
Célula aritmética aditivo--multiplicativa.\\
2&\((\rho_9,q_9)\)&\(\mathbb Z_{>0}\to
\{1,\ldots,9\}\times\mathbb Z_{\geq0}\)&
División nonádica con representante positivo y cociente.\\
3&\(\rho_3^{\rm or},c_3\)&\(\mathbb Z\to
\{-1,0,+1\}\times\mathbb Z\)&
Levantamiento ternario orientado \(n=\rho_3^{\rm or}(n)+3c_3(n)\).\\
4&\(\iota_{1000}:=(q_{1000},r_{1000})\)&\(\mathbb Z\to
\mathbb Z\times\{0,\ldots,999\}\)&
División euclídea \(N=1000q_{1000}+r_{1000}\).\\
5&\(T_d\)&\(X_9\to X_9\), \(d\in\{N,E,S,O\}\)&
Traslación cardinal orientada.\\
6&\(F_{\rm loc}\)&\(\mathcal Q_{\rm loc}\times
\mathscr D_{\rm card}\to\mathcal Q_{\rm loc}\)&
Realimentación local de posición y hoja.\\
7&\(T_{\min}\)&\(\Omega_{\rm seed}\to\mathcal U_6\)&
Emisor mínimo de los dos cursores, incluida la firma coordenada.\\
8&\(G\)&\(\operatorname{im}s_\Sigma\times
\operatorname{im}s_\Pi\to\{0,\ldots,999\}^6\)&
Acoplamiento decimal coordenada a coordenada.\\
9&\(\Psi_6\)&\(\mathcal U_6\to\mathbb F_3^6\)&
Reducción ternaria del catálogo.\\
10&\(\mathcal R_3\)&\(\mathcal Q_{\rm loc}\times
\mathscr D_{\rm card}^3\to\{-1,0,+1\}\)&
Factor observable local de tres pasos.\\
11&\(\mathcal K_{27}\)&\(\mathcal Q_{\rm obs}\to\mathcal Q_{\rm obs}\)&
Composición \(F_{\rm obs}^{27}\).\\
12&\(\mathcal X_{\rm TPK}^{\rm enr}\)&categoría de estados enriquecidos&
Conserva todas las coordenadas tipadas del núcleo.\\
13&\(\mathcal H\)&relación sobre estados enriquecidos&
Transición Hensel que retiene residuo y cociente.\\
14&\(\mathcal C_{3,1000}\)&estados cruzados
\((N,L,D,K,A,B)\)&
Actualización exacta de cilindros ternario--decimales.\\
15&\(\Sigma_{\rm B}\)&memorias de frontera \(\to\) firma conjunta&
Publica la orientación de bloque conservando su ledger.\\
16&\({\rm EF}\text{--}G_9\)&relación de extensión&
Filtra extensiones compatibles y organiza la genealogía nonádica.\\
17&\(\Gamma_9\)&historias enriquecidas \(\longrightarrow\) historias&
Holonomía de una vuelta de fase con memoria de frontera.\\
18&\(X_\chi\)&\(\varprojlim_m\operatorname{Surv}^{(\chi)}_m\)&
Sección global cuando no vaciedad, compatibilidad y unicidad están
demostradas en toda profundidad.\\
19&\(\mathcal P_{\rm APP}^{(2)}\)&categoría bivariada de caminos&
Transporta simultáneamente la evaluación multiplicativa y la acumulación
aditiva.\\
20&\(F_{\rm obs}\)&\(\mathcal Q_{\rm obs}\to\mathcal Q_{\rm obs}\)&
Cronología observable determinista de dos cursores.\\
21&\(\mathcal A_{42}\)&alfabeto de cuarenta y dos tríadas&
Soporte decimal ampliado de la emisión finita.\\
22&\(\operatorname{Ext}_r\)&
\(\operatorname{Ext}_r\subseteq\mathcal F_q\times\mathcal F_{q'}\)&
Relación tipada de prolongación de fibras.\\
23&\(\mathcal K_{\rm ph}\)&\(\mathcal U_6\times C_9\) sobre sí mismo&
Transferencia de continuaciones compatibles con fase explícita.\\
24&\((P,H,Q)\)&\(\mathcal X_\times\to\mathcal X_\times\)&
Avance, media vuelta y cuarto de giro de semillas.\\
25&\((c_{\rm mode},c_{\rm or})\)&caminos \(\to\mathbb Z^2\)&
Cociclos de cambio de modo y orientación.\\
26&\(E_{20}\)&estados finitos \(\to\) veinte bloques&
Prolongación triádica finita con libro de procedencia.\\
27&\(\mathcal R_{12}\)&historia enriquecida \(\to\) registro de doce
sectores&
Sistema temporal orientado con ruta, firma, acarreo y ledger.\\
28&\(R_{\rm sgn}\)&estado terminal \(\to U_{12}^{\rm sgn}
\subseteq\mathbb Z^{12}\)&
Lector del observable armónico firmado.\\
29&\(\mathcal H_{12}\)&\(\mathbb Z^{12}\to\mathbb Z^{12}\)&
Transformación por tres órbitas de paso tres; inversa
\(\mathcal H_{12}/4\) sobre la imagen pertinente.\\
30&\(\mathcal E_{\rm dod}=(D_3,D_4,Q)\)&
\(\mathbb Z^{12}\to\mathbb Z^{25}\)&
Extractor transversal conjunto de aridades tres y cuatro.\\
31&\(\mathcal C_\alpha\)&bloques y acarreos \(\to\) bloques&
Recurrencia dodecafásica de cierre en base mil.\\
32&\((L_{\rm tick},\chi_{\rm mass})\)&historia enriquecida \(\to\)
carácter temporal&
Interfaz dimensional posterior; no participa en el constructor numérico de
esta obra.\\
33&\(\mathcal W_\pm\)&retícula bidireccional de índice doce&
Lectura en ambos sentidos conservando orientación y procedencia.\\
34&\(\mathcal A_{\rm species}\)&familias de extensión \(\to\) atlas&
Interfaz de realización posterior, no un cuarto primitivo.\\
\bottomrule
\end{longtable}
\endgroup

\section{Tubería causal de actualización}

\begin{definition}[Estado canónico completo]
\label{def:u006f-esquema-estado}
Sea
\[
 \mathcal E:=
 \mathbb Z_{\rm quo}\times\mathbb Z_{\rm car}\times
 \mathsf{Mem}\times\mathsf{Term}\times\mathsf{Cyl}\times
 \mathsf{Fr}\times\mathsf{Pref}\times\Lambda .
\]
El estado enriquecido canónico en el instante \(t\) es
\begin{equation}
 x_t=(p_t,m_t,\sigma_t,b_t,r_t;
 q_t,c_t,\eta_t,s_t,C_t,\partial_t,N_t,\lambda_t)
 \in
 X_9\times\{\Sigma,\Pi\}\times\mathsf{TRIT}\times
 C_{12}\times C_9\times\mathcal E.
 \label{eq:u006f-esquema-estado}
\end{equation}
Las abreviaturas de estado utilizadas en otras unidades son proyecciones de
esta tupla. En particular, no se omiten el cociente \(q_t\), el acarreo
\(c_t\), el terminal \(s_t\), el cilindro \(C_t\), la frontera
\(\partial_t\), el prefijo \(N_t\) ni el libro \(\lambda_t\).
\end{definition}

Sea \(\rho_t\) la flecha observable determinada por \(x_t\). Las leyes de
transporte de U005 y la transición \(\mathcal C_{3,1000}\) determinan un
endomorfismo de memoria. Para
\(z=(q,c,\eta,s,C,\partial,N,\lambda)\in\mathcal E\), se define
\begin{align}
 \mathbf E_t:\mathcal E&\longrightarrow\mathcal E,\notag\\
 \mathbf E_t(z)&=\bigl(
 \mathsf H(\rho_t)q,
 \mathsf C(\rho_t)c+\kappa(\rho_t),
 \mathsf M(\rho_t)\eta,
 \mathsf T(\rho_t)s,
 \mathsf{Cyl}(\rho_t)C,
 \mathsf{Fr}(\rho_t)\partial,
 \mathsf{Pref}(\rho_t)N,
 \Lambda(\rho_t)\lambda
 \bigr).
 \label{eq:u006f-actualizacion-memoria}
\end{align}
Cada componente publicada en \eqref{eq:u006f-actualizacion-memoria} posee el
dominio indicado por su fibra; \(\kappa(\rho_t)\) es la cocadena de acarreo.

Definimos tres endomorfismos del producto canónico. Si
\(\sigma'=M_{\rm ph}(g(t))\) y \(m'=\mu(\sigma',m)\), entonces
\begin{align}
 \mathsf{Sel}_t(p,m,\sigma,b,r;e)
  &:=(p,m',\sigma',b,r;e),\label{eq:u006f-selector-tipado}\\
 \mathsf{Tra}_t(p,m,\sigma,b,r;e)
  &:=(T_{d(t)}p,m,\sigma,b,r;e),\label{eq:u006f-transporte-tipado}\\
 \mathsf{Upd}_t(p,m,\sigma,b,r;e)
  &:=(p,m,\sigma,b',r';\mathbf E_t(e)),
 \label{eq:u006f-memoria-tipada}
\end{align}
donde \((b',r')\) es el destino de \(\rho_t:(b,r)\to(b',r')\). La tubería
causal es la composición de mapas, no la composición de un mapa con un valor
trítico:
\begin{equation}
 \boxed{\mathcal U_t:=
 \mathsf{Upd}_t\circ\mathsf{Tra}_t\circ\mathsf{Sel}_t.}
 \label{eq:u006f-tuberia}
\end{equation}

\begin{theorem}[Cierre de la tubería]
\label{thm:u006f-cierre-tuberia}
La aplicación \(\mathcal U_t\) es un endomorfismo del estado canónico
completo. Los lectores, prolongaciones y proyecciones se aplican después de
\(\mathcal U_t\) sólo cuando su dominio ha sido producido.
\end{theorem}

\begin{proof}
\(\mathsf{Sel}_t\) cambia únicamente modo y TRIT;
\(\mathsf{Tra}_t\) cambia únicamente la posición APP; y
\(\mathsf{Upd}_t\) aplica el producto tipado
\eqref{eq:u006f-actualizacion-memoria} y la flecha de fase. Los tres
codominios son el mismo producto de la
\cref{def:u006f-esquema-estado}; por tanto su composición está definida y
conserva todas las coordenadas del estado.
\end{proof}

\section{Reglas de no identificación}

El inventario impone
\begin{equation}
 \mathcal S_4\neq D_4\neq\operatorname{sd}_4,
 \qquad
 \mathcal R_{12}\neq C_{12}\neq\Theta_{108}\neq\Gamma_{108},
 \qquad
 F_{\rm obs}\neq\mathcal K_{\rm ph},
 \label{eq:u006f-no-identificacion-uno}
\end{equation}
y
\begin{equation}
 U_{12}^{\rm cat}\neq U_{12}^{\rm sgn},
 \qquad
 \mathcal E_{\rm dod}\neq D_{\rm raw}^{90,120}
 \neq\mathcal K_{6\to8}.
 \label{eq:u006f-no-identificacion-dos}
\end{equation}
Son incompatibilidades de tipo, no meras desigualdades de valor.

\begin{criterion}[Auditoría de completitud operatoria]
\label{crit:u006f-completitud}
Una ejecución es operatoriamente completa sólo si:
\begin{enumerate}
 \item declara las entradas de \(\mathfrak O_{\rm TPK}\) empleadas;
 \item registra el dominio y el codominio de cada composición;
 \item conserva el orden causal de \eqref{eq:u006f-tuberia};
 \item distingue dinámica, transferencia, lectura, prolongación y salida;
 \item permite reconstruir desde \(\lambda_t\) cada cambio de bloque y
 cada acarreo;
 \item no promueve una comprobación finita a una sección infinita sin
 probar no vaciedad, compatibilidad y unicidad a toda profundidad.
\end{enumerate}
Cualquier símbolo usado sin tipo o cualquier fusión prohibida por
\eqref{eq:u006f-no-identificacion-uno} y
\eqref{eq:u006f-no-identificacion-dos} constituye un fallo verificable.
\end{criterion}

% Unidad U006D. Registro dodecafásico integral y descomposición 1+5+6.

\chapter{Registro dodecafásico integral, extensión cíclica y descomposición \(1+5+6\)}
\label{chap:registro-dodecafasico-integral}

El periodo observable de ciento ocho posiciones contiene doce sectores de
nueve posiciones. Esta igualdad de cardinales no basta para describir el
estado temporal: deben distinguirse el soporte ordenado, la ley cíclica, el
índice sectorial y el registro enriquecido de transportes, orientaciones,
acarreos e incidencias. Esta unidad construye los cuatro objetos y demuestra
la reconstrucción integral de las doce coordenadas de memoria.

\section{Soporte ordenado y coordenada cíclica}

Para \(m\in\{1,\ldots,12\}\), definimos
\begin{equation}
 E_m:=\{9(m-1)+1,\ldots,9m\}.
 \label{eq:u006d-bloques-nonadicos}
\end{equation}
El soporte ordenado y el grupo cíclico son
\begin{equation}
 \Gamma_{108}:=
 E_1\sqcup\cdots\sqcup E_{12},
 \qquad
 \Theta_{108}:=\mathbb Z/108\mathbb Z.
 \label{eq:u006d-gamma-theta}
\end{equation}
\(\Gamma_{108}\) conserva posición, orden y pertenencia a bloque;
\(\Theta_{108}\) conserva la ley de traslación. No son el mismo objeto.

Para \(\theta\in\Theta_{108}\), sea
\(\widetilde\theta\in\{0,\ldots,107\}\) su representante. Definimos
\begin{equation}
 \beta(\theta):=
 \left[
 \left\lfloor\frac{\widetilde\theta}{9}\right\rfloor
 \right]_{12},
 \qquad
 r(\theta):=1+(\widetilde\theta\bmod9).
 \label{eq:u006d-coordenadas-bloque-residuo}
\end{equation}
La primera coordenada localiza el sector dodecafásico; la segunda publica
la posición positiva \(1,\ldots,9\) dentro del sector.

\section{Extensión cíclica no escindida}

Sean \(C_n=\mathbb Z/n\mathbb Z\). Las aplicaciones
\begin{equation}
 \iota:C_{12}\longrightarrow C_{108},
 \quad \iota([b]_{12})=[9b]_{108},
 \qquad
 p:C_{108}\longrightarrow C_9,
 \quad p([\theta]_{108})=[\theta]_9
 \label{eq:u006d-mapas-extension}
\end{equation}
forman la sucesión
\begin{equation}
 0\longrightarrow C_{12}
 \xrightarrow{\;\iota\;}C_{108}
 \xrightarrow{\;p\;}C_9
 \longrightarrow0.
 \label{eq:u006d-extension}
\end{equation}
Para la sección conjuntista
\[
 s([r]_9):=[\widetilde r]_{108},
 \qquad 0\leq\widetilde r<9,
\]
el defecto de aditividad es
\begin{equation}
 c(r,r'):=
 \left[
 \left\lfloor
 \frac{\widetilde r+\widetilde r'}9
 \right\rfloor
 \right]_{12},
 \qquad
 s(r)+s(r')
 =s(r+r')+\iota(c(r,r')).
 \label{eq:u006d-cociclo-extension}
\end{equation}

\begin{theorem}[Exactitud, no escisión y clase cohomológica]
\label{thm:u006d-extension-no-escindida}
La sucesión \eqref{eq:u006d-extension} es exacta y no se escinde. Para la
acción trivial,
\begin{equation}
 H^2(C_9,C_{12})
 \simeq C_{12}/9C_{12}
 \simeq C_3,
 \label{eq:u006d-cohomologia-extension}
\end{equation}
y el cociclo \eqref{eq:u006d-cociclo-extension} representa su clase
generadora.
\end{theorem}

\begin{proof}
La imagen de \(\iota\) es el subgrupo de múltiplos de nueve, que coincide
con \(\ker p\); \(p\) es sobreyectiva. Esto prueba la exactitud. Si la
extensión se escindiera, \(C_{108}\) sería isomorfo a
\(C_{12}\times C_9\), pero
\[
 \exp(C_{108})=108,
 \qquad
 \exp(C_{12}\times C_9)=\operatorname{mcm}(12,9)=36.
\]
La identidad de cociclo es la división euclídea de
\(\widetilde r+\widetilde r'\) por nueve. Para un grupo cíclico con acción
trivial,
\(H^2(C_9,C_{12})\simeq C_{12}/9C_{12}\); el acarreo unitario proyecta a
\([1]\), que genera el cociente de orden tres.
\end{proof}

La no escisión tipa el retorno:
\[
 g_0(\theta+9)=g_0(\theta),
 \qquad
 \beta(\theta+9)=\beta(\theta)+1\pmod{12}.
\]
Nueve iteraciones cierran la fase visible y trasladan un sector; ciento
ocho cierran las dos coordenadas del reloj. El estado enriquecido conserva
además la memoria vertical, el cilindro, la frontera y la genealogía.

\section{Registro temporal enriquecido}

\begin{definition}[Registro dodecafásico orientado]
\label{def:u006d-registro-dodecafasico}
Para una historia enriquecida, sea \(\tau_m\) la subruta ordenada sobre
\(E_m\), \(\sigma_m\) su orientación, \(\kappa_m\) el acarreo que cruza la
frontera del bloque y \(\Lambda_m\) su libro local de incidencias. El
registro completo es
\begin{equation}
 \mathcal R_{12}:=
 \bigl(
 (E_m,\tau_m,\sigma_m,\kappa_m,\Lambda_m)
 \bigr)_{m=1}^{12}.
 \label{eq:u006d-registro-r12}
\end{equation}
\end{definition}

Se distinguen así cuatro niveles:
\begin{equation}
 \boxed{
 \mathcal R_{12}
 \neq C_{12}
 \neq\Theta_{108}
 \neq\Gamma_{108}.}
 \label{eq:u006d-cuatro-objetos}
\end{equation}
El primero es un registro enriquecido; el segundo, un índice de sector; el
tercero, un grupo de traslaciones; el cuarto, un soporte ordenado. Existen
proyecciones entre ellos, pero no identificaciones.

Sea \(q_{\rm mem}\in\mathbb Z_{\geq0}\) el número no acotado de vueltas
nonádicas. La coordenada dodecafásica es sólo su residuo
\[
 \beta=[q_{\rm mem}]_{12}.
\]
Tras ciento ocho pasos,
\[
 \beta\longmapsto\beta,
 \qquad
 q_{\rm mem}\longmapsto q_{\rm mem}+12.
\]
El calendario retorna; la memoria vertical no se reinicia.

\section{Extractor firmado de eventos}

Sea \((d_t)_{t=1}^{108}\) el registro digital sobre
\(\Gamma_{108}\). Definimos
\begin{equation}
 \mathcal D_{\rm evt}:=\{2,4,6,8\},
 \qquad
 \Sigma_{12}:=(+,+,+,-,-,-,+,+,+,-,-,-).
 \label{eq:u006d-datos-eventos}
\end{equation}
Para \(m=0,\ldots,11\), sea
\begin{equation}
 e_m:=
 \Sigma_{12}(m+1)\,
 \#\{t\in E_{m+1}:d_t\in\mathcal D_{\rm evt}\}
 \in\mathbb Z.
 \label{eq:u006d-eventos-firmados}
\end{equation}
La aplicación conserva el censo local y la orientación sectorial como
factores separados.

Emparejamos las dos semicoronas mediante
\begin{equation}
 p_i:=e_i+e_{i+6},
 \qquad
 a_i:=e_i-e_{i+6},
 \qquad 0\leq i\leq5.
 \label{eq:u006d-pares-antipodales}
\end{equation}
Definimos el censo total, las cinco diferencias respecto del sexto par y
las seis antisimetrías:
\begin{equation}
 s_C:=\sum_{i=0}^{5}p_i,
 \qquad
 M_i:=p_i-p_5\quad(0\leq i\leq4),
 \qquad
 A_6:=(a_0,\ldots,a_5).
 \label{eq:u006d-coordenadas-uno-cinco-seis}
\end{equation}

\begin{definition}[Lector integral \(1+5+6\)]
\label{def:u006d-lector-integral}
El lector lineal es
\begin{equation}
 \mathcal L_{12}:\mathbb Z^{12}\longrightarrow\mathbb Z^{12},
 \qquad
 \mathcal L_{12}(e)
 :=
 (s_C,M_0,\ldots,M_4,a_0,\ldots,a_5).
 \label{eq:u006d-lector-l12}
\end{equation}
La distribución de coordenadas es
\begin{equation}
 \underbrace{1}_{s_C}
 +\underbrace{5}_{M_0,\ldots,M_4}
 +\underbrace{6}_{a_0,\ldots,a_5}
 =12.
 \label{eq:u006d-uno-cinco-seis}
\end{equation}
\end{definition}

\begin{theorem}[Reconstrucción exacta e imagen integral]
\label{thm:u006d-inversa-l12}
El lector \(\mathcal L_{12}\) es inyectivo. Sobre su imagen,
\begin{align}
 p_5&=
 \frac{s_C-\sum_{i=0}^{4}M_i}{6},
 &
 p_i&=p_5+M_i\quad(0\leq i\leq4),
 \label{eq:u006d-inversa-p}\\
 e_i&=\frac{p_i+a_i}{2},
 &
 e_{i+6}&=\frac{p_i-a_i}{2}
 \quad(0\leq i\leq5).
 \label{eq:u006d-inversa-e}
\end{align}
Una tupla de salida pertenece a \(\operatorname{im}\mathcal L_{12}\) si y
sólo si
\begin{equation}
 s_C-\sum_{i=0}^{4}M_i\equiv0\pmod6
 \label{eq:u006d-condicion-divisibilidad}
\end{equation}
y, para los \(p_i\) reconstruidos,
\begin{equation}
 p_i\equiv a_i\pmod2
 \qquad(0\leq i\leq5).
 \label{eq:u006d-condicion-paridad}
\end{equation}
\end{theorem}

\begin{proof}
De \(M_i=p_i-p_5\) se obtiene
\[
 s_C=\sum_{i=0}^{4}(p_5+M_i)+p_5
 =6p_5+\sum_{i=0}^{4}M_i,
\]
que da la primera ecuación de \eqref{eq:u006d-inversa-p}; las demás son
inmediatas. Sumar y restar
\(p_i=e_i+e_{i+6}\) y \(a_i=e_i-e_{i+6}\) produce
\eqref{eq:u006d-inversa-e}. La divisibilidad por seis y las seis
congruencias de paridad son precisamente las condiciones necesarias y
suficientes para que todas las coordenadas reconstruidas sean enteras.
\end{proof}

La división por seis aparece después del censo orientado, no durante él. Las
condiciones integrales de esta unidad son exactamente
\eqref{eq:u006d-condicion-divisibilidad} y
\eqref{eq:u006d-condicion-paridad}; no se atribuye aquí una forma normal de
Smith a una matriz que no haya sido previamente definida.

\section{Dos palabras dodecafásicas de distinta procedencia}

Debe conservarse la distinción
\begin{equation}
 U_{12}^{\rm cat}:=U_6\Vert U_6,
 \qquad
 U_{12}^{\rm sgn}:=\mathcal H_{12}(K).
 \label{eq:u006d-dos-palabras-doce}
\end{equation}
La primera concatena dos emisiones catalogales de longitud seis. La segunda
es una coordenada firmada ligada reversiblemente a \(K\) por el operador
\(\mathcal H_{12}\). Su longitud común no implica una misma procedencia,
semántica o ley de transporte.

\begin{criterion}[Criterios de refutación]
\label{crit:u006d-refutacion}
La construcción queda refutada si la sucesión
\eqref{eq:u006d-extension} no es exacta, si su clase cohomológica es
trivial, si una tupla de la imagen viola
\eqref{eq:u006d-condicion-divisibilidad} o
\eqref{eq:u006d-condicion-paridad}, si las fórmulas inversas no recuperan
los doce enteros iniciales, o si el retorno de \(\beta\) se interpreta como
reinicio de \(q_{\rm mem}\).
\end{criterion}

% Unidad U016. Registro firmado, lectura terminal y reconstruccion de K.
% Redaccion autonoma desde la autoridad material 1063 y su sucesor V2.

\chapter{Registro dodecafásico firmado y reconstrucción integral del sello}
\label{chap:registro-dodecafasico-hadamard-k}

La construcción del sello entero no comienza con \(\alpha\), con su inversa
ni con una tabla metrológica. Su dominio es el estado terminal enriquecido
del Topological Phase Kernel obtenido después de la cadena ya construida
\begin{equation}
 \boxed{
 \mathrm{APP}\longrightarrow\mathrm{TRIT}\longrightarrow\mathrm{TPK}
 \longrightarrow\Omega_{\rm seed}
 \xrightarrow{T_{\min}}\mathcal U_6^{\rm cat}
 \xrightarrow{\rho_3}W_{243}
 \longrightarrow w_{30}\longrightarrow R_{36}
 \longrightarrow G_9\longrightarrow x_{\rm term}.}
 \label{eq:cadena-upstream-registro-dodecafasico}
\end{equation}
Las flechas anteriores fueron definidas con sus respectivos dominios:
pares de semillas, firmas, emisiones, palabras ternarias, fronteras,
historias compatibles y holonomía nonádica. Esta unidad no vuelve a escoger
ninguno de esos datos. Construye las cartas enteras del mismo estado
terminal y demuestra su equivalencia.

\section{Ventanas de la traza orientada de longitud ciento ocho}

Sea
\[
 \Theta_{108}=\mathbb Z/108\mathbb Z
\]
el calendario observable. Sus doce ventanas nonádicas son
\[
 E_m=\{9m+1,\ldots,9m+9\},
 \qquad 0\leq m<12.
\]
La partición es disjunta y exhaustiva:
\[
 \{1,\ldots,108\}=\bigsqcup_{m=0}^{11}E_m.
\]
El conjunto de códigos de dirección que publican un evento se fija como
\[
 D_{\rm evt}=\{2,4,6,8\}.
\]
La orientación de las doce ventanas es
\begin{equation}
 \Sigma_{12}=(+,+,+,-,-,-,+,+,+,-,-,-).
 \label{eq:signatura-doce-ventanas}
\end{equation}
Si \(d_t\) es la dirección codificada en el paso \(t\), el evento firmado
de la ventana \(m\) es
\begin{equation}
 e_m=\Sigma_{12}(m)\,
 \#\{t\in E_m:d_t\in D_{\rm evt}\}.
 \label{eq:evento-firmado-ventana}
\end{equation}
La fórmula distingue el número no negativo de eventos y el signo de la
orientación. No reemplaza una dirección cardinal por un signo: las dos
coordenadas siguen presentes en el estado TPK.

\begin{definition}[Registro enriquecido dodecafásico]
\label{def:registro-enriquecido-dodecafasico}
El registro transportado por la traza es
\[
 \mathcal R_{12}(x)=
 \bigl((E_m,\tau_m,\sigma_m,\kappa_m,\Lambda_m)\bigr)_{m=0}^{11},
\]
donde \(\tau_m\) es la orientación TRIT, \(\sigma_m\) la firma de
frontera, \(\kappa_m\) el acarreo, y \(\Lambda_m\) el segmento de ledger
que conserva la ruta. Existe una cadena de proyecciones
\[
 \mathcal X_{\rm TPK}^{\rm enr}
 \longrightarrow\mathcal R_{12}
 \longrightarrow\Theta_{108},
\]
pero ninguna de las dos flechas es una identificación.
\end{definition}

En particular, \(\Theta_{108}\) sólo conserva la posición del calendario.
El registro \(\mathcal R_{12}\) conserva además orientación, firma,
acarreo y procedencia. El estado completo conserva todavía cilindro,
frontera, cociente, terminal y memoria vertical.

\section{Carta integral de la oposición dodecafásica}

Para \(0\leq i<6\), las posiciones opuestas determinan
\begin{equation}
 p_i=e_i+e_{i+6},
 \qquad
 a_i=e_i-e_{i+6}.
 \label{eq:pares-opuestos-registro}
\end{equation}
La carga central y los cinco márgenes relativos son
\begin{equation}
 s_C=\sum_{i=0}^{5}p_i,
 \qquad
 M_i=p_i-p_5\quad(0\leq i<5).
 \label{eq:carga-margenes-registro}
\end{equation}

\begin{definition}[Carta integral \(1+5+6\)]
\label{def:carta-integral-156}
Se define
\begin{equation}
 \mathcal L_{12}(e_0,\ldots,e_{11})
 =\bigl(s_C,M_0,\ldots,M_4,a_0,\ldots,a_5\bigr).
 \label{eq:carta-integral-156}
\end{equation}
La partición de coordenadas es
\[
 1+5+6=12.
\]
\end{definition}

\begin{theorem}[Inversión exacta de la carta integral]
\label{thm:inversion-carta-integral-156}
Una tupla
\[
 \ell=(s_C,M_0,\ldots,M_4,a_0,\ldots,a_5)\in\mathbb Z^{12}
\]
pertenece a la imagen de \(\mathcal L_{12}\) si y sólo si
\begin{align}
 s_C-\sum_{i=0}^{4}M_i&\equiv0\pmod6,
 \label{eq:congruencia-seis-carta}\\
 p_i&\equiv a_i\pmod2
 \quad(0\leq i<6),
 \label{eq:congruencia-paridad-carta}
\end{align}
donde
\begin{align}
 p_5&=\frac{s_C-\sum_{i=0}^{4}M_i}{6},
 \label{eq:inversa-p5-carta}\\
 p_i&=p_5+M_i\quad(0\leq i<5).
 \label{eq:inversa-pi-carta}
\end{align}
En ese caso, la única preimagen entera es
\begin{equation}
 e_i=\frac{p_i+a_i}{2},
 \qquad
 e_{i+6}=\frac{p_i-a_i}{2}.
 \label{eq:inversa-eventos-carta}
\end{equation}
\end{theorem}

\begin{proof}
Sumando \(M_i=p_i-p_5\) para \(0\leq i<5\), obtenemos
\[
 \sum_{i=0}^{4}M_i
 =\sum_{i=0}^{4}p_i-5p_5
 =s_C-6p_5,
\]
lo que fuerza \eqref{eq:inversa-p5-carta} y la congruencia módulo seis.
Una vez reconstruidos los seis \(p_i\), el sistema
\[
 p_i=e_i+e_{i+6},
 \qquad a_i=e_i-e_{i+6}
\]
tiene la única solución \eqref{eq:inversa-eventos-carta}. Esa solución es
entera exactamente cuando \(p_i\) y \(a_i\) tienen la misma paridad.
Las dos familias de congruencias son, por tanto, necesarias y suficientes.
\end{proof}

Este teorema es anterior al cierre numérico de \(\alpha\). Explica cómo la
traza orientada produce una carta integral reversible, y por qué no basta
con conservar el calendario reducido.

\section{Selección del estado terminal común}

Las tres secciones compatibles llegan a profundidad terminal con catálogos
de cardinalidades
\[
 |\mathcal T_\pi|=196,
 \qquad
 |\mathcal T_e|=197,
 \qquad
 |\mathcal T_\varphi|=196.
\]
El producto contiene
\begin{equation}
 196\cdot197\cdot196=7\,567\,952
 \label{eq:cardinal-producto-terminal}
\end{equation}
ternas ordenadas.

El margen de columnas del lector terminal es
\[
 q_\partial=(1,2,2,1,2,0)\in\mathbb F_3^6,
\]
y la orientación de filas del canal conjunto es
\[
 \sigma=(1,1,-1)=(1,1,2)\in\mathbb F_3^3.
\]
El primer dato actúa sobre columnas; el segundo sobre filas. No se deduce
uno del otro.

Las firmas locales de Gram, norma, peso y distancia reducen
\eqref{eq:cardinal-producto-terminal} a \(331\) ternas. El margen
\(q_\partial\) deja dos índices:
\[
 (120,197,157),
 \qquad
 (129,197,148).
\]
Sus matrices son
\[
 B_0=
 \begin{pmatrix}
 0&2&0&2&2&0\\
 0&0&1&2&2&0\\
 1&0&1&0&1&0
 \end{pmatrix},
 \qquad
 B_1=
 \begin{pmatrix}
 0&2&1&0&2&0\\
 0&0&1&2&2&0\\
 1&0&0&2&1&0
 \end{pmatrix}.
\]
Las sumas de filas son
\[
 (0,2,0),
 \qquad
 (2,2,1).
\]
Como
\[
 \langle\sigma\rangle
 =\{(0,0,0),(1,1,2),(2,2,1)\},
\]
sólo \(B_1\) satisface la condición axial.

\begin{theorem}[Unicidad de la publicación terminal visible]
\label{thm:unicidad-publicacion-terminal-visible}
La selección conjunta determina
\begin{equation}
 B_{\rm term}=
 \begin{pmatrix}
 021020\\001220\\100210
 \end{pmatrix},
 \label{eq:matriz-terminal-visible}
\end{equation}
correspondiente a la terna \((129,197,148)\). Ninguna evaluación decimal de
\(\alpha\) interviene en el censo ni en las dos condiciones de selección.
\end{theorem}

\begin{proof}
El censo y el margen dejan exactamente \(B_0,B_1\). La carga de \(B_0\)
no pertenece a \(\langle\sigma\rangle\), mientras la de \(B_1\) es
\(2\sigma\). La segunda candidata es, por tanto, la única que satisface
simultáneamente ambos requisitos tipados.
\end{proof}

\section{Dos publicaciones del mismo estado, no una conversión entre ellas}

Sea
\[
 x_{\rm term}\in
 \mathcal X_{\rm TPK}^{\rm term,enr}
 \subseteq\mathcal X_{\rm TPK}^{[108]}.
\]
La lectura visible es
\[
 \pi_{\rm term}(x_{\rm term})=B_{\rm term}.
\]
La carta firmada conserva otra coordenada del mismo estado.

\begin{definition}[Lector firmado producido por el registro terminal]
\label{def:subespacio-firmado-terminal}
El estado de entrada del lector es exclusivamente el estado terminal
enriquecido ya producido por la tubería causal
\(\mathcal U_t=\mathsf{Upd}_t\circ\mathsf{Tra}_t\circ\mathsf{Sel}_t\):
\begin{equation}
 x_{\rm term}
 =\bigl(\mathcal U_{108}\circ\cdots\circ\mathcal U_1\bigr)
   (x_{\rm seed})
 \in\mathcal X_{\rm TPK}^{\rm term,enr}.
 \label{eq:subespacio-firmado-terminal}
\end{equation}
La composición del registro orientado con sus lectores de canal \(90/120\)
extrae del libro \(\lambda_{108}\), sin añadir coordenadas externas,
\begin{equation}
 \mathcal E_{108}^{90,120}\!\left(\mathcal R_{12}(x_{\rm term})\right)
 =\bigl(b^{90},b^{120},Q_{\rm TPK}\bigr),
 \label{eq:canales-ledger-previos-u}
\end{equation}
donde
\begin{align}
 b^{90}={}&(-495,-116,-684,108,601,-90,
             -173,-88,313,560,-397,461),\notag\\
 b^{120}={}&(-425,-281,-481,671,-255,30,
              475,-543,680,251,6,-128),\notag\\
 Q_{\rm TPK}={}&6263.\notag
\end{align}
Estas veinticinco coordenadas son salidas del registro de la
\cref{def:registro-enriquecido-dodecafasico}, no datos añadidos al dominio.

Sobre una salida compatible \((b^{90},b^{120},Q)\), definimos el lector
armónico integral \(\Pi_H\) de forma material. Para \(r=1,2,3\), sean
\begin{equation}
 B_r=\sum_{j=0}^{3}b^{120}_{r+3j},
 \quad
 A_1=\frac{Q+2B_1+B_2}{3},
 \quad A_2=A_1-B_1,
 \quad A_3=A_2-B_2,
 \label{eq:lector-armonico-a-desde-ledger}
\end{equation}
y \(d_{r,j}=b^{90}_{r+3j}\). El bloque orbital \(r\) es
\begin{equation}
 \Pi_H^{(r)}(b^{90},b^{120},Q)
 =\bigl(A_r,d_{r,1}-d_{r,3},
              d_{r,0}+d_{r,2},d_{r,0}-d_{r,2}\bigr).
 \label{eq:lector-armonico-bloque-desde-ledger}
\end{equation}
Intercalando por columnas los tres bloques se obtiene \(\Pi_H\), y la
lectura firmada queda construida como la composición
\begin{equation}
 \boxed{
 R_{\rm sgn}:=
 \Pi_H\circ\mathcal E_{108}^{90,120}\circ\mathcal R_{12}:
 \mathcal X_{\rm TPK}^{\rm term,enr}
 \longrightarrow\mathcal U_{12}^{\rm sgn}:=
 \operatorname{im}R_{\rm sgn}\subseteq\mathbb Z^{12}.}
 \label{eq:lector-firmado-terminal}
\end{equation}
Ni \(U\) ni \(K\) pertenecen al dominio: \(U\) es la salida de
\eqref{eq:lector-firmado-terminal}, y \(K\) se reconstruye después mediante
la inversa integral de Hadamard.
\end{definition}

La escritura conjunta correcta es
\begin{equation}
 x_{\rm term}
 \xrightarrow{(\pi_{\rm term},R_{\rm sgn})}
 (B_{\rm term},U_{12}^{\rm sgn}).
 \label{eq:doble-publicacion-terminal}
\end{equation}
No se interpone una flecha
\(B_{\rm term}\to U_{12}^{\rm sgn}\): la matriz visible ha olvidado
precisamente las coordenadas de orientación, oposición, cociente, acarreo y
ledger utilizadas por la carta firmada. La fuente de ambas publicaciones es
el estado enriquecido producido por \eqref{eq:cadena-upstream-registro-dodecafasico}.

La evaluación del registro firmado da
\begin{equation}
 \boxed{
 U=(2378,1406,2479,-452,998,-551,
 -668,-204,-371,-322,-28,-997).}
 \label{eq:vector-u-firmado-autonomo}
\end{equation}

\section{Naturalidad respecto de la profundidad}

Para \(N'\geq N\geq108\), el truncamiento
\begin{equation}
 \tau_{N',N}:
 \mathcal X_{\rm TPK}^{[N'],\rm enr}
 \longrightarrow\mathcal X_{\rm TPK}^{[N],\rm enr}
 \label{eq:truncamiento-preserva-ku-autonomo}
\end{equation}
actúa sobre la historia enriquecida y su libro, no sobre una pareja
\((K,U)\) añadida al estado.

\begin{theorem}[Cuadrado de naturalidad de la lectura firmada]
\label{thm:naturalidad-lector-firmado-autonomo}
Se cumple
\begin{equation}
 \boxed{
 R_{{\rm sgn},N}\circ\tau_{N',N}
 =R_{{\rm sgn},N'}.}
 \label{eq:cuadrado-naturalidad-lector-firmado}
\end{equation}
\end{theorem}

\begin{proof}
Cada actualizador \(\mathcal U_t\) añade la incidencia siguiente al libro
sin reescribir las ciento ocho ventanas ya generadas. Por ello
\(\mathcal R_{12}\), \(\mathcal E_{108}^{90,120}\) y \(\Pi_H\) leen las
mismas coordenadas al truncar cualquier prolongación posterior. La
conmutatividad es una igualdad de aplicaciones sobre estados producidos, no
la conservación tautológica de una salida incluida en la entrada.
\end{proof}

\section{Órbitas de paso tres y transformada de Hadamard}

El orden cíclico del sello se distribuye en las tres órbitas
\begin{align}
 K^{(1)}&=(K_1,K_4,K_7,K_{10}),
 \label{eq:orbita-k-1}\\
 K^{(2)}&=(K_2,K_5,K_8,K_{11}),
 \label{eq:orbita-k-2}\\
 K^{(3)}&=(K_3,K_6,K_9,K_{12}).
 \label{eq:orbita-k-3}
\end{align}
Sea
\begin{equation}
 H_4=
 \begin{pmatrix}
 1&1&1&1\\
 1&1&-1&-1\\
 1&-1&1&-1\\
 1&-1&-1&1
 \end{pmatrix}.
 \label{eq:matriz-hadamard-cuatro-autonoma}
\end{equation}
Si \(P_{\rm orb}\) reordena las doce coordenadas de acuerdo con
\eqref{eq:orbita-k-1}--\eqref{eq:orbita-k-3}, la transformación global es
\begin{equation}
 \mathcal H_{12}
 =P_{\rm orb}^{-1}
 (H_4\oplus H_4\oplus H_4)P_{\rm orb}.
 \label{eq:hadamard-global-doce}
\end{equation}

Los tres bloques de \eqref{eq:vector-u-firmado-autonomo}, leídos por
órbitas, son
\begin{align*}
 U^{(1)}&=(2378,-452,-668,-322),\\
 U^{(2)}&=(1406,998,-204,-28),\\
 U^{(3)}&=(2479,-551,-371,-997).
\end{align*}

\begin{lemma}[Cuarto de vuelta algebraico de Hadamard]
\label{lem:cuarto-inversa-hadamard}
La matriz \(H_4\) satisface
\begin{equation}
 H_4^2=4I_4,
 \qquad
 H_4^{-1}=\frac14H_4,
 \qquad
 \det H_4=-16.
 \label{eq:identidades-hadamard-cuatro}
\end{equation}
\end{lemma}

\begin{proof}
Las filas de \(H_4\) son ortogonales y cada una posee norma cuadrada cuatro.
Como la matriz es simétrica, \(H_4^{\mathsf T}H_4=H_4^2=4I_4\). La
fórmula de la inversa y el determinante se deducen inmediatamente.
\end{proof}

\begin{theorem}[Reconstrucción entera única del sello]
\label{thm:reconstruccion-entera-k-hadamard}
La inversión
\begin{equation}
 K^{(r)}=\frac14H_4U^{(r)}
 \qquad(r=1,2,3)
 \label{eq:inversion-bloques-hadamard}
\end{equation}
produce
\begin{align*}
 K^{(1)}&=(234,729,621,794),\\
 K^{(2)}&=(543,659,058,146),\\
 K^{(3)}&=(140,824,914,601).
\end{align*}
Al deshacer el orden orbital se obtiene
\begin{equation}
 \boxed{
 K=(234,543,140,729,659,824,621,058,914,794,146,601).}
 \label{eq:sello-k-doce-autonomo}
\end{equation}
Es la única preimagen racional de \(U\), y pertenece a
\(\mathbb Z^{12}\).
\end{theorem}

\begin{proof}
Por \eqref{eq:identidades-hadamard-cuatro}, la preimagen racional es única.
Las multiplicaciones explícitas son
\begin{align*}
 H_4U^{(1)}&=(936,2916,2484,3176),\\
 H_4U^{(2)}&=(2172,2636,232,584),\\
 H_4U^{(3)}&=(560,3296,3656,2404).
\end{align*}
Todas las coordenadas son divisibles por cuatro. Dividir y reintercalar
produce \eqref{eq:sello-k-doce-autonomo}.
\end{proof}

El factor \(1/4\) no se escoge para ajustar una salida: es la inversa
forzada por \(H_4^2=4I_4\). Tampoco es la raíz cuadrada de \(-1\). La
estructura compleja interna \(A_W^2=-I\) aparecerá después de construir la
matriz de Paley; son dos hechos de tipos diferentes.

\section{Estructura integral y forma normal de Smith}

\begin{theorem}[Imagen integral de \(H_4\)]
\label{thm:smith-hadamard-autonomo}
La forma normal de Smith es
\begin{equation}
 \operatorname{SNF}(H_4)=\operatorname{diag}(1,2,2,4).
 \label{eq:snf-hadamard-autonoma}
\end{equation}
En consecuencia,
\begin{equation}
 \operatorname{coker}\mathcal H_{12}
 \cong(\mathbb Z/2\mathbb Z)^6
 \oplus(\mathbb Z/4\mathbb Z)^3,
 \label{eq:cociente-hadamard-doce}
\end{equation}
y la imagen de \(\mathcal H_{12}\) posee índice
\[
 16^3=4096
\]
en \(\mathbb Z^{12}\). Para \(u\in\mathbb Z^4\),
\begin{equation}
 u\in H_4\mathbb Z^4
 \quad\Longleftrightarrow\quad
 H_4u\equiv0\pmod4.
 \label{eq:criterio-imagen-hadamard}
\end{equation}
\end{theorem}

\begin{proof}
El máximo común divisor de las entradas es uno. Los menores de orden dos son
pares y alguno tiene valor absoluto dos; los de orden tres son múltiplos de
cuatro y alguno tiene valor absoluto cuatro; el determinante tiene módulo
dieciséis. Los divisores determinantes son \(1,2,4,16\), de donde se obtienen
los factores invariantes \(1,2,2,4\). La suma directa triple repite los
factores y multiplica los índices. Finalmente,
\(H_4^{-1}=H_4/4\), por lo que la preimagen es entera exactamente bajo la
congruencia \eqref{eq:criterio-imagen-hadamard}.
\end{proof}

\section{Certificación en el segundo sistema de coordenadas: diferencias de pasos tres y cuatro}

Con índices cíclicos módulo doce, definimos
\begin{equation}
 (D_3K)_m=K_m-K_{m+3},
 \qquad
 (D_4K)_m=K_m-K_{m+4},
 \qquad
 Q(K)=\sum_{m=1}^{12}K_m.
 \label{eq:extractor-diferencias-k}
\end{equation}
Aplicados al sello ya reconstruido por Hadamard, estos operadores deben
recuperar exactamente las coordenadas de canal que el registro TPK produjo
antes de \(U\). Para \eqref{eq:sello-k-doce-autonomo},
\begin{align*}
 D_3K={}&(-495,-116,-684,108,601,-90,\\
         &\qquad-173,-88,313,560,-397,461),\\
 D_4K={}&(-425,-281,-481,671,-255,30,\\
         &\qquad475,-543,680,251,6,-128),\\
 Q(K)={}&6263.
\end{align*}
Por comparación con \eqref{eq:canales-ledger-previos-u}, queda probada la
igualdad de salidas
\begin{equation}
 \boxed{
 D_3K=b^{90},\qquad D_4K=b^{120},\qquad Q(K)=Q_{\rm TPK}.}
 \label{eq:certificacion-canales-ledger-desde-k}
\end{equation}
El sistema diferencial certifica y reconstruye el registro previo; no
selecciona \(K\), \(U\) ni ninguna constante.

\begin{theorem}[Completitud del extractor transversal]
\label{thm:completitud-extractor-transversal-k}
El operador
\[
 \mathcal E_{\rm dod}=(D_3,D_4,Q):
 \mathbb Z^{12}\longrightarrow\mathbb Z^{25}
\]
tiene rango doce y forma normal de Smith
\begin{equation}
 \operatorname{SNF}(\mathcal E_{\rm dod})
 =\operatorname{diag}(\underbrace{1,\ldots,1}_{11},12).
 \label{eq:snf-extractor-transversal-k}
\end{equation}
En particular, \((D_3K,D_4K,Q(K))\) determina un único \(K\).
\end{theorem}

\begin{proof}
Si \(D_3v=D_4v=0\), entonces \(v\) es constante a lo largo de los pasos
tres y cuatro. Como \(\gcd(3,4)=1\), esos pasos generan todo
\(\mathbb Z/12\mathbb Z\), de modo que
\[
 \ker D_3\cap\ker D_4=\langle\mathbf1\rangle.
\]
La carga total no se anula en esa recta, pues \(Q(\mathbf1)=12\). El rango
es doce. Las filas de \((D_3,D_4)\) son incidencias orientadas de un grafo
de Cayley conexo; un árbol generador aporta once factores invariantes
unitarios. Todo menor máximo debe contener la fila \(Q\), y las operaciones
unimodulares reducen su última contribución a doce. Existe un menor máximo de
módulo doce; por tanto el último factor es exactamente doce.
\end{proof}

La reconstrucción de \(U\) desde estas coordenadas no necesita volver a
multiplicar por \(H_4\). Para \(r=1,2,3\), sea
\begin{equation}
 B_r=\sum_{j=0}^{3}(D_4K)_{r+3j}.
 \label{eq:br-diferencias-k}
\end{equation}
Entonces
\[
 (B_1,B_2,B_3)=(972,-1073,101).
\]
Las sumas orbitales son
\begin{align}
 A_1&=\frac{Q+2B_1+B_2}{3}=2378,
 \label{eq:a1-reconstruccion-u}\\
 A_2&=A_1-B_1=1406,
 \label{eq:a2-reconstruccion-u}\\
 A_3&=A_2-B_2=2479.
 \label{eq:a3-reconstruccion-u}
\end{align}
Si \(d_{r,j}=(D_3K)_{r+3j}\), el bloque firmado de la órbita \(r\) es
\begin{equation}
 \bigl(A_r,
 d_{r,1}-d_{r,3},
 d_{r,0}+d_{r,2},
 d_{r,0}-d_{r,2}\bigr),
 \label{eq:bloque-firmado-desde-diferencias}
\end{equation}
que reproduce los tres bloques \(U^{(r)}\).

\section{Tétrada de umbral del sello}

La completación decimal local satisface
\[
 1000=729+271.
\]
El soporte de las coordenadas del sello que alcanzan el umbral \(729\) es
\begin{equation}
 \boxed{
 B_K:=\{m:K_m\geq729\}=\{4,6,9,10\}.}
 \label{eq:tetrada-umbral-k}
\end{equation}
Esta tétrada se obtiene después de construir \(K\). En esta unidad no se la
denomina todavía estrella de Witt: tal objeto sólo estará definido después
de construir el código ternario, sus hexadas y el sistema de Steiner.

\section{Independencia causal respecto de la salida física}

\begin{theorem}[Aislamiento de \(U\) y \(K\)]
\label{thm:aislamiento-u-k-alpha-codata}
La composición
\begin{equation}
 \Omega_{\rm seed}\longrightarrow x_{\rm term}
 \xrightarrow{\mathcal R_{12}}
 \mathcal R_{12}(x_{\rm term})
 \xrightarrow{\mathcal E_{108}^{90,120}}
 (b^{90},b^{120},Q_{\rm TPK})
 \xrightarrow{\Pi_H}U
 \xrightarrow{\mathcal H_{12}^{-1}}K
 \label{eq:composicion-upstream-u-k}
\end{equation}
no recibe como entrada \(\alpha\), \(\alpha^{-1}\), CODATA ni una cadena
decimal objetivo. El sello \(K\) queda determinado antes de definir la
recurrencia que producirá \(\alpha\).
\end{theorem}

\begin{proof}
La primera parte de la composición utiliza únicamente semillas, emisiones,
lectores ternarios, operadores \(L_0,L_1\), relaciones de supervivencia,
transporte de fase y las cartas del estado TPK. El registro
\(\mathcal R_{12}\) y el lector de sus canales producen
\((b^{90},b^{120},Q_{\rm TPK})\); las fórmulas
\eqref{eq:lector-armonico-a-desde-ledger} y
\eqref{eq:lector-armonico-bloque-desde-ledger} producen \(U\); y sólo
entonces actúa la transformación lineal
\(\mathcal H_{12}^{-1}=\mathcal H_{12}/4\) sobre su imagen integral.
Ninguna de las fórmulas contiene una variable física o una tabla
metrológica. La recurrencia en base mil se introduce sólo en la unidad
siguiente, después de haber fijado \(K\).
\end{proof}

\section{Criterios de refutación}

\begin{criterion}[Refutación del registro y del sello]
La construcción queda refutada si se verifica cualquiera de los hechos
siguientes:
\begin{enumerate}
 \item las ventanas \(E_m\) no particionan los ciento ocho pasos;
 \item la carta \(\mathcal L_{12}\) incumple su inversa o sus congruencias;
 \item el censo terminal no produce \(331\) ternas, dos matrices y una única
       carga axial;
 \item se intenta obtener \(U\) de \(B_{\rm term}\) después de olvidar las
       coordenadas firmadas;
 \item falla el cuadrado de naturalidad
       \eqref{eq:cuadrado-naturalidad-lector-firmado};
 \item \(H_4^2\neq4I_4\), alguna preimagen no es entera o el vector
       reconstruido difiere de \eqref{eq:sello-k-doce-autonomo};
 \item la forma normal de Smith difiere de
       \eqref{eq:snf-hadamard-autonoma};
 \item el extractor \((D_3,D_4,Q)\) pierde rango o no reproduce \(U\);
 \item \(B_K\) se introduce antes de construir \(K\), o se le atribuye una
       estructura de Witt antes de definir \(W_{12}\);
 \item una cifra de \(\alpha\) o una referencia metrológica interviene en
       cualquiera de las flechas de \eqref{eq:composicion-upstream-u-k}.
\end{enumerate}
\end{criterion}

\end{HMTScientificOwnerSubordinated}

\HMTFlushCurrentChapter
\HMTSuccessorChapterInput
  {Generación coinductiva de \texorpdfstring{\(\pi\)}{pi}: carácter de clausura, \(5\) regiones orientadas, multiplicidad \texorpdfstring{\(432+4\cdot144=1008\)}{432+4x144=1008} y precisión arbitraria}
  {../colaboracion/parte_iii/tex/capitulos_finales/28_generacion_pi.tex}

\HMTFlushCurrentChapter
\HMTSuccessorChapterInput
  {Generación coinductiva de \texorpdfstring{\(e\)}{e}: carácter exponencial, semigrupo factorial y retorno de fase con memoria de estado}
  {../colaboracion/parte_iii/tex/capitulos_finales/29_generacion_e.tex}

\HMTFlushCurrentChapter
\HMTSuccessorChapterInput
  {Generación coinductiva de \texorpdfstring{\(\varphi\)}{phi}: incidencia modal, rayo positivo de Perron y unicidad de la autoescala}
  {../colaboracion/parte_iii/tex/capitulos_finales/30_generacion_phi.tex}

% Propietario dodecafásico exacto y reunión transversal exacta. El primero
% deriva la coordenada firmada y el sello; el segundo hace visible, antes del
% capítulo de publicación, la bifurcación numérica/excepcional del mismo
% estado. Ninguno recibe \alpha ni valores convencionales como entrada.
\begin{HMTScientificOwnerSubordinated}
\chapter{El cierre dodecafásico de \texorpdfstring{\(\alpha\)}{alpha}}
\label{chap:alpha}
\index{alpha@\(\alpha\)}
\index{ciclo dodecafásico}

Esta sección establece el cierre dodecafásico mediante aritmética entera en
base mil. La construcción explicita las doce recurrencias, los cocientes
euclídeos propagados y las condiciones de frontera que determinan la salida.

\section{El selector terminal de los tres canales}
\label{sec:selector-terminal-alpha-rev6}

El cierre reúne tres generaciones ya construidas, pero no las identifica por
su nombre decimal. A profundidad terminal, los catálogos de colas compatibles
poseen cardinales
\[
 |\mathcal T_\pi|=196,\qquad
 |\mathcal T_e|=197,\qquad
 |\mathcal T_\varphi|=196.
\]
El objeto de selección es, por tanto,
\(\mathcal T_\pi\times\mathcal T_e\times\mathcal T_\varphi\), con
\[
 196\cdot197\cdot196=7\,567\,952
\]
ternas posibles.

La frontera de coordenadas fijada por el lector terminal es
\[
 q=(1,2,2,1,2,0)\in\mathbb F_3^6.
\]
La ley de canal \(\pi+e-\varphi\) posee vector de signos
\[
 \sigma=(1,1,-1)=(1,1,2)\in\mathbb F_3^3
\]
y recta dual no nula
\[
 \langle2\sigma\rangle
 =\langle(2,2,1)\rangle.
\]
El primer dato actúa sobre las columnas de la frontera; el segundo, sobre
sus filas. No son dos rasgos intercambiables.

\begin{theorem}[Selector matricial terminal de \(\alpha\)]
\label{thm:selector-terminal-alpha-rev6}
En el catálogo terminal anterior, las firmas locales de Gram, norma, peso y
distancia de Hamming reducen las \(7\,567\,952\) ternas a \(331\). Al exigir
el margen de columnas \(q\) quedan dos candidatas:
\[
 (120,197,157),\qquad(129,197,148).
\]
La condición de que la carga de filas pertenezca a la recta del canal
\(\langle\sigma\rangle\) selecciona exactamente la segunda. Sus bloques son
\[
 \boxed{
 b_\pi=021020,\qquad
 b_e=001220,\qquad
 b_\varphi=100210.}
\]
\end{theorem}

\begin{proof}
Las dos matrices terminales que sobreviven al margen de columnas son
\[
B_0=
\begin{pmatrix}
0&2&0&2&2&0\\
0&0&1&2&2&0\\
1&0&1&0&1&0
\end{pmatrix},
\qquad
B_1=
\begin{pmatrix}
0&2&1&0&2&0\\
0&0&1&2&2&0\\
1&0&0&2&1&0
\end{pmatrix}.
\]
Sus sumas de filas en \(\mathbb F_3\) son, respectivamente,
\[
 (0,2,0),\qquad(2,2,1).
\]
La primera no pertenece a
\(\langle(1,1,2)\rangle=\{0,(1,1,2),(2,2,1)\}\); la segunda sí, pues
\((2,2,1)=2\sigma\). Por tanto sólo \(B_1\) satisface simultáneamente la
frontera de coordenadas y la carga del canal. Sus tres filas son los bloques
mostrados y corresponden a la terna \((129,197,148)\).
\end{proof}

El teorema no obtiene \(q\) de la carga de filas: \(q\) es el dato
coordenado del lector terminal y \(\sigma\) es la orientación del canal.
Lo demostrado es que, una vez declaradas ambas fronteras del mismo objeto,
la rama real es única. Esta separación impide convertir una firma encontrada
a posteriori en una ley supuestamente derivada.

\begin{definition}[Dos lecturas del estado terminal enriquecido]
Sea
\[
 x_{\rm term}\in\mathcal X_{\rm TPK}^{\rm term,enr}
 \subseteq\mathcal X_{\rm TPK}^{[108]}.
\]
Este estado es la salida
\[
 x_{\rm term}
 =(\mathcal U_{108}\circ\cdots\circ\mathcal U_1)(x_{\rm seed})
\]
de la tubería completa de selección, transporte y actualización del TPK.
El selector anterior es la lectura
\[
 \pi_{\rm term}(x_{\rm term})
 =B_{\rm term}
 :=\begin{pmatrix}021020\\001220\\100210\end{pmatrix}.
\]
El registro orientado y el lector de los canales de \(90^\circ\) y
\(120^\circ\) producen, desde el libro del mismo estado,
\begin{equation}\label{eq:alpha-subespacio-firmado}
 \mathcal E_{108}^{90,120}\!
 \left(\mathcal R_{12}(x_{\rm term})\right)
 =(b^{90},b^{120},Q_{\rm TPK}),
\end{equation}
con
\[
\begin{aligned}
 b^{90}={}&(-495,-116,-684,108,601,-90,
             -173,-88,313,560,-397,461),\\
 b^{120}={}&(-425,-281,-481,671,-255,30,
              475,-543,680,251,6,-128),\\
 Q_{\rm TPK}={}&6263.
\end{aligned}
\]
Para \(r=1,2,3\), sean
\[
 B_r=\sum_{j=0}^{3}b^{120}_{r+3j},\qquad
 A_1=\frac{Q_{\rm TPK}+2B_1+B_2}{3},\qquad
 A_2=A_1-B_1,\qquad A_3=A_2-B_2,
\]
y \(d_{r,j}=b^{90}_{r+3j}\). El lector armónico \(\Pi_H\) forma los
bloques
\begin{equation}\label{eq:alpha-modulo-imagen-H12}
 \Pi_H^{(r)}
 =\bigl(A_r,d_{r,1}-d_{r,3},
              d_{r,0}+d_{r,2},d_{r,0}-d_{r,2}\bigr)
\end{equation}
y los intercala por columnas. La lectura firmada es, por tanto, la
composición producida
\begin{equation}\label{eq:alpha-Rsgn-definido}
 \boxed{
 R_{\rm sgn}:=
 \Pi_H\circ\mathcal E_{108}^{90,120}\circ\mathcal R_{12}:
 \mathcal X_{\rm TPK}^{\rm term,enr}
 \longrightarrow
 \mathcal U_{12}^{\rm sgn}:=\operatorname{im}R_{\rm sgn}.}
\end{equation}
Ni \(U\) ni \(K\) son coordenadas añadidas al dominio: \(U\) es la salida
de \eqref{eq:alpha-Rsgn-definido}, y \(K\) será su única preimagen integral
por Hadamard.
No se afirma una factorización
\(B_{\rm term}\to U_{12}^{\rm sgn}\): ambas salidas conservan componentes
distintas del mismo estado terminal.

Para \(N'\geq N\geq108\), los truncamientos actúan sólo sobre las historias
enriquecidas:
\begin{equation}\label{eq:alpha-truncamiento-firmado-preserva-KU}
 \tau_{N',N}:
 \mathcal X_{\rm TPK}^{[N'],\rm enr}
 \longrightarrow\mathcal X_{\rm TPK}^{[N],\rm enr}.
\end{equation}
Si \(R_{\rm sgn,N}\) denota la misma composición en profundidad \(N\), la
naturalidad es el cuadrado concreto
\begin{equation}\label{eq:alpha-cuadrado-naturalidad-Rsgn}
 \boxed{
 R_{\rm sgn,N}\circ\tau_{N',N}
 =R_{\rm sgn,N'}.}
\end{equation}
\end{definition}

\begin{theorem}[Construcción reversible de la coordenada firmada]
\label{thm:alpha-Rsgn-construido}
El compuesto \eqref{eq:alpha-Rsgn-definido}, aplicado al estado terminal
producido, da exactamente
\begin{equation}\label{eq:alpha-Rsgn-evaluacion}
 U=R_{\rm sgn}(x_{\rm term})
 =(2378,1406,2479,-452,998,-551,-668,-204,-371,-322,-28,-997).
\end{equation}
Su inversión por tres bloques de Hadamard de orden cuatro produce después
\[
 \boxed{
 K=(234,543,140,729,659,824,621,058,914,794,146,601).}
\]
La transformación es invertible; satisface el cuadrado
\eqref{eq:alpha-cuadrado-naturalidad-Rsgn} bajo los truncamientos
\eqref{eq:alpha-truncamiento-firmado-preserva-KU}; y es independiente de
\(\alpha\), CODATA y de toda cifra metrológica.
\end{theorem}

\begin{proof}
Las fórmulas \eqref{eq:alpha-subespacio-firmado} y
\eqref{eq:alpha-modulo-imagen-H12} dan
\[
 (B_1,B_2,B_3)=(972,-1073,101),\qquad
 (A_1,A_2,A_3)=(2378,1406,2479),
\]
y producen los tres bloques firmados
\[
\begin{aligned}
 U^{(1)}&=(2378,-452,-668,-322),\\
 U^{(2)}&=(1406,998,-204,-28),\\
 U^{(3)}&=(2479,-551,-371,-997).
\end{aligned}
\]
Su intercalado es \eqref{eq:alpha-Rsgn-evaluacion}. Como
\(H_4^2=4I_4\) y
\(\det(H_4\oplus H_4\oplus H_4)=(-16)^3=-4096\), se tiene
\begin{equation}\label{eq:alpha-Rsgn-inversa}
 K=\mathcal H_{12}^{-1}U
 =\frac14\mathcal H_{12}U
\end{equation}
después de deshacer la permutación de órbitas. En efecto,
\[
 H_4U^{(1)}=(936,2916,2484,3176),\quad
 H_4U^{(2)}=(2172,2636,232,584),
\]
y
\[
 H_4U^{(3)}=(560,3296,3656,2404).
\]
La división forzada por cuatro da las órbitas
\((234,729,621,794)\), \((543,659,058,146)\) y
\((140,824,914,601)\), cuyo intercalado es el sello declarado.

Como control independiente, los operadores cíclicos
\((D_3K)_m=K_m-K_{m+3}\) y
\((D_4K)_m=K_m-K_{m+4}\), junto con \(Q(K)=\sum_mK_m\), satisfacen
\begin{equation}\label{eq:alpha-extractor-90-120}
 \operatorname{rank}(D_3,D_4)=11,
 \qquad
 \operatorname{rank}(D_3,D_4,Q)=12.
\end{equation}
Por ello \((D_3K,D_4K,Q(K))\) reconstruye un único sello; además, el
cálculo recupera exactamente
\[
 D_3K=b^{90},\qquad D_4K=b^{120},\qquad Q(K)=Q_{\rm TPK},
\]
lo cual certifica el lector previo del registro, sin convertir sus salidas
en entradas. El truncamiento
\eqref{eq:alpha-truncamiento-firmado-preserva-KU} no reescribe las primeras
ciento ocho ventanas del libro; por ello ambos miembros de
\eqref{eq:alpha-cuadrado-naturalidad-Rsgn} leen el mismo registro y
producen \(U\). Ninguna de las matrices utilizadas contiene \(\alpha\) ni
un dato experimental.
\end{proof}

El calendario observable de 108 pasos es una proyección estrictamente más
pobre: reiniciar sus cursores cada 9 pasos borra precisamente las
coordenadas bidireccionales que el registro enriquecido conserva. Por
eso el registro crudo no sustituye al estado enriquecido y no se usa para
reconstruir \(R_{\rm sgn}\). El teorema de selección anterior construye
\(B_{\rm term}\); el teorema \ref{thm:alpha-Rsgn-construido} construye, de
manera separada, la segunda publicación del mismo estado.

\section{Estado dodecafásico y aplicaciones reversibles}

El teorema de esta sección recibe del
teorema~\ref{thm:alpha-Rsgn-construido} el vector firmado
\[
\Usgn=(2378,1406,2479,-452,998,-551,-668,-204,-371,-322,-28,-997).
\]
La ruta tipada demostrada a partir de esa entrada es
\[
\Usgn\longleftrightarrow K,
\]
\[
K\longleftrightarrow(D_3K,D_4K,Q),
\]
\[
(P,E,\Phi,K,c_{13}=0)
\longleftrightarrow A=\alpha_{12}.
\]
La evaluación directa no recibe \(\alpha\) ni un valor metrológico como
entrada. Recibe \(\Usgn\), las tres representaciones arquimedianas y la
condición de frontera. El
\cref{thm:selector-terminal-alpha-rev6} construye antes la selección de los
tres bloques terminales; la presente capa comienza con el vector firmado que
los transporta al cierre dodecafásico. La traza reducida de una sola
coordenada posicional no interviene en esta determinación.

\begin{remark}[Dos representaciones dodecafásicas de tipos distintos]
\[
U_{12}^{\rm cat}=U_6\Vert U_6
\]
es una palabra periódica del catálogo, mientras
\[
\Usgn=(2378,1406,2479,-452,998,-551,-668,-204,-371,-322,-28,-997)
\]
es un vector entero firmado del espacio completo de estados de TPK\@. El primero pertenece a una
representación de bloques y tiene periodo seis. El segundo conserva orientación, no es
una palabra base mil y no tiene ese periodo. No son intercambiables.
\end{remark}

\section{El vector canónico y las tres órbitas}

El vector dodecafásico es
\[
\boxed{
K=(234,543,140,729,659,824,621,058,914,794,146,601).}
\]
Su suma es
\[
Q(K)=6263.
\]
El paso tres divide \(\Cdoce\) en tres órbitas de cuatro sectores:
\[
\begin{aligned}
K^{(1)}&=(K_1,K_4,K_7,K_{10})=(234,729,621,794),\\
K^{(2)}&=(K_2,K_5,K_8,K_{11})=(543,659,058,146),\\
K^{(3)}&=(K_3,K_6,K_9,K_{12})=(140,824,914,601).
\end{aligned}
\]
Cada órbita recorre cuatro sectores separados por \(90^\circ\). Las etiquetas
cardinales norte--este--sur--oeste aparecen después de fijar origen y orientación; el orden
primario es \((m,m+3,m+6,m+9)\).

\section{Primera representación: transformada de Hadamard}

Sea
\[
H_4=
\begin{pmatrix}
1&1&1&1\\
1&1&-1&-1\\
1&-1&1&-1\\
1&-1&-1&1
\end{pmatrix}.
\]
La acción sobre las tres órbitas da
\[
\begin{aligned}
H_4K^{(1)}&=(2378,-452,-668,-322),\\
H_4K^{(2)}&=(1406,998,-204,-28),\\
H_4K^{(3)}&=(2479,-551,-371,-997).
\end{aligned}
\]
Al intercalar por columnas recuperamos
\[
\mathcal H_{12}(K)=\Usgn.
\]
Como
\[
H_4^2=4I_4,
\qquad
\det H_4=-16,
\]
la transformación global es invertible sobre \(\mathbb Q^{12}\):
\[
K^{(r)}=\frac14H_4U^{(r)},
\qquad
\det\mathcal H_{12}=(-16)^3=-4096\ne0.
\]
En la representación canónica, la inversa produce un único vector entero de
\([0,999]^{12}\).

\begin{theorem}[Estructura integral de la transformación de Hadamard]
\label{thm:smith-hadamard}
La forma normal de Smith del bloque \(H_4\), denotada por
\(\operatorname{SNF}(H_4)\), es
\[
 \operatorname{SNF}(H_4)=\operatorname{diag}(1,2,2,4).
\]
En consecuencia, para la suma directa de los tres bloques,
\[
 \operatorname{coker}\mathcal H_{12}
 \cong(\mathbb Z/2\mathbb Z)^6\oplus(\mathbb Z/4\mathbb Z)^3,
\]
y la imagen de \(\mathbb Z^{12}\) posee índice \(4096\). Un vector
\(u\in\mathbb Z^4\) pertenece a \(H_4\mathbb Z^4\) exactamente cuando
\[
 H_4u\equiv0\pmod 4.
\]
\end{theorem}

\begin{proof}
El máximo común divisor de las entradas de \(H_4\) es uno. Los menores de
orden dos son pares y alguno tiene valor absoluto dos; los menores de orden
tres son múltiplos de cuatro y alguno tiene valor absoluto cuatro; por
último, \(\lvert\det H_4\rvert=16\). Los divisores determinantes son,
por tanto, \(1,2,4,16\), de donde se obtienen los factores invariantes
\(1,2,2,4\). La suma directa de tres bloques repite estos factores y tiene
índice \(16^3=4096\). Puesto que \(H_4^{-1}=H_4/4\), la última
congruencia equivale a la integralidad de la preimagen.
\end{proof}

\section{Transportador ciclotómico relativo entre los sectores tres y cuatro}
\label{sec:transportador-ciclotomico-a34-rev8}

Sea \(\mathsf R\) la rotación regular del ciclo \(C_{12}\) sobre
\(\mathbb Q^{12}\), y sean
\[
 W_d^{\rm cyc}:=\ker\Phi_d(\mathsf R)
\]
sus sectores ciclotómicos. Los proyectores racionales sobre los sectores de
órdenes tres y cuatro pueden escribirse
\[
 \Pi_3=
 \frac14(I+\mathsf R^3+\mathsf R^6+\mathsf R^9)
 -\frac1{12}\sum_{k=0}^{11}\mathsf R^k,
\]
\[
 \Pi_4=
 \frac16(I-\mathsf R^2+\mathsf R^4-\mathsf R^6
              +\mathsf R^8-\mathsf R^{10}).
\]
Para fijar el orden de coordenadas, sea
\[
 \mathsf P_{\rm orb}(z_1,\ldots,z_{12})
 =\bigl((z_1,z_4,z_7,z_{10}),
        (z_2,z_5,z_8,z_{11}),
        (z_3,z_6,z_9,z_{12})\bigr).
\]
La transformada global usada arriba es, en el orden natural de \(C_{12}\),
\[
 \mathcal H_{12}
 :=\mathsf P_{\rm orb}^{-1}
   (H_4\oplus H_4\oplus H_4)\mathsf P_{\rm orb}.
\]
Es decir, \(H_4\oplus H_4\oplus H_4\) actúa en el orden de las tres
órbitas de paso tres, no sobre tres cuaternas contiguas del orden natural.
Definimos entonces el transportador relativo
\[
 A_{34}:=\left.\Pi_4\mathcal H_{12}\Pi_3
          \right|_{W_3^{\rm cyc}}:
 W_3^{\rm cyc}\longrightarrow W_4^{\rm cyc}.
\]

\begin{theorem}[Transportador ciclotómico dodecafásico]
\label{thm:transportador-ciclotomico-a34-rev8}
En las bases
\[
 u_1=(1,-1,0)^{\times4},\qquad
 u_2=(0,1,-1)^{\times4}
\]
de \(W_3^{\rm cyc}\), y
\[
 v_1=(1,0,-1,0)^{\times3},\qquad
 v_2=(0,1,0,-1)^{\times3}
\]
de \(W_4^{\rm cyc}\), se cumple
\[
 A_{34}u_1=\frac23(v_1-v_2),\qquad
 A_{34}u_2=\frac23(v_1+v_2),
\]
y, por tanto,
\[
 \boxed{\det A_{34}=\frac89}.
\]
En particular, \(A_{34}\) es un isomorfismo racional entre ambos sectores.
\end{theorem}

\begin{proof}
Se sustituyen las cuatro sucesiones periódicas anteriores en las expresiones
de \(\Pi_3,\Pi_4\) y en los tres bloques de Hadamard. La matriz de \(A_{34}\)
en las bases declaradas es
\[
 \begin{pmatrix}
  2/3&2/3\\
 -2/3&2/3
 \end{pmatrix},
\]
cuyo determinante es \(8/9\).
\end{proof}

La involución racional que intercambia \(W_3^{\rm cyc}\) y
\(W_4^{\rm cyc}\) mediante \(A_{34}\) y \(A_{34}^{-1}\), y fija
\(W_{12}^{\rm cyc}\), transporta la suma
\(W_4^{\rm cyc}\oplus W_{12}^{\rm cyc}\) a
\(W_3^{\rm cyc}\oplus W_{12}^{\rm cyc}\). En esta última suma, el polinomio
ciclotómico de la rotación es
\[
 \Phi_3(x)\Phi_{12}(x),
\]
que coincide con el polinomio del Coxeter de tipo \(E_6\). Ésta es una firma
espectral exacta y prepara el corredor geométrico posterior.

\begin{statusbox}[title={Transporte racional y continuación integral}]
El operador \(A_{34}\) realiza el transporte ciclotómico racional entre los
sectores de órdenes tres y cuatro. Al fijar \(W_{12}^{\rm cyc}\), sitúa la
dodecafase en el sector espectral exacto
\(\Phi_3\Phi_{12}\) de tipo \(E_6\). La continuación integral se construye
mediante dos fuentes posteriores y compatibles: el
\cref{thm:entrelazador-dodecafase-witt} proporciona la isometría
\(M_{12}\)-equivariante hacia el módulo de Witt, y el
\cref{thm:identificacion-isometrica-tres-proyectores} encadena los tres
proyectores de incidencia. Sobre esa realización integral, el
\cref{p12-83-c20-s6:thm:grafo-27-rectas-nuevo} construye el grafo de las
veintisiete rectas y el
\cref{p12-83-c20-s6:chap:holonomia-schlafli-e6-nuevo} identifica su acción de
Weyl. Así, \(A_{34}\) fija la etapa racional de la genealogía y los
entrelazadores de Witt--Schläfli realizan su prolongación reticular e
incidencial.
\end{statusbox}

La separación de signo es
\[
U_{12}^{+}
=(2378,1406,2479,0,998,0,0,0,0,0,0,0),
\]
\[
U_{12}^{-}
=(0,0,0,452,0,551,668,204,371,322,28,997),
\]
y
\[
\Usgn=U_{12}^{+}-U_{12}^{-}.
\]

\section{Segunda representación: diferencias transversales}

Con índices en \(\{1,\ldots,12\}\) y la convención cíclica
\[
 K_{m+r}:=K_{1+((m+r-1)\bmod12)},
\]
definimos
\[
(D_3K)_m=K_m-K_{m+3},
\qquad
(D_4K)_m=K_m-K_{m+4}.
\]
El extractor transversal dodecafásico es el operador
\[
 \mathcal E_{\mathrm{dod}}
 :=(D_3,D_4,Q).
\]
El subíndice ``dod'' declara su dominio. Los desplazamientos tres y cuatro
corresponden a separaciones angulares de \(90^\circ\) y \(120^\circ\) en el
ciclo orientado \(C_{12}\), pero \(\mathcal E_{\mathrm{dod}}\) no es el operador de
incidencia hexada--octada ni el núcleo angular utilizado posteriormente en
las torres.
Los datos son
\[
\begin{aligned}
D_3K={}&(-495,-116,-684,108,601,-90,\\
       &\qquad-173,-88,313,560,-397,461),\\[1mm]
D_4K={}&(-425,-281,-481,671,-255,30,\\
       &\qquad475,-543,680,251,6,-128),\\[1mm]
Q={}&6263.
\end{aligned}
\]
Si \(D_3v=D_4v=0\), el vector es constante en los pasos tres y cuatro. Como
\(\gcd(3,4)=1\), esos pasos conectan todo el ciclo y
\[
\ker D_3\cap\ker D_4=\langle\mathbf1\rangle.
\]
Por tanto
\[
\operatorname{rank}(D_3,D_4)=11.
\]
La carga total elimina la traslación constante, porque
\(Q(\mathbf1)=12\ne0\). Luego
\[
\boxed{\operatorname{rank}(D_3,D_4,Q)=12,}
\]
y
\[
K\longleftrightarrow(D_3K,D_4K,Q)
\]
es un sistema de coordenadas exacto.

\begin{theorem}[Forma normal integral del sistema transversal]
\label{thm:smith-dodecafase}
Sea
\[
 \mathcal D=(D_3,D_4,Q):
 \mathbb Z^{12}\longrightarrow\mathbb Z^{25}.
\]
Sus factores invariantes no nulos son
\[
 \operatorname{SNF}(\mathcal D)
 =\operatorname{diag}(\underbrace{1,\ldots,1}_{11},12),
\]
y, por consiguiente,
\[
 \operatorname{coker}\mathcal D
 \cong\mathbb Z^{13}\oplus\mathbb Z/12\mathbb Z.
\]
\end{theorem}

\begin{proof}
Las filas de \((D_3,D_4)\) son incidencias orientadas del grafo
\[
 \operatorname{Cay}
 \bigl(\mathbb Z/12\mathbb Z,\{\pm3,\pm4\}\bigr).
\]
Como los pasos tres y cuatro generan \(\mathbb Z/12\mathbb Z\), el grafo
es conexo.
Las incidencias de un árbol generador proporcionan once menores reducidos
unimodulares; los once primeros factores invariantes son, pues, iguales a
uno.

Todo menor de orden doce contiene la fila \(Q\), ya que el bloque de
incidencias tiene rango once. Al sumar las once primeras columnas a la
última, operación unimodular, la última columna se anula en las filas de
incidencia y vale doce en la fila \(Q\). Todos los menores máximos son
divisibles por doce. La elección de las once incidencias de un árbol da un
menor restante de valor absoluto uno y, por tanto, un determinante máximo
de valor absoluto doce. El último factor invariante es exactamente doce.
\end{proof}

La componente \(\mathbb Z/12\mathbb Z\) registra la condición de
compatibilidad integral de la carga total con el potencial reconstruido.
En particular, la unicidad racional probada por el rango se refina aquí a
una afirmación sobre la red integral y su única obstrucción torsional.

El observable firmado también puede reconstruirse desde estas coordenadas.
Para \(r=1,2,3\), sea
\[
B_r=\sum_{j=0}^{3}(D_4K)_{r+3j}.
\]
Se obtiene
\[
(B_1,B_2,B_3)=(972,-1073,101).
\]
Las sumas orbitales son
\[
A_1=\frac{Q+2B_1+B_2}{3}=2378,
\quad
A_2=A_1-B_1=1406,
\quad
A_3=A_2-B_2=2479.
\]
Si
\(d_{r,j}=(D_3K)_{r+3j}\) para \(j=0,1,2,3\), con la misma convención
cíclica, el bloque firmado de la órbita \(r\) es
\[
(A_r,d_{r,1}-d_{r,3},d_{r,0}+d_{r,2},d_{r,0}-d_{r,2}),
\]
exactamente el bloque Hadamard correspondiente.

\begin{remark}[Codominios de \(\mathcal H_{12}\) y
\(\mathcal E_{\mathrm{dod}}\)]
La escritura correcta es
\[
\mathcal H_{12}(K)=\Usgn,
\qquad
\mathcal E_{\mathrm{dod}}(K)=(D_3K,D_4K,Q).
\]
Ambas representaciones son compatibles porque reconstruyen el mismo vector,
pero no se ha definido aquí un diagrama entre sus codominios. Sus fórmulas y
codominios son distintos. Compartir los pasos \(3/4\) o los índices \(90/120\) no
autoriza a identificarlas.
\end{remark}

\section{Las tres coordenadas arquimedianas}

Los vectores de doce tríadas son
\[
\begin{aligned}
P={}&(141,592,653,589,793,238,462,643,383,279,502,884),\\
E={}&(718,281,828,459,045,235,360,287,471,352,662,497),\\
\Phi={}&(618,033,988,749,894,848,204,586,834,365,638,117).
\end{aligned}
\]
En forma concatenada:
\[
\begin{aligned}
\iota(P)&=141592653589793238462643383279502884,\\
\iota(E)&=718281828459045235360287471352662497,\\
\iota(\Phi)&=618033988749894848204586834365638117.
\end{aligned}
\]
\(P\) es la coordenada decimal compartida por las cinco regiones de
\(\pi\); el cálculo no afirma que esas regiones sean idénticas.

\section{Recurrencia de acarreo en base mil}

Fijamos
\[
c_{13}=0
\]
y calculamos de derecha a izquierda, \(m=12,11,\ldots,1\):
\[
s_m=P_m+E_m-\Phi_m-K_m+c_{m+1},
\]
\[
A_m=s_m\bmod1000,
\qquad
c_m=\frac{s_m-A_m}{1000},
\qquad0\le A_m<1000.
\]
La tabla se imprime en orden creciente de índice, pero cada acarreo procede de la fila
inferior.

\begingroup
\scriptsize
\setlength{\tabcolsep}{3.3pt}
\begin{longtable}{@{}r|rrrr|rr|rr@{}}
\caption{Cierre dodecafásico completo en base mil.}\label{tab:acarreo-alpha}\\
\toprule
\(m\)&\(P_m\)&\(E_m\)&\(\Phi_m\)&\(K_m\)&\(c_{m+1}\)&\(s_m\)&\(A_m\)&\(c_m\)\\
\midrule
\endfirsthead
\toprule
\(m\)&\(P_m\)&\(E_m\)&\(\Phi_m\)&\(K_m\)&\(c_{m+1}\)&\(s_m\)&\(A_m\)&\(c_m\)\\
\midrule
\endhead
1&141&718&618&234&0&7&007&0\\
2&592&281&033&543&0&297&297&0\\
3&653&828&988&140&\(-1\)&352&352&0\\
4&589&459&749&729&\(-1\)&\(-431\)&569&\(-1\)\\
5&793&045&894&659&\(-2\)&\(-717\)&283&\(-1\)\\
6&238&235&848&824&\(-1\)&\(-1200\)&800&\(-2\)\\
7&462&360&204&621&0&\(-3\)&997&\(-1\)\\
8&643&287&586&058&\(-1\)&285&285&0\\
9&383&471&834&914&\(-1\)&\(-895\)&105&\(-1\)\\
10&279&352&365&794&0&\(-528\)&472&\(-1\)\\
11&502&662&638&146&0&380&380&0\\
12&884&497&117&601&0&663&663&0\\
\bottomrule
\end{longtable}
\endgroup

Los acarreos son
\[
(c_1,\ldots,c_{13})
=(0,0,0,-1,-1,-2,-1,0,-1,-1,0,0,0).
\]
Los valores negativos de \(s_m\) son préstamos ordinarios en base mil. La
división euclídea fija resto y cociente de forma única.

La salida es
\[
\boxed{
A=(007,297,352,569,283,800,997,285,105,472,380,663),}
\]
y por tanto
\[
\boxed{
\begin{aligned}
\alpha_{12}
&:=\frac{7297352569283800997285105472380663}{10^{36}},\\
&=0.007297352569283800997285105472380663.
\end{aligned}}
\]
Este racional es la ventana dodecafásica inicial de la sección compatible
generada por el transporte TPK. Si
\(\rho_{N',N}:C_{N'}^\alpha\to C_N^\alpha\) son sus truncamientos, la
salida completa es
\[
 \boxed{
 \alpha_{\mathrm{HMT}}
 :=\varprojlim_N\alpha_N,\qquad
 \rho_{N',N}(\alpha_{N'})=\alpha_N,\qquad
 \alpha_{36}=\alpha_{12}.}
\]
El retorno de fase prolonga la memoria y estrecha los cilindros
supervivientes a profundidad arbitraria; no reinicia el estado ni termina la
construcción en treinta y seis cifras.

\section{La identidad afín completa}

Coordenada a coordenada,
\[
\boxed{
P_m+E_m-\Phi_m-K_m+c_{m+1}
=A_m+1000c_m.}
\]
Si
\[
C=(c_1,\ldots,c_{12}),
\qquad
C^+=(c_2,\ldots,c_{13}),
\]
la forma vectorial correcta es
\[
\boxed{P+E-\Phi-K+C^+-1000C=A.}
\]
Para evitar identificar las colas finitas con las constantes completas,
definimos
\[
 \tau_{36}(x)=10^{-36}\left\lfloor10^{36}
 \bigl(x-\lfloor x\rfloor\bigr)\right\rfloor,
 \qquad
 \kappa(K)=10^{-36}\iota(K).
\]
La notación decimal exacta de la identidad concatenada es entonces
\[
 \tau_{36}(\pi)+\tau_{36}(e)-\tau_{36}(\varphi)-\kappa(K)
 =\alpha_{12}.
\]
Sin el truncamiento explícito, la inmersión \(\kappa\) y el cociclo
\(C^+-1000C\), la igualdad coordenada sería falsa.

Definimos
\[
\iota(v_1,\ldots,v_{12})
=\sum_{m=1}^{12}v_m1000^{12-m}.
\]
Al multiplicar cada recurrencia por \(1000^{12-m}\) y sumar, los acarreos
telescopan:
\[
\iota(P)+\iota(E)-\iota(\Phi)-\iota(K)-\iota(A)
=1000^{12}c_1-c_{13}.
\]
Como \(c_1=c_{13}=0\):
\[
\boxed{
\iota(P)+\iota(E)-\iota(\Phi)-\iota(K)=\iota(A).}
\]
Los dos enteros restantes son
\[
\iota(K)=234543140729659824621058914794146601,
\]
\[
\iota(A)=007297352569283800997285105472380663.
\]

\section{Inversión y equivalencia de las tres representaciones}

La recurrencia se invierte por
\[
K_m=P_m+E_m-\Phi_m+c_{m+1}-A_m-1000c_m,
\]
y la representación entera concatenada por
\[
\iota(K)=\iota(P)+\iota(E)-\iota(\Phi)-\iota(A).
\]
El diagrama completo es
\[
\boxed{
K\longleftrightarrow\Usgn,
\qquad
K\longleftrightarrow(D_3K,D_4K,Q),
\qquad
K\longleftrightarrow A\mid(P,E,\Phi,c_{13}=0).}
\]
Las dos primeras son representaciones lineales; la tercera, una aplicación
afín reversible en la fibra
condicionada por \(P,E,\Phi\) y la frontera.

\begin{remark}[Conmutatividad del cierre]
La aplicación directa obtiene \(K\) desde la representación firmada y produce
\(A\). La fórmula inversa demuestra que la salida y la frontera recuperan el
mismo vector sin libertad residual. La invertibilidad no constituye una
evidencia estadística independiente; la conmutatividad de las tres
representaciones es una propiedad algebraica del cierre.
\end{remark}

La unicidad se expresa en tres niveles de la misma construcción:
\begin{enumerate}
  \item fijados el vector firmado, \(P,E,\Phi\) y la frontera, el vector y la
  salida son únicos;
  \item las representaciones son reversibles y conmutativas;
  \item la genealogía completa
  \(\mathrm{APP}\to\mathrm{TRIT}\to\mathrm{TPK}\) produce esas primitivas
  antes del cierre local, y sus truncamientos compatibles prolongan la
  identidad a la sección coinductiva \(\alpha_{\mathrm{HMT}}\).
\end{enumerate}

\section{Condición de frontera dodecafásica}

La primera ventana dodecafásica usa \(c_{13}=0\) y termina en
\[
\ldots|380|663.
\]
Esta frontera forma parte del dominio de la aplicación afín reversible. La
prolongación no introduce otro sistema: itera el mismo transporte graduado
del TPK sobre
\[
 w_6\to w_{12}\to w_{18}\to w_{24}\to w_{30}\to R_{36}\to G_9
\]
y construye las ventanas sucesivas con retorno de fase, avance de memoria y
acarreo terminal compatible. El
\cref{sec:l9s102:alpha-dos-vias} demuestra que los cuadrados de truncamiento
de esta vía y del lector analítico conmutan y que sus cilindros
supervivientes poseen intersección unitaria.

\subsection{Clase integral del acarreo}

La condición terminal puede compararse con el problema periódico. Sea
\(S\) el desplazamiento cíclico sobre \(\mathbb Z^{12}\) y, para una base
entera \(b\geq2\), definamos
\[
 \partial_b=S-bI.
\]
La identidad afín periódica adopta la forma
\[
 A=P+E-\Phi-K+\partial_bC.
\]
Para todo \(h\in\mathbb Z^{12}\), la transformación
\[
 K\longmapsto K+\partial_bh,
 \qquad
 C\longmapsto C+h
\]
conserva \(A\). Así, la elección de acarreos es una elección de
representante dentro de una clase integral.

\begin{proposition}[Cociente periódico del acarreo]
\label{prop:cociente-periodico-acarreo}
En un ciclo de longitud doce,
\[
 \operatorname{coker}(S-bI)\cong
 \mathbb Z/(b^{12}-1)\mathbb Z.
\]
La concatenación en base \(b\), tomada módulo \(b^{12}-1\), determina la
clase periódica.
\end{proposition}

\begin{proof}
El módulo cíclico se identifica con
\(\mathbb Z[t]/(t^{12}-1)\). Imponer la relación \(t=b\) produce
\(\mathbb Z/(b^{12}-1)\mathbb Z\). La evaluación del polinomio de
coeficientes es la concatenación en base \(b\), con la orientación fijada
por el orden de los doce sectores.
\end{proof}

El cierre de \(\alpha\) no emplea el cociente periódico: utiliza una cadena
orientada con extremo derecho y exige \(c_{13}=0\). La división euclídea
selecciona entonces un único representante. Esta diferencia explica por
qué la frontera abierta forma parte de la definición de la aplicación y no
puede reemplazarse por una identificación cíclica de los acarreos.

\section{Comparación con la referencia metrológica}

Después de la generación, la aritmética entera de la primera ventana
verifica el reconocimiento
\[
\boxed{
\alpha_{12}
=10^{-36}
\left\lfloor\frac{10^{36}}{137.035999084}\right\rfloor.}
\]
Su inversa es
\[
\alpha_{12}^{-1}
=137.03599908400000000000000000000001066588\ldots,
\]
y reproduce exactamente la cadena central publicada en el ajuste de 2018 del
Committee on Data for Science and Technology (CODATA) \cite{Tiesinga2021CODATA}:
\[
\boxed{137.035999084.}
\]
No es una coincidencia de coma flotante. Las treinta y seis cifras de
\(\alpha_{12}\) constituyen el truncamiento exacto del recíproco de la
cadena central impresa. Esta comprobación finita no define
\(\alpha_{\mathrm{HMT}}\) ni limita su profundidad: la sección coinductiva
ya producida determina todas sus ventanas compatibles, mientras CODATA sólo
reconoce después el prefijo que su resolución experimental permite
contrastar.

El reconocimiento metrológico no forma parte de la recurrencia. El resultado
matemático es la identidad exacta de las dos vías
\[
 \alpha_A=\alpha_B=\alpha_{\mathrm{HMT}},
\]
cuyas ventanas dodecafásicas son determinadas por el vector firmado, las
tres palabras generadas y las fronteras compatibles. La tabla de CODATA
sólo compara después una salida ya fijada.

La comparación con ajustes metrológicos posteriores es un control externo
separado. La salida HMT queda congelada y, por ello, puede ser contrastada sin
modificar el generador.

\section{Independencia de las coordenadas y rigidez}

Las pruebas del cierre distinguen modificaciones profundas de perturbaciones
locales:
\begin{center}
\small
\begin{tabularx}{\textwidth}{@{}lYY@{}}
\toprule
Variante & Resultado & Qué controla \\
\midrule
protocolo canónico derecha--izquierda & igualdad exacta & referencia \\
usar \(U_6\Vert U_6\) como \(K\) & falla macroscópicamente & distinción de tipos \\
rotar \(K\) & rompe el cierre & origen dodecafásico \\
reflejar \(K\) & invierte el borde & orientación \\
acarreo izquierda--derecha & cambia la palabra & sentido del transporte \\
perturbar \(K_i\) en una unidad & falla la palabra exacta & rigidez profunda \\
\bottomrule
\end{tabularx}
\end{center}
Una perturbación local puede alterar sólo una cifra profunda; sigue siendo una
refutación de la igualdad exacta. Las variantes estructurales producen fallos
mayores y muestran que el cierre no es invariante bajo órdenes arbitrarios.

\begin{remark}[Estructura del cierre dodecafásico]
El cierre de \(\alpha\) comprende un ciclo orientado, dos representaciones
exactas del vector \(K\), doce entradas tipadas, una recurrencia de derecha a
izquierda, trece acarreos, una condición de frontera y una identidad entera
telescópica. Estos datos determinan la salida de manera reproducible.
\end{remark}

\end{HMTScientificOwnerSubordinated}
\begingroup
\let\HMTRecoveredOriginalPart\part
\let\part\section
\begin{HMTScientificOwnerSubordinated}
% Adaptador local de aristas para el propietario exacto
% propietarios_exactos/alpha/03a_alpha_y_dependencias_transversales.tex.
% El propietario conserva sus once llamadas \HMTInputLive. En este ensamblado
% sus cuerpos tienen residencias exactas o finales propias; la llamada verifica
% que la arista esté registrada y la deja en el log, pero no vuelve a incluir el
% cuerpo. La existencia y el orden de ambas puntas se auditan contra el grafo de
% inputs. Así no se anteponen lectores ni se publica la copia histórica débil de
% 08a.

\makeatletter
\newcommand{\HMTRegisterLiveInputEdge}[3]{%
  \expandafter\def\csname hmt@liveinput@residence@\detokenize{#1}\endcsname{#2}%
  \expandafter\def\csname hmt@liveinput@order@\detokenize{#1}\endcsname{#3}%
}
\newcommand{\HMTInputLive}[1]{%
  \ifcsname hmt@liveinput@residence@\detokenize{#1}\endcsname
    \edef\HMTLiveInputResidence{%
      \csname hmt@liveinput@residence@\detokenize{#1}\endcsname}%
    \edef\HMTLiveInputOrder{%
      \csname hmt@liveinput@order@\detokenize{#1}\endcsname}%
    \typeout{HMT-INPUT-LIVE-EDGE: #1.tex =>
      \HMTLiveInputResidence\space [\HMTLiveInputOrder]; body not reloaded}%
  \else
    \PackageError{hmt-alpha-live-adapter}
      {Arista sin residencia efectiva: #1.tex}
      {Registre un propietario exacto o final; no se admite una copia reductiva.}%
  \fi
}

% Todas las aristas del propietario 03a, en orden de aparición.
\HMTRegisterLiveInputEdge
  {sections/hmt/08_dodecafase_alpha}
  {sections/hmt/08_dodecafase_alpha.tex}
  {UPSTREAM-CH31-PREOWNER}
\HMTRegisterLiveInputEdge
  {sections/hmt/08a_alpha_dos_vias_t1_rev8}
  {../colaboracion/parte_iii/source/public_final/alpha/08a_alpha_dos_vias_t1_rev8_body.tex}
  {DOWNSTREAM-CH31-FINAL}
\HMTRegisterLiveInputEdge
  {sections/hmt/08b_alpha_t1_10000_rev9}
  {../colaboracion/parte_iii/source/public_final/alpha/08b_alpha_t1_10000_rev9.tex}
  {DOWNSTREAM-CH31-FINAL}
\HMTRegisterLiveInputEdge
  {sections/ampliacion_20260818/03_extension_holografica_euler}
  {sucesor_102/deltas_ley9/c57_sistema_operatorio_euler_hmt.tex}
  {DOWNSTREAM-CH57-SUCCESSOR}
\HMTRegisterLiveInputEdge
  {sections/hmt/07b_realizacion_analitica_contraangulo}
  {sucesor_102/wrappers/capitulo_38_jerarquia_unica.tex}
  {DOWNSTREAM-CH38-FINAL}
\HMTRegisterLiveInputEdge
  {sections/hmt/16c_transportador_positivo_accion}
  {../colaboracion/parte_iii/source/public_final/ascensos_20260826/16d_medida_accion_positiva.tex + sucesor_102/wrappers/capitulo_38_jerarquia_unica.tex}
  {DOWNSTREAM-CH35-AND-CH38-FINAL}
\HMTRegisterLiveInputEdge
  {sections/hmt/delta_297/16c_entrelazador_dodecafase_witt}
  {sections/hmt/delta_297/16c_entrelazador_dodecafase_witt.tex}
  {UPSTREAM-PART-I-CH20}
\HMTRegisterLiveInputEdge
  {sections/hmt/delta_297/16d_medida_accion_positiva}
  {../colaboracion/parte_iii/source/public_final/ascensos_20260826/16d_medida_accion_positiva.tex}
  {DOWNSTREAM-CH35-FINAL}
\HMTRegisterLiveInputEdge
  {sections/md/04b_banda_particula_virtual}
  {integracion_83/parte_iv/chapters/75_definicion_particula.tex}
  {DOWNSTREAM-PART-IV-CH94}
\HMTRegisterLiveInputEdge
  {sections/md/06f_biografia_electron_rev6}
  {integracion_83/parte_iv/body/B001_ch48_electron_primera_realizacion_a3f11d97b332_06f_biografia_electron_rev6.tex}
  {DOWNSTREAM-PART-IV-CH58}
\HMTRegisterLiveInputEdge
  {sections/ampliacion_20260818/07_transduccion_metrologica_estados_genealogicos}
  {integracion_83/continuo_constantes/tex/transiciones/umbral_09_metrologia_y_programa_cuerpo.tex}
  {DOWNSTREAM-TRANSITION-III-IV}

% Al no reincluir cuerpos, las antiguas operaciones de descenso/restauración
% deben ser neutras. El grupo exterior impide que estas definiciones se filtren.
\renewcommand{\HMTDemoteHeadings}{}
\renewcommand{\HMTRestoreHeadings}{}

% El propietario conserva su cabecera primaria. Sus otras dos \chapter* eran
% rótulos de cuerpos que ahora residen después; se suprimen junto con las dos
% entradas de índice asociadas, evitando capítulos vacíos sin tocar el dueño.
\let\HMTAlphaOwnerChapter\chapter
\let\HMTAlphaOwnerAddContentsLine\addcontentsline
\newif\ifHMTAlphaOwnerPrimaryChapterSeen
\newif\ifHMTAlphaOwnerSuppressNextContentsLine
\HMTAlphaOwnerPrimaryChapterSeenfalse
\HMTAlphaOwnerSuppressNextContentsLinefalse
\long\def\chapter{%
  \@ifstar{\HMTAlphaOwnerStarChapter}{%
    \@ifnextchar[{\HMTAlphaOwnerOptionalChapter}{\HMTAlphaOwnerPlainChapter}}%
}
\long\def\HMTAlphaOwnerOptionalChapter[#1]#2{%
  \HMTAlphaOwnerPrimaryChapterSeentrue
  \HMTAlphaOwnerChapter[#1]{#2}%
}
\long\def\HMTAlphaOwnerPlainChapter#1{%
  \HMTAlphaOwnerPrimaryChapterSeentrue
  \HMTAlphaOwnerChapter{#1}%
}
\long\def\HMTAlphaOwnerStarChapter#1{%
  \ifHMTAlphaOwnerPrimaryChapterSeen
    \HMTAlphaOwnerSuppressNextContentsLinetrue
    \typeout{HMT-ALPHA-OWNER-EMPTY-HEADING-SKIP: \detokenize{#1}}%
  \else
    \HMTAlphaOwnerPrimaryChapterSeentrue
    \HMTAlphaOwnerChapter*{#1}%
  \fi
}
\renewcommand{\addcontentsline}[3]{%
  \ifHMTAlphaOwnerSuppressNextContentsLine
    \HMTAlphaOwnerSuppressNextContentsLinefalse
    \typeout{HMT-ALPHA-OWNER-EMPTY-TOC-SKIP: #1/#2}%
  \else
    \HMTAlphaOwnerAddContentsLine{#1}{#2}{#3}%
  \fi
}
\makeatother

\part{Cierre dodecafásico y dependencias transversales del estado}

\chapter[Derivación interna de la constante de estructura fina]{Derivación
interna de \texorpdfstring{\(\alpha_{\rm HMT}\)}{alpha HMT}: genealogía de
supervivencia, cierre dodecafásico, dos construcciones compatibles y
bifurcación de Witt}
\label{chap:alpha-rev7}
\HMTClaim{HMT-R-ALPHA-DODECAPHASE}
\HMTClaim{HMT-R-DODEC-EXTRACTOR}

\paragraph{Cuatro estratos de \(\alpha_{\rm HMT}\).}
\begin{enumerate}
 \item El constructor genealógico es el estado terminal enriquecido con sus
 dos publicaciones paralelas \(B_{\rm term}\) y \(U_{12}^{\rm sgn}\).
 \item La caracterización analítica es la cadena reversible
 \(U_{12}^{\rm sgn}\leftrightarrow K\to A\to\alpha_{12}\), seguida por su
 prolongación coinductiva.
 \item Al ser adimensional, su transducción común no introduce una unidad:
 fija la carta angular \(A_{\rm deg}=1000\alpha_{\rm HMT}\) utilizada por
 los canales posteriores.
 \item La comparación con CODATA se abre únicamente después de obtener
 \(\alpha_{\rm HMT}\); no selecciona \(K\), el acarreo ni la prolongación.
\end{enumerate}
La genealogía completa conserva explícitamente la cadena de entrada y alcanza
un estado terminal enriquecido. Sus dos publicaciones son paralelas:
\[
\begin{aligned}
 w_6\longrightarrow w_{30}\longrightarrow R_{36}\longrightarrow G_9
 \longrightarrow x_{\rm term}
 &\xrightarrow{\ \pi_{\rm term}\ }B_{\rm term},\\
 x_{\rm term}
 &\xrightarrow{\ R_{\rm sgn}\ }U_{12}^{\rm sgn}
 \xleftrightarrow{\ \mathcal H_{12}\ }K
 \longrightarrow A\longrightarrow\alpha_{12}.
\end{aligned}
\]
\(B_{\rm term}\) no genera \(U_{12}^{\rm sgn}\): ambos conservan
coordenadas diferentes del mismo estado. La transformación
\(\mathcal H_{12}\) es invertible, y el extractor
\((D_3K,D_4K,Q(K))\) tiene rango \(12\). En este dominio de primitivas
terminales explícitas, la cadena no recibe \(\alpha\), CODATA, una masa ni
otra coordenada metrológica como entrada.
La evaluación arquimediana forma \((P,E,\Phi,K,c_{13}=0)\), aplica el cociclo
de acarreo y deriva \(\alpha_{\rm HMT}\). La evaluación de incidencia
transporta las mismas palabras y el mismo sello a
\(B_K\subset H_\alpha^-\), que será la entrada de la rama Witt--Golay. La
construcción E21/22 y la prolongación remota
\(\mathsf T_{\alpha,\infty}\) ---denominada \(T1\) en la notación
original--- conservan dominios y mapas propios como desarrollos de la
genealogía dodecafásica principal.

\HMTDemoteHeadings
\HMTInputLive{sections/hmt/08_dodecafase_alpha}
\HMTInputLive{sections/hmt/08a_alpha_dos_vias_t1_rev8}
\HMTInputLive{sections/hmt/08b_alpha_t1_10000_rev9}
\HMTRestoreHeadings

\HMTInputLive{sections/ampliacion_20260818/03_extension_holografica_euler}

\HMTInputLive{sections/hmt/07b_realizacion_analitica_contraangulo}
\HMTInputLive{sections/hmt/16c_transportador_positivo_accion}
\HMTInputLive{sections/hmt/delta_297/16c_entrelazador_dodecafase_witt}
\begingroup
\def\HMTInternalTraceability{1}
\HMTInputLive{sections/hmt/delta_297/16d_medida_accion_positiva}
\endgroup

\HMTInputLive{sections/md/04b_banda_particula_virtual}

\chapter*{Construcción holográfica completa del electrón como dependencia del lector}
\addcontentsline{toc}{chapter}{Construcción holográfica completa del electrón como dependencia del lector}
\HMTInputLive{sections/md/06f_biografia_electron_rev6}

\chapter*{Transducción metrológica de estados genealógicos}
\addcontentsline{toc}{chapter}{Transducción metrológica de estados genealógicos}
\HMTInputLive{sections/ampliacion_20260818/07_transduccion_metrologica_estados_genealogicos}

\end{HMTScientificOwnerSubordinated}
\let\part\HMTRecoveredOriginalPart
\endgroup

\HMTFlushCurrentChapter
\HMTSuccessorChapterInput
  {Clausura dodecafásica de \texorpdfstring{\(\alpha_{\mathrm{HMT}}\)}{alpha HMT} y lector analítico \texorpdfstring{\(E_{21/22}^{(9)}\)}{E21/22 de orden 9}: generación reversible, raíz simple y comparador tipado}
  {../colaboracion/parte_iii/tex/capitulos_finales/31_dos_derivaciones_alpha.tex}
% Insercion editor-ready para el capitulo 31.

\section{Dos lectores internos de \texorpdfstring{\(\alpha\)}{alpha}: clausura dodecafásica y núcleo analítico nonádico de orden \texorpdfstring{\(9\)}{9}}
\label{sec:l9s102:alpha-dos-vias}

Esta sección desarrolla la coordenada \(\alpha\) de la
\hyperref[sec:tesis-inaugural-helicoidal-hmt-md]{tesis inaugural}. Ambos
lectores pertenecen a la misma genealogía APP--TRIT--TPK y preceden a la
comparación metrológica. La vía A opera mediante la clausura dodecafásica
reversible del estado firmado y publica la sección coinductiva
\(\alpha_{\mathrm{HMT}}\). La vía B construye, sin consultar esa salida, el
núcleo analítico nonádico exacto de orden nueve
\(E_{21/22}^{(9)}:=P_9-\log_{10}(10/9)\) y publica su raíz positiva simple y
única \(\xi_9\). Sus dominios, operaciones y certificados se mantienen
diferenciados:
\begin{equation}
 X_{12}^{\mathrm{sig}}
 \xrightarrow{\ \Gamma_{12}^{\alpha}\ }
 \alpha_{12},
 \qquad
 E_{21/22}^{(9)}:I_\alpha\longrightarrow\mathbb R,
 \quad E_{21/22}^{(9)}(\xi_9)=0,
 \quad (E_{21/22}^{(9)})'(\xi_9)\ne0.
 \label{eq:l9s102:alpha-dos-mapas}
\end{equation}
La clausura entera publica primero la ventana finita \(\alpha_{12}\) y su
familia compatible de prolongaciones; su límite inverso es
\(\alpha_{\mathrm{HMT}}\). El segundo lector publica \(\xi_9\) como testigo
finito de orden nueve y, por iteración del mismo transporte graduado del TPK,
construye el sistema compatible completo
\(E_{21/22}^{(6+3m)}\). Los dos sistemas inversos determinan en cada
profundidad el mismo cilindro superviviente \(C_N^\alpha\); sus diámetros
tienden a cero. Por tanto, la intersección es unitaria y el teorema de las
dos vías establece
\begin{equation}
 \alpha_A:=\varprojlim_N\rho_N(\alpha_{12}),
 \qquad
 E_{21/22}:=\varprojlim_m E_{21/22}^{(6+3m)},
 \qquad
 \alpha_B:=\operatorname*{arg\,zero}_{x\in I_\alpha}E_{21/22}(x),
 \qquad
 \boxed{\alpha_A=\alpha_B=\alpha_{\mathrm{HMT}}}.
 \label{eq:l9s102:alpha-identidad-limite-dos-vias}
\end{equation}
Ninguna vía introduce un valor objetivo como dato de entrada; la
identificación metrológica pertenece únicamente al reconocimiento
convencional posterior. La igualdad
\(\operatorname{Tr}_9(\xi_9^{-1})=
\operatorname{Tr}_9(\alpha_{12}^{-1})=137.035999084\) es un certificado
finito de la identidad anterior, nunca su definición ni su alcance máximo.
La ley de prolongación no se añade como dato: es la iteración del transporte
graduado ya producido por APP--TRIT--TPK, con retorno de fase, avance de
memoria y truncamientos compatibles.

\subsection{Núcleo analítico \texorpdfstring{\(P_6\)}{P6} y prolongación exacta de orden \texorpdfstring{\(9\)}{9}}
\label{sec:l9s102:alpha-p6}

La caracterización de vacancias emplea
\begin{equation}
 \begin{aligned}
 P_6(x)={}&2\pi_{\mathrm{HMT}}x-\frac74x^2
 +\frac{x^3}{2\pi_{\mathrm{HMT}}}
 +\frac{x^4}{20}-\frac{2x^5}{21}-\frac{x^6}{46},\\
 P_6(x)={}&\log_{10}\frac{10}{9}.
 \end{aligned}
 \label{eq:l9s102:p6}
\end{equation}
La raíz positiva certificada de este jet satisface
\begin{equation}
 x_6^{-1}=137.03599908400002702009\ldots .
 \label{eq:l9s102:raíz-p6}
\end{equation}
El polinomio \(P_6\) proporciona el jet inicial de la caracterización
analítica. La holonomía nonádica prolonga exactamente sus grados \(7,8,9\) y
la misma regla graduada se itera sobre los grados sucesivos. El cociente de
jets es el primer certificado finito de esa prolongación coinductiva; no
selecciona ni sustituye su sección límite.

\subsection{Genealogía APP--TRIT--TPK de los invariantes \texorpdfstring{\(120,45,11\)}{120, 45, 11} y \texorpdfstring{\(495\)}{495}}
\label{sec:l9s102:alpha-genealogia-7-9}

La semicorona de fase del TPK construye
\begin{equation}
 |\mathcal F_8^{\mathrm{ph}}|=120,
 \qquad |L_{\mathrm{ph}}|=45.
 \label{eq:l9s102:alpha-120-45}
\end{equation}
APP aporta el bloque central
\begin{equation}
 B_c=\begin{pmatrix}7&2\\2&7\end{pmatrix},
 \qquad
 \det B_c=45,
 \qquad
 \operatorname{Spec}(B_c)=\{9,5\}.
 \label{eq:l9s102:alpha-bc}
\end{equation}
El extractor transversal dodecafásico proporciona
\begin{equation}
 r_{\mathrm{dod}}=\operatorname{rank}(D_3,D_4)=11,
 \qquad
 (D_3K)_1=-495,
 \qquad
 495=45\cdot11=\det(B_c)\,r_{\mathrm{dod}}.
 \label{eq:l9s102:alpha-495}
\end{equation}
Los coeficientes del jet prolongado proceden, por tanto, de tres
invariantes previos de APP--TRIT--TPK:
\begin{equation}
 -\frac1{120}\longmapsto\frac1{45}
 \longmapsto\frac2{45\cdot11}.
 \label{eq:l9s102:alpha-coeficientes}
\end{equation}

\subsection{Resolvente nilpotente \texorpdfstring{\((I-N)^{-1}\)}{(I-N)^-1} y construcción exacta de la cola \texorpdfstring{\(T_{7:9}\)}{T7:9}}
\label{sec:l9s102:alpha-resolvente}

Sobre la base \(e_7,e_8,e_9\), definimos
\begin{equation}
 Ne_7=-\frac83e_8,
 \qquad
 Ne_8=\frac2{11}e_9,
 \qquad
 Ne_9=0,
 \qquad
 v_7=-\frac1{120}e_7.
 \label{eq:l9s102:alpha-nilpotente}
\end{equation}
Como \(N^3=0\), su resolvente finito actua exactamente por
\begin{equation}
 (I-N)^{-1}v_7
 =v_7+Nv_7+N^2v_7
 =-\frac{e_7}{120}+\frac{e_8}{45}+\frac{2e_9}{495}.
 \label{eq:l9s102:alpha-resolvente-finito}
\end{equation}
La evaluación \(e_n\mapsto x^n\) produce
\begin{equation}
 \boxed{
 T_{7:9}(x)=-\frac{x^7}{120}+\frac{x^8}{45}
 +\frac{2x^9}{495}.}
 \label{eq:l9s102:alpha-t79}
\end{equation}

\subsection{Jet nonádico \texorpdfstring{\(P_9\)}{P9}, raíz prolongada y escala de residuos certificados}
\label{sec:l9s102:alpha-jet-p9}

El jet nonádico de orden \(9\) queda definido por
\begin{equation}
 P_9(x)=P_6(x)+T_{7:9}(x),
 \qquad
 P_9(x)=\log_{10}\frac{10}{9}.
 \label{eq:l9s102:alpha-p9}
\end{equation}
\begin{theorem}[Existencia, simplicidad y unicidad de la segunda salida analítica]
\label{thm:l9s102:raiz-p9-unica}
El operador \(E_{21/22}^{(9)}\) posee una única raíz positiva \(\xi_9\) en
\(I_\alpha=(0,10^{-2})\). Está encerrada por el intervalo racional exacto
\begin{equation}
 \frac{729735256928380099}{10^{20}}
 <\xi_9<
 \frac{729735256928380100}{10^{20}},
 \label{eq:l9s102:encierro-raiz-p9}
\end{equation}
y su inversa satisface
\begin{equation}
 \xi_9^{-1}
 =137.0359990840000000000000111926291900\ldots .
 \label{eq:l9s102:alpha-raíz-p9}
\end{equation}
\end{theorem}

\begin{proof}
Todos los coeficientes de \(P_9\) se extraen de los invariantes anteriores a
la resolución. Para \(0\leq x\leq10^{-2}\), usando
\(3.14<\pi_{\rm HMT}<3.15\), se obtiene
\[
 (E_{21/22}^{(9)})'(x)
 \geq 2\pi_{\rm HMT}-\frac72\,10^{-2}
 -\frac{10}{21}\,10^{-8}-\frac6{46}\,10^{-10}
 -\frac7{120}\,10^{-12}>6.24.
\]
Las evaluaciones dirigidas en los extremos de
\cref{eq:l9s102:encierro-raiz-p9} tienen signos opuestos. El teorema del
valor intermedio proporciona la existencia; la cota derivativa proporciona
la unicidad y la simplicidad. La división intervalar produce la expansión de
\(\xi_9^{-1}\), que constituye el testigo de orden nueve del sistema
compatible completo.
\end{proof}

\begin{table}[htbp]
\centering
\caption{Descenso certificado del residuo de la caracterización analítica}
\label{tab:l9s102:alpha-residuos}
\begin{tabular}{@{}lc@{}}
\toprule
Núcleo evaluado en \(\alpha_{12}\) & Residuo \\
\midrule
\(P_6\) & \(9.0038886\times10^{-18}\) \\
\(P_6-x^7/120\) & \(-1.7893115\times10^{-19}\) \\
\(P_6-x^7/120+x^8/45\) & \(-2.3708601\times10^{-22}\) \\
\(P_9\) & \(3.7297132\times10^{-27}\) \\
\bottomrule
\end{tabular}
\end{table}

\subsection{Operador dodecafásico alternativo \texorpdfstring{\(A_{\mathrm{dod}}\)}{A_dod} y control de unicidad del lector tipado}
\label{sec:l9s102:alpha-control}

El operador de control
\begin{equation}
 A_{\mathrm{dod}}
 =I-\frac14(D_3^*D_3+D_4^*D_4)
 \label{eq:l9s102:alpha-control-alternativo}
\end{equation}
produce una sucesión diferente de coeficientes. Este control demuestra que
la simetría dodecafásica aislada resulta insuficiente para seleccionar la
prolongación \(T_{7:9}\); la selección pertenece a la composición tipada del
lector \(\Gamma_{E21}^{(9)}\) con el estado enriquecido. La diferencia entre
ambos operadores certifica la necesidad de conservar dominio, orientación,
graduación y memoria, y prueba la especificidad del lector finito. La cola
analítica pertenece a la iteración del lector tipado completo, no a la
simetría aislada de este operador de control.

\subsection{Alcance exacto del operador nilpotente y exclusión de colas ajenas al transporte TPK}
\label{sec:l9s102:alpha-alcance}

La aritmética de \(120,45,11,495\), el operador nilpotente, el resolvente y
la identidad \cref{eq:l9s102:alpha-t79} poseen exactitud interna. La raíz
positiva simple y única de \(E_{21/22}^{(9)}\) constituye el testigo de orden
nueve de la segunda lectura interna. La concordancia de \(P_9\) con la
ventana dodecafásica certifica el primer cuadrado finito del sistema
compatible. La unicidad canónica de la prolongación \(7{:}9\) queda tipada
por el morfismo
\begin{equation}
 \Gamma_{E21}^{(9)}:
 X_{\mathrm{APP\text{-}TRIT\text{-}TPK}}
 \longrightarrow \mathbb Q(\pi_{\mathrm{HMT}})[x]/(x^{10}),
 \label{eq:l9s102:alpha-gamma-e21}
\end{equation}
cuya acción determina conjuntamente la semilla orientada, las dos
transferencias, la graduación \(7,8,9\), el cociente de jets
\(F^7/F^{10}\) y su transporte bajo el retorno de fase.

Sea \(K=\mathbb Q(\pi_{\rm HMT},\log_{10}(10/9))\) y sea
\(j_9:K[[x]]\to K[x]/(x^{10})\) el truncamiento. Entonces
\begin{equation}
 \ker j_9=x^{10}K[[x]],
 \qquad
 E_h(x):=E_{21/22}^{(9)}(x)+x^{10}h(x),
 \qquad h\in K[[x]].
\label{eq:l9s102:alpha-limite-jets-e21}
\end{equation}
Todos los \(E_h\) poseen el mismo jet certificado. Sin embargo,
\(E_0\) y \(E_{1/729}\) tienen raíces distintas en \(I_\alpha\), pues
\(E_{1/729}(\xi_9)=\xi_9^{10}/729>0\). Esta observación demuestra únicamente
que un jet aislado no determina una cola formal. No afecta a la segunda vía
HMT, porque ésta no permite elegir libremente \(h\): la recurrencia del
transporte TPK construye la única familia compatible
\(E_{21/22}^{(6+3m)}\), y sus truncamientos conmutan con los de la vía
dodecafásica. Elegir \(h\) mediante el valor objetivo sería circular y queda
excluido; elegirlo mediante el transporte generado es precisamente la
operación ya demostrada.

El falsador separa así tres afirmaciones: invalida el certificado finito si
fallan la monotonía, el cambio de signo o el aislamiento de \(\xi_9\);
invalida la prolongación \(7{:}9\) si sus coeficientes dejan de factorizar por
\(\Gamma_{E21}^{(9)}\); e invalida la identidad completa si falla la
compatibilidad de algún cuadrado del sistema inverso o si los diámetros de
los cilindros supervivientes no tienden a cero. Ninguna de estas condiciones
introduce un morfismo pendiente: son los falsadores materiales de la
construcción que establece
\(\alpha_A=\alpha_B=\alpha_{\mathrm{HMT}}\).


\HMTFlushCurrentChapter
\HMTSuccessorChapterInput
  {Construcción interna de \texorpdfstring{\(i=\sqrt{-1}\)}{i=sqrt(-1)} mediante el operador de rotación \texorpdfstring{\(J\)}{J}: identidad \texorpdfstring{\(J^2=-I\)}{J2=-I} y geometría elíptica de fase}
  {../colaboracion/parte_iii/tex/capitulos_finales/32_raiz_menos_uno.tex}

\HMTFlushCurrentChapter
\HMTSuccessorChapterInput
  {Derivación Hensel--Witt de \texorpdfstring{\(-1/12\)}{-1/12}: prolongación compatible, regularización y memoria modular}
  {../colaboracion/parte_iii/tex/capitulos_finales/33_menos_un_doceavo.tex}

\HMTFlushCurrentChapter
\HMTSuccessorChapterInput
  {Derivación determinantal de la escala absoluta de acción: forma normal de Smith, conúcleo de orden \(16\) y valoración \texorpdfstring{\(10^{-34}\)}{10^-34}}
  {../colaboracion/parte_iii/tex/capitulos_finales/34_escala_accion.tex}
% Distribución causal exacta del delta Ley 9; contenido conservado sin síntesis.
\Needspace{28\baselineskip}
\section{Bifurcación causal posterior a \texorpdfstring{\(\alpha_{\mathrm{HMT}}\)}{alpha HMT}}
\label{sec:l9s102:bifurcacion-accion-angular}

La genealogía de la acción y del par angular forma un diagrama bifurcado.
La rama regional produce la mantisa reducida \(\eta_{\mathrm{ret}}\); la rama
determinantal produce el orden decimal \(\varepsilon_H=-34\). Ambas
confluyen en la recta de acción. El par angular se construye en una rama
hermana a partir de \(\alpha_{\mathrm{HMT}}\) y
\(\eta_{\mathrm{ret}}\):
\begin{equation}
\begin{tikzcd}[column sep=small,row sep=normal]
\alpha_{\mathrm{HMT}},\pi,\varphi
  \arrow[r] \arrow[d]
& H_5,\mathscr C_\pi
  \arrow[d]
&& \alpha_{\mathrm{HMT}},\varphi,H_4
  \arrow[d] & {}\\
x_A,D_A,\rho_\pi
  \arrow[r]
& \eta_{\mathrm{ret}}
  \arrow[r] \arrow[d]
& \hbar_{\mathrm{ret}}\in\mathcal L_S
& \Sigma_H,\varepsilon_H=-34
  \arrow[l] & {}\\
& (A_{\deg},C^*_{\deg})
  \arrow[r]
& (A_{\mathrm{rad}},C^*_{\mathrm{rad}})
  \arrow[r]
& (q_+,q_-)
  \arrow[r]
& V_{90,120}.
\end{tikzcd}
\label{diag:l9s102:dag-accion-angular}
\end{equation}
El orden expositivo situa la derivación de Planck antes de la publicación
angular completa; el diagrama conserva simultaneamente la dependencia
matemática exacta de cada salida.

\section{Conúcleo local y derivación del orden absoluto \texorpdfstring{\(10^{-34}\)}{10^-34}}
\label{sec:l9s102:orden-accion}

La forma normal de Smith del bloque dodecafásico local es
\begin{equation}
 \operatorname{SNF}(H_4)=\operatorname{diag}(1,2,2,4),
 \qquad
 \operatorname{coker}(H_4)
 \cong\mathbb Z/2\mathbb Z\oplus\mathbb Z/2\mathbb Z
 \oplus\mathbb Z/4\mathbb Z.
 \label{eq:l9s102:snf-h4}
\end{equation}
El conúcleo local posee orden \(16\). Para la sección constante
\(s_\alpha(g)=\alpha_{\mathrm{HMT}}\), la norma local y una actuación de
\(\varphi\) producen
\begin{equation}
 \boxed{
 \Sigma_H=\varphi\mathcal N_{G_H}(s_\alpha)
 =\varphi\alpha_{\mathrm{HMT}}^{16}.}
 \label{eq:l9s102:sigma-h}
\end{equation}
Las desigualdades exactas
\begin{equation}
 \alpha_{\mathrm{HMT}}^{16}<10^{-34}
 <\varphi\alpha_{\mathrm{HMT}}^{16}<10^{-33}
 \label{eq:l9s102:orden-34}
\end{equation}
determinan
\begin{equation}
 \boxed{
 \varepsilon_H=\left\lfloor\log_{10}\Sigma_H\right\rfloor=-34.}
 \label{eq:l9s102:epsilon-h}
\end{equation}


\HMTFlushCurrentChapter
\HMTSuccessorChapterInput
  {Carácter regional de grado \(5\), construcción de \texorpdfstring{\(\eta_{\mathrm{ret}}\)}{eta ret} y medida positiva de acción sobre el cociente de emisiones}
  {../colaboracion/parte_iii/tex/capitulos_finales/35_medida_accion_positiva.tex}
% Distribución causal exacta del delta Ley 9; contenido conservado sin síntesis.
\section{Carácter regional de grado \texorpdfstring{\(5\)}{5} y mantisa reducida de acción}
\label{sec:l9s102:eta-ret}

El bloque central
\begin{equation}
 \mathcal B_c=\begin{pmatrix}7&2\\2&7\end{pmatrix}
 \label{eq:l9s102:bc-accion}
\end{equation}
posee autovalores \(9\) y \(5\). La longitud \(4\) de las órbitas de paso
\(3\), el rango \(12\) del ciclo dodecafásico y
\(104976/1944=54\) fijan el carácter
\begin{equation}
 \boxed{
 H_5(\alpha,\varphi)
 =\frac12\sqrt{\frac{1000\alpha}{\varphi}}-\alpha
 +\frac9{16}\alpha^2-\frac59\alpha^3
 +\frac7{48}\alpha^4-\frac1{54}\alpha^5.}
 \label{eq:l9s102:h5}
\end{equation}
La completación nonádica construye la coordenada analítica y el determinante
regional
\begin{equation}
 x_A=\frac{50\pi}{9}\alpha_{\mathrm{HMT}},
 \qquad
 D_A=e^{-2x_A}
 =\exp\!\left(-\frac{100\pi}{9}\alpha_{\mathrm{HMT}}\right).
 \label{eq:l9s102:xa-da}
\end{equation}
La microfibra aporta \(\rho_\pi=169\), y su serie absolutamente
convergente es
\begin{equation}
 \mathscr C_\pi(\alpha)
 =169\alpha^6+\frac{D_A\alpha^7}{1-D_A\alpha}.
 \label{eq:l9s102:cpi}
\end{equation}
La mantisa reducida posterior al retorno se construye entonces por
\begin{equation}
 \boxed{
 \eta_{\mathrm{ret}}
 =H_5-\frac{90}{\pi}\mathscr C_\pi(\alpha_{\mathrm{HMT}})
 \in\mathbb R_{>0}.}
 \label{eq:l9s102:eta-ret}
\end{equation}
Su miembro derecho contiene exclusivamente objetos generados con anterioridad
por APP--TRIT--TPK y por la clausura de \(\alpha_{\mathrm{HMT}}\).



\HMTFlushCurrentChapter
\HMTSuccessorChapterInput
  {Derivación autodual de la constante de Planck y de su recta de acción: periodicidad temporal de \(108\) fases, cierre de Compton y covariancia bidireccional}
  {../colaboracion/parte_iii/tex/capitulos_finales/36_planck_autodual.tex}
% Distribución causal exacta del delta Ley 9; contenido conservado sin síntesis.
\section{Recta de acción y constante de Planck}
\label{sec:l9s102:recta-planck}

Sea \(\mathcal L_S\) la recta real de acción con base positiva
\(\mathcal U_S\). La transducción
\begin{equation}
 \tau_S:\mathbb R_{>0}\longrightarrow\mathcal L_S,
 \qquad u\longmapsto u\,10^{\varepsilon_H}\mathcal U_S
 \label{eq:l9s102:transductor-accion}
\end{equation}
publica las secciones previa y posterior al retorno:
\begin{equation}
 \hbar_{\mathrm{pre}}=H_5\,10^{-34}\mathcal U_S,
 \qquad
 \boxed{
 \hbar_{\mathrm{ret}}=\eta_{\mathrm{ret}}\,10^{-34}\mathcal U_S.}
 \label{eq:l9s102:hbar-pre-ret}
\end{equation}
Las acciones de ciclo completo son
\begin{equation}
 h_{\mathrm{pre}}=2\pi\hbar_{\mathrm{pre}},
 \qquad
 h_{\mathrm{ret}}=2\pi\hbar_{\mathrm{ret}}.
 \label{eq:l9s102:h-pre-ret}
\end{equation}
Para una fase \(\theta\), las cartas angular y de acción conmutan:
\begin{equation}
 S(\theta)
 =\hbar_{\mathrm{ret}}\theta_{\mathrm{rad}}
 =\hbar_{\mathrm{ret}}\frac{\pi}{180}\theta_{\deg}
 =h_{\mathrm{ret}}\frac{\theta_{\deg}}{360}.
 \label{eq:l9s102:accion-fase}
\end{equation}



% Transición visible de precedencia causal. Reúne propietarios ya activos sin
% repetir sus pruebas y mantiene tipadas las ramas dimensional y compleja.
% Reunión expositiva de precedencia causal entre resultados ya establecidos.
% Los propietarios íntegros continúan gobernando cada flecha.
% Propietarios reunidos por el grafo activo:
% - sucesor_102/deltas_ley9/c27_teorema_generacion_coinductiva_profundidad_arbitraria.tex
% - sucesor_102/deltas_ley9/c34_diagrama_bifurcado_y_decada_accion.tex
% - sucesor_102/deltas_ley9/c36_recta_accion_planck.tex
% - colaboracion/parte_iii/tex/capitulos_finales/32_raiz_menos_uno.tex
% - incorporaciones/parte_iv_ampliaciones_20260824/sources/editor_ready/03_pareja_planck_respuesta.tex
% - integracion_83/parte_iv/chapters/48_electron_primera_realizacion.tex

\section{Genealogía operatoria de la acción, la geometría angular y la masa}
\label{sec:l9s102:transicion-constantes-accion-geometria-masas}

La dinámica coinductiva APP--TRIT--TPK se prolonga a profundidad arbitraria
sobre un mismo estado enriquecido y sobre la estructura discreta conjunta del
continuo. En cada prolongación conserva valor, genealogía, fase, acarreo,
orientación, vacancia, ruta, frontera, memoria y supervivencia. Sus
publicaciones distinguidas son primero las tres coordenadas internas y,
mediante el sello dodecafásico, la coordenada \(\alpha_{\mathrm{HMT}}\):
\[
 \mathrm{APP}\longrightarrow\mathrm{TRIT}\longrightarrow\mathrm{TPK}
 \longrightarrow\mathcal S_{\infty}^{\mathrm{enr}}
 \longrightarrow\mathcal C_{\mathrm{cont}}^{\mathrm{disc}}
 \longrightarrow(\pi_{\mathrm{HMT}},\varphi_{\mathrm{HMT}},e_{\mathrm{HMT}})
 \longrightarrow\alpha_{\mathrm{HMT}}.
\]
La última flecha representa la publicación firmada de
\(\alpha_{\mathrm{HMT}}\) dentro de esa misma dinámica; no recibe valores
convencionales ni convierte a \(\pi,\varphi,e\) convencionales en generadores.

Sólo después se aplica el lector determinantal de acción. Éste publica la
valoración \(10^{-34}\mathcal U_S\), las dos secciones
\(\hbar_{\mathrm{pre}},\hbar_{\mathrm{ret}}\) y sus acciones de ciclo
completo
\[
 h_{\mathrm{pre}}=2\pi_{\mathrm{HMT}}\hbar_{\mathrm{pre}},
 \qquad
 h_{\mathrm{ret}}=2\pi_{\mathrm{HMT}}\hbar_{\mathrm{ret}}.
\]
A partir de este punto causal permanecen separadas dos ramas que no se
seleccionan entre sí. La rama dimensional--gravitatoria prolonga la acción y
la sección causal hasta las formas ya demostradas de la pareja de Planck,
\[
 \ell^{2}=G\hbar c^{-3},
 \qquad
 \frac{\ell}{Q}=Gc^{-4};
\]
la rama compleja publica, desde el operador interno de rotación,
\[
 u^2+1=0,
 \qquad u\longmapsto J_{+1},
 \qquad J_{+1}^{2}=-I.
\]
Ésta es la realización operatoria interna de \(i\), no la introducción de un
símbolo complejo sin antecedente. Las ramas dimensional--gravitatoria y
compleja son paralelas sobre el estado enriquecido común: la primera no
convierte a \(i\) en entrada de la escala de acción y, recíprocamente,
\(h\), \(\hbar\), \(c^{-3}\) y \(c^{-4}\) no generan a \(i\). Las identidades
de la pareja de Planck permanecen así separadas de la realización compleja.

Escribiendo \(\prec_{\mathrm{HMT}}\) para precedencia causal con conservación
de tipo ---y no para afirmar que cada objeto de la derecha sea función de
todos los de la izquierda---, se obtiene
\[
\begin{aligned}
 (\pi,\varphi,e,\alpha)_{\mathrm{HMT}}
 &\prec_{\mathrm{HMT}}(10^{-34}\mathcal U_S,\hbar,h)\\
 &\prec_{\mathrm{HMT}}
 \bigl\{G\hbar c^{-3},Gc^{-4};\,J_{+1}^{2}=-I, i\bigr\}\\
 &\prec_{\mathrm{HMT}}(A_{\mathrm{deg}},C^{*}_{\mathrm{deg}})\\
 &\prec_{\mathrm{HMT}}(\text{electrón},\text{ley de masas}).
\end{aligned}
\]
La construcción angular conserva sus entradas efectivas
\(\alpha_{\mathrm{HMT}}\) y \(\eta_{\mathrm{ret}}\): las dos ramas paralelas
precedentes fijan la disponibilidad conjunta y el tipado, pero no se promueven
retrospectivamente a selectores del ángulo o del contraángulo.

En la ruta electrónica central, los lectores bidireccionales conservan dos
etapas distintas de una sola genealogía:
\[
 \begin{aligned}
 E_e^{(0)}
   &=0.5109989224880660725\ldots\ \mathrm{MeV},
   &&\text{etapa basal},\\
 E_e^{\beta}
   &=E_e^{(0)}R_{\mathrm{act}}^{\,80-54\alpha_{\mathrm{HMT}}+6\Delta_4}
     =0.510998950690000808608\ldots\ \mathrm{MeV},
   &&\text{salida beta rectora},
 \end{aligned}
\]
con
\[
 K_D(\Gamma_e)=K_{\Omega}(\Gamma_e)=(80,54,6).
\]
La coincidencia del triple en la fibra central no identifica los dominios ni
las operaciones de \(K_D\) y \(K_{\Omega}\). En el registro canónico
histórico, \texttt{MD-ELECTRON-001} designa el candidato calibrado de la etapa
basal. Esa identificación conserva su función de procedencia, pero no gobierna
la salida beta ni selecciona la ruta; el basal y la clausura beta permanecen
como etapas tipadas de los lectores bidireccionales.

\paragraph{Criterio de refutación.}
La genealogía anterior queda refutada si un valor convencional o la salida
beta se usa para seleccionar una ruta, si la etapa basal se presenta como
cierre rector, si se identifica cualquiera de las dos ramas paralelas sin un
mapa tipado, o si el ángulo, el contraángulo, el electrón o una masa se colocan
como entradas del generador de constantes o de acción.


\HMTFlushCurrentChapter
\HMTSuccessorChapterInput
  {Operador de fase anterior a la transducción dimensional: descomposición en \(27\) sectores, radio espectral y factorización de la doble proyección}
  {../colaboracion/parte_iii/tex/capitulos_finales/37_operador_angular.tex}

% La introducción heredada del capítulo 38 despejaba primero C*/2-alpha.
% Se conserva como antecedente en la cantera, pero la residencia pública
% presenta la construcción directa eta_ret -> hbar_ret -> (A,C*).
\HMTFlushCurrentChapter
\chapter{Equivalencia exacta de las cartas global y analítica del par angular: grados, radianes, fases \texorpdfstring{\(q_{\pm}\)}{q+/-} y orientación complementaria}
\label{chap:p3-final:38:angulo_contraangulo}
La componente adimensional posterior al retorno se construye antes del
contraángulo mediante
\(\eta_{\mathrm{ret}}=H_5-(90/\pi)C_\pi(\alpha_{\mathrm{HMT}})\).
En la rama determinantal paralela se obtiene la valoración \(10^{-34}\), y
ambas ramas confluyen en la magnitud
\(\hbar_{\mathrm{ret}}=\eta_{\mathrm{ret}}10^{-34}U_S\).
Sólo después se publican
\(A_{\mathrm{deg}}=10^3\alpha_{\mathrm{HMT}}\) y
\(C^*_{\mathrm{deg}}=2(\eta_{\mathrm{ret}}+\alpha_{\mathrm{HMT}})\);
el despeje \(C^*/2-\alpha_{\mathrm{HMT}}\) queda como verificación inversa.
% Adaptador único del capítulo 38 del sucesor 102.
% Sustituye exclusivamente la jerarquía de encabezados. Los cuerpos de los
% dos propietarios científicos y del delta Ley 9 se conservan línea a línea,
% reordenados conforme al DAG causal aprobado. Véase MAPA_SEGMENTOS_CAPITULO_38.tsv.

% HMT_SOURCE_SEGMENT_BEGIN id=B00 source=colaboracion/parte_iii/source/public_final/base_83/c30_body.tex body_lines=1-2
\label{chap:p3-83-30-angulo_contraangulo}

% HMT_SOURCE_SEGMENT_END id=B00
% HMT_SOURCE_SEGMENT_BEGIN id=A00 source=colaboracion/parte_iii/source/public_final/ascensos_20260826/16f_realizacion_analitica_contraangulo.tex body_lines=1-4
\label{p3-20260826:chap:realizacion-analitica-contraangulo}
\index{contraángulo}
\index{microfibra de pi@microfibra de \(\pi\)!carácter analítico}
\index{carácter determinantal}
% HMT_SOURCE_SEGMENT_END id=A00
% HMT_SOURCE_SEGMENT_BEGIN id=D00 source=manuscrito/sucesor_102/deltas_ley9/c38_par_angular_cartas_y_canales.tex body_lines=1
% Distribución causal exacta del delta Ley 9; contenido conservado sin síntesis.
% HMT_SOURCE_SEGMENT_END id=D00

\section{Fase angular intrínseca \texorpdfstring{\(\alpha_{\mathrm{HMT}}\)}{alpha HMT}, elevación nonádica \texorpdfstring{\(10^3\)}{10^3} y precursor \texorpdfstring{\(x_A\)}{x_A}}
\label{sec:s102:c38:1}

\subsection{Construcción del precursor analítico \texorpdfstring{\(x_A\)}{x_A}, de la contracción \texorpdfstring{\(D_A\)}{D_A} y de la corrección regional \texorpdfstring{\(\mathscr C_\pi(\alpha)\)}{C_pi(alpha)}}
\label{sec:s102:c38:1:1}

% HMT_SOURCE_SEGMENT_BEGIN id=B01 source=colaboracion/parte_iii/source/public_final/base_83/c30_body.tex body_lines=4-32
\label{p3-83-c30-s1:chap:realizacion-analitica-contraangulo}
\index{coordenada angular conjugada \(C^*\)}
\index{microfibra de pi@microfibra de \(\pi\)!carácter analítico}
\index{carácter determinantal}

El catálogo finito y una evaluación analítica responden a preguntas de tipos
distintos. El catálogo determina qué regiones existen, cuáles comparten un
retorno y qué información conserva cada una. Una evaluación analítica debe
decidir, además, cómo se transforma la longitud de una palabra en grado, cómo
se prolonga una rama después de su retorno y con qué carácter escalar se lee
esa prolongación. Este capítulo construye esa segunda operación de manera
explícita.

La construcción parte de cuatro datos ya obtenidos en los capítulos
anteriores: la microfibra de cinco regiones de \(\pi\), la longitud seis de
cada palabra regional, el cierre dodecafásico de \(\alpha\) y el transporte
bicapa asociado a la coordenada común
\(A=10^3\alpha\). Ningún valor de la coordenada angular conjugada \(C_*^{\rm HMT}\), de la componente reducida de
acción, del funcional de Barbero--Immirzi o de una constante metrológica se
emplea en la evaluación. El resultado es un carácter regional convergente,
una recurrencia exacta y una fase orientada cuya carta sexagesimal reproduce
la coordenada angular conjugada publicada.

La separación entre la parte finita y la parte analítica es esencial. El
entero \(169\) será un teorema del catálogo. La serie que lo prolonga será el
teorema de un lector analítico cuyas reglas se declaran antes de evaluarlo.
Así, las premisas adicionales no quedan ocultas dentro de una coincidencia
decimal.

% HMT_SOURCE_SEGMENT_END id=B01
\section{Construcción preangular de \texorpdfstring{\(\eta_{\mathrm{ret}}\)}{eta_ret} mediante \texorpdfstring{\(H_5\)}{H5}, \texorpdfstring{\(D_A\)}{D_A} y \texorpdfstring{\(\mathscr C_\pi(\alpha)\)}{C_pi(alpha)}}
\label{sec:s102:c38:2}

\subsection{Derivación independiente de \texorpdfstring{\(\eta_{\mathrm{ret}}=H_5-(90/\pi)\mathscr C_\pi(\alpha)\)}{eta_ret=H5-(90/pi)C_pi(alpha)} mediante el lector analítico regional}
\label{sec:s102:c38:2:1}

% HMT_SOURCE_SEGMENT_BEGIN id=A01 source=colaboracion/parte_iii/source/public_final/ascensos_20260826/16f_realizacion_analitica_contraangulo.tex body_lines=5-38,40-47

El catálogo finito y una evaluación analítica responden a preguntas de tipos
distintos. El catálogo determina qué regiones existen, cuáles comparten un
retorno y qué información conserva cada una. Una evaluación analítica debe
decidir, además, cómo se transforma la longitud de una palabra en grado, cómo
se prolonga una rama después de su retorno y con qué carácter escalar se lee
esa prolongación. Esta sección construye esa segunda operación de manera
explícita.

La construcción parte de la microfibra de cinco regiones de \(\pi\), la
longitud seis de cada palabra regional, el cierre dodecafásico de \(\alpha\)
y la inserción angular declarada
\[
 \iota_\angle(x):=10^3x\,u_\circ,\qquad
 \Adeg:=10^3\alpha .
\]
Aquí \(u_\circ\) es la base sexagesimal cuya vuelta completa tiene
coordenada \(360\).
El factor \(10^3\) pertenece a esa normalización de carta y no se deduce del
cierre aritmético de \(\alpha\). El transporte bicapa se construye después
a partir de \(\Adeg\). La función que evalúa el lector no recibe como
argumento el contraángulo, la coordenada posterior al retorno, el funcional
hexada--octada ni una constante metrológica. El resultado es un carácter
regional convergente, una recurrencia exacta y una fase orientada en la carta
sexagesimal declarada.

La separación entre la parte finita y la parte analítica es esencial. El
entero \(169\) es un teorema del catálogo. La serie que lo prolonga pertenece
a un lector analítico cuyas reglas se declaran antes de evaluarlo.

La fase se denota por \(y_*\) y, cuando se enfatiza su origen regional, por
\(y_{\rm reg}\); ambos símbolos designan el mismo lector construido a
continuación.


La evaluación recibe el catálogo de \(468\) emisiones, el valor exacto de
\(\alpha_{\rm HMT}\), \(\varphi=(1+\sqrt5)/2\), la fórmula fijada de \(H_5\),
la carta angular y las cinco reglas del lector regional. Con esos datos
calcula \(169\), \(D_A\), la recurrencia y \(C^*\). El teorema es exacto en
esta clase declarada de lectores; no usa \(C^*\) como entrada ni afirma que
las reglas de \(H_5\) sean la única clase natural posible.

% HMT_SOURCE_SEGMENT_END id=A01
\subsection{Persistencia del retorno en las \texorpdfstring{\(5\)}{5} regiones de la microfibra}
\label{sec:s102:c38:2:2}

% HMT_SOURCE_SEGMENT_BEGIN id=B02 source=colaboracion/parte_iii/source/public_final/base_83/c30_body.tex body_lines=34-155

Sea
\[
\mathcal F_\pi=\Psi^{-1}(010211)\subset\mathcal U_6
\]
la microfibra relativa al catálogo. Recordemos sus cinco elementos, ahora
separados en bloque inicial y bloque terminal:
\[
\begin{aligned}
 &(501,614,498\mid169,272,272),\\
 &(810,923,870\mid169,272,272),\\
 &(810,983,810\mid169,272,272),\\
 &(870,923,810\mid169,272,272),\\
 &(870,923,810\mid418,674,521).
\end{aligned}
\]
Definimos la proyección de retorno
\[
p_{\rm ret}:\mathbb Z^6\longrightarrow\mathbb Z^3,
\qquad
p_{\rm ret}(u_1,\ldots,u_6)=(u_4,u_5,u_6).
\]

\begin{proposition}[Bloque terminal persistente]
\label{p3-83-c30-s1:prop:bloque-terminal-persistente}
El multiconjunto \(p_{\rm ret}(\mathcal F_\pi)\) posee un único elemento de
multiplicidad máxima:
\[
b_\pi=(169,272,272),
\qquad
\operatorname{mult}_{\mathcal F_\pi}(b_\pi)=4.
\]
El único bloque terminal restante es
\[
b_\pi^-=(418,674,521)
\]
y tiene multiplicidad uno.
\end{proposition}

\begin{proof}
La afirmación se obtiene por enumeración exacta de las cinco filas del
catálogo que satisfacen \(\Psi(U)=010211\). Cuatro filas poseen bloque
terminal \((169,272,272)\); la quinta, que porta el retorno mutado, posee
\((418,674,521)\). La multiplicidad máxima es cuatro y se alcanza una sola
vez.
\end{proof}

La diferencia orientada entre el retorno mutado y el persistente es
\[
 \omega_\pi=b_\pi^- - b_\pi=(249,402,249),
 \qquad
 \langle(1,1,1),\omega_\pi\rangle=900.
\]
Esta identidad separa dos objetos que no deben confundirse. El entero
\(169\) es una coordenada de un dato terminal, es decir, un dato de grado
cero. El vector \(\omega_\pi\) puede leerse como una cocadena de grado uno
después de orientar la arista que une ambos bloques. Llamarlo cociclo exigiría,
además, fijar un complejo de cocadenas y verificar que su coborde se anula.
La prolongación analítica construida más adelante no presupone esa
identificación.

La selección de \(169\) no exige privilegiar la primera posición del bloque.
El grupo simétrico \(\mathfrak S_3\) actúa permutando sus tres coordenadas.
El estabilizador de \(b_\pi\) intercambia las dos apariciones de \(272\) y
fija la única coordenada de multiplicidad uno. Denotemos por
\(\operatorname{sing}(b)\) esa coordenada cuando el bloque tiene tipo
\((r,s,s)\), con \(r\ne s\). Equivalentemente,
\[
\operatorname{sing}(b)
=\sum_{j=1}^3b_j-2\operatorname{moda}(b).
\]
En particular,
\[
\boxed{\operatorname{sing}(b_\pi)=169.}
\]

\begin{definition}[Selector residual de la microfibra]
\label{p3-83-c30-s1:def:selector-residual-pi}
El selector \(\rho_\pi\) aplica sucesivamente dos operaciones:
\begin{enumerate}
  \item selecciona el bloque terminal de multiplicidad máxima dentro de
  \(p_{\rm ret}(\mathcal F_\pi)\);
  \item retiene la coordenada fija por el estabilizador de ese bloque.
\end{enumerate}
Por la \cref{p3-83-c30-s1:prop:bloque-terminal-persistente},
\[
\boxed{\rho_\pi=169.}
\]
\end{definition}

\begin{theorem}[Unicidad equivariante del retorno]
\label{p3-83-c30-s1:thm:unicidad-retorno-169}
Entre los selectores que
\begin{enumerate}
  \item son invariantes bajo permutaciones de las cinco regiones;
  \item son equivariantes bajo permutaciones de las tres coordenadas
  terminales;
  \item seleccionan el bloque de multiplicidad máxima;
  \item retienen una coordenada fija por su estabilizador,
\end{enumerate}
el valor de retorno está determinado de manera única y es \(169\).
\end{theorem}

\begin{proof}
La invariancia bajo \(\mathfrak S_5\) impide usar el orden de las filas. La
tercera condición selecciona \(b_\pi\), que es el único modo por la
\cref{p3-83-c30-s1:prop:bloque-terminal-persistente}. Su estabilizador en
\(\mathfrak S_3\) posee dos órbitas sobre las coordenadas: la pareja de
valores \(272\) y el valor unipuntual \(169\). La cuarta condición selecciona la
segunda órbita. La equivarianza conserva su valor al desplazar su posición.
\end{proof}

Este argumento mejora la lectura posicional del retorno. El \(169\) no es
«la primera coordenada» de una fila escogida: es un invariante del
multiconjunto regional y de la acción de sus simetrías de reordenación.
Tampoco se lo selecciona por rareza global. En las \(468\) emisiones del
catálogo, \(169\) aparece \(84\) veces, uniformemente distribuido con
\(14\) apariciones en cada una de las seis posiciones. La unicidad del
\cref{p3-83-c30-s1:thm:unicidad-retorno-169} es, por tanto, relacional y regional: reside
en la combinación de microfibra, persistencia modal y estabilizador, no en la
frecuencia aislada del entero.

% HMT_SOURCE_SEGMENT_END id=B02
\subsection{Construcción del lector analítico regional y de su dominio de convergencia}
\label{sec:s102:c38:2:3}

% HMT_SOURCE_SEGMENT_BEGIN id=B05 source=colaboracion/parte_iii/source/public_final/base_83/c30_body.tex body_lines=261-301

La extensión analítica se fija mediante las reglas siguientes.

\begin{definition}[Lector regional determinantal]
\label{p3-83-c30-s1:def:lector-regional-determinantal}
El lector regional determinantal satisface:
\begin{enumerate}
  \item \emph{compatibilidad de grado}: la longitud cilíndrica es el grado
  analítico;
  \item \emph{descomposición de renovación}: el impulso finito de retorno y
  la continuación estricta pertenecen a sumandos aditivos distintos;
  \item \emph{normalización de continuación}: después de registrar una sola
  vez el residuo \(\rho_\pi\), la rama superviviente comienza en la unidad de
  la línea determinante;
  \item \emph{covariancia por concatenación}: cada prolongación elemental
  actúa mediante \(\bigwedge^2\mathcal T=D_A\);
  \item \emph{orientación}: en la convención cardinal fijada para APP, el
  retorno regional pertenece al sector compensador y se sustrae de la fase
  de acción de quinto grado.
\end{enumerate}
\end{definition}

La tercera regla tiene contenido matemático. El valor \(169\) se registra
una vez como amplitud del retorno; no se multiplica de nuevo en cada
prolongación. La continuación estricta cuenta la única línea que prosigue y
la normaliza por su unidad. Esta es la diferencia entre una recurrencia de
renovación y la iteración homogénea del impulso inicial.

La necesidad de declarar esa normalización puede probarse. Para todo
\(\lambda\in\mathbb R\), la familia
\[
F_\lambda(z)
=169z^6+\lambda\frac{D_Az^7}{1-D_Az}
\]
conserva el mismo retorno finito y la misma razón determinantal entre
coeficientes consecutivos. El catálogo y la razón por sí solos no distinguen
\(\lambda\). La unidad de la línea determinante fija
\(\lambda=1\) antes de evaluar \(z=\alpha\). De este modo, la normalización
no se obtiene ajustando el valor final; forma parte de la definición del
carácter.

% HMT_SOURCE_SEGMENT_END id=B05
\subsection{Resolución analítica de la recurrencia de retorno}
\label{sec:s102:c38:2:4}

% HMT_SOURCE_SEGMENT_BEGIN id=B06 source=colaboracion/parte_iii/source/public_final/base_83/c30_body.tex body_lines=303-364

Definimos los coeficientes \(\kappa_n\) por
\[
\kappa_6=169,
\qquad
\kappa_7=D_A,
\qquad
\kappa_{n+1}=D_A\kappa_n\quad(n\ge7).
\]
La primera igualdad es el impulso regional; las dos siguientes describen la
continuación normalizada.

\begin{theorem}[Realización analítica graduada]
\label{p3-83-c30-s1:thm:realizacion-analitica-graduada}
El germen analítico compatible con las reglas de la
\cref{p3-83-c30-s1:def:lector-regional-determinantal} es único y vale
\[
\boxed{
\mathscr C_\pi(z)
=169z^6+\frac{D_Az^7}{1-D_Az}.}
\]
Sus coeficientes satisfacen
\[
\kappa_n=D_A^{n-6}\qquad(n\ge7).
\]
\end{theorem}

\begin{proof}
Escribamos
\[
F(z)=169z^6+\sum_{k\ge1}c_kz^{6+k}.
\]
La normalización de continuación y la covariancia determinantal dan
\[
c_1=D_A,
\qquad
c_{k+1}=D_Ac_k.
\]
Por inducción, \(c_k=D_A^k\). La suma geométrica produce
\[
\sum_{k\ge1}D_A^kz^{6+k}
=\frac{D_Az^7}{1-D_Az}.
\]
No queda ningún coeficiente libre.
\end{proof}

Como \(0<\alpha<1\) y \(0<D_A<1\),
\[
0<D_A\alpha<1.
\]
En consecuencia, \(\mathscr C_\pi(\alpha)\) converge absolutamente y su
radio de convergencia es \(D_A^{-1}>1\). Numéricamente,
\[
D_A=0.7751291192471564301888840964802\ldots,
\]
\[
D_A\alpha=0.00565639046986492673885079095469\ldots.
\]
La rapidez de esta contracción explica por qué la suma completa y su
truncamiento de grado siete poseen las mismas quince primeras cifras en la
carta final.

% HMT_SOURCE_SEGMENT_END id=B06
\subsection{Persistencia del retorno en las \texorpdfstring{\(5\)}{5} regiones y compatibilidad de la componente \texorpdfstring{\(\eta_{\mathrm{ret}}\)}{eta_ret}}
\label{sec:s102:c38:2:5}

% HMT_SOURCE_SEGMENT_BEGIN id=A02 source=colaboracion/parte_iii/source/public_final/ascensos_20260826/16f_realizacion_analitica_contraangulo.tex body_lines=49-170

Sea
\[
\mathcal F_\pi=\Psi^{-1}(010211)\subset\mathcal U_6
\]
la microfibra relativa al catálogo. Recordemos sus cinco elementos, ahora
separados en bloque inicial y bloque terminal:
\[
\begin{aligned}
 &(501,614,498\mid169,272,272),\\
 &(810,923,870\mid169,272,272),\\
 &(810,983,810\mid169,272,272),\\
 &(870,923,810\mid169,272,272),\\
 &(870,923,810\mid418,674,521).
\end{aligned}
\]
Definimos la proyección de retorno
\[
p_{\rm ret}:\mathbb Z^6\longrightarrow\mathbb Z^3,
\qquad
p_{\rm ret}(u_1,\ldots,u_6)=(u_4,u_5,u_6).
\]

\begin{proposition}[Bloque terminal persistente]
\label{p3-20260826:prop:bloque-terminal-persistente}
El multiconjunto \(p_{\rm ret}(\mathcal F_\pi)\) posee un único elemento de
multiplicidad máxima:
\[
b_\pi=(169,272,272),
\qquad
\operatorname{mult}_{\mathcal F_\pi}(b_\pi)=4.
\]
El único bloque terminal restante es
\[
b_\pi^-=(418,674,521)
\]
y tiene multiplicidad uno.
\end{proposition}

\begin{proof}
La afirmación se obtiene por enumeración exacta de las cinco filas del
catálogo que satisfacen \(\Psi(U)=010211\). Cuatro filas poseen bloque
terminal \((169,272,272)\); la quinta, que porta el retorno mutado, posee
\((418,674,521)\). La multiplicidad máxima es cuatro y se alcanza una sola
vez.
\end{proof}

La diferencia orientada entre el retorno mutado y el persistente es
\[
 \omega_\pi=b_\pi^- - b_\pi=(249,402,249),
 \qquad
 \langle(1,1,1),\omega_\pi\rangle=900.
\]
Esta identidad separa dos objetos que no deben confundirse. El entero
\(169\) es una coordenada de un dato terminal, es decir, un dato de grado
cero. El vector \(\omega_\pi\) puede leerse como una cocadena de grado uno
después de orientar la arista que une ambos bloques. Llamarlo cociclo exigiría,
además, fijar un complejo de cocadenas y verificar que su coborde se anula.
La prolongación analítica construida más adelante no presupone esa
identificación.

La selección de \(169\) no exige privilegiar la primera posición del bloque.
El grupo simétrico \(\mathfrak S_3\) actúa permutando sus tres coordenadas.
El estabilizador de \(b_\pi\) intercambia las dos apariciones de \(272\) y
fija la única coordenada de multiplicidad uno. Denotemos por
\(\operatorname{sing}(b)\) esa coordenada cuando el bloque tiene tipo
\((r,s,s)\), con \(r\ne s\). Equivalentemente,
\[
\operatorname{sing}(b)
=\sum_{j=1}^3b_j-2\operatorname{moda}(b).
\]
En particular,
\[
\boxed{\operatorname{sing}(b_\pi)=169.}
\]

\begin{definition}[Selector residual de la microfibra]
\label{p3-20260826:def:selector-residual-pi}
El selector \(\rho_\pi\) aplica sucesivamente dos operaciones:
\begin{enumerate}
  \item selecciona el bloque terminal de multiplicidad máxima dentro de
  \(p_{\rm ret}(\mathcal F_\pi)\);
  \item retiene la coordenada fija por el estabilizador de ese bloque.
\end{enumerate}
Por la \cref{p3-20260826:prop:bloque-terminal-persistente},
\[
\boxed{\rho_\pi=169.}
\]
\end{definition}

\begin{theorem}[Unicidad equivariante del retorno]
\label{p3-20260826:thm:unicidad-retorno-169}
Entre los selectores que
\begin{enumerate}
  \item son invariantes bajo permutaciones de las cinco regiones;
  \item son equivariantes bajo permutaciones de las tres coordenadas
  terminales;
  \item seleccionan el bloque de multiplicidad máxima;
  \item retienen una coordenada fija por su estabilizador,
\end{enumerate}
el valor de retorno está determinado de manera única y es \(169\).
\end{theorem}

\begin{proof}
La invariancia bajo \(\mathfrak S_5\) impide usar el orden de las filas. La
tercera condición selecciona \(b_\pi\), que es el único modo por la
\cref{p3-20260826:prop:bloque-terminal-persistente}. Su estabilizador en
\(\mathfrak S_3\) posee dos órbitas sobre las coordenadas: la pareja de
valores \(272\) y el valor unipuntual \(169\). La cuarta condición selecciona la
segunda órbita. La equivarianza conserva su valor al desplazar su posición.
\end{proof}

Este argumento mejora la lectura posicional del retorno. El \(169\) no es
«la primera coordenada» de una fila escogida: es un invariante del
multiconjunto regional y de la acción de sus simetrías de reordenación.
Tampoco se lo selecciona por rareza global. En las \(468\) emisiones del
catálogo, \(169\) aparece \(84\) veces, uniformemente distribuido con
\(14\) apariciones en cada una de las seis posiciones. La unicidad del
\cref{p3-20260826:thm:unicidad-retorno-169} es, por tanto, relacional y regional: reside
en la combinación de microfibra, persistencia modal y estabilizador, no en la
frecuencia aislada del entero.

% HMT_SOURCE_SEGMENT_END id=A02
\subsection{Aplicación del lector regional al par angular conjugado}
\label{sec:s102:c38:2:6}

% HMT_SOURCE_SEGMENT_BEGIN id=A05 source=colaboracion/parte_iii/source/public_final/ascensos_20260826/16f_realizacion_analitica_contraangulo.tex body_lines=277-317

La extensión analítica se fija mediante las reglas siguientes.

\begin{definition}[Lector regional determinantal]
\label{p3-20260826:def:lector-regional-determinantal}
El lector regional determinantal satisface:
\begin{enumerate}
  \item \emph{compatibilidad de grado}: la longitud cilíndrica es el grado
  analítico;
  \item \emph{descomposición de renovación}: el impulso finito de retorno y
  la continuación estricta pertenecen a sumandos aditivos distintos;
  \item \emph{normalización de continuación}: después de registrar una sola
  vez el residuo \(\rho_\pi\), la rama superviviente comienza en la unidad de
  la línea determinante;
  \item \emph{covariancia por concatenación}: cada prolongación elemental
  actúa mediante \(\bigwedge^2\mathcal T=D_A\);
  \item \emph{orientación}: en la convención cardinal fijada para APP, el
  retorno regional pertenece al sector compensador y se sustrae de la fase
  de acción de quinto grado.
\end{enumerate}
\end{definition}

La tercera regla tiene contenido matemático. El valor \(169\) se registra
una vez como amplitud del retorno; no se multiplica de nuevo en cada
prolongación. La continuación estricta cuenta la única línea que prosigue y
la normaliza por su unidad. Esta es la diferencia entre una recurrencia de
renovación y la iteración homogénea del impulso inicial.

La necesidad de declarar esa normalización puede probarse. Para todo
\(\lambda\in\mathbb R\), la familia
\[
F_\lambda(z)
=169z^6+\lambda\frac{D_Az^7}{1-D_Az}
\]
conserva el mismo retorno finito y la misma razón determinantal entre
coeficientes consecutivos. El catálogo y la razón por sí solos no distinguen
\(\lambda\). La unidad de la línea determinante fija
\(\lambda=1\) antes de evaluar \(z=\alpha\). De este modo, la normalización
no se obtiene ajustando el valor final; forma parte de la definición del
carácter.

% HMT_SOURCE_SEGMENT_END id=A05
\subsection{Suma cerrada y publicación de la componente posterior al retorno}
\label{sec:s102:c38:2:7}

% HMT_SOURCE_SEGMENT_BEGIN id=A06 source=colaboracion/parte_iii/source/public_final/ascensos_20260826/16f_realizacion_analitica_contraangulo.tex body_lines=319-380

Definimos los coeficientes \(\kappa_n\) por
\[
\kappa_6=169,
\qquad
\kappa_7=D_A,
\qquad
\kappa_{n+1}=D_A\kappa_n\quad(n\ge7).
\]
La primera igualdad es el impulso regional; las dos siguientes describen la
continuación normalizada.

\begin{theorem}[Realización analítica graduada]
\label{p3-20260826:thm:realizacion-analitica-graduada}
El germen analítico compatible con las reglas de la
\cref{p3-20260826:def:lector-regional-determinantal} es único y vale
\[
\boxed{
\mathscr C_\pi(z)
=169z^6+\frac{D_Az^7}{1-D_Az}.}
\]
Sus coeficientes satisfacen
\[
\kappa_n=D_A^{n-6}\qquad(n\ge7).
\]
\end{theorem}

\begin{proof}
Escribamos
\[
F(z)=169z^6+\sum_{k\ge1}c_kz^{6+k}.
\]
La normalización de continuación y la covariancia determinantal dan
\[
c_1=D_A,
\qquad
c_{k+1}=D_Ac_k.
\]
Por inducción, \(c_k=D_A^k\). La suma geométrica produce
\[
\sum_{k\ge1}D_A^kz^{6+k}
=\frac{D_Az^7}{1-D_Az}.
\]
No queda ningún coeficiente libre.
\end{proof}

Como \(0<\alpha<1\) y \(0<D_A<1\),
\[
0<D_A\alpha<1.
\]
En consecuencia, \(\mathscr C_\pi(\alpha)\) converge absolutamente y su
radio de convergencia es \(D_A^{-1}>1\). Numéricamente,
\[
D_A=0.7751291192471564301888840964802\ldots,
\]
\[
D_A\alpha=0.00565639046986492673885079095469\ldots.
\]
La rapidez de esta contracción explica por qué la suma completa y su
truncamiento de grado siete poseen las mismas quince primeras cifras en la
carta final.

% HMT_SOURCE_SEGMENT_END id=A06
\section{Confluencia de \texorpdfstring{\(\eta_{\mathrm{ret}}\)}{eta_ret} con la valoración \texorpdfstring{\(\varepsilon_H=-34\)}{epsilon_H=-34} en la magnitud \texorpdfstring{\(\hbar_{\mathrm{ret}}\)}{hbar_ret}}
\label{sec:s102:c38:3}

\subsection{Derivación del carácter determinantal del transporte parangular}
\label{sec:s102:c38:3:1}

% HMT_SOURCE_SEGMENT_BEGIN id=B04 source=colaboracion/parte_iii/source/public_final/base_83/c30_body.tex body_lines=196-259

El cierre dodecafásico proporciona \(\alpha\). La completación de las tres
parejas nonádicas proporciona la escala \(10^3\), de modo que
\[
A=10^3\alpha.
\]
Elegida la carta angular arquimediana, la coordenada común en radianes es
\[
x=\frac{\pi A}{180}=\frac{50\pi}{9}\alpha.
\]
Sea \(R\) una involución real de traza nula en dimensión dos,
\[
R^2=I_2,
\qquad
\operatorname{tr}R=0,
\]
y consideremos el transporte formal bicapa
\[
\mathcal T(x,y)=\exp(-xI_2-yR).
\]
La variable orientada \(y\) todavía no se ha evaluado. Sin embargo, la acción
de \(\mathcal T\) sobre la línea determinante es independiente de ella:
\[
\begin{aligned}
\det\mathcal T(x,y)
&=\exp\!\left(\operatorname{tr}(-xI_2-yR)\right)\\
&=e^{-2x}.
\end{aligned}
\]
Definimos
\[
\boxed{
D_A=e^{-2x}=e^{-100\pi\alpha/9}.}
\]
El subíndice recuerda que se trata del determinante de la contracción común
asociada a \(A\), no de una función de \(C_*^{\rm HMT}\).

\begin{proposition}[Carácter de la línea determinante]
\label{p3-83-c30-s1:prop:caracter-determinante}
La aplicación
\[
\delta_A:(\mathbb N_0,+)\longrightarrow(\mathbb R_{>0},\cdot),
\qquad
\delta_A(k)=D_A^k,
\]
es el único carácter multiplicativo cuyo valor sobre una prolongación
elemental es \(D_A\).
\end{proposition}

\begin{proof}
La multiplicatividad da
\[
\delta_A(k)
=\delta_A(\underbrace{1+\cdots+1}_{k\text{ veces}})
=\delta_A(1)^k=D_A^k.
\]
La misma identidad prueba existencia y unicidad.
\end{proof}

El determinante permanece invariante bajo conjugación de
\(\mathcal T\) y bajo el intercambio de hojas \(R\mapsto-R\). Por tanto, la
cola construida a partir de \(D_A\) pertenece al modo común. La orientación
aparecerá en el signo con el que el carácter regional corrige la fase.

% HMT_SOURCE_SEGMENT_END id=B04
\subsection{Carácter de acción de grado \texorpdfstring{\(5\)}{5} en la construcción regional}
\label{sec:s102:c38:3:2}

% HMT_SOURCE_SEGMENT_BEGIN id=B07 source=colaboracion/parte_iii/source/public_final/base_83/c30_body.tex body_lines=366-405

El bloque central
\[
\mathcal B_c=\begin{pmatrix}7&2\\2&7\end{pmatrix}
\]
posee autovalores \(9\) y \(5\). El ciclo dodecafásico aporta el rango
\(12\); sus órbitas de paso tres poseen longitud cuatro; y el censo completo
aporta el cociente exacto
\[
\frac{104976}{1944}=54.
\]
A partir de estos invariantes se fija el carácter de acción de quinto grado
\[
\boxed{
H_5(\alpha,\varphi)
=\frac12\sqrt{\frac{1000\alpha}{\varphi}}-\alpha
+\frac9{16}\alpha^2-\frac59\alpha^3
+\frac7{48}\alpha^4-\frac1{54}\alpha^5.}
\]
Los cuatro coeficientes racionales pueden escribirse como
\[
\frac9{16}=\frac{\lambda_+}{4^2},
\qquad
\frac59=\frac{\lambda_-}{\lambda_+},
\qquad
\frac7{48}=\frac{\tfrac12\operatorname{tr}\mathcal B_c}{4\cdot12},
\qquad
\frac1{54}=\frac{1944}{104976},
\]
con \((\lambda_+,\lambda_-)=(9,5)\). Estas identidades certifican la
procedencia de los números utilizados. La asignación sucesiva de esos
invariantes a los grados dos, tres, cuatro y cinco forma parte del carácter
analítico; no se deduce de la mera existencia del catálogo finito.

Esta precisión elimina una ambigüedad habitual. Una combinación de
invariantes de la construcción es una definición interna y reproducible; su
naturalidad requiere las reglas del lector que especifican orden, signos y
normalización. Aquí esas reglas son visibles y quedan sometidas a los mismos
controles negativos que el resto de la construcción.

% HMT_SOURCE_SEGMENT_END id=B07
\subsection{Descomposición exacta de la recta de acción en secciones pre-retorno y post-retorno}
\label{sec:s102:c38:3:3}

% HMT_SOURCE_SEGMENT_BEGIN id=B12 source=colaboracion/parte_iii/source/public_final/base_83/c30_body.tex body_lines=571-617

El carácter de quinto grado y la componente de acción posterior al retorno
son dos valores distintos:
\[
H_5
=1.05457181764615446962582182943306\ldots,
\]
\[
\bar h_{\rm HMT}^{\rm ret}
=1.05457181691503906113369058472430\ldots.
\]
El primero es la acción local antes de incorporar el carácter regional; el
segundo es la acción orientada que queda después del retorno. La corrección
que genera la coordenada angular conjugada impide identificarlos.

La comparación metrológica debe respetar esta distinción. El primer valor da
la mantisa
\[
2\pi H_5
=6.62607014999998792600111034989947\ldots,
\]
mientras que la acción posterior da
\[
2\pi\bar h_{\rm HMT}^{\rm ret}
=6.62607014540625433351074984759288\ldots.
\]
En el Sistema Internacional, el valor
\[
h=6.62607015\times10^{-34}\;\mathrm{J\,s}
\]
es exacto por definición \cite{BIPMSI}. Por ello,
\[
2\pi H_5-6.62607015
=-1.2073998889650101\ldots\times10^{-14},
\]
y no es una igualdad exacta. La proximidad del valor anterior al retorno es
un control numérico notable; no autoriza a sustituir la constante definitoria
del SI por la componente posterior ni a confundir las dos acciones HMT\@.

Este resultado determina con precisión la situación matemática. El lector
regional genera \(C_*^{\rm HMT}\) y, desde él, el funcional
\(\Gamma_{6\to8}\). El mismo cálculo separa la aproximación de Planck
producida por \(H_5\) de la acción reducida que interviene en el
la coordenada angular conjugada. La separación no debilita la construcción: elimina una
identificación incompatible con las cifras exactas y convierte cada
comparación en una afirmación falsable de tipo bien definido.

% HMT_SOURCE_SEGMENT_END id=B12
\subsection{Factorización del carácter determinantal en las \texorpdfstring{\(2\)}{2} cartas de fase}
\label{sec:s102:c38:3:4}

% HMT_SOURCE_SEGMENT_BEGIN id=A04 source=colaboracion/parte_iii/source/public_final/ascensos_20260826/16f_realizacion_analitica_contraangulo.tex body_lines=211-275

El cierre dodecafásico proporciona \(\alpha\). La inserción angular declarada
emplea la escala \(10^3\), compatible con la agrupación de tres parejas
nonádicas, de modo que
\[
\Adeg=10^3\alpha\,{}^\circ.
\]
Elegida la carta angular arquimediana, la coordenada común en radianes es
\[
x=: \Arad=\frac{\pi\Adeg}{180}=\frac{50\pi}{9}\alpha.
\]
Sea \(R\) una involución real de traza nula en dimensión dos,
\[
R^2=I_2,
\qquad
\operatorname{tr}R=0,
\]
y consideremos el transporte formal bicapa
\[
\mathcal T(x,y)=\exp(-xI_2-yR).
\]
La variable orientada \(y\) todavía no se ha evaluado. Sin embargo, la acción
de \(\mathcal T\) sobre la línea determinante es independiente de ella:
\[
\begin{aligned}
\det\mathcal T(x,y)
&=\exp\!\left(\operatorname{tr}(-xI_2-yR)\right)\\
&=e^{-2x}.
\end{aligned}
\]
Definimos
\[
\boxed{
D_A=e^{-2x}=e^{-100\pi\alpha/9}.}
\]
El subíndice recuerda que se trata del determinante de la contracción común
asociada a \(A\), no de una función del contraángulo.

\begin{proposition}[Carácter de la línea determinante]
\label{p3-20260826:prop:caracter-determinante}
La aplicación
\[
\delta_A:(\mathbb N_0,+)\longrightarrow(\mathbb R_{>0},\cdot),
\qquad
\delta_A(k)=D_A^k,
\]
es el único carácter multiplicativo cuyo valor sobre una prolongación
elemental es \(D_A\).
\end{proposition}

\begin{proof}
La multiplicatividad da
\[
\delta_A(k)
=\delta_A(\underbrace{1+\cdots+1}_{k\text{ veces}})
=\delta_A(1)^k=D_A^k.
\]
La misma identidad prueba existencia y unicidad.
\end{proof}

El determinante permanece invariante bajo conjugación de
\(\mathcal T\) y bajo el intercambio de hojas \(R\mapsto-R\). Por tanto, la
cola construida a partir de \(D_A\) pertenece al modo común. La orientación
aparecerá en el signo con el que el carácter regional corrige la fase.

% HMT_SOURCE_SEGMENT_END id=A04
\subsection{Publicación del carácter de acción de grado \texorpdfstring{\(5\)}{5} en la carta parangular}
\label{sec:s102:c38:3:5}

% HMT_SOURCE_SEGMENT_BEGIN id=A07 source=colaboracion/parte_iii/source/public_final/ascensos_20260826/16f_realizacion_analitica_contraangulo.tex body_lines=382-418

El bloque central
\[
\mathcal B_c=\begin{pmatrix}7&2\\2&7\end{pmatrix}
\]
posee autovalores \(9\) y \(5\). El ciclo dodecafásico aporta el rango
\(12\); sus órbitas de paso tres poseen longitud cuatro; y el censo completo
aporta el cociente exacto
\[
\frac{104976}{1944}=54.
\]
El lector vigente declara el carácter de acción de quinto grado
\[
\boxed{
H_5(\alpha,\varphi)
=\frac12\sqrt{\frac{1000\alpha}{\varphi}}-\alpha
+\frac9{16}\alpha^2-\frac59\alpha^3
+\frac7{48}\alpha^4-\frac1{54}\alpha^5.}
\]
Los cuatro coeficientes racionales pueden escribirse como
\[
\frac9{16}=\frac{\lambda_+}{4^2},
\qquad
\frac59=\frac{\lambda_-}{\lambda_+},
\qquad
\frac7{48}=\frac{\tfrac12\operatorname{tr}\mathcal B_c}{4\cdot12},
\qquad
\frac1{54}=\frac{1944}{104976},
\]
con \((\lambda_+,\lambda_-)=(9,5)\). Estas identidades fijan la
representabilidad aritmética interna de los coeficientes del carácter. Su
orden, sus signos y su normalización pertenecen a la regla graduada del lector
regional y se aplican antes de toda carta de unidades. Las ablaciones
certifican que cada término interviene efectivamente en la fase resultante.
En particular, ni la mantisa SI de la acción ni el valor sexagesimal del
contraángulo forman parte del dominio de la función generadora.

% HMT_SOURCE_SEGMENT_END id=A07
\section{Aplicación parangular \texorpdfstring{\(P(\alpha,\eta_{\mathrm{ret}})\)}{P(alpha,eta_ret)} y clausura bicapa \texorpdfstring{\(C^*=2(\eta_{\mathrm{ret}}+\alpha)\)}{C*=2(eta_ret+alpha)}}
\label{sec:s102:c38:4}

\subsection{Fase orientada y coordenada angular conjugada}
\label{sec:s102:c38:4:1}

% HMT_SOURCE_SEGMENT_BEGIN id=B08 source=colaboracion/parte_iii/source/public_final/base_83/c30_body.tex body_lines=407-472

La fase de quinto grado anterior al retorno es
\[
y_5=\frac\pi{90}\bigl(H_5+\alpha\bigr).
\]
El carácter regional completo produce la corrección
\(\mathscr C_\pi(\alpha)\). Con la orientación fijada en la
\cref{p3-83-c30-s1:def:lector-regional-determinantal}, definimos
\[
\boxed{
y_*=\frac\pi{90}\bigl(H_5+\alpha\bigr)
-169\alpha^6
-\frac{D_A\alpha^7}{1-D_A\alpha}.}
\]
Ésta es la fórmula primaria. La unidad sexagesimal no interviene hasta aplicar la carta
\[
C_*^{\rm HMT}=\frac{180}{\pi}y_*.
\]

\begin{theorem}[Coordenada angular conjugada del lector regional determinantal]
\label{p3-83-c30-s1:thm:contraangulo-regional}
Fijados el cierre dodecafásico de \(\alpha\),
\(\varphi=(1+\sqrt5)/2\) y las reglas del lector regional determinantal, la
fase \(y_*\) y su coordenada angular conjugada quedan determinados sin introducir el valor
buscado. Su evaluación es
\[
\boxed{
y_*=0.03706622646583826398493779409230638\ldots,}
\]
\[
\boxed{
C_*^{\rm HMT}
=2.12373833896864572426195138039336\ldots{}^\circ.}
\]
Por tanto, a quince cifras decimales,
\[
\boxed{C_*^{\rm HMT}=2.123738338968646^\circ.}
\]
\end{theorem}

\begin{proof}
El \cref{p3-83-c30-s1:thm:unicidad-retorno-169} determina \(169\). El cierre de
\(\alpha\) determina \(A\), \(x\) y
\(D_A=e^{-100\pi\alpha/9}\). El \cref{p3-83-c30-s1:thm:realizacion-analitica-graduada}
determina la serie completa. La fórmula de \(H_5\) contiene únicamente
\(\alpha\), \(\varphi\) e invariantes estructurales declarados. Sustituir
estos datos en la expresión de \(y_*\) da el valor impreso. El valor decimal
de \(C_*^{\rm HMT}\) no aparece en ninguno de los miembros derechos.
\end{proof}

La componente reducida de acción posterior al retorno queda ahora como
salida:
\[
\boxed{
\bar h_{\rm HMT}^{\rm ret}
=\frac{C_*^{\rm HMT}}2-\alpha
=H_5-\frac{90}{\pi}\mathscr C_\pi(\alpha).}
\]
Numéricamente,
\[
\bar h_{\rm HMT}^{\rm ret}
=1.05457181691503906113369058472430\ldots.
\]
La identidad inversa no se ha usado para introducir \(C_*^{\rm HMT}\); se
obtiene después de generar \(y_*\).

% HMT_SOURCE_SEGMENT_END id=B08
\subsection{Coordenada conjugada \texorpdfstring{\(C^*\)}{C*} y reversión de su orientación bajo la involución bicapa}
\label{sec:s102:c38:4:2}

% HMT_SOURCE_SEGMENT_BEGIN id=A08 source=colaboracion/parte_iii/source/public_final/ascensos_20260826/16f_realizacion_analitica_contraangulo.tex body_lines=420-503

La fase de quinto grado anterior al retorno es
\[
y_5=\frac\pi{90}\bigl(H_5+\alpha\bigr).
\]
El carácter regional completo produce la corrección
\(\mathscr C_\pi(\alpha)\). Con la orientación fijada en la
\cref{p3-20260826:def:lector-regional-determinantal}, definimos
\[
\boxed{
\Crad:=y_*=\frac\pi{90}\bigl(H_5+\alpha\bigr)
-169\alpha^6
-\frac{D_A\alpha^7}{1-D_A\alpha}.}
\]
En adelante escribimos también
\[
 \boxed{y_{\rm reg}:=y_*}
\]
cuando sea necesario distinguir esta realización de la fase espectral de
\Gtr. Las dos notaciones designan aquí el mismo objeto regional canónico.
Ésta es la fórmula primaria. La unidad sexagesimal no interviene hasta aplicar la carta
\[
\boxed{\Cdeg:=\frac{180}{\pi}\Crad.}
\]
Los símbolos \(\Crad\) y \(\Cdeg\) no designan dos contraángulos: son las
coordenadas radial y sexagesimal del mismo elemento angular.

\begin{theorem}[Contraángulo del lector regional determinantal]
\label{p3-20260826:thm:contraangulo-regional}
Fijados el cierre dodecafásico de \(\alpha\),
\(\varphi=(1+\sqrt5)/2\), el germen \(H_5\) y las reglas del lector regional
determinantal, la fase \(\Crad\) y su contraángulo quedan determinados sin
introducir el valor buscado en la función de evaluación. Su evaluación es
\[
\boxed{
\Crad=0.03706622646583826398493779409230638\ldots,}
\]
\[
\boxed{
\Cdeg
=2.12373833896864572426195138039336\ldots{}^\circ.}
\]
Por tanto, a quince cifras decimales,
\[
\boxed{\Cdeg=2.123738338968646^\circ.}
\]
\end{theorem}

\begin{proof}
El \cref{p3-20260826:thm:unicidad-retorno-169} determina \(169\). El cierre de
\(\alpha\) determina \(\Adeg\), \(\Arad\) y
\(D_A=e^{-100\pi\alpha/9}\). El \cref{p3-20260826:thm:realizacion-analitica-graduada}
determina la serie completa. La fórmula fijada de \(H_5\) contiene únicamente
\(\alpha\), \(\varphi\) e invariantes estructurales declarados. Sustituir
estos datos en la expresión de \(\Crad\) da el valor impreso. El valor decimal
del contraángulo no aparece en ninguno de los miembros derechos de esta
evaluación. La comparación con una mantisa metrológica de acción se efectúa
únicamente después de obtener la fase y no modifica la función generadora.
\end{proof}

Sea \(u_\circ\) la base sexagesimal cuya vuelta completa tiene coordenada
\(360\), e introduzcamos la coordenada angular
\(\alpha_{\angle,{\rm deg}}:=\alpha_{\rm HMT}\). La salida posterior al
retorno queda entonces tipada por
\[
\boxed{
\etaRet
=\frac{\Cdeg}2-\alpha_{\angle,{\rm deg}}
=H_5-\frac{90}{\pi}\mathscr C_\pi(\alpha).}
\]
Su coordenada radial es
\[
 \etaRetrad=\frac{\pi}{180}\etaRet
 =\frac{\Crad}{2}-\frac{\pi}{180}\alpha_{\rm HMT}.
\]
Numéricamente,
\[
\etaRet
=1.05457181691503906113369058472430\ldots.
\]
La identidad inversa no se ha usado para introducir \(\Cdeg\); se obtiene
después de generar \(\Crad\). El símbolo \(\hbar\) queda reservado para el
elemento de la recta de acción construido en la sección siguiente.

% HMT_SOURCE_SEGMENT_END id=A08
\subsection{Teorema directo \texorpdfstring{\(C^*=2(\eta_{\mathrm{ret}}+\alpha)\)}{C*=2(eta_ret+alpha)} y carácter inverso de la identidad \texorpdfstring{\(\eta_{\mathrm{ret}}=C^*/2-\alpha\)}{eta_ret=C*/2-alpha}}
\label{sec:s102:c38:4:3}

% HMT_SOURCE_SEGMENT_BEGIN id=A13 source=colaboracion/parte_iii/source/public_final/ascensos_20260826/16f_realizacion_analitica_contraangulo.tex body_lines=750-797

El carácter de quinto grado y la coordenada posterior al retorno
son dos valores distintos:
\[
H_5
=1.05457181764615446962582182943306\ldots,
\]
\[
\etaRet
=1.05457181691503906113369058472430\ldots.
\]
El primero es el carácter local antes de incorporar el carácter regional; el
segundo es la coordenada orientada que queda después del retorno. La corrección
que genera el contraángulo impide identificarlos.

La comparación metrológica debe respetar esta distinción. El primer valor da
la mantisa
\[
2\pi H_5
=6.62607014999998792600111034989947\ldots,
\]
mientras que la coordenada posterior da
\[
2\pi\etaRet
=6.62607014540625433351074984759288\ldots.
\]
Como control externo únicamente, retenemos la mantisa \(6.62607015\) del
valor definitorio del Sistema Internacional \cite{BIPMSI}. No se incorpora
todavía su exponente ni su unidad a la recta HMT: la década se deriva en el
\cref{chap:orden-determinantal-accion} y la unidad pertenece a una carta
metrológica posterior. La comparación de mantisas da
\[
2\pi H_5-6.62607015
=-1.2073998889650101\ldots\times10^{-14},
\]
y no es una igualdad exacta. La proximidad del valor anterior al retorno es
un control numérico notable; no autoriza a sustituir la constante definitoria
del SI por la coordenada posterior ni a confundir una coordenada angular con
el elemento de acción que se construirá al incorporar la década.

Este resultado determina con precisión la situación matemática. El lector
regional genera \((\Crad,\Cdeg)\) y, desde ese único elemento angular, el funcional
\(\Gamma_{6\to8}\). El mismo cálculo separa la proximidad externa de la
mantisa producida por \(H_5\) de la coordenada de retorno que interviene en el
contraángulo. La separación no debilita la construcción: elimina una
identificación incompatible con las cifras exactas y convierte cada
comparación en una afirmación falsable de tipo bien definido.

% HMT_SOURCE_SEGMENT_END id=A13
% HMT_SOURCE_SEGMENT_BEGIN id=D01 source=manuscrito/sucesor_102/deltas_ley9/c38_par_angular_cartas_y_canales.tex body_lines=3-20
\label{sec:l9s102:par-angular}

La completación cubica
\begin{equation}
 10^3=(9+1)^3=9^3+3\cdot9^2+3\cdot9+1
 =729+243+27+1
 \label{eq:l9s102:base-1000}
\end{equation}
determina la carta global de la coordenada directa. La clausura bicapa aporta
el factor \(2\) de la coordenada conjugada:
\begin{equation}
 \boxed{
 A_{\deg}=10^3\alpha_{\mathrm{HMT}},
 \qquad
 C^*_{\deg}=2(\eta_{\mathrm{ret}}+\alpha_{\mathrm{HMT}})
 \quad\text{en }\mathbb R/(360\mathbb Z).}
 \label{eq:l9s102:a-cstar-deg}
\end{equation}
% HMT_SOURCE_SEGMENT_END id=D01
% HMT_SOURCE_SEGMENT_BEGIN id=D04 source=manuscrito/sucesor_102/deltas_ley9/c38_par_angular_cartas_y_canales.tex body_lines=60-66
La identidad inversa
\begin{equation}
 \eta_{\mathrm{ret}}=\frac12C^*_{\deg}-\alpha_{\mathrm{HMT}}
 \label{eq:l9s102:eta-corolario-inverso}
\end{equation}
queda establecida como corolario algebraico de la construcción directa.

% HMT_SOURCE_SEGMENT_END id=D04
\section{Conmutatividad exacta de las cartas en vueltas, grados y radianes}
\label{sec:s102:c38:5}

\subsection{Derivación de la longitud cilíndrica y del grado analítico del precursor parangular}
\label{sec:s102:c38:5:1}

% HMT_SOURCE_SEGMENT_BEGIN id=B03 source=colaboracion/parte_iii/source/public_final/base_83/c30_body.tex body_lines=157-194

Sea \(\mathcal P\) el módulo libre generado por las palabras regionales y
sea \(F^n\mathcal P\) su filtración por longitud. La graduación asociada
registra el primer nivel en el que aparece cada palabra. Para un parámetro
\(z\), el carácter ordinario de longitud es
\[
\chi_z([w])=z^{|w|}.
\]
La multiplicatividad
\[
\chi_z([uv])=\chi_z([u])\chi_z([v])
\]
expresa que la concatenación suma longitudes. Al evaluar en
\(z=\alpha\), una palabra de longitud \(n\) recibe grado \(\alpha^n\).

\begin{definition}[Compatibilidad de grado]
\label{p3-83-c30-s1:def:compatibilidad-grado}
El lector analítico regional es compatible con la filtración cilíndrica
cuando envía el componente homogéneo de longitud \(n\) al grado analítico
\(n\):
\[
\operatorname{gr}_n\mathcal P\longrightarrow
\alpha^n\mathbb R.
\]
\end{definition}

Las cinco regiones pertenecen a \(\mathcal U_6\); por ello, el dato de
retorno aparece en grado seis:
\[
\rho_\pi[U_6]\longmapsto169\alpha^6.
\]
Este seis es la longitud de la palabra regional. No es la longitud de los
paseos cerrados del \cref{chap:cinco-ciclos-pi}. Aquellos paseos recorren a
la vez los cuatro anclajes
\(U_\pi,U_e,U_\varphi,U_{\rm APP}\) y poseen longitud mínima ocho. Una
palabra, una arista del grafo de bloques y un paseo cerrado son objetos de
tipos distintos; la graduación utilizada aquí pertenece al primero.

% HMT_SOURCE_SEGMENT_END id=B03
\subsection{Publicación del grado analítico en la carta global de fase}
\label{sec:s102:c38:5:2}

% HMT_SOURCE_SEGMENT_BEGIN id=A03 source=colaboracion/parte_iii/source/public_final/ascensos_20260826/16f_realizacion_analitica_contraangulo.tex body_lines=172-209

Sea \(\mathcal P\) el módulo libre generado por las palabras regionales y
sea \(F^n\mathcal P\) su filtración por longitud. La graduación asociada
registra el primer nivel en el que aparece cada palabra. Para un parámetro
\(z\), el carácter ordinario de longitud es
\[
\chi_z([w])=z^{|w|}.
\]
La multiplicatividad
\[
\chi_z([uv])=\chi_z([u])\chi_z([v])
\]
expresa que la concatenación suma longitudes. Al evaluar en
\(z=\alpha\), una palabra de longitud \(n\) recibe grado \(\alpha^n\).

\begin{definition}[Compatibilidad de grado]
\label{p3-20260826:def:compatibilidad-grado}
El lector analítico regional es compatible con la filtración cilíndrica
cuando envía el componente homogéneo de longitud \(n\) al grado analítico
\(n\):
\[
\operatorname{gr}_n\mathcal P\longrightarrow
\alpha^n\mathbb R.
\]
\end{definition}

Las cinco regiones pertenecen a \(\mathcal U_6\); por ello, el dato de
retorno aparece en grado seis:
\[
\rho_\pi[U_6]\longmapsto169\alpha^6.
\]
Este seis es la longitud de la palabra regional. No es la longitud de los
paseos cerrados del \cref{chap:cinco-ciclos-pi}. Aquellos paseos recorren a
la vez los cuatro anclajes
\(U_\pi,U_e,U_\varphi,U_{\rm APP}\) y poseen longitud mínima ocho. Una
palabra, una arista del grafo de bloques y un paseo cerrado son objetos de
tipos distintos; la graduación utilizada aquí pertenece al primero.

% HMT_SOURCE_SEGMENT_END id=A03
% HMT_SOURCE_SEGMENT_BEGIN id=D02 source=manuscrito/sucesor_102/deltas_ley9/c38_par_angular_cartas_y_canales.tex body_lines=21-34
El calendario global exhibe las tres factorizaciones
\begin{equation}
 360=12\cdot30=4\cdot90=3\cdot120.
 \label{eq:l9s102:360-factorizaciones}
\end{equation}
La carta analítica se obtiene mediante
\begin{equation}
 A_{\mathrm{rad}}=\frac\pi{180}A_{\deg}
 =\frac{50\pi}{9}\alpha_{\mathrm{HMT}},
 \qquad
 C^*_{\mathrm{rad}}=\frac\pi{180}C^*_{\deg}
 =\frac\pi{90}(\eta_{\mathrm{ret}}+\alpha_{\mathrm{HMT}}).
 \label{eq:l9s102:a-cstar-rad}
\end{equation}
% HMT_SOURCE_SEGMENT_END id=D02
% HMT_SOURCE_SEGMENT_BEGIN id=D03 source=manuscrito/sucesor_102/deltas_ley9/c38_par_angular_cartas_y_canales.tex body_lines=35-59

\begin{theorem}[Conmutatividad de las factorizaciones regional y de acción]
\label{thm:l9s102:conmutatividad-cstar}
Las dos rutas hacia la fase conjugada producen la misma salida:
\begin{equation}
 \boxed{
 \frac\pi{180}C^*_{\deg}
 =\frac\pi{90}(\eta_{\mathrm{ret}}+\alpha_{\mathrm{HMT}})
 =\frac\pi{90}(H_5+\alpha_{\mathrm{HMT}})
 -\mathscr C_\pi(\alpha_{\mathrm{HMT}})=y_*.}
 \label{eq:l9s102:cuadrado-cstar}
\end{equation}
\end{theorem}

\begin{proof}
Sustituyendo \cref{eq:l9s102:eta-ret} en el segundo miembro se obtiene
\[
 \frac\pi{90}
 \left(H_5-\frac{90}{\pi}\mathscr C_\pi+\alpha_{\mathrm{HMT}}\right)
 =\frac\pi{90}(H_5+\alpha_{\mathrm{HMT}})-\mathscr C_\pi.
\]
La última expresion es la definicion regional de \(y_*\), mientras la
primera igualdad procede de \cref{eq:l9s102:a-cstar-deg}.
\end{proof}

% HMT_SOURCE_SEGMENT_END id=D03
\section{Descomposición espectral \texorpdfstring{\(q_\pm\)}{q+/-} y cálculo funcional sobre los canales \texorpdfstring{\(90/120\)}{90/120} previamente generados}
\label{sec:s102:c38:6}

\subsection{Reconstrucción espectral del par angular y transición hexada--octada}
\label{sec:s102:c38:6:1}

% HMT_SOURCE_SEGMENT_BEGIN id=B11 source=colaboracion/parte_iii/source/public_final/base_83/c30_body.tex body_lines=523-569

La fase común y la fase orientada determinan los dos canales
\[
q_-=e^{-x+y_*},
\qquad
q_+=e^{-x-y_*}.
\]
Se recuperan exactamente los invariantes
\[
q_-q_+=D_A,
\qquad
\alpha=-\frac9{100\pi}\log D_A,
\qquad
C_*^{\rm HMT}=\frac{90}{\pi}\log\frac{q_-}{q_+}.
\]
La primera igualdad lee la contracción común; la tercera, la separación
orientada de las hojas. El intercambio \(R\mapsto-R\) permuta
\(q_+\) y \(q_-\), conserva \(D_A\) e invierte el signo de \(y_*\).

Para
\[
S_{90}(q)=\frac{q^{90}}{1-q^{270}},
\qquad
S_{120}(q)=\frac{q^{120}}{1-q^{360}},
\]
el funcional de transición hexada--octada se escribe íntegramente en la
carta de fases:
\[
\Gamma_{6\to8}
=\frac{180}{\pi}xy_*
+12\bigl(S_{90}(q_-)-S_{90}(q_+)\bigr)
-\bigl(S_{120}(q_-)-S_{120}(q_+)\bigr).
\]
La evaluación producida por la coordenada \(C_*^{\rm HMT}\) del
\cref{p3-83-c30-s1:thm:contraangulo-regional} es
\[
\boxed{
\Gamma_{6\to8}
=0.274007672694459274747255260774\ldots.}
\]
Así, el valor interno reconocido en Mecánica Dimensional como funcional de
Barbero--Immirzi deja de recibir \(C_*^{\rm HMT}\) como literal independiente:
ambos descienden del mismo par de fases generado en este capítulo. La
interpretación geométrica y física de \(\Gamma_{6\to8}\) pertenece a la obra
de Mecánica Dimensional; su evaluación matemática pertenece ya a la interfaz
previamente construida.

% HMT_SOURCE_SEGMENT_END id=B11
\subsection{Transportador relativo y disjunción espectral de los bloques}
\label{sec:s102:c38:6:2}

% HMT_SOURCE_SEGMENT_BEGIN id=B13 source=colaboracion/parte_iii/source/public_final/base_83/c30_body.tex body_lines=619-675
\label{p3-83-c30-s1:sec:transportador-relativo-split}

En la base ordenada de canales, sea
\[
 R=\begin{pmatrix}0&1\\1&0\end{pmatrix},
 \qquad P_\pm=\frac{I\pm R}{2},
 \qquad
 B_0=7I+2R,
 \qquad B_1=\Adeg I+\Cdeg R.
\]
La descomposición escindida
\[
 xI+yR=\rho(x,y)e^{\eta(x,y)R},
 \quad \rho(x,y)=\sqrt{x^2-y^2},
 \quad \eta(x,y)=\operatorname{artanh}(y/x)
\]
construye el transportador único
\[
 T_{\rm rel}
 =\frac{\sqrt{(\Adeg)^2-(\Cdeg)^2}}{\sqrt{45}}
   \exp\!\left[
   \left(\operatorname{artanh}\frac{\Cdeg}{\Adeg}
   -\operatorname{artanh}\frac27\right)R\right],
 \qquad T_{\rm rel}B_0=B_1.
\]
En la base espectral actúa por
\(9\mapsto\Adeg+\Cdeg\) y
\(5\mapsto\Adeg-\Cdeg\).

\begin{theorem}[Disjunción espectral de las realizaciones local y angular]
\label{p3-83-c30-s1:thm:no-go-entrelazador-bloques-split}
Los espectros
\[
 \operatorname{Spec}(B_0)=\{9,5\},\qquad
 \operatorname{Spec}(B_1)=\{\Adeg+\Cdeg,\Adeg-\Cdeg\}
\]
son disjuntos. En consecuencia,
\[
 \operatorname{Hom}_{B_0,B_1}
 :=\{\Phi:\Phi B_0=B_1\Phi\}=\{0\}.
\]
\end{theorem}

\begin{proof}
Las cotas racionales construidas para las coordenadas angulares dan
\(7.29<\Adeg<7.30\) y \(2.12<\Cdeg<2.13\); por tanto,
\(9.41<\Adeg+\Cdeg<9.43\) y
\(5.16<\Adeg-\Cdeg<5.18\). En la base que diagonaliza \(R\), cada entrada
de un entrelazador queda multiplicada por la diferencia entre un autovalor
de \(B_0\) y uno de \(B_1\); la disjunción obliga a que todas se anulen.
\end{proof}

El transportador relativo compara dos bloques dentro del álgebra escindida;
los entrelazadores dodecafásico--Witt y de incidencia actúan entre
representaciones equivalentes en sus residencias propias. Esta separación
tipa positivamente ambas operaciones y conserva sus dominios.

% HMT_SOURCE_SEGMENT_END id=B13
\subsection{Evaluación espectral \texorpdfstring{\(V_{90,120}\)}{V_90,120} de los canales de incidencia previamente generados por el TPK}
\label{sec:s102:c38:6:3}

% HMT_SOURCE_SEGMENT_BEGIN id=A11 source=colaboracion/parte_iii/source/public_final/ascensos_20260826/16f_realizacion_analitica_contraangulo.tex body_lines=554-572

La fase común y la fase orientada determinan los dos canales
\[
q_-=e^{-\Arad+\Crad},
\qquad
q_+=e^{-\Arad-\Crad}.
\]
Se recuperan exactamente los invariantes
\[
q_-q_+=D_A,
\qquad
\alpha=-\frac9{100\pi}\log D_A,
\qquad
\Cdeg=\frac{90}{\pi}\log\frac{q_-}{q_+}.
\]
La primera igualdad lee la contracción común; la tercera, la separación
orientada de las hojas. El intercambio \(R\mapsto-R\) permuta
\(q_+\) y \(q_-\), conserva \(D_A\) e invierte el signo de \(\Crad\).

% HMT_SOURCE_SEGMENT_END id=A11
% HMT_SOURCE_SEGMENT_BEGIN id=D05 source=manuscrito/sucesor_102/deltas_ley9/c38_par_angular_cartas_y_canales.tex body_lines=68-97
\label{sec:l9s102:qpm-posteriores}

Sea \(H^2=I\) y \(P_\pm=(I\pm H)/2\). El bloque angular
\begin{equation}
 B=A_{\mathrm{rad}}I+C^*_{\mathrm{rad}}H
 =(A_{\mathrm{rad}}+C^*_{\mathrm{rad}})P_+
 +(A_{\mathrm{rad}}-C^*_{\mathrm{rad}})P_-
 \label{eq:l9s102:b-angular}
\end{equation}
se publica multiplicativamente por
\begin{equation}
 e^{-B}=q_+P_++q_-P_-,
 \qquad
 q_\pm=\exp[-(A_{\mathrm{rad}}\pm C^*_{\mathrm{rad}})].
 \label{eq:l9s102:qpm-angular}
\end{equation}
La involución \(C^*\mapsto-C^*\) permuta \(q_+\) y \(q_-\), y conserva
\begin{equation}
 q_+q_-=e^{-2A_{\mathrm{rad}}}=D_A.
 \label{eq:l9s102:det-angular}
\end{equation}
Los cardinales \(90/120\), construidos previamente por la incidencia del
TPK, reciben ahora su evaluación espectral:
\begin{equation}
 V_{90,120}
 =(I-T^{90})(I-T^{120})^{-1}
 =\sum_{\sigma\in\{+,-\}}
 \frac{1-q_\sigma^{90}}{1-q_\sigma^{120}}P_\sigma.
 \label{eq:l9s102:v90120-posterior}
\end{equation}
% HMT_SOURCE_SEGMENT_END id=D05
\subsubsection{Transportador relativo escindido}
\label{sec:s102:c38:6:3:1}

% HMT_SOURCE_SEGMENT_BEGIN id=A12 source=colaboracion/parte_iii/source/public_final/ascensos_20260826/16f_realizacion_analitica_contraangulo.tex body_lines=574-748
\label{p3-20260826:sec:transportador-relativo-split}

La comparación con el bloque central de APP puede efectuarse sin identificar
representaciones distintas. Fijemos la base ordenada de canales y la
involución
\[
 R=\begin{pmatrix}0&1\\1&0\end{pmatrix},
 \qquad R^2=I,
 \qquad
 P_\pm=\frac{I\pm R}{2}.
\]
En el álgebra conmutativa
\[
 \mathscr B_R:=\{xI+yR:x,y\in\mathbb R\}
\]
consideremos, dentro de la carta sexagesimal ya declarada,
\[
 B_0:=7I+2R,
 \qquad
 B_1:=\Adeg I+\Cdeg R.
\]
El cambio ortogonal de base
\[
 S=\frac1{\sqrt2}
 \begin{pmatrix}1&1\\1&-1\end{pmatrix}
\]
diagonaliza simultáneamente la involución y ambos bloques. Para \(x>|y|\),
definamos
\[
 \rho(x,y):=\sqrt{x^2-y^2},
 \qquad
 \eta(x,y):=\operatorname{artanh}\!\left(\frac yx\right).
\]
Entonces
\[
 xI+yR=\rho(x,y)e^{\eta(x,y)R}.
\]
Si
\[
 \rho_0=\sqrt{45},\quad
 \eta_0=\operatorname{artanh}(2/7),\quad
 \rho_1=\rho(\Adeg,\Cdeg),\quad
 \eta_1=\eta(\Adeg,\Cdeg),
\]
la única multiplicación izquierda que lleva \(B_0\) a \(B_1\) es
\[
 \boxed{
 T_{\rm rel}
 :=\frac{\rho_1}{\sqrt{45}}
 e^{(\eta_1-\eta_0)R},
 \qquad
 T_{\rm rel}B_0=B_1.}
\]

\begin{proposition}[Transportador relativo de los dos bloques]
\label{p3-20260826:prop:transportador-relativo-split}
Una vez fijados \(R\), el orden de los canales, la rama positiva y la carta
sexagesimal, \(T_{\rm rel}\) es el único elemento de
\(\mathscr B_R\) que satisface \(TB_0=B_1\). En la base espectral actúa por
\[
 9\longmapsto\Adeg+\Cdeg,
 \qquad
 5\longmapsto\Adeg-\Cdeg.
\]
\end{proposition}

\begin{proof}
La descomposición polar escindida da la fórmula impresa y prueba la existencia.
Como \(\det B_0=45\ne0\), toda solución matricial de \(TB_0=B_1\) satisface
\(T=B_1B_0^{-1}\); por tanto es única. La diagonalización mediante \(S\)
produce los dos cocientes espectrales
\((\Adeg+\Cdeg)/9\) y \((\Adeg-\Cdeg)/5\), equivalentes a la expresión
exponencial de \(T_{\rm rel}\).
\end{proof}

Este transportador no es un entrelazador representacional. En efecto, un
entrelazador \(\Phi\) entre los endomorfismos \(B_0\) y \(B_1\) debería
satisfacer
\[
 \Phi B_0=B_1\Phi.
\]

\begin{theorem}[Obstrucción espectral al entrelazamiento]
\label{p3-20260826:thm:no-go-entrelazador-bloques-split}
Para los valores de \(\Adeg\) y \(\Cdeg\) construidos en esta sección,
\[
 \boxed{\Phi B_0=B_1\Phi\quad\Longrightarrow\quad\Phi=0.}
\]
\end{theorem}

\begin{proof}
Los espectros son
\[
 \operatorname{Spec}(B_0)=\{9,5\},
 \qquad
 \operatorname{Spec}(B_1)
 =\{\Adeg+\Cdeg,\Adeg-\Cdeg\}.
\]
Las cotas racionales previamente demostradas proporcionan los intervalos
disjuntos
\[
 7.29<\Adeg<7.30,
 \qquad
 2.12<\Cdeg<2.13,
\]
y, por consiguiente,
\[
 9.41<\Adeg+\Cdeg<9.43,
 \qquad
 5.16<\Adeg-\Cdeg<5.18.
\]
En la base que diagonaliza \(R\), cada entrada \(\phi_{ij}\) de \(\Phi\)
queda multiplicada por la diferencia entre un autovalor de \(B_0\) y uno de
\(B_1\). Todas esas diferencias son no nulas; luego todas las entradas de
\(\Phi\) se anulan.
\end{proof}

Numéricamente,
\[
 \rho_1=6.9814819335172378048\ldots,
 \qquad
 \frac{\rho_1}{\sqrt{45}}=1.0407378791354140779\ldots,
\]
\[
 \eta_1-\eta_0
 =0.005796351470571295590\ldots.
\]
La fórmula conserva el álgebra escindida, los proyectores \(P_\pm\) y la forma
lorentziana hasta el factor conforme \(\rho_1^2/45\). Depende, sin embargo,
de \(B_1\) ya construido: no constituye una derivación anterior de
\(\Adeg\) o \(\Cdeg\), ni debe confundirse con los entrelazadores
equivariantes entre realizaciones de un mismo módulo.

Para
\[
S_{90}(q)=\frac{q^{90}}{1-q^{270}},
\qquad
S_{120}(q)=\frac{q^{120}}{1-q^{360}},
\]
el funcional de transición hexada--octada se escribe íntegramente en la
carta de fases:
\[
\Gamma_{6\to8}
=\frac{180}{\pi}\Arad\Crad
+12\bigl(S_{90}(q_-)-S_{90}(q_+)\bigr)
-\bigl(S_{120}(q_-)-S_{120}(q_+)\bigr).
\]
La evaluación producida por el contraángulo del
\cref{p3-20260826:thm:contraangulo-regional} es
\[
\boxed{
\Gamma_{6\to8}
=0.274007672694459274747255260774\ldots.}
\]
Las formas equivalentes del funcional matemático interno son
\[
\boxed{
\Gamma_{6\to8}
=\Arad\Cdeg+K_{6\to8}
=\frac{180}{\pi}\Arad\Crad+K_{6\to8}
=\frac{\pi\Adeg\Cdeg}{180}+K_{6\to8}.}
\]
Estas identidades evalúan \(\Gamma_{6\to8}\) después de generar
\((\Crad,\Cdeg)\); no se utiliza en sentido inverso para seleccionar el
contraángulo. La lectura física en la conexión de Ashtekar--Barbero se
representa después mediante
\[
 \Gamma_{6\to8}\dashrightarrow\gamma_{\rm BI}.
\]
La flecha no introduce un tercer funcional ni retroactúa sobre \(C^*\). Las
series \(S_{90}\), \(S_{120}\), los coeficientes \(12\) y \(-1\) y su
asignación al paso \(6\to8\) pertenecen a la construcción matemática
declarada. La identificación física exige, además, la normalización propia
de la conexión.

% HMT_SOURCE_SEGMENT_END id=A12
\section{Naturalidad bajo refinamiento y separación respecto del reconocimiento metrológico}
\label{sec:s102:c38:7}

\subsection{Estimación exacta del resto de la prolongación analítica}
\label{sec:s102:c38:7:1}

% HMT_SOURCE_SEGMENT_BEGIN id=B09 source=colaboracion/parte_iii/source/public_final/base_83/c30_body.tex body_lines=474-499

La expresión hallada inicialmente,
\[
y_7=\frac\pi{90}(H_5+\alpha)-169\alpha^6-D_A\alpha^7,
\]
es el truncamiento de grado siete del carácter completo. Produce
\[
C_7=2.12373833896864600265403827103919\ldots{}^\circ.
\]
El resto omitido no es un término indeterminado; es la cola geométrica
exacta
\[
\mathscr C_\pi(\alpha)
-169\alpha^6-D_A\alpha^7
=\frac{D_A^2\alpha^8}{1-D_A\alpha}.
\]
Por tanto,
\[
\left|C_7-C_*^{\rm HMT}\right|
=\frac{180}{\pi}\frac{D_A^2\alpha^8}{1-D_A\alpha}
=2.7839208689064583\ldots\times10^{-16}\;{}^\circ.
\]
La recurrencia transforma así una aproximación de séptimo grado en una
realización analítica completa y proporciona una cota cerrada para cada
truncamiento posterior.

% HMT_SOURCE_SEGMENT_END id=B09
\subsection{Dependencia funcional del selector regional y de la continuación analítica}
\label{sec:s102:c38:7:2}

% HMT_SOURCE_SEGMENT_BEGIN id=B10 source=colaboracion/parte_iii/source/public_final/base_83/c30_body.tex body_lines=501-521

La construcción cambia al retirar cualquiera de sus componentes activos.
Los siguientes valores se obtienen sin recalibrar los demás términos:
\[
\begin{array}{@{}lc@{}}
\toprule
\text{variante}&\text{coordenada angular conjugada en grados}\\
\midrule
\rho_\pi=168&2.1237383389772976863904487\ldots\\
\rho_\pi=169&2.1237383389686457242619514\ldots\\
\rho_\pi=170&2.1237383389599937621334541\ldots\\
\text{sin cola determinantal}&2.1237383389686949415301675\ldots\\
D_A\text{ sustituido por }1&2.1237383389686313409965826\ldots\\
\bottomrule
\end{array}
\]
Estas ablaciones no constituyen por sí solas una estimación probabilística.
Demuestran que el selector equivariante y el carácter determinantal son
componentes efectivas de la salida y que no pueden suprimirse manteniendo la
misma realización.

% HMT_SOURCE_SEGMENT_END id=B10
\subsection{Interpretación operatoria de la coordenada angular conjugada}
\label{sec:s102:c38:7:3}

% HMT_SOURCE_SEGMENT_BEGIN id=B14 source=colaboracion/parte_iii/source/public_final/base_83/c30_body.tex body_lines=677-715

El procedimiento completo puede resumirse como una cadena de aplicaciones:
\[
\begin{aligned}
\mathcal F_\pi
&\longmapsto b_\pi
\longmapsto\rho_\pi=169,\\
\alpha
&\longmapsto A
\longmapsto x
\longmapsto D_A,\\
(\rho_\pi,D_A)
&\longmapsto\mathscr C_\pi(\alpha),\\
(H_5,\mathscr C_\pi(\alpha))
&\longmapsto y_*
\longmapsto C_*^{\rm HMT},\\
(x,y_*)
&\longmapsto(q_+,q_-)
\longmapsto\Gamma_{6\to8}.
\end{aligned}
\]
La primera línea pertenece al catálogo finito; la segunda, al cierre común de
\(\alpha\); la tercera, a la prolongación determinantal; la cuarta produce la
coordenada orientada; la quinta la transporta hacia la interfaz dimensional.

En esta organización, \(\alpha\) mide la contracción común del transporte y
\(C_*^{\rm HMT}\) mide la separación orientada de sus dos hojas. El retorno
regional no añade una constante externa: selecciona una amplitud finita y la
prolonga mediante el determinante ya fijado por \(\alpha\). La carta en
grados es posterior. El objeto primario es el par de fases \((x,y_*)\).

Una verificación exhaustiva reconstruye las cinco regiones directamente
desde el catálogo, recorre las acciones de
\(\mathfrak S_5\) y \(\mathfrak S_3\), comprueba la recurrencia y suma la
serie. La entrada de \(\alpha\) es la salida exacta de su cierre
dodecafásico, no una cifra metrológica transcrita. La comprobación mantiene
separados los teoremas del catálogo, las reglas del lector y el contraste
externo; por ello ninguna comparación experimental interviene en la función
generadora.
% HMT_SOURCE_SEGMENT_END id=B14
\subsection{Certificación de convergencia del truncamiento parangular}
\label{sec:s102:c38:7:4}

% HMT_SOURCE_SEGMENT_BEGIN id=A09 source=colaboracion/parte_iii/source/public_final/ascensos_20260826/16f_realizacion_analitica_contraangulo.tex body_lines=505-530

La expresión hallada inicialmente,
\[
y_7=\frac\pi{90}(H_5+\alpha)-169\alpha^6-D_A\alpha^7,
\]
es el truncamiento de grado siete del carácter completo. Produce
\[
C_7=2.12373833896864600265403827103919\ldots{}^\circ.
\]
El resto omitido no es un término indeterminado; es la cola geométrica
exacta
\[
\mathscr C_\pi(\alpha)
-169\alpha^6-D_A\alpha^7
=\frac{D_A^2\alpha^8}{1-D_A\alpha}.
\]
Por tanto,
\[
\left|C_7-\Cdeg\right|
=\frac{180}{\pi}\frac{D_A^2\alpha^8}{1-D_A\alpha}
=2.7839208689064583\ldots\times10^{-16}\;{}^\circ.
\]
La recurrencia transforma así una aproximación de séptimo grado en una
realización analítica completa y proporciona una cota cerrada para cada
truncamiento posterior.

% HMT_SOURCE_SEGMENT_END id=A09
\subsection{Naturalidad del selector regional bajo refinamiento y cambio de carta}
\label{sec:s102:c38:7:5}

% HMT_SOURCE_SEGMENT_BEGIN id=A10 source=colaboracion/parte_iii/source/public_final/ascensos_20260826/16f_realizacion_analitica_contraangulo.tex body_lines=532-552

La construcción cambia al retirar cualquiera de sus componentes activos.
Los siguientes valores se obtienen sin recalibrar los demás términos:
\[
\begin{array}{@{}lc@{}}
\toprule
\text{variante}&\text{contraángulo en grados}\\
\midrule
\rho_\pi=168&2.1237383389772976863904487\ldots\\
\rho_\pi=169&2.1237383389686457242619514\ldots\\
\rho_\pi=170&2.1237383389599937621334541\ldots\\
\text{sin cola determinantal}&2.1237383389686949415301675\ldots\\
D_A\text{ sustituido por }1&2.1237383389686313409965826\ldots\\
\bottomrule
\end{array}
\]
Estas ablaciones no constituyen por sí solas una estimación probabilística.
Demuestran que el selector equivariante y el carácter determinantal son
componentes efectivas de la salida y que no pueden suprimirse manteniendo la
misma realización.

% HMT_SOURCE_SEGMENT_END id=A10
\subsection{Independencia metrológica de la construcción parangular}
\label{sec:s102:c38:7:6}

% HMT_SOURCE_SEGMENT_BEGIN id=A14 source=colaboracion/parte_iii/source/public_final/ascensos_20260826/16f_realizacion_analitica_contraangulo.tex body_lines=799-831

El procedimiento completo puede resumirse como una cadena de aplicaciones:
\[
\begin{aligned}
\mathcal F_\pi
&\longmapsto b_\pi
\longmapsto\rho_\pi=169,\\
\alpha
&\longmapsto \Adeg
\longmapsto \Arad
\longmapsto D_A,\\
(\rho_\pi,D_A)
&\longmapsto\mathscr C_\pi(\alpha),\\
(H_5,\mathscr C_\pi(\alpha))
&\longmapsto \Crad
\longmapsto \Cdeg,\\
(\Arad,\Crad)
&\longmapsto(q_+,q_-)
\longmapsto\Gamma_{6\to8}.
\end{aligned}
\]
La primera línea pertenece al catálogo finito; la segunda, al cierre común de
\(\alpha\); la tercera, a la prolongación determinantal; la cuarta produce la
coordenada orientada; la quinta la transporta hacia la interfaz dimensional.

En esta organización, \(\alpha\) mide la contracción común del transporte y
el contraángulo mide la separación orientada de sus dos hojas. En la
evaluación vigente, el retorno regional selecciona una amplitud finita y la
prolonga mediante el determinante ya fijado por \(\alpha\), sin insertar una
constante externa adicional. La carta en grados es posterior. El objeto
primario es el par de fases \((\Arad,\Crad)\). La recurrencia, la suma
regional y las ablaciones dependen del catálogo y de
\(\alpha_{\rm HMT}\), pero no reciben una constante metrológica adicional.
% HMT_SOURCE_SEGMENT_END id=A14



\HMTFlushCurrentChapter
\HMTSuccessorChapterInput
  {Operador constitutivo positivo del vacío: completación \(3\to4\), reconstrucción espectral y derivación conjunta de \texorpdfstring{\(\varepsilon,\mu,Z\)}{epsilon, mu, Z} y \(c\)}
  {../colaboracion/parte_iii/tex/capitulos_finales/39_vacio_constitutivo.tex}

\HMTFlushCurrentChapter
\HMTSuccessorChapterInput
  {Derivación de la carga y de los cuantos eléctricos: resistencia de von Klitzing, conductancia, constante de Josephson y flujo magnético}
  {../colaboracion/parte_iii/tex/capitulos_finales/40_cuantos_electricos.tex}

\HMTFlushCurrentChapter
\HMTSuccessorChapterInput
  {Caracteres cuadráticos y constantes espectrales: Catalán, Apéry, Euler--Mascheroni y jerarquía de Stieltjes}
  {../colaboracion/parte_iii/tex/capitulos_finales/41_catalan_apery_euler.tex}

\HMTFlushCurrentChapter
\HMTSuccessorChapterInput
  {Constantes de Euler--Kronecker y distribuciones primas: caracteres dodecafásicos, Meissel--Mertens y series de primos gemelos}
  {../colaboracion/parte_iii/tex/capitulos_finales/42_euler_kronecker_primos.tex}

\HMTFlushCurrentChapter
\HMTSuccessorChapterInput
  {Familia Bendersky--Glaisher--Kinkelin: productos regulares, derivadas de la función zeta y jerarquías de normalización}
  {../colaboracion/parte_iii/tex/capitulos_finales/43_bendersky_glaisher.tex}

\HMTFlushCurrentChapter
\HMTSuccessorChapterInput
  {Dinámica de Gauss y constantes de Khinchin, Lévy y Lochs: cofinalidad de cilindros, operador de transferencia y espectro de Gauss--Kuzmin--Wirsing}
  {../colaboracion/parte_iii/tex/capitulos_finales/44_gauss_khinchin_levy_lochs.tex}

% L9-059. Propietario analítico íntegro de los funcionales espectrales
% triádicos (SHA-256 c53c5efa08ecacfed1a5cd01f20eaa372a29a56adab4873bcb442709963496f4).
% Se consume una sola vez y se subordina únicamente su jerarquía visible.
\begin{HMTScientificOwnerSubordinated}
\chapter{Construcción analítica de los funcionales de la torre triádica}
\label{chap:funcionales-espectrales-triadicos}

La síntesis anterior ha situado las constantes espectrales dentro de la
torre triádica. Este bloque contiene su construcción analítica primaria:
sucesión de operadores autoadjuntos sobre ciclos de orden \(3^m\), límite de
trazas de calor y transformada de Mellin. Las constantes clásicas aparecen
como evaluaciones, partes finitas o derivadas del funcional meromorfo
resultante, no como prefijos reconocidos en un catálogo.

La normalización del operador $D_m$ contiene explícitamente $\pi$ en la escala
que relaciona el laplaciano discreto con el núcleo de calor gaussiano. El
apartado constituye, por tanto, una extensión analítica de la rama
arquimediana ya construida. Las constantes espectrales posteriores se obtienen
simultáneamente del funcional resultante, con esa normalización como hipótesis
de partida.

\section{Límites inductivos y escala de refinamiento}

La prolongación espectral exige distinguir dos niveles analíticos. El
primero corresponde a familias acotadas y compatibles; el segundo, a
operadores no acotados, para los cuales el entrelazamiento algebraico no
resuelve por sí solo los problemas de dominio.

\begin{theorem}[Límite inductivo acotado]
\label{pf84:E03-06}
Sean \(J_m:\mathcal H_m\to\mathcal H_{m+1}\) isometrías y
\(A_m\in\mathcal B(\mathcal H_m)\) operadores autoadjuntos tales que
\[
 A_{m+1}J_m=J_mA_m,
 \qquad
 \sup_m\lVert A_m\rVert<\infty.
\]
Existe un único operador acotado autoadjunto \(A\) sobre
\(\mathcal H_\infty=\varinjlim(\mathcal H_m,J_m)\) cuya restricción a cada
nivel es \(A_m\).
\end{theorem}

\begin{proof}
En la unión algebraica densa se define
\[
 A[J_{m,\infty}\xi]:=J_{m,\infty}A_m\xi.
\]
El entrelazamiento hace que la definición no dependa del representante. La
cota uniforme permite extender \(A\) por continuidad a
\(\mathcal H_\infty\); la simetría en el subespacio denso y la acotación
implican su autoadjunción.
\end{proof}

\begin{remark}[Control para operadores no acotados]
\label{pf84:E03-07}
Si los \(A_m\) no son acotados, la identidad
\(A_{m+1}J_m=J_mA_m\) sólo define una acción en una unión de dominios.
Para obtener un operador cerrado o autoadjunto deben declararse los dominios
transportados, probarse cerrabilidad y controlarse los índices de
deficiencia, o bien demostrarse esencial autoadjunción sobre un núcleo común
compatible. Omitir estas obligaciones constituye un fallo de tipo, no una
abreviación de la prueba.
\end{remark}

Sea ahora \(\mathcal S_\infty\) un conjunto no vacío de secciones
supervivientes compatibles y sea \(\mathcal P_m\) el conjunto finito de sus
prefijos de longitud \(m\). Escribimos
\[
 \mathcal H_m:=\ell^2(\mathcal P_m),
 \qquad
 E_m(\gamma,\gamma')=
 \begin{cases}
 1,&\gamma'\text{ prolonga }\gamma,\\
 0,&\text{en otro caso},
 \end{cases}
\]
y
\[
 d_m(\gamma):=
 \#\{\gamma'\in\mathcal P_{m+1}:\gamma'\mapsto\gamma\},
 \qquad
 D_m^{\deg}:=E_mE_m^*=\operatorname{diag}(d_m(\gamma)).
\]
Nos restringimos al subespacio vivo \(d_m(\gamma)>0\), y cada prefijo de
nivel \(m+1\) posee un único padre.

\begin{theorem}[Isometría normalizada de prolongación]
\label{pf84:C01}
El operador
\[
 J_m:=E_m^*(D_m^{\deg})^{-1/2}:
 \mathcal H_m\longrightarrow\mathcal H_{m+1}
\]
es una isometría, y
\[
 J_me_\gamma=
 \frac1{\sqrt{d_m(\gamma)}}
 \sum_{\gamma'\mapsto\gamma}e_{\gamma'}.
\]
\end{theorem}

\begin{proof}
La unicidad del padre da
\[
 J_m^*J_m
 =(D_m^{\deg})^{-1/2}E_mE_m^*(D_m^{\deg})^{-1/2}=I.
\]
La fórmula sobre \(e_\gamma\) se obtiene desarrollando la columna
correspondiente de \(E_m^*\).
\end{proof}

Puesto
\[
 \mathcal K_m:=\mathcal H_{m+1}\ominus J_m\mathcal H_m,
 \qquad
 \mathcal H_{\mathrm{surv}}
 \cong\mathcal H_0\oplus\bigoplus_{m\ge0}\mathcal K_m,
\]
se define el operador de refinamiento
\[
 \mathcal R|_{\mathcal H_0}=0,
 \qquad
 \mathcal R|_{\mathcal K_m}=3^{m+1}I_{\mathcal K_m},
\]
con dominio maximal
\[
 \Dom(\mathcal R)=
 \mathcal H_0\oplus
 \left\{(\xi_m)_{m\ge0}:
 \sum_{m\ge0}3^{2m+2}\lVert\xi_m\rVert^2<\infty\right\}.
\]

\begin{theorem}[Operador de refinamiento]
\label{pf84:C02}
Si cada \(\mathcal K_m\) es finito dimensional, \(\mathcal R\) es positivo,
autoadjunto y de resolvente compacto. Su espectro mide profundidad de
refinamiento; no constituye un espectro de primos ni de ceros.
\end{theorem}

\begin{proof}
La descomposición ortogonal reduce \(\mathcal R\) a bloques escalares reales
sobre su dominio maximal, lo que prueba positividad y autoadjunción. Los
autovalores \(3^{m+1}\) escapan a infinito y sus multiplicidades son finitas;
por tanto, el resolvente es compacto. Un bloque infinito dimensional a
autovalor acotado falsaría esta última conclusión.
\end{proof}

\section{Operadores entero, primo y orbital}

Sea \(\mathscr W_{\mathrm{bal}}\) el conjunto de palabras ternarias
balanceadas reducido mediante el teorema
\cref{pf84:C03}. Sobre \(\ell^2(\mathscr W_{\mathrm{bal}})\), el operador
\[
 Be_w=b(w)e_w
\]
se toma en su dominio maximal. Su restricción al subespacio de los enteros
positivos se denota por \(Q\).

\begin{theorem}[Operadores entero y positivo]
\label{pf84:C04}
El operador \(B\) es autoadjunto, posee espectro simple \(\ZZ\) y
resolvente compacto. El operador \(Q\) posee espectro simple
\(\NN_{\ge1}\) y resolvente compacto.
\end{theorem}

\begin{proof}
La biyección de \cref{pf84:C03} enumera cada entero exactamente una vez. Una
diagonal real sobre su dominio maximal es autoadjunta; como sus entradas
escapan a infinito fuera de conjuntos finitos, su resolvente es compacto. La
restricción positiva conserva estas propiedades.
\end{proof}

La prueba clásica omitida de inversión theta, continuación y ecuación
funcional se fusiona una sola vez con la realización triádica ya construida
en \cref{pf84:C05}. Allí se fijan literalmente
\(\Theta_B=1+2\Theta_+\) y \(Z=\mathcal Z_3\); por tanto, este bloque no
define un segundo origen de \(\zeta\).

Sea \(\mathbb P\) el conjunto de primos ordinarios,
\(\mathfrak h_{\mathbb P}:=\ell^2(\mathbb P)\) y
\[
 H_{\mathbb P}e_p=(\log p)e_p.
\]
En la Fock bosónica \(\Gamma_s(\mathfrak h_{\mathbb P})\), las ocupaciones
\((k_p)_p\) tienen soporte finito.

\begin{theorem}[Representación de Fock de la factorización]
\label{pf84:C08}
La aplicación sobre bases
\[
 U_F\lvert(k_p)_p\rangle=e_{\prod_pp^{k_p}}
\]
se extiende a un unitario
\[
 U_F:\Gamma_s(\mathfrak h_{\mathbb P})
 \longrightarrow\ell^2(\NN_{\ge1})
\]
y satisface
\[
 U_F\,d\Gamma(H_{\mathbb P})\,U_F^*=\log Q.
\]
\end{theorem}

\begin{proof}
La factorización única da una biyección entre ocupaciones finitas y enteros
positivos, luego una identificación unitaria de las bases ortonormales.
Además,
\[
 \sum_pk_p\log p=\log\!\left(\prod_pp^{k_p}\right),
\]
lo que prueba el entrelazamiento. La estructura de partida explícita es la
elección de los modos primos y de la Fock simétrica.
\end{proof}

Sobre
\[
 \mathcal H_{\mathrm{orb}}
 :=\ell^2\{(p,k,\varepsilon):
 p\in\mathbb P,\ k\ge1,\ \varepsilon=\pm1\},
\]
definimos
\[
 D_{\mathrm{orb}}e_{p,k,\varepsilon}
 =\varepsilon k\log p\,e_{p,k,\varepsilon},
 \qquad
 W_{\log}e_{p,k,\varepsilon}
 =(\log p)e_{p,k,\varepsilon}.
\]

\begin{proposition}[Traza orbital]
\label{pf84:C09}
El operador \(D_{\mathrm{orb}}\) es autoadjunto y de resolvente compacto. Si
\(f\in C_c(\RR)\) es acotada, entonces
\[
 \Tr\!\left(
 W_{\log}e^{-|D_{\mathrm{orb}}|/2}f(D_{\mathrm{orb}})
 \right)
 =
 \sum_{p,k}(\log p)p^{-k/2}
 \bigl[f(k\log p)+f(-k\log p)\bigr].
\]
\end{proposition}

\begin{proof}
La diagonal es real y sus longitudes escapan de todo compacto con
multiplicidad finita. El soporte compacto de \(f\) hace que el producto
regulado tenga rango finito; su traza es la suma de las entradas diagonales.
El índice \(p\) recorre aquí los primos ordinarios. El resultado realiza el
lado primo clásico, pero no identifica \(D_{\mathrm{orb}}\) con un operador
dual de los ceros.
\end{proof}

\section{Forma completada de Guinand--Weil}

Fijamos
\[
 \mathcal D_{\mathrm W}:=C_c^\infty(\RR),
 \qquad
 \widehat\psi(t)=\int_{\RR}\psi(x)e^{-itx}\,\dd x,
 \qquad
 \psi(x)=\frac1{2\pi}\int_{\RR}\widehat\psi(t)e^{itx}\,\dd t,
\]
y, para \(\psi,\phi\in\mathcal D_{\mathrm W}\),
\[
 h_{\psi,\phi}(z)
 :=\widehat\psi(z)\,
 \overline{\widehat\phi(\overline z)}.
\]
En la recta real \(h_{\psi,\psi}=|\widehat\psi|^2\).

\begin{definition}[Clase de prueba]
\label{pf84:E01}
La clase de prueba es \(\mathcal D_{\mathrm W}\), estable por derivación,
reflexión y traslación. Se utiliza la inversa
\[
 \check h(x)=\frac1{2\pi}\int_{\RR}h(t)e^{itx}\,\dd t.
\]
La suavidad y el soporte compacto proporcionan decaimiento rápido en bandas
y soporte compacto para \(\check h\), condiciones que garantizan la
convergencia de los bloques siguientes.
\end{definition}

Sea \(\psi_\Gamma=\Gamma'/\Gamma\). Definimos
\[
 \Irr_\#:=
 \left\{\operatorname{ev}_9(w):
 w\in\mathscr W_9^+,\ 
 P_{\Irr,\#}e_w=e_w\right\}.
\]
Por \cref{pf84:C07}, esta imagen coincide con los primos ordinarios; aquí se
usa únicamente para tipar el índice del bloque orbital, sin repetir aquella
prueba.
\begin{align}
 \mathcal P(h)&:=h(i/2)+h(-i/2),\\
 \mathcal G(h)&:=-\log\pi\,\check h(0)
 +\frac1{2\pi}\int_{\RR}h(t)
 \Re\psi_\Gamma\!\left(\frac14+\frac{it}{2}\right)\dd t,\\
 \mathcal O_\#(\check h)&:=
 \sum_{p\in\Irr_\#}\sum_{k\ge1}
 (\log p)p^{-k/2}
 \bigl[\check h(k\log p)+\check h(-k\log p)\bigr].
\end{align}

\begin{definition}[Funcional completado]
\label{pf84:E02}
Para \(h=h_{\psi,\phi}\), se fija
\[
 \mathcal W(h):=\mathcal P(h)+\mathcal G(h)
 -\mathcal O_\#(\check h).
\]
El bloque orbital procede del selector irreducible del
\cref{chap:producto-nativo}; los bloques polar y gamma pertenecen a la
completación clásica de zeta. Esta diferencia de procedencia se conserva en
la notación y en la prueba.
\end{definition}

\begin{theorem}[Fórmula explícita de Guinand--Weil]
\label{pf84:E03}
Para \(\psi,\phi\in C_c^\infty(\RR)\) y \(h=h_{\psi,\phi}\),
\[
 \lim_{T\to\infty}
 \sum_{\substack{\rho:\,\zeta(\rho)=0\\|\Im\rho|\le T}}
 h\!\left(\frac{\rho-1/2}{i}\right)
 =\mathcal W(h),
\]
donde los ceros no triviales se cuentan con multiplicidad y el límite se
toma con la simetrización clásica.
\end{theorem}

\begin{proof}
Paley--Wiener coloca \(h_{\psi,\phi}\) en la clase admisible de la fórmula
explícita. Los polos de
\(\Lambda_\zeta(s)=\pi^{-s/2}\Gamma(s/2)\zeta(s)\) producen
\(\mathcal P\), el factor arquimediano produce \(\mathcal G\) y el producto
de Euler produce la suma sobre potencias primas. La identificación de esta
última suma con \(\mathcal O_\#\) se realiza únicamente mediante
\cref{pf84:E04}; así no se duplica la prueba de primalidad nativa.
\end{proof}

\begin{proposition}[Reindexación del bloque orbital]
\label{pf84:E04}
Si los irreducibles del producto nativo coinciden exactamente con los primos
ordinarios y se conservan los pesos
\((\log p)p^{-k/2}\) y las dos orientaciones \(\pm k\log p\), entonces
\[
 \mathcal O_\#(\check h)=\mathcal O(\check h).
\]
\end{proposition}

\begin{proof}
La biyección entre irreducibles nativos y primos reindexa la suma sin
modificar índices, pesos ni argumentos. Un irreducible no primo, un primo
ausente o un cambio de peso falsaría la identidad.
\end{proof}

\begin{remark}[Separación meromorfa--entera]
\label{pf84:E05}
La función \(\Lambda_\zeta\) posee polos simples en \(0\) y \(1\), que
originan el bloque polar. En cambio,
\[
 \xi(s)=\frac12s(s-1)\Lambda_\zeta(s)
\]
es entera: los factores \(s(s-1)\) cancelan precisamente esos polos. No se
atribuyen, por tanto, polos en \(0\) o \(1\) a \(\xi\).
\end{remark}

\begin{definition}[Forma hermitiana y preoperador]
\label{pf84:E06}
Se definen
\[
 \mathcal B_{\mathrm W}(\psi,\phi)
 :=\mathcal W(h_{\psi,\phi}),
 \qquad
 \mathcal V_{\mathrm W}:=
 \mathcal D_{\mathrm W}/\operatorname{Rad}\mathcal B_{\mathrm W},
\]
y, sobre el cociente algebraico, la regla candidata
\[
 D_{\mathrm W,0}[\psi]:=[-i\psi'].
\]
No se presupone que \(\mathcal B_{\mathrm W}\) sea positiva ni que
\(\mathcal V_{\mathrm W}\) sea ya un espacio de Hilbert. La regla sólo queda
bien definida después de probar la invariancia del radical.
\end{definition}

\begin{theorem}[Simetría algebraica e invariancias]
\label{pf84:E07}
El radical es invariante bajo \(-i\,\dd/\dd x\);
\(D_{\mathrm W,0}\) está bien definido y es
\(\mathcal B_{\mathrm W}\)-simétrico. Las traslaciones
\(T_a\psi(x)=\psi(x-a)\) preservan la forma y la reflexión
\(J\psi(x)=\psi(-x)\) satisface
\[
 JD_{\mathrm W,0}J=-D_{\mathrm W,0}.
\]
\end{theorem}

\begin{proof}
La convención de Fourier da
\[
 h_{-i\psi',\phi}(z)=h_{\psi,-i\phi'}(z).
\]
La identidad vale antes de aplicar cualquiera de los tres bloques de
\(\mathcal W\), de modo que la derivada preserva el radical y el
preoperador es simétrico. Las fases opuestas de las traslaciones se cancelan
en \(h\); la reflexión cambia el signo de la derivada y deja invariantes los
bloques polar, gamma y orbital. Esta simetría algebraica no equivale todavía
a autoadjunción en un Hilbert positivo.
\end{proof}

\section{Geometría local de la recta crítica}

\begin{lemma}[Coordenada de Möbius de la recta crítica]
\label{pf84:RH-01}
Para \(s\in\CC\setminus\{0,1\}\), sea
\[
 \tau(s):=i\frac{s-1}{s}.
\]
Con \(F(s)=1-s\) y \(C(s)=\overline s\),
\[
 \tau(Fs)=-\frac1{\tau(s)},
 \qquad
 \tau(Cs)=-\overline{\tau(s)},
\]
y
\[
 \tau(Fs)=\tau(Cs)
 \iff |\tau(s)|=1
 \iff \Re s=\frac12.
\]
\end{lemma}

\begin{proof}
Las dos identidades se obtienen por sustitución directa. Su igualdad
equivale a \(\tau(s)^{-1}=\overline{\tau(s)}\), es decir,
\(|\tau(s)|=1\); desarrollar esta última condición da
\(\Re s=1/2\). Para los ceros de zeta se requieren además las simetrías
funcional y conjugada. La equivalencia parametriza la recta; no sitúa por sí
sola los ceros en ella.
\end{proof}

\begin{lemma}[Signo local de Hermite--Biehler]
\label{pf84:RH-04}
Sea \(A\) una función real entera con un cero real simple \(a\). Para
\(z=a+iy\), \(y>0\) suficientemente pequeño,
\[
|A(z)-iA'(z)|^2-|A(z)+iA'(z)|^2
 =-4A'(a)^2y+O(y^2)<0.
\]
La convención local asociada a \(A\) es
\[
 E_A:=A+iA'.
\]
Sólo al especializar \(A=\Xi\) y suponer además que \(a\) es un cero real
simple de \(\Xi\) se obtiene
\[
 E_\Xi=\Xi+i\Xi'.
\]
\end{lemma}

\begin{proof}
Se desarrolla \(A\) y \(A'\) en Taylor alrededor de \(a\), usando
\(A(a)=0\) y \(A'(a)\ne0\), y se retienen los términos de primer orden en
\(y\). El resultado es estrictamente local: no demuestra que \(E_A\) posea la
propiedad global de Hermite--Biehler.
\end{proof}

\begin{proposition}[Obstrucción por leyes de conteo]
\label{pf84:E14}
Sea
\[
 N_A(T):=\Tr\mathbf1_{[-T,T]}(A).
\]
Los operadores \(Q\), \(B\), \(D_{\mathrm{orb}}\) y \(\mathcal R\) no
poseen, por sí solos, la ley de conteo exigida a un eventual operador dual
de los ceros.
\end{proposition}

\begin{proof}
Para \(T\ge1\),
\[
 N_Q(T)=\lfloor T\rfloor,
 \qquad
 N_B(T)=2\lfloor T\rfloor+1.
\]
El teorema de los números primos da
\[
 N_{D_{\mathrm{orb}}}(T)
 =2\sum_{p\le e^T}\left\lfloor\frac{T}{\log p}\right\rfloor
 \ge2\pi(e^T)\sim\frac{2e^T}{T}.
\]
Finalmente, \(N_{\mathcal R}(T)\) es constante en cada intervalo
multiplicativo \([3^m,3^{m+1})\), pues \(\mathcal R\) sólo tiene niveles
\(3^{m+1}\) con multiplicidades finitas.

Para
\(\Xi(z)=\xi(1/2+iz)\), la ecuación funcional implica
\(\Xi(-z)=\Xi(z)\); los ceros \(\pm\gamma\) tienen la misma multiplicidad y
\[
 N_{\mathrm{sym}}(T)=2N_+(T)+O(1).
\]
En consecuencia, la convención bilateral requiere
\[
 N_{\mathrm{sym}}(T)\sim\frac{T}{\pi}\log T,
\]
mientras la unilateral equivalente, bajo esta simetría y emparejamiento de
multiplicidades, requiere
\[
 N_+(T)\sim\frac{T}{2\pi}\log T.
\]
Ninguna de las cuatro funciones anteriores tiene esa asintótica.
\end{proof}

\section{Laplacianos sobre los ciclos de orden \texorpdfstring{$3^m$}{potencias de tres}}

Para $m\geq 1$, sean

\[
 N_m=3^m,
 \qquad
 C_m=\ZZ/N_m\ZZ,
\]

y considérese el laplaciano combinatorio

\[
 (\Delta_m f)(j)=2f(j)-f(j+1)-f(j-1),
 \qquad f\in\ell^2(C_m).
\]

Los caracteres

\[
 e_{m,n}(j)=N_m^{-1/2}
 \exp\!\left(\frac{2\pi i n j}{N_m}\right),
 \qquad n\in C_m,
\]

forman una base ortonormal y satisfacen

\[
 \Delta_m e_{m,n}
 =4\sin^2\!\left(\frac{\pi n}{N_m}\right)e_{m,n}.
\]

Como $N_m$ es impar, cada clase de $C_m$ posee un único representante
$\widetilde n$ con

\[
 -\frac{N_m-1}{2}\leq\widetilde n\leq\frac{N_m-1}{2}.
\]

Definimos el operador autoadjunto $D_m$ por

\begin{equation}
 D_m e_{m,n}
 =\frac{N_m}{\sqrt\pi}
 \sin\!\left(\frac{\pi\widetilde n}{N_m}\right)e_{m,n}.
 \label{eq:operador-espectral-Dm}
\end{equation}

Entonces

\begin{equation}
 D_m^2=\frac{N_m^2}{4\pi}\,\Delta_m.
 \label{eq:Dm-laplaciano}
\end{equation}

Sea $P_0$ la proyección ortogonal sobre el carácter constante. Para toda
función par acotada $F$ sobre el espectro de $D_m$, introducimos la traza
positiva

\begin{equation}
 \operatorname{Tr}^{+}_m\!\bigl(F(D_m)\bigr)
 :=\frac12\operatorname{Tr}
 \!\left(F(D_m)P_0^{\perp}\right).
 \label{eq:traza-positiva}
\end{equation}

El factor $1/2$ selecciona una contribución por cada par espectral
$\{n,-n\}$; la eliminación de $P_0$ excluye el autovalor nulo.

\begin{theorem}[Límite theta de las trazas de calor]
\label{thm:limite-theta-triadico}
Para todo $x>0$,
\[
 \lim_{m\to\infty}
 \operatorname{Tr}^{+}_m\!\left(e^{-xD_m^2}\right)
 =\Theta_+(x)
 :=\sum_{n\geq1}e^{-\pi n^2x}.
\]
\end{theorem}

\begin{proof}
La traza bilateral, una vez retirado el modo constante, es
\[
 \sum_{\substack{n\in C_m\\n\neq0}}
 \exp\!\left[
 -\frac{N_m^2x}{\pi}
 \sin^2\!\left(\frac{\pi\widetilde n}{N_m}\right)
 \right].
\]
Para cada entero fijo $n$, el sumando converge a $e^{-\pi n^2x}$. Además,
para $|\widetilde n|\leq N_m/2$, la desigualdad
$\sin y\geq 2y/\pi$, válida en $0\leq y\leq\pi/2$, proporciona una
mayorante gaussiana sumable e independiente de $m$. La convergencia
dominada aplicada a la suma discreta implica
\[
 \lim_{m\to\infty}
 \operatorname{Tr}\!\left(e^{-xD_m^2}P_0^\perp\right)
 =\sum_{n\in\ZZ\setminus\{0\}}e^{-\pi n^2x}.
\]
La división entre dos establecida en \eqref{eq:traza-positiva} conserva una
contribución de cada par $\{n,-n\}$ y da $\Theta_+(x)$.
\end{proof}

\section{Transformada theta--Mellin}

Para $\Re s>1$, definimos

\begin{equation}
 \mathcal Z_{3}(s)
 :=
 \frac{\displaystyle
 \int_0^\infty\Theta_+(x)x^{s/2-1}\,\dd x}
 {\pi^{-s/2}\Gamma(s/2)}.
 \label{eq:funcional-theta-mellin}
\end{equation}

La convergencia absoluta permite integrar término a término. Como

\[
 \int_0^\infty e^{-\pi n^2x}x^{s/2-1}\,\dd x
 =\Gamma(s/2)(\pi n^2)^{-s/2},
\]

se obtiene

\begin{equation}
 \mathcal Z_3(s)=\sum_{n\geq1}n^{-s}=\zeta(s),
 \qquad \Re s>1.
 \label{eq:Z3-zeta}
\end{equation}

La igualdad se prolonga meromórficamente por la continuación de la función
zeta de Riemann \cite{Apostol1976}. En lo sucesivo, todas las evaluaciones
proceden de \eqref{eq:funcional-theta-mellin}, de sus derivadas o de la
descomposición residual de su serie de Dirichlet.

\begin{theorem}[Inversión theta y ecuación funcional]
\label{pf84:C05}
Con la notación operatorial precedente,
\[
 \Theta_B(x):=\Tr(e^{-\pi xB^2})
 =1+2\Theta_+(x),
 \qquad
 Z(s):=\Tr(Q^{-s})=\mathcal Z_3(s)
 \quad(\Re s>1).
\]
La sumación de Poisson y la transformada de Mellin producen
\[
 \Theta_B(x)=x^{-1/2}\Theta_B(1/x)
\]
y la continuación de
\[
 \xi(s)=\frac12s(s-1)\pi^{-s/2}\Gamma(s/2)Z(s),
 \qquad
 \xi(s)=\xi(1-s).
\]
\end{theorem}

\begin{proof}
La igualdad \(\Theta_B=1+2\Theta_+\) separa el modo cero y empareja
\(\pm n\); la identidad \(Z=\mathcal Z_3\) es
\eqref{eq:Z3-zeta}. La sumación de Poisson aplicada a la gaussiana da la
inversión de \(\Theta_B\). En el semiplano de convergencia, la integración
término a término proporciona
\[
 \pi^{-s/2}\Gamma(s/2)Z(s)
 =\int_0^\infty x^{s/2-1}\Theta_+(x)\,\dd x.
\]
Al dividir la integral en \((0,1)\) y \((1,\infty)\), aplicar la inversión
theta al primer tramo y separar los dos polos clásicos, se obtiene la
continuación y la ecuación funcional. Esta única prueba clásica incorporada
no localiza los ceros de \(\zeta\).
\end{proof}

\section{Descomposición residual módulo tres}

Para $r\in\{0,1,2\}$, sea

\begin{equation}
 Z_r(s)=
 \sum_{\substack{n\geq1\\n\equiv r\pmod3}}n^{-s}.
 \label{eq:canales-residuales-Zr}
\end{equation}

En términos de la zeta de Hurwitz,

\[
 Z_0(s)=3^{-s}\zeta(s),
 \qquad
 Z_1(s)=3^{-s}\zeta\!\left(s,\frac13\right),
 \qquad
 Z_2(s)=3^{-s}\zeta\!\left(s,\frac23\right).
\]

Los tres términos tienen residuo $1/3$ en $s=1$. Escribimos sus
desarrollos de Laurent en la forma

\begin{equation}
 Z_r(1+z)-\frac1{3z}=\sum_{n\geq0}c_{r,n}\,z^n,
 \qquad
 \zeta(1+z)-\frac1z=\sum_{n\geq0}a_n\,z^n.
 \label{eq:coeficientes-Laurent-residuales}
\end{equation}

\begin{theorem}[Descomposición de los coeficientes de Laurent]
\label{thm:descomposicion-laurent-triadica}
Para todo $n\geq0$,
\[
 a_n=c_{0,n}+c_{1,n}+c_{2,n},
 \qquad
 \gamma_n=(-1)^n n!\,a_n.
\]
Si $\chi_{-3}$ denota el carácter real no principal módulo tres, entonces
\[
 \frac{L^{(n)}(1,\chi_{-3})}{n!}=c_{1,n}-c_{2,n}.
\]
Además, con \(\lambda=\log3\),
\[
 c_{0,n}=\frac13\left[
 \frac{(-\lambda)^{n+1}}{(n+1)!}
 +\sum_{k=0}^{n}a_k\frac{(-\lambda)^{n-k}}{(n-k)!}
 \right].
\]
Estas tres identidades determinan de manera única
\((c_{0,n},c_{1,n},c_{2,n})\) a partir de
\((a_n,L^{(n)}(1,\chi_{-3}))\) y de los coeficientes agregados de orden
menor.
\end{theorem}

\begin{proof}
La partición de los enteros positivos por clases residuales proporciona
$\zeta=Z_0+Z_1+Z_2$. Al sustraer los polos en $s=1$ y comparar los
coeficientes de \eqref{eq:coeficientes-Laurent-residuales} se obtiene la
primera identidad. La segunda coincide con la convención usual para las
constantes de Stieltjes. Además,
$Z_1-Z_2=L(s,\chi_{-3})$; los polos de $Z_1$ y $Z_2$ se cancelan y la
comparación de coeficientes da la tercera relación. Por último,
\[
 Z_0(1+z)=\frac13e^{-\lambda z}\zeta(1+z).
\]
El producto de las dos series de Laurent proporciona la fórmula explícita
para \(c_{0,n}\). La suma y la diferencia recuperan entonces
\(c_{1,n}\) y \(c_{2,n}\), lo que prueba la unicidad afirmada.
\end{proof}

La suma de los tres términos recupera el sector agregado; la diferencia de
las clases $1$ y $2$ conserva la orientación del carácter; la comparación
con la clase divisible por tres registra el cambio de escala. Estas tres
combinaciones son linealmente independientes y determinan los tres
coeficientes $c_{r,n}$.

\section{Euler--Mascheroni y constantes de Stieltjes}

Sea

\[
 C_r:=\operatorname*{FP}_{s=1}Z_r(s).
\]

Las identidades clásicas para la función digamma dan
\cite{Apostol1976,Lagarias2013}

\begin{equation}
 C_0=\frac{\gamma-\log3}{3},
 \label{eq:C0-parte-finita}
\end{equation}

\begin{equation}
 C_1=\frac\gamma3+\frac{\log3}{6}+\frac{\pi}{6\sqrt3},
 \qquad
 C_2=\frac\gamma3+\frac{\log3}{6}-\frac{\pi}{6\sqrt3}.
 \label{eq:C1-C2-partes-finitas}
\end{equation}

\begin{corollary}
\label{cor:gamma-log3-L1}
Las partes finitas satisfacen
\[
 \boxed{\gamma=C_0+C_1+C_2},
\]
\[
 \boxed{L(1,\chi_{-3})=C_1-C_2=\frac{\pi}{3\sqrt3}},
\]
\[
 \boxed{\log3=C_1+C_2-2C_0}.
\]
\end{corollary}

\begin{proof}
En orden cero, el \cref{thm:descomposicion-laurent-triadica} da
$a_0=c_{0,0}+c_{1,0}+c_{2,0}$. Por definición,
$a_0=\gamma$ y $c_{r,0}=C_r$, lo que prueba la primera igualdad. La
diferencia de las dos clases orientadas es
$C_1-C_2=L(1,\chi_{-3})$. Sustituyendo
\eqref{eq:C0-parte-finita}--\eqref{eq:C1-C2-partes-finitas} se obtienen las
otras dos identidades.
\end{proof}

Una única extracción de coeficientes mediante la fórmula integral de
Cauchy aplicada a \eqref{eq:coeficientes-Laurent-residuales} determina
simultáneamente $\gamma_0,\ldots,\gamma_6$. Las discretizaciones con 96 y
192 nodos coinciden hasta orden $10^{-87}$; la identidad de los
coeficientes, no obstante, procede del
\cref{thm:descomposicion-laurent-triadica} y no de esa comparación numérica.

\section{Evaluación de la constante de Apéry}

\begin{proposition}
\label{prop:evaluacion-apery}
El funcional \eqref{eq:funcional-theta-mellin} satisface
\[
 \boxed{
 \mathcal Z_3(3)=\zeta(3)
 =1.202056903159594285399738161511449\ldots}.
\]
\end{proposition}

\begin{proof}
La igualdad exacta es la especialización $s=3$ de \eqref{eq:Z3-zeta}.
Para certificar cualquier prefijo decimal sin modificar el funcional,
considérese $S_N=\sum_{n=1}^{N}n^{-3}$. El criterio integral da
\[
 S_N+\frac1{2(N+1)^2}
 <\zeta(3)<
 S_N+\frac1{2N^2}.
\]
Los intervalos racionales anidados resultantes determinan las cifras
impresas.
\end{proof}

La irracionalidad de $\zeta(3)$ es el teorema clásico de Apéry
\cite{Apery1979}; no se deduce de la mera evaluación del funcional.

\section{Constante de Glaisher--Kinkelin}

La derivada espectral en $s=-1$ define

\begin{equation}
 \boxed{
 A_{\mathrm{GK}}
 =\exp\!\left(\frac1{12}-\mathcal Z_3'(-1)\right)}.
 \label{eq:glaisher-kinkelin-espectral}
\end{equation}

Esta es la fórmula del determinante regularizado asociado al producto
hiperfactorial \cite{Vardi1988}. Una derivación central de cinco puntos,
evaluada a dos escalas, proporciona

\[
 A_{\mathrm{GK}}
 =1.2824271291006226368753425688697917\ldots.
\]

El término $1/12$ pertenece a la regularización del hiperfactorial y no se
introduce como parámetro de ajuste.

\section{Constante de Euler--Kronecker del campo de Eisenstein}

El carácter $\chi_{-3}$ determina el campo cuadrático imaginario

\[
 F_{-3}=\QQ(\sqrt{-3}),
 \qquad
 \zeta_{F_{-3}}(s)=\zeta(s)L(s,\chi_{-3}).
\]

Su constante de Euler--Kronecker es

\begin{equation}
 \boxed{
 \gamma_{F_{-3}}
 =\gamma+\frac{L'(1,\chi_{-3})}{L(1,\chi_{-3})}}.
 \label{eq:euler-kronecker-eisenstein}
\end{equation}

Los coeficientes residuales de
\eqref{eq:coeficientes-Laurent-residuales} dan

\[
 L(1,\chi_{-3})
 =0.6045997880780726168646927525473852\ldots,
\]
\[
 L'(1,\chi_{-3})
 =0.2226629869686015094866602627647443\ldots,
\]

y, por \eqref{eq:euler-kronecker-eisenstein},

\[
 \gamma_{F_{-3}}
 =0.9454972808716807032397499941581890\ldots.
\]

La elección de $F_{-3}$ procede del carácter residual módulo tres; no de
una búsqueda sobre los valores decimales de constantes de campos
cuadráticos.

\section{Producto local asociado a los primos gemelos}

Para un patrón finito $H=\{h_1,\ldots,h_k\}\subset\ZZ$, definimos

\[
 \nu_p(H):=\left|\{-h_1,\ldots,-h_k\}\bmod p\right|,
 \qquad
 \sigma_p(H)=
 \frac{1-\nu_p(H)/p}{(1-1/p)^k}.
\]

Para $H=\{0,2\}$, los factores locales conducen a

\begin{equation}
 C_2^{\mathrm{twin}}
 =\prod_{p>2}\frac{p(p-2)}{(p-1)^2}.
 \label{eq:constante-primos-gemelos}
\end{equation}

Sea \(P_X\) el producto restringido a los primos \(2<p\leq X\). Para
\(p>X\), todos los factores pertenecen a \((0,1)\); al ampliar el producto
de la cola desde los primos a todos los enteros se obtiene la cota
telescópica
\[
 \prod_{n>X}\left(1-\frac1{(n-1)^2}\right)
 =\frac{X-1}{X}
 \leq
 \prod_{p>X}\left(1-\frac1{(p-1)^2}\right)<1.
\]
El producto \(P_{2\times10^6}\), evaluado con redondeo decimal dirigido, y
esta desigualdad racional para la cola dan

\begin{equation}
 0.6601615071224244826403933551
 <C_2^{\mathrm{twin}}<
 0.6601618372035081248537566591.
 \label{eq:intervalo-primos-gemelos}
\end{equation}

El recuento directo hasta $10^6$ contiene 8169 pares de primos gemelos. La
convergencia del producto local y este censo finito no implican la
existencia de infinitos pares; esta última es una proposición aritmética
distinta.

\section{Constante de Meissel--Mertens y Hamiltoniano primo}

La constante de Meissel--Mertens se define por

\begin{equation}
 B_1=\lim_{x\to\infty}
 \left(\sum_{p\leq x}\frac1p-\log\log x\right)
 \label{eq:meissel-mertens-def}
\end{equation}

y satisface

\begin{equation}
 B_1=\gamma+
 \sum_p\left(\log\!\left(1-\frac1p\right)+\frac1p\right).
 \label{eq:meissel-mertens-euler}
\end{equation}

Con la notación de \cref{pf84:C08}, ponemos
\[
 \mathcal F_P:=\Gamma_s(\mathfrak h_{\mathbb P}),
 \qquad
 H_P:=d\Gamma(H_{\mathbb P}).
\]
No se redefine aquí la Fock bosónica de modos primos ni su Hamiltoniano. El
operador ya construido tiene función de partición

\begin{equation}
 \operatorname{Tr}_{\mathcal F_P}(e^{-sH_P})
 =\prod_p\frac1{1-p^{-s}}
 =\zeta(s),
 \qquad \Re s>1.
 \label{eq:hamiltoniano-primo-zeta}
\end{equation}

Sea \(\mathcal H_P^{(1)}\) el subespacio de una partícula, generado por
\(\{|p\rangle:p\text{ primo}\}\). Entonces
\[
 \operatorname{Tr}_{\mathcal H_P^{(1)}}\!\left(
 e^{-H_P}\mathbf 1_{(-\infty,\log x]}(H_P)\right)
 =\sum_{p\leq x}\frac1p.
\]
Por tanto, \(B_1\) es la parte finita, cuando \(x\to\infty\), de esta traza
espectral truncada después de sustraer \(\log\log x\). La igualdad
\eqref{eq:hamiltoniano-primo-zeta} presupone la descomposición de Fock sobre
los primos; por ello constituye una representación operatorial del producto
de Euler, no una deducción de la primalidad a partir del laplaciano
triádico.

\section{Dependencia respecto del operador}

Las evaluaciones anteriores comparten el operador
\eqref{eq:operador-espectral-Dm}, la traza positiva
\eqref{eq:traza-positiva} y la transformada
\eqref{eq:funcional-theta-mellin}. Otras constantes requieren dinámicas
distintas. La constante de Catalan está asociada al carácter módulo cuatro
$\chi_{-4}$, mientras que la constante de Khinchin pertenece al operador de
Gauss y a su medida invariante. Un prefijo de cualquiera de ellas puede
aparecer en un catálogo finito sin que el funcional triádico la seleccione.
Esta dependencia operatorial distingue la evaluación espectral de la
capacidad puramente combinatoria de codificación.

\end{HMTScientificOwnerSubordinated}

\HMTFlushCurrentChapter
\HMTSuccessorChapterInput
  {Criticidad unimodal de la familia dodecafásica: perfil analítico exacto, aproximantes de Feigenbaum y certificado hiperbólico de renormalización}
  {../colaboracion/parte_iii/tex/capitulos_finales/45_feigenbaum.tex}

\HMTFlushCurrentChapter
\HMTSuccessorChapterInput
  {Aritmética de la familia modular dodecafásica de periodo \(30\): identidades de Rogers--Ramanujan, normalizador \(899\) y levantamiento nonádico de Pell}
  {../colaboracion/parte_iii/tex/capitulos_finales/46_modo_30_pell.tex}
\begin{HMTSpiralChapterBody}
% Corte mecánico íntegro: parte_ii.tex:323--419.
% Residencia material del ensamblaje sucesor de 102 capítulos.
% Residencia causal: capítulo 46.

\chapter{Tornillo espectral de Pell y suspensión logarítmica}

\section{Ley de tornillo}

El corpus contiene la identidad exacta
\begin{equation}
 V^{-1}\mathscr S V=(2+\sqrt3)\ii\,\mathscr S.
 \label{espiral:eq:tornillo-pell}
\end{equation}
Sea \(\varepsilon=2+\sqrt3\). La órbita transversal
\[
 z_n=z_0(\varepsilon\ii)^n
\]
satisface
\[
 r_n=r_0\varepsilon^n,
 \qquad
 \theta_n=\theta_0+\frac{\pi n}{2}.
\]
Eliminando \(n\),
\begin{equation}
 r(\theta)=r_0
 \exp\!\left[
 \frac{2\log(2+\sqrt3)}{\pi}
 (\theta-\theta_0)
 \right].
 \label{espiral:eq:espiral-log-pell}
\end{equation}
La espiral logarítmica es una salida del operador: no se eligió su ecuación
para ajustar después \(\varepsilon\).

\begin{proof}[Derivación de \cref{espiral:eq:espiral-log-pell}]
De \(\theta_n-\theta_0=n\pi/2\) se sigue
\(n=2(\theta_n-\theta_0)/\pi\). Sustituyendo en
\(r_n=r_0\varepsilon^n\),
\[
 r_n=r_0\exp\left(
 \frac{2\log\varepsilon}{\pi}(\theta_n-\theta_0)
 \right).
\]
La interpolación continua es la ecuación indicada.
\end{proof}

\section{Involución conjugada de contracción}

La extensión bilateral \(n\in\Z\) produce
\[
 n\to+\infty:\ r_n\to\infty,
 \qquad
 n\to-\infty:\ r_n\to0.
\]
La inversión \(n\mapsto-n\) intercambia expansión y contracción. Como
\(\varepsilon^{-1}=2-\sqrt3\), ambas ramas pertenecen a la misma unidad de
Pell. No son dos espirales independientes, sino dos orientaciones del mismo
tornillo espectral.

\begin{figure}[htbp]
\centering
\begin{tikzpicture}[scale=.82]
 \draw[->,HMTGray] (-5.1,0)--(5.1,0) node[right] {$x$};
 \draw[->,HMTGray] (0,-3.1)--(0,3.1) node[above] {$y$};
 \draw[HMTBlue,very thick,domain=-7:5.2,samples=240,smooth,variable=\t]
   plot ({.27*exp(.22*\t)*cos(deg(\t))},{.27*exp(.22*\t)*sin(deg(\t))});
 \draw[HMTRed,thick,domain=-7:5.2,samples=240,smooth,variable=\t]
   plot ({.27*exp(.22*\t)*cos(deg(-\t))},{.27*exp(.22*\t)*sin(deg(-\t))});
 \foreach \k in {-4,-3,...,4}{
   \pgfmathsetmacro{\ang}{1.57079632679*\k}
   \fill[HMTNavy]
     ({.27*exp(.22*\ang)*cos(deg(\ang))},
      {.27*exp(.22*\ang)*sin(deg(\ang))}) circle (.8pt);
 }
 \node[HMTBlue,font=\small\sffamily] at (3.5,2.4) {expansión};
 \node[HMTRed,font=\small\sffamily] at (3.5,-2.4) {orientación conjugada};
\end{tikzpicture}
\caption{Proyección transversal del tornillo de Pell. Las marcas representan cuartos de giro.}
\label{espiral:fig:tornillo-pell}
\end{figure}

\section{Covariancia de escala y conservación energética}

La dilatación \(\varepsilon\) no debe interpretarse como crecimiento
gratuito de la amplitud energética de una onda en vacío. Su realización
coherente actúa sobre la escala espacial o espectral:
\[
 (r,\varphi,\lambda)
 \longmapsto
 (\varepsilon r,\varphi+\pi/2,\varepsilon\lambda).
\]
Esta distinción será esencial en la realización electromagnética de la
\cref{espiral:chap:realizacion-em}.

\begin{resultbox}[title={Compatibilidad exacta entre el tornillo espectral de Pell y la suspensión fase--memoria}]
La elevación fase--memoria mínima es una hélice circular. El tornillo
espectral combina un cuarto de giro y una dilatación de Pell; su proyección
es una espiral logarítmica. Ambos proceden del registro HMT, pero no son la
misma periodicidad ni se obtienen el uno del otro por olvido de tipos.
\end{resultbox}

\end{HMTSpiralChapterBody}

\HMTFlushCurrentChapter
\HMTSuccessorChapterInput
  {Teoría aritmética de las vacancias nonádicas: defecto \texorpdfstring{\(10^3-3^6\)}{10^3-3^6}, red radical \(3\to6\to9\), ley potencia--vacancia y torsión de borde}
  {../colaboracion/parte_iii/tex/capitulos_finales/47_vacancias_nonadicas.tex}

\HMTFlushCurrentChapter
\HMTSuccessorChapterInput
  {Constantes térmicas y radiación: Boltzmann, Wien, Stefan--Boltzmann y termodinámica espectral del gas fotónico}
  {../colaboracion/parte_iii/tex/capitulos_finales/48_constantes_termicas.tex}

\HMTFlushCurrentChapter
\HMTSuccessorChapterInput
  {Constantes electrodébiles y mezcla: determinante angular, relación \(W/Z\), ángulo de Weinberg y reflexión racional CKM}
  {../colaboracion/parte_iii/tex/capitulos_finales/49_sector_electrodebil.tex}

\HMTFlushCurrentChapter
\HMTSuccessorChapterInput
  {Derivación intrínseca de \(G\) y reciprocidad contragrediente \texorpdfstring{\(G\leftrightarrow\hbar\)}{G<->hbar}: cierre de fase, escalas de Compton y Planck e invariancia de \(G\hbar\)}
  {../colaboracion/parte_iii/tex/capitulos_finales/50_gravitacion_accion.tex}
% Secciones 1--3 del delta gravitatorio íntegro.
\section{Dominio genealógico de la realización gravitatoria}
\label{sec:l9s102:gravedad-dominio-genealogico}

La realización gravitatoria recibe como objeto de partida una sección del
continuo discreto ya construida por la composición
\begin{equation}
 \mathrm{APP}\longrightarrow\mathrm{TRIT}\longrightarrow\mathrm{TPK}
 \longrightarrow X_{\mathrm{TPK}}^{\mathrm{enr}}
 \longrightarrow\mathcal C_{\mathrm{cont}}^{\mathrm{disc}}.
 \label{eq:l9s102:gravedad-cadena-genealogica}
\end{equation}
APP conserva las evaluaciones aditiva y multiplicativa mediante residuo,
cociente y acarreo; TRIT determina el régimen y la orientación; el TPK
incorpora el selector trítico de fase sin eliminar ninguna de las dos hojas
APP, efectúa el transporte cardinal, conserva la memoria y realiza el retorno
de fase sin reiniciar el estado. El transductor
dimensional actúa después sobre la sección superviviente del continuo:
\begin{equation}
 \mathfrak G_{108}:\mathcal C_{\mathrm{cont}}^{\mathrm{disc}}
 \times\mathcal L_S^+
 \longrightarrow\mathcal G_{\mathrm{HMT\text{-}MD}},
 \qquad
 (x,\hbar_a)\longmapsto
 \bigl(G_a,R_{108},m_{108,a},\ell_{P,a}^{,2}\bigr).
 \label{eq:l9s102:gravedad-mapa-realizacion}
\end{equation}
Los cardinales \(90/120\), la periodicidad \(108\), la acción y el par
angular entran con su genealogía conservada. Ninguna identificación
gravitatoria posterior selecciona retroactivamente estos antecedentes.

\section{Transductor gravitatorio y selector cronológico de periodicidad}
\label{sec:l9s102:gravedad-transductor}

Sean \(c_{\mathrm{int}}>0\), \(t_0>0\),
\(\ell_0=c_{\mathrm{int}}t_0\) y
\(\hbar_a\in\mathcal L_S^+\), con
\(a\in\{\mathrm{pre},\mathrm{ret}\}\). El transductor dimensional es
\begin{equation}
 G_a(\theta)
 =\theta\frac{c_{\mathrm{int}}^5t_0^2}{\hbar_a}
 =\theta\frac{c_{\mathrm{int}}^3\ell_0^2}{\hbar_a},
 \qquad \theta>0.
 \label{eq:l9s102:g-transductor}
\end{equation}
La cronologia ciclica de \(108\) estados determina
\begin{equation}
 T_{108}=108t_0,
 \qquad
 \omega_{108}=\frac{2\pi_{\rm HMT}}{108t_0},
 \qquad
 R_{108}=\frac{c_{\mathrm{int}}}{\omega_{108}}
 =\frac{54}{\pi_{\rm HMT}}\ell_0.
 \label{eq:l9s102:g-reloj108}
\end{equation}
El cuanto cronológico y su masa asociada son
\begin{equation}
 E_{108,a}=\hbar_a\omega_{108},
 \qquad
 m_{108,a}=\frac{E_{108,a}}{c_{\mathrm{int}}^2}.
 \label{eq:l9s102:g-cuanto-masa}
\end{equation}
La longitud de Compton reducida satisface
\begin{equation}
 \bar\lambda_{C,a}
 =\frac{\hbar_a}{m_{108,a}c_{\mathrm{int}}}=R_{108}.
 \label{eq:l9s102:g-compton}
\end{equation}

\begin{theorem}[Unicidad del selector gravitatorio de periodicidad]
\label{thm:l9s102:g-selector}
Con la convencion reducida \(r_g=Gm/c_{\mathrm{int}}^2\), la condición
\(r_{g,a}(\theta)=R_{108}\) posee la solución positiva única
\begin{equation}
 \boxed{
 \theta_{108}^{\mathrm{circ}}=\left(\frac{54}{\pi_{\rm HMT}}\right)^2.}
 \label{eq:l9s102:g-theta}
\end{equation}
La realización correspondiente identifica estructuralmente
\begin{equation}
 R_{108}=\bar\lambda_{C,a}=r_{g,a}=\ell_P^{\mathrm{circ}}.
 \label{eq:l9s102:g-cuatro-longitudes}
\end{equation}
\end{theorem}

\begin{proof}
Por \cref{eq:l9s102:g-transductor,eq:l9s102:g-cuanto-masa},
\[
 r_{g,a}(\theta)
 =\frac{G_a(\theta)m_{108,a}}{c_{\mathrm{int}}^2}
 =\theta\frac{\pi_{\rm HMT}}{54}\ell_0.
\]
La igualdad con \(R_{108}=(54/\pi_{\rm HMT})\ell_0\) impone
\(\theta=(54/\pi_{\rm HMT})^2\). La linealidad en \(\theta\) y la positividad de
los factores determinan la unicidad.
\end{proof}

En la carta en la que \(c_{\mathrm{int}}=(\mu\varepsilon)^{-1/2}\), el
transductor adopta la forma
\begin{equation}
 \boxed{
 G_a=\left(\frac{54}{\pi_{\rm HMT}}\right)^2
 \frac{t_0^2}{\hbar_a(\mu\varepsilon)^{5/2}}.}
 \label{eq:l9s102:g-mu-epsilon}
\end{equation}
El producto constitutivo \(\mu\varepsilon\) gobierna el cono causal y la
escala gravitatoria; el cociente \(Z^2=\mu/\varepsilon\) conserva la
orientación de impedancia.

\section{Covariancia bidireccional de gravedad y acción}
\label{sec:l9s102:g-hbar}

Definimos
\begin{equation}
 R_{\mathrm{act}}=\frac{\hbar_{\mathrm{pre}}}
 {\hbar_{\mathrm{ret}}}.
 \label{eq:l9s102:r-act}
\end{equation}
La masa cronológica y el acoplamiento gravitatorio transforman
contragredientemente:
\begin{equation}
 \frac{m_{108,\mathrm{pre}}}{m_{108,\mathrm{ret}}}=R_{\mathrm{act}},
 \qquad
 \frac{G_{\mathrm{pre}}}{G_{\mathrm{ret}}}=R_{\mathrm{act}}^{-1}.
 \label{eq:l9s102:g-hbar-reciprocidad}
\end{equation}
Su producto permanece fijo:
\begin{equation}
 \boxed{
 G_a\hbar_a
 =\theta_{108}^{\mathrm{circ}}c_{\mathrm{int}}^5t_0^2,
 \qquad
 \ell_P^2=\frac{G_a\hbar_a}{c_{\mathrm{int}}^3}
 =\theta_{108}^{\mathrm{circ}}\ell_0^2.}
 \label{eq:l9s102:g-hbar-invariante}
\end{equation}
La acción conserva la orientación de la hoja; el área de Planck conserva la
geometría común de ambas orientaciones.

\begin{theorem}[Independencia de carta de las secciones \(h\) y \(G\)]
\label{thm:l9s102:covariancia-unidades-g}
Sean \(u_L,u_T,u_S\) bases positivas de las rectas dimensionales de
longitud, tiempo y acción, y cámbiense por
\[
 u_L'=\lambda_Lu_L,\qquad
 u_T'=\lambda_Tu_T,\qquad
 u_S'=\lambda_Su_S,
 \qquad \lambda_L,\lambda_T,\lambda_S>0.
\]
Las coordenadas de las mismas secciones verifican
\begin{equation}
 c_{\rm int}'=\frac{\lambda_T}{\lambda_L}c_{\rm int},\qquad
 t_0'=\frac{t_0}{\lambda_T},\qquad
 \hbar_a'=\frac{\hbar_a}{\lambda_S},\qquad
 G_a'=\frac{\lambda_S\lambda_T^3}{\lambda_L^5}G_a.
 \label{eq:l9s102:cambio-carta-g}
\end{equation}
Para la base inducida
\(u_G=u_L^5u_T^{-3}u_S^{-1}\) se obtiene
\begin{equation}
 \hbar_a'u_S'=\hbar_au_S,
 \qquad G_a'u_G'=G_au_G.
 \label{eq:l9s102:secciones-h-g-invariantes}
\end{equation}
Por consiguiente,
\(h_a=2\pi_{\rm HMT}\hbar_a\),
\(G_a=\theta_{108}^{\rm circ}c_{\rm int}^5t_0^2/\hbar_a\) y
\(G_a\hbar_a=\theta_{108}^{\rm circ}c_{\rm int}^5t_0^2\)
son identidades de secciones, no igualdades dependientes de una elección de
unidades. La carta metrológica publica sus coordenadas numéricas después de
la construcción HMT--MD.
\end{theorem}

\begin{proof}
Las tres primeras transformaciones de
\eqref{eq:l9s102:cambio-carta-g} son la ley contragrediente de coordenadas
en \(L\otimes T^{-1}\), \(T\) y \(S\). Su sustitución en
\eqref{eq:l9s102:g-transductor} produce la cuarta. La base de la recta
gravitatoria se transforma por el factor inverso, lo que prueba
\eqref{eq:l9s102:secciones-h-g-invariantes}. Los escalares
\(\pi_{\rm HMT}\) y \(\theta_{108}^{\rm circ}\) son adimensionales y
permanecen invariantes.
\end{proof}


\HMTFlushCurrentChapter
\HMTSuccessorChapterInput
  {Derivación HMT del parámetro de Barbero--Immirzi y del área mínima: incidencia hexada--octada, canales \(90/120\) y espectro del operador de área}
  {../colaboracion/parte_iii/tex/capitulos_finales/51_barbero_area.tex}
% Secciones 4--7 del delta gravitatorio íntegro.
\section{Incidencia hexada--octada, parámetro de Barbero--Immirzi y espectro de área}
\label{sec:l9s102:barbero-area}

La derivación areal recibe como antecedentes la inclusión de incidencia
\(6\to8\), los cardinales TPK \(90/120\), el defecto normalizado
\(-1/12\) y los canales espectrales \(q_\pm\). El funcional
hexada--octada determina el parámetro de Barbero--Immirzi y la
normalización del operador de área. Esta residencia contiene la derivación
HMT y el espectro; la realización Palatini--Cartan--Holst ocupa una residencia
posterior en Mecánica Dimensional.

Sean \(H\subset O\), con \(|H|=6\) y \(|O|=8\), y consérvese el par
marcado de la hexada en la inclusión
\begin{equation}
 \iota_{6,8}:\mathcal F_6=H\times\binom H2
 \lhook\joinrel\longrightarrow
 \mathcal F_8=O\times\binom H2.
 \label{eq:l9s102:barbero-inclusion-incidencia}
\end{equation}
La cuenta interna da
\begin{equation}
 |\mathcal F_6|=6\binom62=90,
 \qquad |\mathcal F_8|=8\binom62=120,
 \qquad
 \frac{90-120}{24\binom62}=-\frac1{12}.
 \label{eq:l9s102:barbero-conteos}
\end{equation}
Una vez publicados \(\pi_{\rm HMT}\), el ángulo \(A\) y el
contraángulo único \(C^*\), los dos canales son
\begin{equation}
 q_\pm=\exp\!\left[-\frac{\pi_{\rm HMT}}{180}(A\pm C^*)\right],
 \qquad 0<q_+<q_-<1,
 \label{eq:l9s102:barbero-qpm}
\end{equation}
y sus series de retorno orientado son
\begin{equation}
 S_{90}(q)=\frac{q^{90}}{1-q^{270}},\qquad
 S_{120}(q)=\frac{q^{120}}{1-q^{360}},\qquad
 \Delta S_m=S_m(q_-)-S_m(q_+).
 \label{eq:l9s102:barbero-series}
\end{equation}

El orden causal de los coeficientes se fija en el registro temporal del TPK,
antes de evaluar el coeficiente de polo. Sea \(C_{12}=\mathbb Z/12\mathbb Z\) la rueda
dodecafásica y sea \(\partial_{\rm ret}\) su única costura orientada de retorno.
En el módulo libre de sucesos de incidencia, la vuelta fundamental es
\begin{equation}
 c_{12}^{+}
 =\sum_{j\in C_{12}}[j,\mathcal F_6]
  -[\partial_{\rm ret},\mathcal F_8].
 \label{eq:l9s102:barbero-cadena-dodecafasica}
\end{equation}
El signo del segundo término es el signo de frontera. La involución
bidireccional invierte la orientación completa,
\(c_{12}^{-}=-c_{12}^{+}\), sin modificar las multiplicidades absolutas de
los dos tipos de incidencia. Defínase el lector de retorno sobre los
generadores por
\begin{equation}
 \mathscr R([j,\mathcal F_6])=S_{90},
 \qquad
 \mathscr R([\partial_{\rm ret},\mathcal F_8])=S_{120}.
 \label{eq:l9s102:barbero-lector-retorno}
\end{equation}
La covariancia de fase hace constante la primera imagen sobre los doce
sectores; la segunda actúa únicamente en la costura ampliada. Por linealidad,
\begin{equation}
 \boxed{\mathscr R(c_{12}^{+})=12S_{90}-S_{120}.}
 \label{eq:l9s102:barbero-peso-generado}
\end{equation}
Así, el peso \(12:-1\) es la imagen de la vuelta fundamental: doce sectores y
una frontera orientada. No procede de imponer previamente \(1/24\), de una
comparación decimal ni del valor convencional de Barbero--Immirzi.

\Needspace{12\baselineskip}
\begin{theorem}[Determinación dodecafásica del funcional de Barbero--Immirzi]
\label{thm:l9s102:barbero-funcional}
La cadena fundamental \eqref{eq:l9s102:barbero-cadena-dodecafasica} determina
el par orientado \((a,b)=(12,1)\). El coeficiente de \((1-q)^{-1}\) de su imagen es
entonces \(1/24\), como salida exacta del funcional ya construido. Por tanto,
después de evaluar los canales conjugados \(q_\pm\), la salida adimensional es
\begin{equation}
 \boxed{
 \gamma_{\rm HMT}=\Gamma_{6\to8}
 =\frac{\pi_{\rm HMT}AC^*}{180}
   +12\Delta S_{90}-\Delta S_{120}.}
 \label{eq:l9s102:barbero-gamma}
\end{equation}
Esta fórmula recibe el par angular ya generado y no utiliza como entrada
el valor convencional del parámetro de Barbero--Immirzi.
\end{theorem}

\begin{proof}
La suma de \eqref{eq:l9s102:barbero-cadena-dodecafasica} contiene exactamente
los doce sectores de \(C_{12}\), mientras la costura de retorno aparece una
sola vez y con orientación negativa. El lector
\eqref{eq:l9s102:barbero-lector-retorno} produce por ello
\eqref{eq:l9s102:barbero-peso-generado}, antes de todo cálculo de polo. Para
cada \(m>0\),
\[
 p_1(S_m)
 =\lim_{q\to1}(1-q)\frac{q^m}{1-q^{3m}}
 =\frac1{3m}.
\]
En consecuencia,
\[
 p_1\bigl(\mathscr R(c_{12}^{+})\bigr)
 =\frac{12}{270}-\frac1{360}
 =\frac{48-3}{1080}=\frac1{24}.
\]
El coeficiente \(p_1\) está referido a \((1-q)^{-1}\).
En la coordenada compleja \(q-1\),
\[
 \operatorname*{Res}_{q=1}S_m(q)=-\frac1{3m},
 \qquad
 \operatorname*{Res}_{q=1}\mathscr R(c_{12}^{+})(q)=-\frac1{24}.
\]
La ecuación diofántica posterior
\[
 \frac a{270}-\frac b{360}=\frac1{24}
 \quad\Longleftrightarrow\quad 4a-3b=45
\]
posee las soluciones enteras positivas
\((a,b)=(12+3t,1+4t)\), \(t\ge0\); la única solución mínima es
\((12,1)\). Esta cuenta certifica aritméticamente el par que la órbita
fundamental ya había generado: cualquier \(t>0\) altera la multiplicidad de
los sectores o de la costura y deja de representar una vuelta del TPK.
Finalmente, sustituir \(q_\pm\) en la imagen orientada y añadir la componente
angular previamente generada prueba \eqref{eq:l9s102:barbero-gamma}.
\end{proof}

La secuencia causal queda fijada sin intercambio de papeles:
\begin{equation}
\begin{aligned}
 (\mathcal F_6\hookrightarrow\mathcal F_8,C_{12},\partial_{\rm ret})
 &\longrightarrow (90,120,-1/12,c_{12}^{+})\\
 &\longrightarrow \mathscr R(c_{12}^{+})=12S_{90}-S_{120},\\
 \mathscr R(c_{12}^{+})
 &\longrightarrow p_1\bigl(\mathscr R(c_{12}^{+})\bigr)=\frac1{24},\\
 \bigl(\mathscr R(c_{12}^{+}),A,C^*,q_\pm\bigr)
 &\longrightarrow \Gamma_{6\to8}=\gamma_{\rm HMT}
 \longrightarrow \widehat{\mathcal A}_{\rm phys}.
\end{aligned}
 \label{eq:l9s102:barbero-cadena-completa}
\end{equation}
El par angular y sus canales proceden de la construcción anterior.
La evaluación del coeficiente de polo constituye una comprobación
del funcional de retorno; la evaluación conjugada emplea las
coordenadas angulares previamente generadas.
La realización posterior en la acción de Holst reconoce la misma sección
adimensional; no interviene en su selección.

Para los dos sectores de borde, sean
\begin{equation}
 N_m=\sum_{e\in\partial A_m}N_e,
 \qquad m\in\{90,120\},
 \qquad \operatorname{Spec}(N_e)=\{0,1,2\}.
 \label{eq:l9s102:n90-n120}
\end{equation}
El operador físico de área se escribe
\begin{equation}
 \boxed{
 \frac{\widehat{\mathcal A}_{\mathrm{phys}}}{\ell_P^2}
 =\gamma_{\mathrm{HMT}}a_0(N_{90}+N_{120}).}
 \label{eq:l9s102:area-operador}
\end{equation}
Sobre \(36\) enlaces qutrit,
\begin{equation}
 \operatorname{Spec}(N_{90}+N_{120})=\{0,1,\ldots,72\},
 \label{eq:l9s102:area-numero-spectrum}
\end{equation}
y la multiplicidad del autovalor \(n\) es el coeficiente de \(z^n\) en
\((1+z+z^2)^{36}\). De aqui resulta
\begin{equation}
 \operatorname{Spec}\!\left(
 \frac{\widehat{\mathcal A}_{\mathrm{phys}}}{\ell_P^2}\right)
 =\{n\gamma_{\mathrm{HMT}}a_0:0\le n\le72\}.
 \label{eq:l9s102:area-spectrum}
\end{equation}

\section{Geometría total y memoria de procedencia}
\label{sec:l9s102:area-hamiltoniano}

El Hamiltoniano modular del borde es
\begin{equation}
 \boxed{
 K_A=-\log\rho_A
 =cI+\Delta_{\mathrm{mod}}(90N_{90}+120N_{120}).}
 \label{eq:l9s102:hamiltoniano-modular}
\end{equation}
Sean
\begin{equation}
 N=N_{90}+N_{120},
 \qquad
 D=90N_{90}+120N_{120}.
 \label{eq:l9s102:n-d}
\end{equation}
La matriz de paso
\(\left(\begin{smallmatrix}1&1\\90&120\end{smallmatrix}\right)\)
posee determinante \(30\), y su inversa reconstruye los sectores:
\begin{equation}
 \boxed{
 N_{90}=4N-\frac{D}{30},
 \qquad
 N_{120}=\frac{D}{30}-3N.}
 \label{eq:l9s102:reconstruccion-sectores}
\end{equation}
El área publica el número total de excitaciones; el Hamiltoniano modular
publica su sector de procedencia. El par conjunto reconstruye la ocupación
completa de la frontera.

\section{Estado bicapa e información orientada}
\label{sec:l9s102:estado-bicapa}

Sea \(R=R^*=R^{-1}\) la involución de hojas y
\(y=C^*_{\mathrm{rad}}\). El estado normalizado es
\begin{equation}
 \boxed{
 \rho_y=\frac{e^{-yR}}{2\cosh y}.}
 \label{eq:l9s102:rho-y}
\end{equation}
Su Hamiltoniano modular y su entropía son
\begin{equation}
 K_y=-\log\rho_y=\log(2\cosh y)I+yR,
 \label{eq:l9s102:ky}
\end{equation}
\begin{equation}
 S(\rho_y)=\log(2\cosh y)-y\tanh y.
 \label{eq:l9s102:entropia-y}
\end{equation}
La inversión \(y\mapsto-y\) deja fija la entropía y produce la distancia
informacional orientada
\begin{equation}
 \boxed{
 D(\rho_y\Vert\rho_{-y})=2y\tanh y.}
 \label{eq:l9s102:entropia-relativa-hojas}
\end{equation}
La componente par mide la distribucion espectral; la componente impar mide
el coste informacional de invertir la orientación.

\section{Cuanto informacional de energía y celda gravitatoria de media acción}
\label{sec:l9s102:horizonte-media-accion}

La periodicidad de \(108\) estados determina
\begin{equation}
 t_P=\frac{54}{\pi}t_0,
 \qquad
 E_P=\frac{\hbar}{t_P}=\frac{h}{108t_0}.
 \label{eq:l9s102:ep-tp}
\end{equation}
La realización térmica del reloj tritico satisface
\begin{equation}
 \boxed{
 E_P=k_B\Theta_{\mathrm{clk}}\log3.}
 \label{eq:l9s102:planck-trit}
\end{equation}
Para el horizonte de radio \(R\),
\begin{equation}
 E_\Lambda(R)=\frac{c^4R}{2G},
 \qquad
 E_\Lambda(R)=T_H(R)S_H(R).
 \label{eq:l9s102:energia-horizonte}
\end{equation}
En la celda de Planck,
\begin{equation}
 \boxed{
 E_\Lambda
 =T_HS_H
 =\frac{E_P}{2}
 =\frac{k_B\Theta_{\mathrm{clk}}\log3}{2}.}
 \label{eq:l9s102:confluencia-horizonte}
\end{equation}
Por tanto,
\begin{equation}
 \boxed{
 E_\Lambda t_P=\frac{\hbar}{2},
 \qquad
 E_\Lambda T_{108}=\frac h2.}
 \label{eq:l9s102:media-accion}
\end{equation}
La energía de vacío, la entropía areal, la temperatura cronológica y la
acción comparten una misma celda de confluencia con dominios tipados.

Cuando \(S_H/k_B=\pi\), la multiplicidad efectiva satisface
\begin{equation}
 \Omega_{\mathrm{eff}}=e^{S_H/k_B}=e^\pi.
 \label{eq:l9s102:omega-e-pi}
\end{equation}
El operador hiperbólico de Euler posee
\begin{equation}
 \operatorname{Spec}(e^{\pi H})=\{e^\pi,e^{-\pi}\}.
 \label{eq:l9s102:e-pi-spectrum}
\end{equation}
La igualdad de los escalares establece la confluencia espectral; el
entrelazador entre multiplicidad microscópica de horizonte y dinamica de
Perron pertenece a su residencia física especifica.


\HMTFlushCurrentChapter
% Capítulo 52 del sucesor: integración pública del propietario íntegro del radión.
% No carga portada, main, style ni C_estatutos_procedencia del módulo autónomo.

% Adaptador de compatibilidad del propietario íntegro del radión.
%
% El tratado rector ya carga tipografía, geometría, hipervínculos, teoremas,
% tablas, TikZ y tcolorbox.  Este archivo aporta únicamente los nombres y
% entornos exclusivos del propietario, sin importar su preámbulo autónomo ni
% alterar la jerarquía tipográfica general.

\makeatletter
\@ifundefined{color@RadNavy}{\definecolor{RadNavy}{HTML}{102A43}}{}
\@ifundefined{color@RadBlue}{\definecolor{RadBlue}{HTML}{1F5E7A}}{}
\@ifundefined{color@RadTeal}{\definecolor{RadTeal}{HTML}{168C8C}}{}
\@ifundefined{color@RadCopper}{\definecolor{RadCopper}{HTML}{B66A3C}}{}
\@ifundefined{color@RadGold}{\definecolor{RadGold}{HTML}{D4A73A}}{}
\@ifundefined{color@RadInk}{\definecolor{RadInk}{HTML}{1A202C}}{}
\@ifundefined{color@RadMist}{\definecolor{RadMist}{HTML}{EDF4F7}}{}
\@ifundefined{color@RadCream}{\definecolor{RadCream}{HTML}{FAF7F0}}{}
\@ifundefined{color@RadRose}{\definecolor{RadRose}{HTML}{8D3F55}}{}
\@ifundefined{color@RadGreen}{\definecolor{RadGreen}{HTML}{2F7D5D}}{}
\@ifundefined{color@RadGray}{\definecolor{RadGray}{HTML}{66788A}}{}

\providecommand{\TRIT}{\ensuremath{\mathrm{TRIT}}}
\providecommand{\HMT}{\ensuremath{\mathrm{HMT}}}
\providecommand{\Rstate}{\mathcal R_{12}}
\providecommand{\epsR}{\varepsilon_R}
\providecommand{\epsone}{\varepsilon_R^{(1)}}
\providecommand{\pih}{\pi_{\!\HMT}}
\providecommand{\alphah}{\alpha_{\!\HMT}}
\providecommand{\hret}{\hbar_{\mathrm{ret}}}
\providecommand{\hfull}{h_{\mathrm{ret}}}
\providecommand{\diff}{\mathop{}\!\mathrm d}
\providecommand{\e}{\mathrm e}
\providecommand{\R}{\mathbb R}
\providecommand{\C}{\mathbb C}
\providecommand{\F}{\mathbb F}
\providecommand{\Z}{\mathbb Z}
\providecommand{\dd}{\,\mathrm d}
\providecommand{\status}[1]{\textsf{\bfseries #1}}
\providecommand{\sourcepath}[1]{\nolinkurl{#1}}

\@ifundefined{typebox}{%
  \newtcolorbox{typebox}[1][]{enhanced,breakable,colback=white,
    colframe=RadBlue,boxrule=.7pt,arc=1.5mm,left=4mm,right=4mm,
    top=3mm,bottom=3mm,title={Recibo de tipos},fonttitle=\bfseries,
    coltitle=RadNavy,#1}%
}{}
\@ifundefined{narrativebox}{%
  \newtcolorbox{narrativebox}[1][]{enhanced,breakable,colback=white,
    colframe=RadGold,borderline west={2.4pt}{0pt}{RadGold},
    boxrule=.25pt,arc=0mm,left=5mm,right=4mm,top=3mm,bottom=3mm,#1}%
}{}

% Las once portadillas internas del propietario autónomo se convierten en
% divisiones científicas ordinarias dentro del capítulo 52.  Se evita así
% introducir partes autónomas, gigantismo tipográfico o páginas vacías.
\providecommand{\partopener}[3]{\section{#2}\noindent\emph{#3}\par\medskip}

\providecommand{\BigRadionEquation}{%
  \begin{equation*}
  \boxed{\displaystyle
  \frac{180^{\circ}}{\pih}
  =50^{\circ}+A^{\circ}-\frac{20}{81}\,\Delta_4^{\circ}+\epsR^{\circ}}
  \end{equation*}}
\makeatother


\chapter{El radión angular HMT--MD: empalme con acarreo de las secciones directa y torsional y cierre exacto de la unidad de fase, acción y memoria}
\label{chap:sucesor-102-radion}

\section{Composición APP--TRIT--TPK del estado de fase: generación de la unidad orientada con memoria}
% Fragmento público generado desde propietarios íntegros.
% Residencia: 52.1.
% Fuente: ../incoming/radion_monografia_20260827/manuscrito/sections/01_resumen_y_guia.tex:1-128
\paragraph{Síntesis matemática del objeto y de su cadena causal}

Esta monografía reconstruye el radián desde la cadena generativa
\[
\APP\longrightarrow\TRIT\longrightarrow\TPK
\longrightarrow X_{\mathrm{enr}}
\longrightarrow\text{estructura discreta del continuo}
\longrightarrow\MD,
\]
sin tomar como entrada ni el valor metrológico de \(180/\pi\), ni el pequeño
residuo angular que se desea explicar. El resultado es una unidad enriquecida,
el \emph{radión}, definida como la unidad orientada que recorre un radián de
fase y barre exactamente una acción \(\hret\) en la elipse positiva asociada
al mismo estado.

El cierre rector es
\BigRadionEquation
con
\[
S_E=2\pih\left(\frac5{36}+\frac{25}{9}\alphah\right),
\qquad
\Delta_4=
\frac{\pih-e_{\HMT}\log\pih}{270}
\left(1+\frac{\pih}{729}\right),
\]
\[
\rho_{\mathrm{tw}}=\frac{2\cdot90}{729}=\frac{20}{81},
\qquad
\epsone=\frac{20}{81}\Delta_4-(S_E-1).
\]
La división APP prueba \(1<S_E<2\), de modo que el cociente es exactamente
uno y el resto es \(D_E=S_E-1\). La igualdad
\[
1=S_E-\frac{20}{81}\Delta_4+\epsone
\]
es entonces una identidad interna. La carta angular posterior produce
\[
\epsR^\circ
=5.13830167585752820716851646226459474899085744\ldots
\times10^{-8}\,{}^\circ.
\]

El texto demuestra además que la elipse de acción
\[
Q=a_S\cos\theta,\qquad P=b_S\sin\theta,\qquad
a_Sb_S=2\hret
\]
satisface
\[
\operatorname{Area}(0,\theta)=\hret\theta.
\]
Por tanto, un radión barre \(\hret\) y la vuelta \(2\pi\) encierra
\(h_{\mathrm{ret}}=2\pi\hret\). La misma elipse tiene borde simpléctico
finito e interior de área infinita cuando se equipa con la métrica
Hilbert--Beltrami--Klein. Ésta es la forma exacta de la tesis física que guía
la obra: la superficie visible cierra la acción; la profundidad interior
conserva un desarrollo no agotable por la proyección plana.

La exposición mantiene tipadas cuatro realizaciones del mismo estado:

\begin{enumerate}[label=\textbf{\arabic*.}]
\item el círculo \(S^1\), que publica la fase;
\item la elipse positiva, que publica la acción y la anisotropía;
\item el interior de Klein, que mide la profundidad mediante razón doble;
\item la hélice nonádica, que conserva retorno, memoria y contracción.
\end{enumerate}

No son cuatro geometrías superpuestas sin criterio. Son cuatro lectores con
dominio, codominio y dato conservado diferentes, aplicados a un único estado
enriquecido. El sector \(90/120\), la dodecafase, el acarreo del radión, la cola
de Lambert, el núcleo de Barbero y las torres globales se presentan asimismo
como objetos emparentados pero no intercambiables.

\begin{thesisbox}
El radión es la unidad orientada que barre \(\hret\) en la elipse de acción.
El círculo publica su fase, Klein mide su profundidad, la hélice conserva su
memoria y el acarreo coinductivo empalma exactamente sus secciones directa y
torsional.
\end{thesisbox}

\paragraph{Orden expositivo y dependencias de lectura}

El orden no es ornamental. Las Partes I y II producen el objeto y su cierre
desde APP--TRIT--TPK. Sólo después se introducen las realizaciones en lenguaje
de Hilbert, LC, Maxwell, Minkowski, Hodge, Lorentz, modularidad y órbitas. Cada
capítulo distingue tres niveles:

\begin{description}[style=nextline,leftmargin=4.2cm,labelwidth=3.9cm]
\item[Resultado interno.] Igualdad o construcción demostrada en el dominio
HMT con sus primitivas declaradas.
\item[Formalización reunida.] Mapa nuevo que articula resultados ya existentes
sin usar el blanco como selector.
\item[Interfaz física.] Traducción dimensional posterior que reconoce una
estructura convencional sin hacerla origen de la selección.
\end{description}

Los recuadros de \emph{estatuto y falsador} indican qué observación concreta
invalidaría cada promoción fuerte. Esta disciplina no debilita la tesis; la
hace auditable. El objetivo de la narración extensa es que ninguna palabra
como \emph{círculo}, \emph{elipse}, \emph{torsión}, \emph{carga} o
\emph{memoria} oculte un cambio de tipo.

\paragraph{Resultados rectores y alcance probatorio}

\begin{enumerate}[label=\textbf{C\arabic*.}]
\item Cierre independiente del valor final del radión por dos publicaciones del mismo estado y
resta con acarreo APP.
\item Derivación del coeficiente \(20/81\) como dos hojas del sector activo
\(90\), normalizadas por la fibra \(3^6=729\).
\item Igualador tipado entre la costura \(\Re s=1/2\), la cara APP de cien
marcas y la fase dodecafásica \(\theta_2=50^\circ\).
\item Teorema de área del radión: \(\operatorname{Area}(\theta)=\hret\theta\).
\item Estudio focal exacto de la elipse y de sus invariantes euclídeos,
hiperbólicos, simplécticos y modulares.
\item Teorema borde--interior: acción total finita \(h\) e interior de Klein
de área hiperbólica infinita.
\item Realización exacta mediante covarianza mínima y oscilador LC, seguida de
la publicación electromagnética \(90/120\).
\item Atlas de no-identificaciones para todos los objetos \(90/120\), incluida
la separación entre \(\epsR\), cola espectral y núcleo Barbero \(12:-1\).
\item Caracterización estructural de \(\pi\) en la carta del radión, distinta
del generador coinductivo primario.
\item Síntesis ontológica: dimensión como rango mínimo de memorias que una
proyección de fase ya no puede conservar.
\end{enumerate}
% Fuente: ../incoming/radion_monografia_20260827/manuscrito/sections/01b_mapa_resultados.tex:1-133
\paragraph{Correspondencia entre resultados, tipos y residencias causales}

La correspondencia siguiente ordena las afirmaciones rectoras y sus demostraciones,
manteniendo visible la jerarquía que impide que
la riqueza física borre el cierre matemático que la sostiene.

% HMT_RESULT_BEGIN:RDN-01-GENEALOGIA:theorem
\begin{exactbox}[title={Genealogía APP--TRIT--TPK del radión}]
El radión nace de la composición efectiva
\(\APP\to\TRIT\to\TPK\to X_{\mathrm{enr}}\to\MD\); no se obtiene
redefiniendo retrospectivamente el radián convencional.
\end{exactbox}
% HMT_RESULT_END:RDN-01-GENEALOGIA:theorem

% HMT_RESULT_BEGIN:RDN-02-CIERRE:theorem
\begin{exactbox}[title={Cierre exacto e independencia prospectiva}]
El acarreo APP produce internamente
\[
1=S_E-\frac{20}{81}\Delta_4+\varepsilon_R^{(1)},
\qquad
\varepsilon_R^{(1)}=\frac{20}{81}\Delta_4-(S_E-1),
\]
y la carta angular transporta esta misma identidad a
\(180^\circ/\pi_{\!\HMT}\) sin introducir el residuo objetivo.
\end{exactbox}
% HMT_RESULT_END:RDN-02-CIERRE:theorem

% HMT_RESULT_BEGIN:RDN-03-BISAGRA-50:theorem
\begin{exactbox}[title={Coordenada crítica de \(50^\circ\)}]
La coordenada crítica \(\Re s=1/2\), publicada en la cara APP de cien
marcas y en la rueda dodecafásica, selecciona la fase
\(\theta_2=50^\circ\). Es una igualdad de lectores tipados del mismo estado.
\end{exactbox}
% HMT_RESULT_END:RDN-03-BISAGRA-50:theorem

% HMT_RESULT_BEGIN:RDN-04-COEFICIENTE-2081:theorem
\begin{exactbox}[title={Densidad bidireccional nonádica \(20/81\)}]
El coeficiente torsional no se ajusta en la red \(1/81\): se construye como
\(\rho_{\mathrm{tw}}=2N_6/3^6=2\cdot90/729=20/81\), conservando hoja,
orientación y normalización de la fibra nonádica.
\end{exactbox}
% HMT_RESULT_END:RDN-04-COEFICIENTE-2081:theorem

% HMT_RESULT_BEGIN:RDN-05-AREA-ACCION:theorem
\begin{exactbox}[title={Ley areal del radión}]
Para la elipse positiva
\(Q=a_S\cos\theta\), \(P=b_S\sin\theta\),
\(a_Sb_S=2\hbar_{\mathrm{ret}}\), el área orientada barrida es
\(\mathcal A(\theta)=\hbar_{\mathrm{ret}}\theta\). Un radión barre exactamente
\(\hbar_{\mathrm{ret}}\), y la vuelta \(2\pi\) encierra
\(h_{\mathrm{ret}}=2\pi\hbar_{\mathrm{ret}}\).
\end{exactbox}
% HMT_RESULT_END:RDN-05-AREA-ACCION:theorem

% HMT_RESULT_BEGIN:RDN-06-INVARIANTES-FOCALES:theorem
\begin{exactbox}[title={Invariantes focales del operador positivo}]
El par publicado \(A\pm C^\ast\) fija los semiejes, la excentricidad, la
separación focal y el parámetro de compresión simpléctica. En particular,
\((d_1/d_0)^2=C^\ast/2\), identidad adimensional que separa forma de escala.
\end{exactbox}
% HMT_RESULT_END:RDN-06-INVARIANTES-FOCALES:theorem

% HMT_RESULT_BEGIN:RDN-07-BORDE-INTERIOR:theorem
\begin{exactbox}[title={Finitud simpléctica del borde y divergencia del área hiperbólica interior}]
El borde simpléctico de la elipse encierra la acción finita
\(h_{\mathrm{ret}}\), mientras el área Hilbert--Beltrami--Klein del interior
diverge al alcanzar la frontera. Fase visible y profundidad no son la misma
medida.
\end{exactbox}
% HMT_RESULT_END:RDN-07-BORDE-INTERIOR:theorem

% HMT_RESULT_BEGIN:RDN-08-HELICE-MEMORIA:theorem
\begin{exactbox}[title={Elevación helicoidal nonádica con retorno de fase y avance de memoria}]
La monodromía nonádica devuelve la fase sin reiniciar el estado. La hélice
levanta el círculo conservando acarreo y memoria; el cierre \(2\pi/4\pi\)
distingue retorno proyectivo, inversión de orientación y retorno completo.
\end{exactbox}
% HMT_RESULT_END:RDN-08-HELICE-MEMORIA:theorem

% HMT_RESULT_BEGIN:RDN-09-OSCILADOR-EM:development
\begin{narrativebox}[title={Realización oscilatoria y constitutiva electromagnética}]
La misma anisotropía se realiza mediante una covarianza mínima y un oscilador
LC: una coordenada almacena el aspecto eléctrico-elástico y su conjugada el
aspecto magnético-rotacional. La frecuencia y la impedancia quedan tipadas por
separado.
\end{narrativebox}
% HMT_RESULT_END:RDN-09-OSCILADOR-EM:development

% HMT_RESULT_BEGIN:RDN-10-TIPOS-90120:theorem
\begin{exactbox}[title={Separación tipada de los operadores \(90/120\)}]
Extractor dodecafásico, acarreo del radión, cola de Lambert, funcional
hexada--octada, núcleo de Barbero y semigrupo global de torres son objetos
distintos. El peso \(12:-1\) procede de los doce sectores de la cadena
fundamental dodecafásica y de su única costura orientada. El coeficiente de
\((1-q)^{-1}\), igual a \(1/24\), y el contratérmino entero mínimo verifican
posteriormente el par producido.
\end{exactbox}
% HMT_RESULT_END:RDN-10-TIPOS-90120:theorem

% HMT_RESULT_BEGIN:RDN-11-GEOMETRIAS-CAMPO:development
\begin{narrativebox}[title={Realizaciones compleja, lorentziana y dual de Hodge}]
Los regímenes elíptico, parabólico e hiperbólico del TRIT se publican como
rotación, umbral y transformación hiperbólica. La estrella de Hodge, las constitutivas y la fuerza
de Lorentz actúan después sobre el estado ya orientado; no seleccionan sus
constantes.
\end{narrativebox}
% HMT_RESULT_END:RDN-11-GEOMETRIAS-CAMPO:development

% HMT_RESULT_BEGIN:RDN-12-MODULAR-ORBITAL:development
\begin{narrativebox}[title={Caracteres modular, orbital y aritmético de frontera}]
La involución modular, las órbitas elípticas, el avance-retardo y la aritmética
3--6--9 reconocen distintas memorias del mismo estado. La función modular y
las vacancias se usan como lectores posteriores, nunca como sustitutos del
generador APP--TRIT--TPK.
\end{narrativebox}
% HMT_RESULT_END:RDN-12-MODULAR-ORBITAL:development

% HMT_RESULT_BEGIN:RDN-13-SINTESIS-ONTOLOGICA:conclusion
\begin{thesisbox}[title={Síntesis operatoria de fase, acción, profundidad y memoria}]
El círculo es la proyección de fase; la elipse es la sección de acción; Klein
mide la profundidad; la hélice conserva la historia. El radión es la unidad
orientada que mantiene acopladas esas publicaciones sin borrar el residuo que
las distingue.
\end{thesisbox}
% HMT_RESULT_END:RDN-13-SINTESIS-ONTOLOGICA:conclusion

% HMT_RESULT_BEGIN:RDN-14-CONTROL-ALCANCE:control
\begin{statusbox}[title={Delimitación de dominios y continuidad causal}]
Las leyes completas de masas, el refinamiento electrónico, CKM/CP y el
corredor excepcional conservan sus demostraciones rectoras propias. Esta monografía los
cita sólo donde comparten antecedente o carta; no traslada al radión sus
selectores específicos ni confunde sus residencias probatorias.
\end{statusbox}
% HMT_RESULT_END:RDN-14-CONTROL-ALCANCE:control


\subsection{Soporte bivariado de APP, orientación ternaria del TRIT y actualización del estado enriquecido mediante el TPK}
% Fragmento público generado desde propietarios íntegros.
% Residencia: 52.1.1.
% Fuente: ../incoming/radion_monografia_20260827/manuscrito/sections/02_genealogia.tex:1-180
\noindent La genealogía operatoria antecede a la medida angular y determina las variables conservadas por cada lector.\par\medskip

\subsubsection{Soporte APP--TRIT--TPK y construcción del estado enriquecido}

\paragraph{Antecedencia del estado enriquecido respecto de la publicación circular}

La definición escolar de radián comienza con un círculo ya dado, un centro ya
dado y una longitud de arco ya mensurable. Si el radio y el arco tienen la
misma longitud, el ángulo subtendido mide un radián. La definición es correcta
dentro de la geometría euclídea, pero no responde a una pregunta anterior:
¿qué estructura hace posible que una vuelta posea fase, orientación, acción y
memoria, y por qué la unidad natural se relaciona con \(2\pi\)?

HMT invierte el orden causal. Primero construye el estado finito, después su
prolongación coinductiva, luego la vuelta y sólo al final su publicación
angular. El radión no corrige el radián convencional; lo levanta a su dominio
genealógico.

\begin{typebox}
\[
\APP\to\TRIT\to\TPK\to X_{\mathrm{enr}}
\to\mathcal C^{\mathrm{disc}}_{\mathrm{cont}}
\to\bigl(\pih,e_{\HMT},\varphi_{\HMT},\alphah,A,C^*,\Delta_4,\hret\bigr)
\to\mathfrak r_\sigma.
\]
La matemática convencional sólo aparece después como lenguaje de realización:
\(S^1\), elipse simpléctica, métrica de Klein, espacio de Hilbert, circuito
LC, pantalla de Minkowski y ecuaciones de campo.
\end{typebox}

\paragraph{Soporte bivariado de APP: hojas, residuos, cocientes y acarreo}

La Plataforma de Aritmética Posicional (APP) conserva simultáneamente dos
realizaciones de los símbolos \(1,\dots,9\): una hoja aditiva y una hoja
multiplicativa. No se reduce un número a su valor publicado. Junto al valor
permanecen hoja, orientación, cociente, resto, acarreo, frontera, ruta y
supervivencia.

El censo básico tiene tres estratos que jamás deben identificarse por mera
coincidencia cardinal:
\[
104,976\longrightarrow468\;U_6\longrightarrow243\;w_6
\subset\F_3^6,\qquad |\F_3^6|=729.
\]
Las \(468\) emisiones decimales, las \(243\) palabras visibles y el ambiente
de \(729\) palabras son objetos distintos. El ambiente será decisivo para la
normalización \(20/81\).

En base cúbica \(B=1000\), la resta orientada de dos palabras de tríadas se
realiza de derecha a izquierda:
\begin{equation}
u_i=t_i-v_i+c_{i+1},\qquad
r_i=u_i\bmod1000,\qquad
c_i=\frac{u_i-r_i}{1000}.
\label{eq:subacarreo}
\end{equation}
Fijados los operandos y el acarreo remoto, cociente y resto son únicos. Éste es
el operador que, más adelante, publicará las cifras de \(\epsone\). No hay
ningún decimal elegido después de ver el radián.

La completación
\[
10^3=729+243+27+1
\]
no es una curiosidad decimal. Expresa la conservación de la unidad al pasar
del ambiente ternario a una cámara de mil posiciones. La diferencia
\(1000-729=271\) es la corona completante; la identidad correcta es
\(999=729+270\), no \(729+271\).

\paragraph{Orientación ternaria y clasificación de los 3 regímenes operatorios del TRIT}

El TRIT asigna a cada estado local uno de tres regímenes orientados. En la
convención temporal activa,
\[
\begin{aligned}
+1&\longleftrightarrow\text{retorno, memoria, pasado},\\
0&\longleftrightarrow\text{umbral, presente},\\
-1&\longleftrightarrow\text{apertura bidireccional, futuro}.
\end{aligned}
\]
La realización cuadrática distingue:
\begin{equation}
J^2=-I,\qquad N^2=0,\qquad R^2=I.
\label{eq:tres-regimenes}
\end{equation}
Así aparecen, respectivamente, la rotación elíptica, el umbral parabólico y
la separación hiperbólica. No se mezclan tres métricas. Se conservan tres
acciones tipadas sobre un antecedente común.

La raíz interna de \(-1\) no se importa suponiendo de antemano el plano
complejo. En el bloque APP, dos proyecciones \(P,Q\) con razón
\(r=3/4=90/120\) producen
\[
J_{\mathrm{comm}}=\frac{[P,Q]}{\sqrt{r(1-r)}},\qquad J_{\mathrm{comm}}^2=-I,
\]
mientras
\[
R_{\mathrm{comm}}=2P-I,\qquad R_{\mathrm{comm}}^2=I,\qquad
R_{\mathrm{comm}}J_{\mathrm{comm}}=-J_{\mathrm{comm}}R_{\mathrm{comm}}.
\]
La pareja genera la forma mínima de \(\mathrm{Cl}_{1,1}\simeq M_2(\R)\):
rotación y apertura emergen juntas, pero siguen siendo operaciones diferentes.

\paragraph{Composición dinámica del TPK: selección, transporte, memoria y retorno}

El Topological Phase Núcleo no es un nombre para una caja de objetos. Es la
composición que selecciona fase, transporta cardinalidad, levanta hojas,
registra memoria y ejecuta el retorno. La dependencia fundamental es
\[
U_t\longrightarrow\Gamma_9\longrightarrow K_{\mathrm{ph}}.
\]
La holonomía nonádica \(\Gamma_9\) actúa sobre el estado enriquecido; la
publicación \(K_{\mathrm{ph}}\) es posterior y conserva memoria reducida. La
identidad \(K_{\mathrm{ph}}^9|_r=E_+E_\times\) no convierte a \(K_{\mathrm{ph}}\) en
generador.

El certificado nonádico conserva la cadena completa
\[
w_6\to w_{12}\to w_{18}\to w_{24}\to w_{30}
\to R_{36}\to G_9,
\]
con \(396\) niveles, \(2376\) trits, \(43\) vueltas, \(387\) pares
\((K,K+9)\) por canal y cinco regiones de \(\pi\). La fase cumple
\[
g(k+9)=g(k),
\]
pero el estado completo no retorna:
\[
B_{k+9}\ne B_k.
\]
Ésta es la semilla matemática de la hélice: la proyección angular cierra; la
memoria avanza.

\paragraph{Extensión temporal no escindida \(C_{12}\to C_{108}\to C_9\)}

La memoria de doce pasos se organiza mediante
\begin{equation}
0\longrightarrow C_{12}\longrightarrow C_{108}
\longrightarrow C_9\longrightarrow0.
\label{eq:extension-108}
\end{equation}
El cociclo
\[
c(r,r')=\left[\left\lfloor\frac{r+r'}9\right\rfloor\right]_{12}
\]
mide el acarreo entre retorno de fase y avance de memoria. Nueve pasos cierran
la fase observable \(C_9\); no reinician el registro dodecafásico.

Conviene separar cuatro objetos:
\begin{center}
\begin{tabularx}{.95\textwidth}{>{\bfseries}l X}
\toprule
\(C_{12}\) & grupo cíclico de doce posiciones;\\
\(\Rstate\) & registro enriquecido \((E_m,\tau_m,\sigma_m,\kappa_m,\Lambda_m)_{m=1}^{12}\);\\
\(\Theta_{108}\) & publicación cruda en el calendario de 108;\\
\(\Gamma_{108}\) & historia u holonomía levantada.\\
\bottomrule
\end{tabularx}
\end{center}
Por ello, la expresión vaga \enquote{post-\(C_{12}\)} se sustituirá por
\enquote{posterior al registro dodecafásico \(\Rstate\), en su lectura
sexádica} cuando ése sea el mapa efectivo.

\paragraph{Emergencia correlativa de las 5 construcciones del continuo desde el estado único}

La estructura discreta del continuo no se obtiene ensamblando cinco teorías
independientes. Sus cinco aspectos internos ---espectral, solenoidal, calibre,
cohomológico y determinantal--- emergen correlativamente del mismo estado
APP--TRIT--TPK. Esta obra no vuelve a demostrar el tratado integral, pero
preserva el hecho causal que necesita el radión: la sección coinductiva que
genera \(\pi,e,\varphi\) y \(\alpha\) no vive en una recta real desnuda; vive
en un objeto con fibras, relaciones, orientación, acarreo, memoria, terminal y
lector.

\begin{statusbox}
\textbf{Falsador de genealogía.} Si una fórmula posterior recibe
\(180/\pi\), \(\epsR\), una masa o una tabla metrológica para elegir una fase,
un coeficiente o una rama, deja de ser una generación HMT-first. El contraste
externo puede reconocer una salida; no puede seleccionarla.
\end{statusbox}


\subsection{Independencia prospectiva del cierre respecto del valor final}
% Fragmento público generado desde propietarios íntegros.
% Residencia: 52.1.2.
% Fuente: ../incoming/radion_monografia_20260827/manuscrito/sections/02_genealogia.tex:181-272

\subsubsection{Genealogía de las constantes y construcción del par angular}

\paragraph{Sección coinductiva y unicidad del valor generado de \(\pi\)}

Sea \(x_\infty\) la sección global del sistema inverso. Los lectores
coinductivos producen
\[
\operatorname{ev}_\pi(x_\infty)=\pih,\qquad
\operatorname{ev}_e(x_\infty)=e_{\HMT},\qquad
\operatorname{ev}_\varphi(x_\infty)=\varphi_{\HMT}.
\]
El subíndice HMT no designa otro número. Conserva la genealogía de la misma
constante. En particular,
\[
\pih=3.14159265358979323846264338327950288419716939937510\ldots.
\]
La vuelta positiva es
\begin{equation}
T_+(x_\infty)=2\pih.
\label{eq:vuelta-plus}
\end{equation}
Su publicación angular de una unidad es
\begin{equation}
R_{\HMT}^{\circ}=\frac{360^\circ}{T_+}=\frac{180^\circ}{\pih}.
\label{eq:radian-publicado}
\end{equation}
La ecuación no define todavía el radión enriquecido; define la carta visible
del radián una vez generada la vuelta.

\paragraph{Bidireccionalidad de rutas y prolongaciones conjugadas}

La bidireccionalidad no se enuncia como eslogan. La involución concreta
\begin{equation}
J_\alpha(T)=\frac{\alphah^2}{T},\qquad
T_-=J_\alpha(T_+),\qquad T_+T_-=\alphah^2
\label{eq:bidireccional-turn}
\end{equation}
conserva un producto y revierte la hoja. Expansión y contracción son dos
prolongaciones conjugadas del mismo estado. La raíz finita de un jet
cuadrático puede aproximar \(T_+\), pero no sustituye a la sección
coinductiva: la diferencia certificada es
\[
T_{P_9}-T_+=-5.1110566290335\ldots\times10^{-25}.
\]
El control finito es un testigo; el fundamento es el límite.

\paragraph{Dependencia causal entre \(\alpha\), acción y publicación angular}

El cierre dodecafásico genera \(\alphah\) desde el estado firmado y sus
cartas reversibles. Después, el lector determinantal local fija la década de
acción mediante
\[
\Sigma_H=\varphi_{\HMT}\alphah^{16},
\qquad \operatorname{ord}_{10}\Sigma_H=-34.
\]
La sección de retorno produce \(\hret\). Ni \(\alpha\) ni \(\hbar\) se
\enquote{miden en grados}. Se publican posteriormente mediante el par angular
\begin{equation}
P(\alpha,\eta_{\mathrm{ret}})=
\left([10^3\alpha]_{360},[2(\eta_{\mathrm{ret}}+\alpha)]_{360}\right)
=(A,C^*).
\label{eq:parangular}
\end{equation}
En la rama canónica,
\[
A=7.297352569283800997285105472380663\ldots{}^\circ,
\]
\[
C^*=2.1237383389686457242619513803933607937\ldots{}^\circ.
\]
Sólo entonces la carta analítica
\[
x=\frac{\pi A}{180},\qquad y=\frac{\pi C^*}{180}
\]
abre los canales \(q_\pm=e^{-x\mp y}\).

El orden causal completo es
\[
\alpha\to(\eta\parallel-34)\to\hbar
\to P\to\kappa_{360}\to q_\pm
\to\text{realizaciones MD}.
\]
La relación \(C^*/2-\alpha\) puede servir de verificación inversa. No se usa
para generar \(\hret\) retroactivamente.

\begin{narrativebox}
La carta angular no convierte una constante adimensional o una acción en un
ángulo por decreto. Publica dos coordenadas del estado en una rueda finita.
El grado es el idioma de esa publicación; no es la sustancia ontológica de
\(\alpha\) ni de \(\hbar\).
\end{narrativebox}
% Fuente: ../incoming/radion_monografia_20260827/manuscrito/sections/03_cierre_exacto.tex:254-285
\paragraph{Independencia prospectiva respecto del valor final}

El verificador del operador de profundidad prohíbe expresamente tres entradas:
\[
\frac{180}{\pi},\qquad \epsR,\qquad
\text{el decimal objetivo}.
\]
Primero construye \(S_E\) y \(S_T\), después ejecuta la resta APP, y sólo al
final compara el cierre. Los controles dan
\[
\texttt{target\_epsilon\_used=False},\qquad
\texttt{target\_180\_over\_pi\_used=False}.
\]
Asimismo:
\begin{itemize}
\item una sola hoja da un defecto de orden \(10^{-5}\);
\item sustituir \(\theta_2\) por \(\theta_1\) o \(\theta_3\) desplaza el
resultado en aproximadamente \(\pm\pi/6\);
\item eliminar el canal \(e\) deja un defecto de orden \(10^{-3}\);
\item la variante histórica \(\Delta_4(A,C^*)\) produce otro número;
\item la cola espectral posterior a \(\Rstate\) y el operador armónico de
cuatro términos tampoco coinciden con \(\epsone\).
\end{itemize}

\begin{statusbox}
\textbf{Estatuto.} El cierre es una \status{formalización nueva} demostrada
con estructura de partida explícita y certificado computacional independiente del valor final.
La igualdad es exacta en el límite coinductivo. El perfil corto de 36 cifras
de \(\alpha\) genera una variante que diverge sólo después de muchas cifras;
se conserva como testigo histórico, no como valor canónico.
\end{statusbox}



\section{Secciones directa y torsional: acarreo APP y cierre exacto del radión}
\subsection{Construcción independiente de las secciones \(S_E\) y \(S_T\) antes de definir el cociclo \(\varepsilon_R=S_T-S_E\)}
% Fragmento público generado desde propietarios íntegros.
% Residencia: 52.2.1.
% Fuente: ../incoming/radion_monografia_20260827/manuscrito/sections/03_cierre_exacto.tex:1-50
\noindent El cierre exacto se obtiene mediante dos secciones construidas antes del cociclo que las relaciona.\par\medskip

\subsubsection{Construcción independiente de las secciones directa y torsional}

\paragraph{Sección directa del sistema temporal dodecafásico}

La rueda dodecafásica tiene centros
\begin{equation}
\theta_m=30m-10^\circ,\qquad m=1,\ldots,12.
\label{eq:rueda-dodecafasica}
\end{equation}
La segunda fase es \(\theta_2=50^\circ\). Junto a la publicación
\(A=10^3\alphah\), forma la sección directa
\begin{align}
S_E
&=\frac{T_+}{360}\,(50+10^3\alphah)\\
&=2\pih\left(\frac5{36}+\frac{25}{9}\alphah\right).
\label{eq:seccion-electrica}
\end{align}
El subíndice \(E\) significa aquí \emph{directa o común}; su realización
eléctrica llegará después del lector constitutivo. No se postula
literalmente \(A=E\).

Los cilindros coinductivos prueban
\begin{equation}
1<S_E<2.
\label{eq:SE-intervalo}
\end{equation}
La división APP fuerza entonces
\begin{equation}
q_E=\lfloor S_E\rfloor=1,
\qquad
D_E=\operatorname{rem}_1(S_E)=S_E-1.
\label{eq:DE}
\end{equation}
El entero uno no es el número de vueltas de la hélice. Es el cociente único
de la sección directa dentro de la celda de unidad.

\paragraph{Sección torsional determinada por el operador \(\Delta_4\)}

La segunda sección utiliza una salida coinductiva anterior:
\begin{equation}
\Delta_4=
\frac{\pih-e_{\HMT}\log\pih}{270}
\left(1+\frac{\pih}{729}\right).
\label{eq:delta4}
\end{equation}
La palabra \emph{torsión} designa aquí el cuanto global de la carta HMT. No
lo identifica automáticamente con la torsión de Cartan
\(T^I=D_\omega e^I\). El puente hacia Cartan exige otra flecha tipada.

\subsection{Densidad orientada del sector activo y determinación única de \(20/81\)}
% Fragmento público generado desde propietarios íntegros.
% Residencia: 52.2.2.
% Fuente: ../incoming/radion_monografia_20260827/manuscrito/sections/03_cierre_exacto.tex:51-103

\paragraph{Densidad orientada \(20/81\) del sector activo \(90/120\)}

El selector trítico de fase separa
\[
H_+=\{1,4,7\},\qquad H_-=\{2,5,8\},\qquad H_0=\{3,6,9\}.
\]
Hay seis fases activas y quince parejas no ordenadas:
\[
H_{\mathrm{act}}=H_+\cup H_-,\qquad
\left|\binom{H_{\mathrm{act}}}2\right|=\binom62=15.
\]
El modo activo y su extensión a ocho posiciones son
\begin{equation}
N_6=6\binom62=90,\qquad
N_8=8\binom62=120.
\label{eq:90-120-incidencia}
\end{equation}
Dos hojas conjugadas del sector activo, normalizadas por la fibra ambiente,
dan
\begin{equation}
\rho_{\mathrm{tw}}
=\frac{2N_6}{|\F_3^6|}
=\frac{2\cdot90}{729}
=\frac{180}{729}
=\boxed{\frac{20}{81}}.
\label{eq:rho-twoway}
\end{equation}

\begin{proposition}[Unicidad relativa de la densidad]
Fijados el sector activo \(N_6=90\), las dos hojas conjugadas y la
normalización por \(|\F_3^6|=729\), la densidad orientada es necesariamente
\(20/81\).
\end{proposition}
\begin{proof}
La cardinalidad transportada es \(2N_6=180\). La normalización por el ambiente
da \(180/729\); dividir numerador y denominador por \(9\) produce \(20/81\).
No interviene ninguna comparación con \(180/\pi\) ni con \(\epsR\).
\end{proof}

La lectura histórica como proyección sobre una red \(1/81\) es un control de
rigidez útil, pero ya no es el fundamento: el coeficiente se obtiene antes
del cierre. La ablación de una hoja produce \(10/81\), y la salida cambia en
orden \(10^{-5}\); por tanto, el factor dos no es decorativo.

La sección torsional es
\begin{equation}
\Phi_T=\rho_{\mathrm{tw}}\Delta_4,
\qquad
S_T=1+\Phi_T.
\label{eq:seccion-torsional}
\end{equation}


\subsection{El acarreo \(\varepsilon_R\) como cociclo de transición entre secciones}
% Fragmento público generado desde propietarios íntegros.
% Residencia: 52.2.3.
% Fuente: ../incoming/radion_monografia_20260827/manuscrito/sections/03_cierre_exacto.tex:104-253
\subsubsection{Cociclo de transición entre las secciones directa y torsional}

\paragraph{Definición operatoria del acarreo \(\varepsilon_R\)}

\begin{definition}[Acarreo real del radión]
La salida unitaria del radión es el cociclo de transición
\begin{equation}
\epsone:=S_T-S_E
=\Phi_T-D_E
=\frac{20}{81}\Delta_4-operatorname{rem}_1(S_E).
\label{eq:epsilon-def}
\end{equation}
\end{definition}

La ecuación contiene una resta, pero no es un ajuste al blanco. Sus dos
operandos \(S_T\) y \(S_E\) se generan separadamente. La independencia que
se demuestra es independencia respecto de \(180/\pi\) y de un
\(\varepsilon\) objetivo; no se afirma independencia algebraica entre una
diferencia y sus sumandos.

\begin{theorem}[Cierre exacto unitario]
Con \(S_E\), \(\Phi_T\) y \(\epsone\) definidos por
\eqref{eq:seccion-electrica}, \eqref{eq:delta4}, \eqref{eq:rho-twoway} y
\eqref{eq:epsilon-def}, se cumple exactamente
\begin{equation}
\boxed{1=S_E-\Phi_T+\epsone.}
\label{eq:cierre-unitario}
\end{equation}
\end{theorem}
\begin{proof}
Por \eqref{eq:SE-intervalo}, \(D_E=S_E-1\). Por definición,
\(\epsone=\Phi_T-D_E\). Entonces
\[
S_E-\Phi_T+\epsone
=S_E-\Phi_T+\Phi_T-(S_E-1)=1.
\]
La fuerza de la prueba no está en esta cancelación elemental, sino en que la
genealogía previa fija los dos operandos sin usar el cierre como entrada.
\end{proof}

\paragraph{Publicación del cierre en la carta angular}

La carta de la vuelta multiplica una fracción unitaria por
\(360/T_+=180/\pih\). En particular,
\[
\epsR^\circ=\frac{360^\circ}{T_+}\epsone.
\]
Multiplicar \eqref{eq:cierre-unitario} por \(360^\circ/T_+\) y usar
\[
\frac{360}{T_+}S_E
=360\left(\frac5{36}+\frac{25}{9}\alphah\right)
=50+1000\alphah=50+A
\]
produce el cierre angular.

\begin{theorem}[Ecuación exacta del radión]
\BigRadionEquation
\end{theorem}

La ecuación dice algo más preciso que \enquote{el radián es
\(57.2957\ldots{}^\circ\)}. Descompone esa publicación en cuatro operaciones
con demostración de referencia:
\begin{center}
\begin{tabularx}{.96\textwidth}{>{\bfseries}l X}
\toprule
\(50^\circ\) & fase segunda de la rueda; bisagra crítica en una carta declarada;\\
\(A^\circ\) & coordenada común de la publicación parangular;\\
\(-\frac{20}{81}\Delta_4^\circ\) & transporte orientado de la torsión global;\\
\(+\epsR^\circ\) & acarreo real que empalma los dos lifts.\\
\bottomrule
\end{tabularx}
\end{center}

\paragraph{Convergencia por profundidad del refinamiento}

A profundidad \(N\), el generador entrega cilindros racionales
\(p_N\to\pih\) y \(e_N\to e_{\HMT}\). Se definen
\[
S_{E,N}=2p_N\left(\frac5{36}+\frac{25}{9}\alphah\right),
\]
\[
\Delta_{4,N}=\frac{p_N-e_N\log p_N}{270}
\left(1+\frac{p_N}{729}\right),\qquad
S_{T,N}=1+\frac{20}{81}\Delta_{4,N},
\]
y
\[
\mathcal E_N=S_{T,N}-S_{E,N}.
\]
La continuidad del funcional en \(3<p<4\), \(2<e<3\), junto a las
derivadas monótonas, transporta cada par de cilindros a un intervalo
certificado. A profundidad \(45\) bloques, su anchura es menor que
\(4.81\times10^{-130}\).

Cada retorno de nueve niveles contiene seis trits por nivel. Por eso la tasa
de refinamiento es
\begin{equation}
729^{-9}=3^{-54}.
\label{eq:contraccion-54}
\end{equation}
Esta cifra controla la anchura de los intervalos; no es el valor de
\(\epsR\).

La resta de tríadas estabiliza
\begin{center}
\ttfamily
000|000|000|896|802|822|044|562|981|999|050|695|020|651|286|413
\end{center}
y el prefijo de carries
\begin{center}
\ttfamily
0,0,0,0,0,-1,0,-1,-1,-1,0,-1,0.
\end{center}
Éste es el mecanismo concreto de memoria decimal: no una cola elegida, sino
la recurrencia \eqref{eq:subacarreo} aplicada a operandos previamente
construidos.

\paragraph{Evaluación de los valores canónicos}

La salida coinductiva es
\begin{align}
\epsone
&=8.9680282204456298199905069502065128641372736205009\ldots
\times10^{-10},\\
\epsR^\circ
&=5.1383016758575282071685164622645947489908574400054\ldots
\times10^{-8}\,{}^\circ.
\label{eq:epsilon-valores}
\end{align}
La publicación completa resulta
\[
\frac{180^\circ}{\pih}
=57.2957795130823208767981548141051703324054724665643\ldots{}^\circ.
\]

\begin{exactbox}
\begin{align*}
50^\circ+A^\circ
&=57.29735256928380099728510547238066299956\ldots{}^\circ,\\
\frac{20}{81}\Delta_4^\circ
&=0.00157310758449687906223272996065728980465\ldots{}^\circ,\\
50^\circ+A^\circ-\frac{20}{81}\Delta_4^\circ
&=57.29577946169930411822287274242000570976\ldots{}^\circ,\\
\epsR^\circ
&=0.00000005138301675857528207168516462264594749\ldots{}^\circ,\\
\text{suma final}
&=57.29577951308232087679815481410517033241\ldots{}^\circ.
\end{align*}
\end{exactbox}



\section{Coordenada crítica y sistema temporal dodecafásico orientado: publicación afín de la fase de \(50^\circ\)}
\subsection{Aplicación afín \(j_{100}\) de la carta APP y unicidad de la fase \(m=2\) en \(j_{100}(1/2)=\theta_m=50^\circ\)}
% Fragmento público generado desde propietarios íntegros.
% Residencia: 52.3.1.
% Fuente: ../incoming/radion_monografia_20260827/manuscrito/sections/04_cincuenta_critica.tex:1-105
\noindent La coordenada crítica enlaza la carta de Cayley, el registro APP y la fase dodecafásica mediante aplicaciones tipadas.\par\medskip

\subsubsection{Dos realizaciones exactas de la coordenada angular \(50^\circ\)}

\paragraph{Segunda fase del sistema temporal dodecafásico}

Por \eqref{eq:rueda-dodecafasica}, la sucesión de centros comienza
\[
20^\circ,\ 50^\circ,\ 80^\circ,\ 110^\circ,\ldots.
\]
Resolver
\[
30m-10=50
\]
da únicamente \(m=2\). Por tanto, \(50^\circ\) no es una redondeo de \(A\),
una media vuelta ni el centro de un círculo: es una fase tipada del calendario
dodecafásico. La fracción de vuelta es
\[
\frac{50}{360}=\frac5{36},
\]
exactamente el primer sumando de \(S_E/T_+\).

\paragraph{Igualación de la coordenada crítica de Cayley}

La involución funcional
\[
\rho\longmapsto1-\bar\rho
\]
fija
\[
\operatorname{Fix}(\rho\mapsto1-\bar\rho)
=\{\rho:\Re\rho=1/2\}.
\]
El número \(1/2\) es la coordenada real autodual del intervalo funcional
\(0\leq\Re s\leq1\). No es \enquote{la mitad de una recta} en sentido
longitudinal, y no contiene la altura espectral \(\gamma\). Para
\[
\rho_\gamma=\frac12+i\gamma,
\]
la transformación de Cayley se puede escribir
\begin{equation}
\tau_\gamma
=\frac{\rho_\gamma-1}{\rho_\gamma}
=\frac{\gamma+i/2}{\gamma-i/2}
\in S^1,
\qquad
\gamma=\frac12\cot\frac{\theta_\gamma}{2}.
\label{eq:cayley-critica}
\end{equation}
Así, \(1/2\) fija la costura, \(\gamma\) determina el carácter angular y
\(\Gamma_9\) conserva la profundidad de sus iteraciones.

\paragraph{Inyección afín \(j_{100}\) del registro APP en la carta angular}

La completación cúbica dispone una cámara de diez marcas por coordenada. Una
cara contiene \(10^2=100\) posiciones. Declaramos la carta afín
\begin{equation}
j_{100}:[0,1]\longrightarrow[0,100^\circ],
\qquad j_{100}(\sigma)=100\sigma^\circ.
\label{eq:j100}
\end{equation}
Entonces
\[
j_{100}\left(\frac12\right)=50^\circ.
\]
Al exigir simultáneamente \(j_{100}(1/2)=\theta_m\), resulta
\[
100\cdot\frac12=30m-10
\quad\Longleftrightarrow\quad m=2.
\]

\begin{theorem}[Igualador crítico--dodecafásico]
En la carta APP de cien marcas, la coordenada fija de la involución crítica
se publica exactamente en la segunda fase de la rueda:
\begin{equation}
\boxed{j_{100}(1/2)=\theta_2=50^\circ.}
\label{eq:igualador-50}
\end{equation}
La fase \(m=2\) es única.
\end{theorem}

El teorema es exacto una vez declarada \eqref{eq:j100}. Su estatuto es
\status{formalización nueva}: el medio crítico, la completación \(10^3\), la
cara de cien marcas y \(\theta_2\) ya estaban construidos; la elección de esa
cara como carta crítica afín reúne por primera vez sus dominios.

\paragraph{Compatibilidad del valor crítico con la carta de Cayley}

La transformación de Cayley no envía \(1/2\) a \(50^\circ\). Si
\(\gamma=0\), \eqref{eq:cayley-critica} da
\[
\tau_0=-1=e^{i\pi},
\]
es decir, \(180^\circ\). Los \(50^\circ\) pertenecen a la carta APP de cien
marcas y al calendario dodecafásico; el ángulo Cayley depende de \(\gamma\).
Esta distinción impide colapsar todos los ceros de la función zeta en una sola
fase.

\begin{falsifier}
La identificación \(1/2\leftrightarrow50^\circ\) falla si se elimina la
carta \eqref{eq:j100}. La igualdad escalar \(100(1/2)=50\) no basta para
identificar por sí sola dos objetos. El teorema depende de la cara APP
distinguida y de la rueda dodecafásica.
\end{falsifier}


\subsection{Compatibilidad de la fase crítica con el retorno nonádico y el avance de memoria}
% Fragmento público generado desde propietarios íntegros.
% Residencia: 52.3.2.
% Fuente: ../incoming/radion_monografia_20260827/manuscrito/sections/04_cincuenta_critica.tex:106-171
\subsubsection{Separación tipada de los 2 objetos de valor \(1/2\)}

APP contiene además un modo singular con
\[
s=\frac12,\qquad s^2=\frac14,\qquad
1-s^2=\frac34=\frac{90}{120},\qquad
s\sqrt{1-s^2}=\frac{\sqrt3}{4}.
\]
Este \(s\) nace de los ángulos principales de las hojas suma y producto. La
costura crítica \(\Re\rho=1/2\) nace de una involución funcional. Que ambas
coordenadas valgan \(1/2\) es una coincidencia escalar con posible contenido,
pero no autoriza a identificarlas sin un cuadrado tipado.

La forma correcta de conservar el hallazgo es un diagrama de igualadores:
\[
\begin{tikzpicture}[baseline=(current bounding box.center),node distance=14mm and 23mm,
 every node/.style={font=\small}]
\node[draw=RadBlue,rounded corners,fill=RadMist] (crit) {costura crítica \(\Re s=1/2\)};
\node[draw=RadTeal,rounded corners,fill=RadCream,right=of crit] (face) {carta \(j_{100}\)};
\node[draw=RadCopper,rounded corners,fill=RadCream,right=of face] (theta) {fase \(\theta_2=50^\circ\)};
\node[draw=RadRose,rounded corners,fill=white,below=of face] (mode) {modo APP \(s=1/2\), \(1-s^2=90/120\)};
\draw[-{Latex},thick,RadBlue] (crit)--node[above]{\(j_{100}\)}(face);
\draw[-{Latex},thick,RadTeal] (face)--node[above]{igualador}(theta);
\draw[-{Latex},dashed,RadRose] (mode)--node[right,align=left]{misma cifra;\\dominio distinto}(face);
\end{tikzpicture}
\]
El diagrama no finge la flecha ausente. Muestra qué está demostrado, qué se
ha formalizado y qué igualdad sigue necesitando un transporte adicional si
se pretende una identificación operatoria global.

\subsubsection{Compatibilidad entre círculo espectral y memoria helicoidal}

El lector Cayley--Pontryagin transporta una fibra espectral mediante
\[
\widetilde U_{\Gamma_9^n}J\psi_\gamma
=\tau_\gamma^nJ\psi_\gamma.
\]
La publicación circular es \(\tau_\gamma^n\in S^1\). El lift
\[
(\tau_\gamma^n,n)\in S^1\times\Z
\]
conserva el número de retornos. En la hoja contractiva aparece
\[
\widetilde M_9^nJ\psi_\gamma
=\left(\frac59\right)^n\tau_\gamma^nJ\psi_\gamma.
\]
Tenemos, por tanto, tres imágenes tipadas:
\[
\text{círculo de fase},\qquad
\text{hélice de memoria},\qquad
\text{espiral contractiva}.
\]

Un cero no es una vuelta, y un nueve no es un cero particular de la función
zeta. El cero fija un carácter \(\tau_\gamma\) que persiste a través de las
vueltas; el retorno nonádico incrementa memoria. El nueve funciona como cero
residual en \(\Z/9\Z\), una operación aritmética distinta de la enumeración
de ceros espectrales.

\begin{narrativebox}
La recta crítica se vuelve círculo cuando olvidamos la altura lineal y
conservamos el carácter unitario. Se vuelve hélice cuando volvemos a insertar
la memoria de cuántas veces el carácter ha sido transportado. La profundidad
no es una tercera coordenada decorativa: es la información que la proyección
circular no puede retener.
\end{narrativebox}


\section{Par angular y acción reducida: construcción de la elipse positiva y del área por radión}
\subsection{Ley areal del radión: \(\mathcal A(\theta)=\hbar_{\mathrm{ret}}\theta\) y \(\mathcal A(2\pi)=h_{\mathrm{ret}}\)}
% Fragmento público generado desde propietarios íntegros.
% Residencia: 52.4.1.
% Fuente: ../incoming/radion_monografia_20260827/manuscrito/sections/05_geometria_accion.tex:1-217
\noindent La geometría espectral determina la forma y la ley areal determina la acción transportada por la fase.\par\medskip

\subsubsection{Equivalencia entre las cartas espectral y geométrica}

\paragraph{Operador espectral positivo y autovalores conjugados}

El par angular produce el operador bicapa
\begin{equation}
B_1=AI+C^*R,\qquad R^2=I,\qquad
P_\pm=\frac{I\pm R}{2}.
\label{eq:B1}
\end{equation}
Sus autovalores son \(A\pm C^*\). Como \(0<C^*<A\), ambos son positivos.
En la carta \(\kappa=\pi/180\), la elipse espectral tiene semiejes
\begin{equation}
a_1^2=\kappa(A+C^*),\qquad
b_1^2=\kappa(A-C^*).
\label{eq:semiejes-espectrales}
\end{equation}
La descomposición es reversible:
\begin{equation}
A=\frac{a_1^2+b_1^2}{2\kappa},\qquad
C^*=\frac{a_1^2-b_1^2}{2\kappa}.
\label{eq:reconstruccion-AC}
\end{equation}
La elipse no se dibuja para ilustrar un número. Reconstruye exactamente el
par que la generó.

Definamos
\begin{equation}
\rho_1=\sqrt{A^2-C^{*2}},\qquad
\eta=\operatorname{artanh}\frac{C^*}{A}
=\frac12\log\frac{A+C^*}{A-C^*}.
\label{eq:eta}
\end{equation}
Entonces
\[
A\pm C^*=\rho_1e^{\pm\eta},
\]
y
\begin{equation}
a_1=\sqrt{\kappa\rho_1}\,e^{\eta/2},
\qquad
b_1=\sqrt{\kappa\rho_1}\,e^{-\eta/2}.
\label{eq:semiejes-eta}
\end{equation}
El producto conserva escala; el cociente conserva forma:
\[
a_1b_1=\kappa\rho_1,
\qquad
\frac{b_1}{a_1}=e^{-\eta}
=\sqrt{\frac{A-C^*}{A+C^*}}.
\]

\paragraph{Invariantes focales del operador positivo}

La semidistancia focal es
\begin{equation}
c_1^2=a_1^2-b_1^2=2\kappa C^*,
\qquad
F_\pm=(\pm c_1,0).
\label{eq:focos}
\end{equation}
La distancia total vale
\begin{equation}
d_1=2c_1=2\sqrt{2\kappa C^*}
=0.544545509325626623406299779866023\ldots.
\label{eq:d1}
\end{equation}
El número aislado no es la tesis. La identidad \(c_1^2=2\kappa C^*\) dice que
\(C^*\) es literalmente la escisión focal cuadrática en la carta
distinguida. Los focos son los extremos conjugados de la apertura espectral;
no son, por sí solos, los polos eléctrico y magnético. Esa lectura requiere
el transductor constitutivo.

La excentricidad es
\begin{equation}
e_f=\frac{c_1}{a_1},\qquad
e_f^2=1-e^{-2\eta}=\frac{2C^*}{A+C^*}.
\label{eq:excentricidad}
\end{equation}
Numéricamente,
\[
\begin{aligned}
\eta&=0.299689683921630799684926562863927\ldots,\\
e_f&=0.671451895550191986916277055864332\ldots.
\end{aligned}
\]

\paragraph{Bloque APP de referencia y transporte normalizado}

El bloque central
\[
B_0=7I+2R=9P_++5P_-
\]
posee
\[
\rho_0=\sqrt{45},\qquad
\eta_0=\frac12\log\frac95,\qquad
e_0=\frac23.
\]
Su distancia focal es
\[
d_0=4\sqrt\kappa
=0.528443639680801468446003835349358\ldots.
\]
La comparación exacta es
\begin{equation}
\boxed{\left(\frac{d_1}{d_0}\right)^2=\frac{C^*}{2}},
\qquad
d_1=d_0\sqrt{\frac{C^*}{2}}.
\label{eq:invariante-focal-relativo}
\end{equation}
Ésta es la genealogía defendible de \(0.544545\ldots\): una escala focal APP
de referencia multiplicada por un invariante relativo. Su cercanía visual a
\(0.54\) no es una identidad con el defecto APP \(54\).

El transporte desde \(B_0\) se separa en escala y anisotropía:
\begin{equation}
T_{\mathrm{rel}}=\frac{\rho_1}{\sqrt{45}}
\exp\bigl((\eta-\eta_0)R\bigr).
\label{eq:transporte-rel}
\end{equation}
Sobre las hojas,
\[
t_+=\frac{A+C^*}{9}=1.0467878786947163\ldots,
\qquad
t_-=\frac{A-C^*}{5}=1.0347228460630311\ldots.
\]
La deformación no es un solo decimal. Su parte isotrópica es
\(\sqrt{t_+t_-}\); su parte hiperbólica es
\(\tfrac12\log(t_+/t_-)\).

\subsubsection{Elipse positiva asociada al par angular}

\paragraph{Construcción mediante semiejes conjugados}

La escala espectral anterior carece de unidades de acción. Mecánica
Dimensional publica la misma forma sobre la recta interna de acción:
\begin{equation}
a_S=\sqrt{2\hret e^\eta},
\qquad
b_S=\sqrt{2\hret e^{-\eta}}.
\label{eq:semiejes-accion}
\end{equation}
Inmediatamente,
\begin{equation}
a_Sb_S=2\hret,\qquad
\frac{a_S}{b_S}=e^\eta.
\label{eq:producto-forma}
\end{equation}
El producto fija la acción; el cociente fija la anisotropía. La elipse se
parametriza por
\begin{equation}
Q(\theta)=a_S\cos\theta,\qquad
P(\theta)=b_S\sin\theta.
\label{eq:elipse-param}
\end{equation}

Las elipses espectral y de acción son homotéticas:
\begin{equation}
\lambda_{\mathrm{act}}^2=\frac{2\hret}{\kappa\rho_1},
\qquad
(a_S,b_S,c_S)=\lambda_{\mathrm{act}}(a_1,b_1,c_1).
\label{eq:homotecia}
\end{equation}
Comparten la forma \(e^{-\eta}\), no la dimensión física.

\paragraph{Ley areal de la acción por radión}

La forma simpléctica canónica es \(\omega=\diff Q\wedge\diff P\). El área
orientada barrida desde \(0\) hasta \(\theta\) es
\begin{align}
\mathcal A(\theta)
&=\frac12\int_0^\theta
\bigl(Q(t)P'(t)-P(t)Q'(t)\bigr)\,\diff t\\
&=\frac12a_Sb_S\int_0^\theta
(\cos^2t+\sin^2t)\,\diff t.
\end{align}
Usando \eqref{eq:producto-forma}, obtenemos el resultado nuclear.

\begin{theorem}[Acción por radión]
Para la elipse \eqref{eq:elipse-param},
\begin{equation}
\boxed{\mathcal A(\theta)=\hret\theta.}
\label{eq:area-theorem}
\end{equation}
En particular,
\begin{equation}
\boxed{\mathcal A(1)=\hret},
\qquad
\boxed{\mathcal A(2\pi)=2\pi\hret=\hfull}.
\label{eq:area-radion-vuelta}
\end{equation}
\end{theorem}

Esta igualdad permite definir el radión sin apelar primero a la longitud de
arco.

\begin{definition}[Radión orientado]
\begin{equation}
\mathfrak r_\sigma=(\theta=1,\sigma),
\qquad \sigma\in\{+1,-1\},
\qquad
\mathcal A(\mathfrak r_\sigma)=\sigma\hret.
\label{eq:radion-oriented}
\end{equation}
\end{definition}

El radián es la proyección que olvida \(\sigma\), la hoja, la memoria y el
área. El radión conserva esa información. De ahí el nombre \emph{rad--ión}:
la orientación precede a la carga física, y la carga aparece más tarde como
\[
Q(x)=n_Q(x)e_{\mathrm{int}},\qquad
e_{\mathrm{int}}=\sqrt{\frac{2\alphah\hfull}{Z_{\mathrm{int}}}}.
\]


\subsection{Semiejes conjugados, invariantes focales y orientación de la anisotropía}
% Fragmento público generado desde propietarios íntegros.
% Residencia: 52.4.2.
% Fuente: ../incoming/radion_monografia_20260827/manuscrito/sections/05_geometria_accion.tex:218-290
\paragraph{Correspondencia entre fase, acción, área y orientación}

\begin{figure}[H]
\centering
\begin{tikzpicture}[scale=1.12]
  \def\aa{4.0}
  \def\bb{2.95}
  \def\cc{2.70}
  \draw[->,RadGray] (-4.7,0)--(4.7,0) node[right]{\(Q\)};
  \draw[->,RadGray] (0,-3.5)--(0,3.5) node[above]{\(P\)};
  \fill[RadTeal!8] (0,0) ellipse[x radius=\aa,y radius=\bb];
  \draw[RadTeal,line width=1.35pt] (0,0) ellipse[x radius=\aa,y radius=\bb];
  \fill[RadCopper] (-\cc,0) circle (2pt) node[below=2mm]{\(F_-\)};
  \fill[RadCopper] (\cc,0) circle (2pt) node[below=2mm]{\(F_+\)};
  \draw[RadGold,line width=1.1pt,-{Latex}] (\aa,0)
     arc[start angle=0,end angle=57.3,x radius=\aa,y radius=\bb];
  \node[RadGold] at (3.7,2.3) {un radión};
  \draw[dashed,RadBlue] (0,0)--(\aa,0) node[midway,below]{\(a_S\)};
  \draw[dashed,RadBlue] (0,0)--(0,\bb) node[midway,left]{\(b_S\)};
  \node[align=center,RadNavy] at (0,-3.25)
    {área total \(\pi a_Sb_S=2\pi\hret=\hfull\)};
\end{tikzpicture}
\caption{La elipse de acción. Un arco paramétrico de un radián barre
\(\hret\); una vuelta encierra \(\hfull\). Los focos codifican la apertura
\(C^*\) en la carta distinguida.}
\label{fig:elipse-accion}
\end{figure}

\paragraph{Relación entre la acción reducida, la incertidumbre y el área}

La relación \(\hbar=h/(2\pi)\) suele presentarse como una conveniencia
notacional: al describir oscilaciones, rotaciones y transformadas de Fourier,
los factores \(2\pi\) desaparecen si se usa \(\hbar\). La elipse de acción
revela una lectura más fuerte. La división por \(2\pi\) es exactamente el
paso de la acción de una vuelta a la acción por unidad de fase:
\[
h=\mathcal A(2\pi),\qquad \hbar=\mathcal A(1).
\]
La constante reducida no es sólo \(h\) con una notación más cómoda; es la
densidad areal de acción respecto a la fase.

Heisenberg, Dirac y Jordan establecieron el lenguaje de variables conjugadas,
conmutadores y transformaciones canónicas. No se les atribuye aquí una
\enquote{elipse ontológica} literal sin fuente. La contribución HMT--MD es
identificar la figura positiva, construir sus semiejes desde \((A,C^*)\) y
demostrar \eqref{eq:area-theorem}. La genealogía histórica abre el lenguaje;
el teorema actual fija el objeto.

\begin{narrativebox}
La acción no se añade después al movimiento. En la elipse del radión, la
acción es el área que el movimiento genera. Una vuelta no transporta una
constante prefabricada: la encierra.
\end{narrativebox}

\Needspace{10\baselineskip}
\paragraph{Controles algebraicos de los invariantes focales}

Las semejanzas decimales se someten a control:
\begin{center}
\begin{tabularx}{.96\textwidth}{l r X}
\toprule
Comparación & Residuo & Dictamen\\
\midrule
\(d_1\) frente a \(0.54\) & \(4.5455\times10^{-3}\) & no igualdad; el \(54\) no se deduce;\\
\(d_1^2\) frente a \(8/27\) & \(2.3352\times10^{-4}\) & proximidad sin mapa;\\
\(\eta\) frente a \(3/10\) & \(-3.1032\times10^{-4}\) & no igualdad;\\
\(e_f\) frente a \(2/3\) & \(4.7852\times10^{-3}\) & \(2/3\) pertenece exactamente a \(B_0\);\\
\(e^{-\eta}\) frente a \(20/27\) & \(3.0740\times10^{-4}\) & no igualdad.\\
\bottomrule
\end{tabularx}
\end{center}
Los invariantes sólidos son \eqref{eq:eta}, \eqref{eq:excentricidad} y
\eqref{eq:invariante-focal-relativo}. El criterio HMT no consiste en premiar
un parecido decimal, sino en exigir objeto, operación y conservación.


\section{Retorno nonádico e interior proyectivo: disco de Klein y profundidad helicoidal de memoria}
\subsection{Normalización al disco e interior hiperbólico de Hilbert--Beltrami--Klein}
% Fragmento público generado desde propietarios íntegros.
% Residencia: 52.5.1.
% Fuente: ../incoming/radion_monografia_20260827/manuscrito/sections/06_interior_helice.tex:1-108
\noindent La geometría proyectiva distingue la acción finita de la profundidad hiperbólica y la memoria del retorno.\par\medskip

\subsubsection{Interior hiperbólico de Hilbert--Beltrami--Klein}

\paragraph{Normalización proyectiva en el disco de Klein}

La aplicación afín
\[
u=\frac{Q}{a_S},\qquad v=\frac{P}{b_S}
\]
lleva la elipse de acción al disco unitario
\[
D=\{(u,v):u^2+v^2<1\}.
\]
Sobre cada cuerda, la distancia de Hilbert se define mediante la razón doble
de sus extremos de frontera. La métrica infinitesimal resultante es la forma
Beltrami--Klein
\begin{equation}
\diff s_K^2=
\frac{(1-r^2)(\diff u^2+\diff v^2)+(u\diff u+v\diff v)^2}
{(1-r^2)^2},
\qquad r^2=u^2+v^2.
\label{eq:klein-metric}
\end{equation}
Con esta normalización, la curvatura gaussiana es constante:
\begin{equation}
K=-1.
\label{eq:klein-curvature}
\end{equation}
Las geodésicas son cuerdas euclídeas, pero su longitud hiperbólica diverge al
acercarse a la frontera.

\paragraph{Finitud de la frontera e infinitud de la profundidad hiperbólica}

El determinante métrico de \eqref{eq:klein-metric} produce
\begin{equation}
\diff A_K=\frac{\diff u\,\diff v}{(1-r^2)^{3/2}}.
\label{eq:klein-area-density}
\end{equation}
El área del subdisco \(r<R<1\) es
\begin{align}
A_K(R)
&=\int_0^{2\pi}\int_0^R
\frac{r\,\diff r\,\diff\theta}{(1-r^2)^{3/2}}\\
&=2\pi\left(\frac1{\sqrt{1-R^2}}-1\right).
\label{eq:klein-area-R}
\end{align}
\Needspace{4\baselineskip}
Por tanto,
\[
\lim_{R\to1^-}A_K(R)=+\infty.
\]
Simultáneamente, la misma frontera encierra el área simpléctica
\[
\pi a_Sb_S=2\pi\hret=\hfull<\infty.
\]

\begin{theorem}[Finitud de acción e infinitud de profundidad]
La elipse del radión posee una acción de vuelta finita \(\hfull\), mientras
su interior, equipado con la métrica Hilbert--Beltrami--Klein, tiene área
hiperbólica infinita:
\begin{equation}
\boxed{
\operatorname{Area}_{\mathrm{simp}}(\partial\mathbb E_S)=\hfull<\infty,
\qquad
\operatorname{Area}_{K}(\mathbb E_S)=+\infty.}
\label{eq:finite-infinite}
\end{equation}
\end{theorem}

No hay contradicción: \(\operatorname{Area}_{\mathrm{simp}}\) y
\(\operatorname{Area}_K\) son medidas diferentes. La primera cuenta acción
orientada en fase; la segunda mide profundidad proyectiva interior. La tesis
\enquote{borde finito, interior infinito} es ahora una igualdad y un límite,
no una metáfora.

\paragraph{Coordenadas focales en el interior de Klein}

En el disco normalizado, los focos se sitúan en \((\pm e_f,0)\). La distancia
del centro a un foco es
\begin{equation}
u_f=\operatorname{artanh}e_f
=\operatorname{arcosh}(e^\eta)
=0.813382331175342333760889731088228\ldots.
\label{eq:klein-focal}
\end{equation}
La distancia foco--foco es
\[
d_K(F_-,F_+)=2u_f
=1.62676466235068466752177946217646\ldots.
\]
Se cumplen
\begin{equation}
e_f=\tanh u_f,\qquad
e^{-\eta}=\operatorname{sech}u_f,\qquad
\eta=\log\cosh u_f.
\label{eq:focal-hyperbolic-chain}
\end{equation}
La región hasta el radio focal tiene área
\[
A_K(e_f)=2\pi(e^\eta-1)
=2.19559620884061869541206115980360\ldots.
\]

Todas las elipses abiertas son proyectivamente isométricas como dominios de
Hilbert. Por eso la curvatura \(-1\) sola no recuerda la excentricidad. La
forma HMT reside en la estructura conjunta: carta afín distinguida, forma
simpléctica, etiquetado de hojas y parametrización de frontera.

\subsection{Diferencia entre retorno de fase y retorno del estado enriquecido}
% Fragmento público generado desde propietarios íntegros.
% Residencia: 52.5.2.
% Fuente: ../incoming/radion_monografia_20260827/manuscrito/sections/06_interior_helice.tex:109-201
\subsubsection{Elevación helicoidal del retorno nonádico}

\paragraph{Descomposición nonádica y cociente de memoria}

La operación elemental
\begin{equation}
n=9q_9(n)+\rho_9(n),
\qquad
H_9(n)=(\rho_9(n),q_9(n))
\label{eq:H9}
\end{equation}
satisface
\begin{equation}
H_9(n+9)=H_9(n)+(0,1).
\label{eq:H9-helix}
\end{equation}
La primera coordenada cierra; la segunda asciende. Esta identidad es la
versión aritmética mínima de un tornillo.

Para una fibra espectral,
\[
n\longmapsto(\tau_\gamma^n,n)
\]
es una hélice sobre \(S^1\). Añadir la contracción \((5/9)^n\) produce una
espiral interior. El radión participa en ambas lecturas: una unidad de fase
en la frontera y una unidad orientada de acción en el lift.

\paragraph{Rango de memoria como lector dimensional}

Supongamos que una proyección visible conserva el valor de fase, pero pierde
\(d\) cociclos de memoria linealmente independientes. Cualquier lift fiel
necesita al menos \(d\) coordenadas adicionales para reconstruirlos. Esto
motiva una lectura precisa de Mecánica Dimensional:
\begin{equation}
\begin{aligned}
\text{dimensión emergente}
&=\text{rango mínimo de memoria independiente}\\
&\phantom{=}\text{no publicable en la hoja visible}.
\end{aligned}
\label{eq:dimension-memory}
\end{equation}
No se afirma que toda dimensión física sea un contador. Se afirma que, en el
dominio TPK, cada memoria independiente que debe sobrevivir fuerza una
coordenada del lift.

\begin{narrativebox}
La dimensión aparece cuando la proyección ya no puede conservar toda la
historia. El círculo es suficiente para decir dónde está la fase; no es
suficiente para decir cuántas veces volvió, qué hoja transportó, qué intervalo
se contrajo o qué acarreo sobrevivió.
\end{narrativebox}

\paragraph{Ensamblaje de cilindros, espirales y trayectorias helicoidales}

Los cilindros coinductivos no son tubos físicos añadidos a una espiral. Son
intervalos encajados que conservan prefijos y estrechan el valor. La espiral
describe contracción transversal; la hélice describe memoria longitudinal;
el cilindro describe el conjunto de valores compatibles a una profundidad
finita. En el límite, la intersección de cilindros selecciona una sección
única.

La imagen espacial reúne, sin confundir:
\[
\begin{array}{ccl}
\text{ángulo} &\leftrightarrow& \text{fase visible},\\
\text{altura} &\leftrightarrow& \text{retorno/memoria},\\
\text{radio de cilindro} &\leftrightarrow& \text{incertidumbre de refinamiento},\\
\text{hoja} &\leftrightarrow& \text{orientación bidireccional}.
\end{array}
\]
La tasa \(3^{-54}\) contrae el cilindro por retorno; no se suma al ángulo y
no reemplaza el acarreo real \(\epsR\).

\subsubsection{Cierres de fase de órdenes \(2\pi\) y \(4\pi\)}

En la proyección circular, \(2\pi\) restaura la fase. En un lift orientado
con dos hojas, una vuelta puede cerrar en la hoja conjugada y necesitar una
segunda vuelta para restaurar la orientación:
\[
U_{2\pi}=-I,\qquad U_{4\pi}=I.
\]
La estructura coincide con la periodicidad espinorial cuando se aplica el
funtor a una representación \(SU(2)\); no se identifica el lift TPK con
\(SU(2)\) sin ese mapa.

La lectura areal es compatible:
\[
\mathcal A(2\pi)=\hfull,\qquad
\mathcal A(4\pi)=2\hfull.
\]
En una banda de Möbius, un circuito visible invierte hoja y dos circuitos
restauran la orientación. El \(2\pi/4\pi\) no es una duplicación accidental;
es la diferencia entre retorno de fase y retorno del estado enriquecido.


\section{Elipse orientada y variables conjugadas: incertidumbre mínima y oscilación electromagnética}
\subsection{Variables simplécticas conjugadas y determinación de la incertidumbre mínima}
% Fragmento público generado desde propietarios íntegros.
% Residencia: 52.6.1.
% Fuente: ../incoming/radion_monografia_20260827/manuscrito/sections/07_incertidumbre_oscilador.tex:1-111
\noindent La forma simpléctica coordina acción, incertidumbre y dinámica oscilatoria sin alterar la genealogía del estado.\par\medskip

\subsubsection{Realización del estado en un espacio de Hilbert}

\paragraph{Matriz de covarianza mínima}

Una vez construida la elipse de acción, puede realizarse en un par de
operadores conjugados
\[
[\widehat Q,\widehat P]=i\hret.
\]
Definimos
\begin{equation}
V_\eta=\frac{\hret}{2}
\begin{pmatrix}
e^\eta&0\\[2pt]0&e^{-\eta}
\end{pmatrix}.
\label{eq:covariance}
\end{equation}
Su determinante es
\begin{equation}
\det V_\eta=\frac{\hret^2}{4}.
\label{eq:cov-det}
\end{equation}
Por consiguiente,
\begin{equation}
(\Delta Q)^2=\frac{\hret}{2}e^\eta,
\qquad
(\Delta P)^2=\frac{\hret}{2}e^{-\eta},
\qquad
\Delta Q\,\Delta P=\frac{\hret}{2}.
\label{eq:uncertainty}
\end{equation}
La desigualdad de Robertson--Schrödinger se satura.

El aniquilador comprimido
\begin{equation}
\widehat a_\eta=
\frac{e^{-\eta/2}\widehat Q+i e^{\eta/2}\widehat P}
{\sqrt{2\hret}}
\label{eq:a-eta}
\end{equation}
satisface \([\widehat a_\eta,\widehat a_\eta^\dagger]=1\). Su vacío tiene
función de onda
\begin{equation}
\psi_\eta(Q)=
(\pi\hret e^\eta)^{-1/4}
\exp\left(-\frac{Q^2}{2\hret e^\eta}\right).
\label{eq:gaussian}
\end{equation}

\paragraph{Elipse como nivel del funcional de Mahalanobis}

El conjunto
\begin{equation}
\mathbb E_{\mathrm{act}}
=\{x\in\R^2:x^{\mathsf T}V_\eta^{-1}x\le4\}
\label{eq:mahalanobis}
\end{equation}
tiene semiejes
\begin{equation}
a_S=2\Delta Q,\qquad b_S=2\Delta P.
\label{eq:two-sigma}
\end{equation}
Su área es
\[
4\pi\sqrt{\det V_\eta}=2\pi\hret=\hfull.
\]
La elipse HMT construida antes coincide exactamente con el contorno de dos
desviaciones del estado gaussiano mínimo.

\begin{proposition}[Cadena única del cociente de forma]
Tras declarar las interfaces espectral, simpléctica y de Hilbert,
\begin{equation}
\boxed{
e^{-\eta}
=\sqrt{\frac{A-C^*}{A+C^*}}
=\frac{b_1}{a_1}
=\frac{b_S}{a_S}
=\frac{\Delta P}{\Delta Q}.}
\label{eq:shape-chain-hilbert}
\end{equation}
\end{proposition}

La elipse de acción corresponde al nivel clásico \(H=\hret\omega\) y al
contorno \(2\sigma\) del vacío comprimido. No debe confundirse con la energía
media del vacío, \(\hret\omega/2\). Esta diferencia de factor dos es un
control físico obligatorio.

\paragraph{Condiciones de invariancia de la elipse como objeto operatorio}

En la formulación convencional, una elipse de incertidumbre puede verse como
un recurso gráfico para representar una matriz de covarianza. Aquí el orden
causal es inverso: la elipse positiva ya existe como publicación de
\((A,C^*,\hret)\); el espacio de Hilbert la reconoce mediante
\eqref{eq:covariance}. La realidad ontológica reclamada no significa que una
partícula clásica recorra físicamente un óvalo en el plano \((Q,P)\). Significa
que el estado conserva una forma positiva, una acción total y un cociente de
orientación antes de que una teoría probabilista los interprete.

La función de onda no crea la forma; la representa. La incertidumbre no crea
\(\hbar\); publica la imposibilidad de comprimir simultáneamente las dos
direcciones sin alterar el producto simpléctico.

\begin{narrativebox}
La lectura probabilista es una realización posible de la elipse, no su causa.
El objeto primario es la conservación conjunta de producto y orientación. La
probabilidad aparece cuando esa geometría se publica como covarianza de un
estado sobre Hilbert.
\end{narrativebox}


\subsection{Intercambio \(\mu\leftrightarrow\varepsilon\), inversión \(Z\leftrightarrow Z^{-1}\) y conservación de la velocidad causal}
% Fragmento público generado desde propietarios íntegros.
% Residencia: 52.6.2.
% Fuente: ../incoming/radion_monografia_20260827/manuscrito/sections/07_incertidumbre_oscilador.tex:156-197
\paragraph{Intercambio energético entre los sectores eléctrico y magnético}

Los dos términos de \eqref{eq:HLC} se convierten en
\begin{equation}
U_E(\theta)=\hret\omega\cos^2\theta,
\qquad
U_B(\theta)=\hret\omega\sin^2\theta.
\label{eq:energy-exchange}
\end{equation}
Por tanto,
\[
U_E+U_B=\hret\omega.
\]
La polaridad eléctrica--magnética no consiste en etiquetar arbitrariamente
los focos. Consiste en dos reservas conjugadas que se intercambian durante la
órbita manteniendo constantes frecuencia y acción total.

El cociente común se prolonga:
\begin{equation}
\boxed{
e^{-\eta}
=\frac{b_S}{a_S}
=\frac{\Delta P}{\Delta Q}
=\frac{Z_\eta}{Z_*}
=\sqrt{\frac{A-C^*}{A+C^*}}.}
\label{eq:shape-chain-LC}
\end{equation}
Los productos conservados son
\[
a_Sb_S=2\hret,\qquad
\det V_\eta=\frac{\hret^2}{4},
\qquad
L_\eta C_\eta=\omega^{-2}.
\]

\begin{exactbox}
El oscilador LC es la realización física mínima de la doble lectura: el
producto conserva acción y frecuencia; el cociente conserva orientación e
impedancia. Un ciclo completo intercambia electricidad y magnetismo sin
perder el área \(\oint P\,\diff Q=\hfull\).
\end{exactbox}


\subsection{Par inductivo--capacitivo \((L_\eta,C_\eta)\): invariancia de \(L_\eta C_\eta=\omega^{-2}\) y orientación de la impedancia}
% Fragmento público generado desde propietarios íntegros.
% Residencia: 52.6.3.
% Fuente: ../incoming/radion_monografia_20260827/manuscrito/sections/07_incertidumbre_oscilador.tex:112-155
\subsubsection{Oscilador inductivo--capacitivo asociado al par angular}

\paragraph{Construcción independiente de los 2 depósitos conjugados}

Dados \(\omega>0\) y una impedancia de referencia \(Z_*>0\), definimos
\begin{equation}
L_\eta=\frac{Z_*}{\omega}e^{-\eta},
\qquad
C_\eta=\frac{1}{Z_*\omega}e^\eta.
\label{eq:LC-def}
\end{equation}
Entonces
\begin{equation}
L_\eta C_\eta=\omega^{-2},
\qquad
\sqrt{\frac{L_\eta}{C_\eta}}=Z_*e^{-\eta}.
\label{eq:LC-invariants}
\end{equation}
La frecuencia permanece fija; la impedancia conserva la orientación de la
forma.

El Hamiltoniano es
\begin{equation}
H_{LC}=\frac{q^2}{2C_\eta}+\frac{\Phi^2}{2L_\eta}.
\label{eq:HLC}
\end{equation}
Con las variables normalizadas
\begin{equation}
Q=\sqrt{Z_*}\,q,\qquad
P=\frac{\Phi}{\sqrt{Z_*}},
\label{eq:LC-map}
\end{equation}
se obtiene
\begin{equation}
H_{LC}=\frac\omega2
\left(e^{-\eta}Q^2+e^\eta P^2\right).
\label{eq:HLC-QP}
\end{equation}
Sobre \eqref{eq:elipse-param},
\begin{equation}
H_{LC}=\hret\omega.
\label{eq:LC-level}
\end{equation}

% Fuente: ../incoming/radion_monografia_20260827/manuscrito/sections/07_incertidumbre_oscilador.tex:198-239
\paragraph{Oscilador LC como unidad dinámica elemental}

El oscilador armónico atraviesa casi toda la física porque linealiza el
entorno de un equilibrio, descompone campos en modos, organiza cuantos y
permite leer propagación como superposición de fases. En el radión, esa
universalidad recibe una genealogía concreta: el modo no empieza en una
ecuación diferencial externa, sino en una forma positiva con acción de vuelta
fija. La dinámica convencional aparece al elegir un reloj \(\omega\) y un
transductor \(Z_*\).

La frecuencia no altera la elipse de acción: cambia la energía del nivel
mediante \(E=\hret\omega\). Esta separación es crucial. La geometría fija la
acción; el reloj convierte acción en energía.

\subsubsection{Reversión temporal y propagaciones conjugadas}

Sea \(H=H^*\) y sea \(J_E\) una estructura compleja con \(J_E^2=-I\).
Supongamos un involutivo \(C\) tal que
\[
CJ_EC^{-1}=-J_E,\qquad CHC^{-1}=H.
\]
La evolución
\begin{equation}
U(t)=\exp\left(-\frac{J_EHt}{\hret}\right)
\label{eq:evolution-J}
\end{equation}
satisface exactamente
\begin{equation}
CU(t)C^{-1}=U(-t).
\label{eq:time-reversal}
\end{equation}
Ésta es la forma operatoria mínima de la bidireccionalidad: la conjugación
intercambia las prolongaciones temporalmente opuestas.

La teoría del absorbedor de Wheeler--Feynman y la lectura de antipartículas
como trayectorias temporalmente conjugadas ofrecen una realización histórica
posterior. La identidad \eqref{eq:time-reversal} está demostrada; convertirla
en una teoría completa de propagadores avanzados y retardados exige declarar
acción, condiciones de frontera, causalidad e interacciones del absorbedor.
La monografía conserva la intuición física sin convertir una correspondencia
en una igualdad de teorías.



\section{3 regímenes del TRIT: polaridades eléctrica--magnética y elástica--rotacional}
\subsection{Coordinación de las realizaciones elíptica, parabólica e hiperbólica}
% Fragmento público generado desde propietarios íntegros.
% Residencia: 52.7.1.
% Fuente: ../incoming/radion_monografia_20260827/manuscrito/sections/08_complejo_electromagnetismo.tex:1-58
\noindent Los 3 regímenes del TRIT admiten realizaciones compleja, electromagnética y mecánica con tipos diferenciados.\par\medskip

\subsubsection{Operador cuadrático de raíz interna de \(-1\)}

\paragraph{Representación exponencial del régimen elíptico \(J^2=-I\)}

Con \(J^2=-I\),
\begin{equation}
T_{\mathrm{ell}}(x,y)=e^{-xI-yJ}
=e^{-x}(\cos y\,I-\sin y\,J).
\label{eq:elliptic-exp}
\end{equation}
Bajo la identificación generada por \(J\),
\[
T_{\mathrm{ell}}(x,y)\longleftrightarrow e^{-x-iy}.
\]
El escalar \(x\) controla dilatación; \(y\) controla giro orientado. La forma
\(a+bi\) se reconoce como elasticidad más rotación, no como dos cantidades
pegadas arbitrariamente.

\paragraph{Representación exponencial del régimen parabólico \(J^2=0\)}

Con \(N^2=0\), la exponencial se trunca:
\begin{equation}
T_{\mathrm{par}}(x,y)=e^{-x}(I-yN).
\label{eq:parabolic-exp}
\end{equation}
El régimen describe el umbral donde no hay ni giro periódico ni apertura
exponencial. En la lectura temporal activa es el presente euclídeo, pero
conserva memoria de las ramas que lo limitan.

\paragraph{Representación exponencial del régimen hiperbólico \(J^2=+I\)}

Con \(R^2=I\) y \(P_\pm=(I\pm R)/2\),
\begin{equation}
T_{\mathrm{hyp}}(x,y)=e^{-xI-yR}
=q_+P_++q_-P_-,
\qquad q_\pm=e^{-x\mp y}.
\label{eq:hyperbolic-exp}
\end{equation}
Los invariantes elementales son
\begin{equation}
q_+q_-=e^{-2x},
\qquad
\frac{q_-}{q_+}=e^{2y}.
\label{eq:q-product-ratio}
\end{equation}
El producto conserva el modo común; el cociente conserva la separación
orientada.

Aplicado al par angular,
\[
x=\frac{\pi A}{180},\qquad y=\frac{\pi C^*}{180}.
\]
Por ello, \(A\) es la coordenada común de la publicación y \(C^*\) su
contraste orientado. Sólo después de aplicar un lector físico se reconocen
como polos eléctricos y magnéticos.


\subsection{Separación tipada de las polaridades constitutivas y mecánicas}
% Fragmento público generado desde propietarios íntegros.
% Residencia: 52.7.2.
% Fuente: ../incoming/radion_monografia_20260827/manuscrito/sections/08_complejo_electromagnetismo.tex:59-158
\subsubsection{Operador constitutivo sobre los canales \(90/120\)}

\paragraph{Descomposición en componentes común y orientada}

Definimos
\begin{align}
\mathcal E
&=\log\bigl[(1-q_+^{90})(1-q_-^{90})\bigr]
-\log\bigl[(1-q_+^{120})(1-q_-^{120})\bigr],
\label{eq:Ecommon}\\
\mathcal B
&=\log\frac{1-q_-^{90}}{1-q_+^{90}}
-\log\frac{1-q_-^{120}}{1-q_+^{120}}.
\label{eq:Borient}
\end{align}
Bajo \(C^*\mapsto-C^*\), se intercambian \(q_+\leftrightarrow q_-\). Por
tanto, \(\mathcal E\) es par y \(\mathcal B\) es impar.

Las constitutivas normalizadas son
\begin{equation}
\widehat\mu=e^{\mathcal E+\mathcal B},
\qquad
\widehat\varepsilon=e^{\mathcal E-\mathcal B},
\qquad
\widehat Z=e^{\mathcal B},
\qquad
\widehat c=e^{-\mathcal E}.
\label{eq:constitutives}
\end{equation}
Se verifican exactamente
\begin{equation}
\widehat Z=\sqrt{\frac{\widehat\mu}{\widehat\varepsilon}},
\qquad
\widehat c=(\widehat\mu\widehat\varepsilon)^{-1/2}.
\label{eq:constitutive-identities}
\end{equation}
La conjugación produce
\begin{equation}
C^*\mapsto-C^*:\qquad
\widehat\mu\leftrightarrow\widehat\varepsilon,
\quad
\widehat Z\leftrightarrow\widehat Z^{-1},
\quad
\widehat c\mapsto\widehat c.
\label{eq:EM-conjugation}
\end{equation}

\begin{theorem}[Polaridad constitutiva]
El lector \(90/120\) convierte el modo común y el contraste orientado del
par \((A,C^*)\) en dos constitutivas conjugadas. La inversión de hoja
intercambia permeabilidad y permitividad, invierte impedancia y conserva la
velocidad normalizada.
\end{theorem}

Ésta es la formulación precisa de la intuición \enquote{parte eléctrica y
parte magnética}. No se escribe \(A=E\) ni \(C^*=B\). Se escribe el mapa
\[
(A,C^*)\to(x,y)\to(q_+,q_-)
\to(\mathcal E,\mathcal B)
\to(\widehat\varepsilon,\widehat\mu,\widehat Z,\widehat c).
\]

\paragraph{Realización elástica--rotacional de la polaridad constitutiva}

La misma información puede publicarse en el régimen complejo mediante
\begin{equation}
Z_{\mathrm{em}}=e^{\mathcal E/2}
\left(\cos\frac{\mathcal B}{2}\,I
-\sin\frac{\mathcal B}{2}\,J\right).
\label{eq:Zem}
\end{equation}
El factor \(e^{\mathcal E/2}\) es una dilatación elástica; el factor generado
por \(J\) es una rotación. La escritura compleja no añade una dimensión
imaginaria misteriosa a un mundo real previamente dado: registra dos modos
operacionales que el TRIT ya separó.

\paragraph{Transducción de la orientación en carga efectiva}

La hoja \(\sigma\) del radión fija una orientación antes de introducir una
unidad de carga. La Mecánica Dimensional dispone luego el transductor
\[
e_{\mathrm{int}}=\sqrt{\frac{2\alphah\hfull}{Z_{\mathrm{int}}}},
\qquad
Q=n_Qe_{\mathrm{int}}.
\]
La cadena causal es
\[
\text{hoja}\to\text{orientación}\to\text{constitutiva}
\to\text{transductor dimensional}\to\text{carga física}.
\]
Llamar \emph{radión} a la unidad no identifica prematuramente la orientación
con el culombio; hace visible el lugar donde el signo de carga puede residir.

\begin{narrativebox}
La electricidad y el magnetismo no se pegan al círculo desde fuera. Emergen
cuando el mismo estado se lee a través de producto y cociente: el producto
retiene el modo común, el cociente retiene la orientación. El oscilador LC
realiza el intercambio; el lector \(90/120\) publica las constitutivas; la
carta dimensional asigna unidades.
\end{narrativebox}


\section{Separación de las publicaciones \(90/120\): incidencia, registro, espectro, núcleo de Barbero y torres}
\subsection{Separación tipada de las publicaciones del sector \(90/120\)}
% Fragmento público generado desde propietarios íntegros.
% Residencia: 52.8.1.
% Fuente: ../incoming/radion_monografia_20260827/manuscrito/sections/09_sector_90120.tex:176-244
\subsubsection{Jerarquía armónica global y semigrupo \(\langle90,120\rangle\)}

Para evitar una colisión de signos, fijamos
\[
R_{\mathrm{tow}}:=-R_{\mathrm{ch}}.
\]
Entonces
\begin{equation}
T=e^{-xI+yR_{\mathrm{tow}}}
=q_-P_+^{\mathrm{tow}}+q_+P_-^{\mathrm{tow}}.
\label{eq:tower-operator}
\end{equation}
Las torres pares e impares son
\begin{align}
c_n&=\operatorname{Tr}(T^n)
=q_-^n+q_+^n
=2e^{-nx}\cosh(ny),
\label{eq:cn}\\
s_n&=\operatorname{Tr}(R_{\mathrm{tow}}T^n)
=q_-^n-q_+^n
=2e^{-nx}\sinh(ny).
\label{eq:sn}
\end{align}
Bajo \(C^*\mapsto-C^*\), \(c_n\) es par y \(s_n\) es impar. Se cumple
\begin{equation}
c_n^2-s_n^2=4e^{-2nx},
\label{eq:tower-invariant}
\end{equation}
y ambas sucesiones satisfacen la recurrencia de orden dos determinada por los
autovalores \(q_\pm\).

El semigrupo
\begin{equation}
\langle90,120\rangle=30\langle3,4\rangle
\label{eq:semigroup}
\end{equation}
es la jerarquía armónica de un único operador bicapa. No es una lista de ocho
correcciones laterales. Su completitud significa que toda razón positiva
admite una expansión convergente; no significa que una expansión obtenida a
partir del blanco sea un selector físico prospectivo.

\subsubsection{Clasificación tipada y separación de las realizaciones \(90/120\)}

\begin{longtable}{>{\bfseries}p{.19\textwidth} p{.24\textwidth} p{.45\textwidth}}
\toprule
Objeto & Operación & Qué conserva / por qué no es otro objeto\\
\midrule
\endfirsthead
\toprule
Objeto & Operación & Qué conserva / por qué no es otro objeto\\
\midrule
\endhead
Incidencia \(N_6,N_8\) & conteo finito sobre fases activas & cardinalidad y orientación; produce \(90,120\), no una serie;\\
\(-1/12\) & diferencia normalizada \eqref{eq:minus-one-twelve} & escalar de incidencia; no peso Barbero;\\
\(E_{\mathrm{dod}}\) & extractor armónico de la rueda & fase dodecafásica; no acarreo real;\\
\(h=C^*/6\) & reconstrucción de seis parejas & apertura por pareja; no normalización metrológica;\\
\(F_{\mathrm{spec}}\) & cola logarítmica \(j\ge2\) & publicación espectral completa; difiere de \(\epsR^\circ\);\\
\(\epsone\) & resta APP entre \(S_T\) y \(S_E\) & empalme real de dos lifts; no serie de Lambert;\\
\(K_{6\to8}\) & \(12\Delta S_{90}-\Delta S_{120}\) & salto marcado y evaluación conjugada; coeficiente de polo \(1/24\) comprobado; salida Barbero;\\
\(12:-1\) & doce sectores y una costura orientada & peso del núcleo; certificado posterior \eqref{eq:pole-condition}; distinto de \(-1/12\);\\
\(\langle90,120\rangle\) & semigrupo de potencias de \(T\) & jerarquía armónica global; no residuo por especie.\\
\bottomrule
\end{longtable}

\begin{thesisbox}
Todos estos objetos descienden del mismo estado APP--TRIT--TPK. Esa genealogía
común explica por qué comparten números y simetrías. Su tipado explica por qué
no pueden reemplazarse mutuamente.
\end{thesisbox}

\subsection{Conservación de incidencia y orientación entre el núcleo trítico y sus realizaciones posteriores}
% Fragmento público generado desde propietarios íntegros.
% Residencia: 52.8.2.
% Fuente: ../incoming/radion_monografia_20260827/manuscrito/sections/09_sector_90120.tex:1-30
\noindent La jerarquía tipada separa incidencia, registro, evaluación espectral, funcional de Barbero y torres armónicas.\par\medskip

\subsubsection{Incidencia trítica, registro dodecafásico y lectura sexádica}

\paragraph{Derivación combinatoria primaria de \(90\) y \(120\)}

Los conteos
\[
90=6\binom62,
\qquad
120=8\binom62
\]
nacen del selector de fase y de sus incidencias activas. La orientación
normalizada satisface
\begin{equation}
\frac{30}{120}-\frac{30}{90}=-\frac1{12}.
\label{eq:minus-one-twelve}
\end{equation}
Este escalar es un resultado de incidencia. No es el peso \(12:-1\) del
núcleo Barbero y no es un acarreo dodecafásico.

El extractor dodecafásico puede publicarse mediante una combinación de
armónicos como
\[
E_{\mathrm{dod}}(\phi)=12\cos(3\phi)-\cos(4\phi).
\]
Su dominio es la fase de la rueda; su salida es un selector armónico. Aunque
aparecen los números \(12\) y \(-1\), no opera sobre series de Lambert ni
produce \(\epsR\).


\subsection{Origen dodecafásico del peso \(12:-1\) y comprobación diofántica de su coeficiente de polo}
% Fragmento público generado desde propietarios íntegros.
% Residencia: 52.8.3.
% Fuente: ../incoming/radion_monografia_20260827/manuscrito/sections/09_sector_90120.tex:115-175
\subsubsection{Funcional hexada--octada de Barbero--Immirzi}

\paragraph{Series orientadas y salto de incidencia}

Definimos
\begin{equation}
S_{90}(q)=\frac{q^{90}}{1-q^{270}},
\qquad
S_{120}(q)=\frac{q^{120}}{1-q^{360}},
\label{eq:barbero-series}
\end{equation}
y
\[
\Delta S_m=S_m(q_-)-S_m(q_+).
\]
El funcional de inclusión marcada es
\begin{equation}
K_{6\to8}=12\Delta S_{90}-\Delta S_{120}.
\label{eq:barbero-núcleo}
\end{equation}

\paragraph{Comprobación diofántica del peso dodecafásico \(12:-1\)}

La cadena fundamental dodecafásica y su lector de retorno, construidos en
\eqref{eq:l9s102:barbero-cadena-dodecafasica}--%
\eqref{eq:l9s102:barbero-peso-generado}, producen doce términos de tipo
\(S_{90}\) y un término de frontera de tipo \(-S_{120}\). El par orientado
queda determinado antes de calcular el coeficiente de polo.

Consideremos \(aS_{90}-bS_{120}\). La comprobación del coeficiente de
\((1-q)^{-1}\) toma la forma
\begin{equation}
\frac{a}{270}-\frac{b}{360}=\frac1{24}.
\label{eq:pole-condition}
\end{equation}
Multiplicar por \(1080\) da
\[
4a-3b=45.
\]
La única costura de la vuelta fundamental tiene multiplicidad absoluta uno.
Su verificación diofántica corresponde al contratérmino entero mínimo
\(|b|=1\). Para \(b=1\), \(a=12\); para \(b=-1\), \(a=21/2\) no es entero.
Por tanto,
\begin{equation}
\boxed{a:b=12:1,\qquad aS_{90}-bS_{120}=12S_{90}-S_{120}.}
\label{eq:12minus1}
\end{equation}

\begin{proposition}[Unicidad de la comprobación diofántica]
El par generado \((12,1)\) es la única solución entera de
\eqref{eq:pole-condition} con canal principal positivo y contratérmino
mínimo \(|b|=1\).
\end{proposition}

La realización hexada--octada conserva los canales \(90/120\)
previamente construidos; la cadena dodecafásica aporta las multiplicidades
y la orientación de frontera; el lector de retorno asigna las series a los
generadores. Estas operaciones conjuntas producen \(12S_{90}-S_{120}\).
La ecuación \eqref{eq:pole-condition}, la positividad y la minimalidad
conservan su función de comprobación aritmética posterior.

La cadena Barbero es posterior:
\[
(A,C^*)\to(x,y)\to q_\pm
\to K_{6\to8}\to\gamma_{\mathrm{BI}},
\]
con evaluación
\[
\gamma_{\mathrm{BI}}=0.274007672694459\ldots.
\]
Barbero--Immirzi no genera retroactivamente \(C^*\), \(\Delta_4\) ni
\(\epsR\).


\subsection{Extracción de la componente finita posterior al sistema dodecafásico orientado \(R_{12}\): \(C^*\mapsto C^*/6\), resto \(j\ge2\) y reconstrucción de \(6\) secciones}
% Fragmento público generado desde propietarios íntegros.
% Residencia: 52.8.4.
% Fuente: ../incoming/radion_monografia_20260827/manuscrito/sections/09_sector_90120.tex:31-54
\paragraph{Extracción de la componente \(C^*/6\) del registro dodecafásico}

Los doce sectores se agrupan en seis parejas opuestas:
\begin{equation}
p_i=e_i+e_{i+6},\qquad
a_i=e_i-e_{i+6},\qquad i=1,\ldots,6.
\label{eq:sexadic-pairs}
\end{equation}
La reconstrucción tiene perfil \(1+5+6=12\). La apertura por pareja es
\begin{equation}
h=\frac{C^*}{6}.
\label{eq:h-C6}
\end{equation}
No es una normalización exterior ni la sexta parte de un blanco: es el índice
interno de la reconstrucción sexádica del registro \(\Rstate\).

Definimos
\begin{equation}
q_{h,-}=\exp\left[-\frac\pi{180}(A-h)\right],
\qquad
q_{h,+}=\exp\left[-\frac\pi{180}(A+h)\right].
\label{eq:q-h}
\end{equation}


\subsection{Descomposición cuantitativa de las publicaciones de acarreo y espectral: \(F_{\mathrm{spec}}=\varepsilon_R^\circ+\chi_R\)}
% Fragmento público generado desde propietarios íntegros.
% Residencia: 52.8.5.
% Fuente: ../incoming/radion_monografia_20260827/manuscrito/sections/09_sector_90120.tex:55-114
\subsubsection{Publicación espectral de la cola de órdenes \(j\ge2\)}

\paragraph{Definición del funcional espectral}

Para \(m\in\{90,120\}\), sea
\begin{equation}
L_m^{(2)}(q)=\sum_{j\ge2}\frac{q^{mj}}{j}
=-\log(1-q^m)-q^m.
\label{eq:L2}
\end{equation}
La lectura de una pareja es
\begin{equation}
T_h^{(2)}=\frac{180}{\pi}
\Bigl[
L_{90}^{(2)}(q_{h,-})-L_{90}^{(2)}(q_{h,+})
-L_{120}^{(2)}(q_{h,-})+L_{120}^{(2)}(q_{h,+})
\Bigr].
\label{eq:T-h}
\end{equation}
La reconstrucción de las seis parejas da
\begin{equation}
F_{\mathrm{spec}}=6T_h^{(2)}
=5.1499235153094574410\ldots\times10^{-8}\,{}^\circ.
\label{eq:Fspec}
\end{equation}

El acarreo del radión es
\[
F_{\mathrm{acarreo}}=\epsR^\circ
=5.1383016758575282072\ldots\times10^{-8}\,{}^\circ.
\]
No coinciden:
\begin{equation}
\chi_R:=F_{\mathrm{spec}}-F_{\mathrm{acarreo}}
=1.16218394519292338\ldots\times10^{-10}\,{}^\circ.
\label{eq:chiR}
\end{equation}

La igualdad
\[
F_{\mathrm{acarreo}}=F_{\mathrm{spec}}-\chi_R
\]
es un cambio de carta válido una vez construidas ambas salidas. No debe
presentarse como generación independiente de \(\epsR^\circ\) si \(\chi_R\) se
define por la propia resta. La cola está cerrada y es significativa; lo falso
es identificarla directamente con el acarreo.

\begin{falsifier}
Tras convertir la cola expresada en grados a la unidad angular interna,
\[
\frac{\pi}{180^\circ}F_{\mathrm{spec}}-\epsone
=2.0283936357433839\ldots\times10^{-12}\ne0.
\]
Por tanto, ni el log completo ni la cola \(j\ge2\), aun después de aplicar la
carta de unidades, son \(\epsone\). Una
segunda factorización espectral autónoma necesitaría generar \(\chi_R\) sin
usar la diferencia como definición; el cierre APP--TPK del radión no depende
de esa promoción.
\end{falsifier}



\section{Holonomía orientada y geometría física: realizaciones posteriores de Minkowski, Hodge, Lorentz y Cartan--Holst}
% Fragmento público generado desde propietarios íntegros.
% Residencia: 52.9.
% Fuente: ../incoming/radion_monografia_20260827/manuscrito/sections/10_minkowski_hodge_lorentz.tex:1-55
\noindent La holonomía orientada admite realizaciones causales y geométricas posteriores en Minkowski, Hodge, Lorentz y Cartan--Holst.\par\medskip

\paragraph{Construcción de la métrica lorentziana desde el bloque APP}

\paragraph{Matriz APP de parámetros \(7\) y \(2\)}

La matriz central
\begin{equation}
B_c=\begin{pmatrix}7&2\\2&7\end{pmatrix}
=9P_++5P_-
\label{eq:Bc}
\end{equation}
se normaliza por su determinante:
\begin{equation}
M_c=\frac{B_c}{\sqrt{45}}
=\exp\left(\frac12\log\frac95\,R\right)
\in SO^+(1,1).
\label{eq:Mc}
\end{equation}
La rapidez interna es
\[
\eta_c=\frac12\log\frac95.
\]
La publicación proyectiva asociada satisface
\begin{equation}
P_B^n=P_++\left(\frac59\right)^nP_-.
\label{eq:PB-contract}
\end{equation}
La misma razón \(5/9\) que aparece como contracción espectral se reconoce aquí
como separación de un modo estable y otro contractivo. La identidad es
operatoria; no convierte todas las apariciones de \(5/9\) en el mismo mapa.

\paragraph{Incidencia \(6\to8\) y cociente de firma causal}

Sea \(M\) la incidencia tipada entre seis caras y ocho vértices. La relación
espectral recuperada es
\begin{equation}
MM^{\mathsf T}=12P_u+4P_-,
\label{eq:MMt}
\end{equation}
donde \(P_u\) proyecta sobre el modo uniforme y \(P_-\) sobre el subespacio
orientado. El cociente superviviente de dimensión cuatro se descompone como
\begin{equation}
Q_4\simeq\R u\oplus E_-.
\label{eq:radion-Q4}
\end{equation}
Declarando negativo el modo uniforme y positivos los tres modos orientados,
se obtiene
\begin{equation}
B_{\mathrm{caus}}=\operatorname{diag}(-1,1,1,1).
\label{eq:minkowski}
\end{equation}
La pantalla de Minkowski no selecciona el estado HMT. Es la realización
causal posterior del cociente que la incidencia ya construyó.


\subsection{Dualidad de Hodge sobre \(2\)-formas lorentzianas}
% Fragmento público generado desde propietarios íntegros.
% Residencia: 52.9.1.
% Fuente: ../incoming/radion_monografia_20260827/manuscrito/sections/10_minkowski_hodge_lorentz.tex:56-100
\subsubsection{Dualidad de Hodge y operador de cuarto de giro}

Fijadas la firma \((-+++ )\) y una orientación, la estrella de Hodge sobre
dos-formas satisface
\begin{equation}
\star:\Lambda^2Q_4^*\to\Lambda^2Q_4^*,
\qquad \star^2=-I.
\label{eq:hodge-star}
\end{equation}
De este modo, el cuarto de giro que en el plano de fase se expresaba mediante
\(J^2=-I\) reaparece en el espacio de campos como estructura compleja sobre
\(\Lambda^2\). Los dominios son diferentes, pero la genealogía de orientación
se conserva mediante el funtor.

Sea \(A_{\mathrm{em}}\) un potencial y
\begin{equation}
F=\diff A_{\mathrm{em}}.
\label{eq:F-dA}
\end{equation}
La descomposición de \(F\) respecto a un observador temporal produce campos
eléctrico y magnético; \(\star F\) intercambia sus papeles con el signo fijado
por la orientación y la firma. La polaridad constitutiva de
\eqref{eq:EM-conjugation} adquiere así una realización diferencial.

\paragraph{Cadena de transducción electromagnética y causal}

El transporte que debe mantenerse visible es
\begin{equation}
\Gamma_9
\longrightarrow\text{hélice de memoria}
\longrightarrow(Q_4,B_{\mathrm{caus}},\star)
\longrightarrow A_{\mathrm{em}}
\longrightarrow F=\diff A_{\mathrm{em}}
\longrightarrow
m\nabla_u u=q(\iota_uF)^\sharp.
\label{eq:lorentz-chain}
\end{equation}
Cada flecha añade estructura: la holonomía conserva retorno; el cociente fija
causalidad; Hodge fija dualidad; el potencial fija el campo; el acoplamiento
con carga produce dinámica.

La ecuación final es la fuerza de Lorentz en forma geométrica. La hélice TPK
es su antecedente de memoria, no una órbita física ya movida por un campo. La
dinámica aparece sólo al introducir \(F\), \(q\) y la conexión \(\nabla\).


\subsection{Separación entre memoria torsional interna y realización geométrica de Cartan--Holst}
% Fragmento público generado desde propietarios íntegros.
% Residencia: 52.9.2.
% Fuente: ../incoming/radion_monografia_20260827/manuscrito/sections/10_minkowski_hodge_lorentz.tex:101-186
\subsubsection{Realización lorentziana de las trayectorias helicoidales}

En un campo magnético uniforme, una partícula de carga \(q\), masa \(m\) y
factor relativista \(\gamma\) posee
\begin{equation}
\omega_c=\frac{|q|B}{\gamma m},
\qquad
r_L=\frac{p_\perp}{|q|B},
\qquad
\Delta z=\frac{2\pi p_\parallel}{|q|B}.
\label{eq:lorentz-helix}
\end{equation}
La fase ciclotrónica cierra mientras el movimiento paralelo avanza: la órbita
es helicoidal. El paralelismo con \eqref{eq:H9-helix} es estructural y puede
formalizarse por un transductor que preserve fase, orientación y paso. No se
identifica el cociente \(q_9\) con la coordenada espacial \(z\); se reconoce
la misma arquitectura de retorno más avance.

La fuerza de Lorentz no se deriva de la imagen de una espiral por semejanza.
Se deriva al completar \eqref{eq:lorentz-chain}. La imagen ayuda a comprender
por qué fase circular y memoria longitudinal son compatibles; la forma de
campo determina la dinámica.

\subsubsection{Confluencia entre acción, área, información y gravedad}

\paragraph{Realización Cartan--Holst de la holonomía}

La promoción gravitatoria se expresa mediante un morfismo de holonomías
\begin{equation}
\operatorname{Hol}_{\mathcal A}
\bigl(\mathfrak R_{\mathrm{grav}}(\gamma)\bigr)
=\rho_C\bigl(\operatorname{Hol}_{\TPK}(\gamma)\bigr).
\label{eq:holonomy-gravity}
\end{equation}
El lado derecho contiene la memoria de ruta; \(\rho_C\) la realiza en una
conexión gravitatoria. La torsión de Cartan
\[
T^I=D_\omega e^I
\]
aparece en el codominio. La torsión global \(\Delta_4\) es una coordenada
upstream que puede alimentar el mapa, pero no se renombra como \(T^I\) sin
\eqref{eq:holonomy-gravity}.

\paragraph{Transducción del área de acción en área informacional}

El radión fija un cuanto de acción por fase. El lector hexada--octada de
Barbero usa posteriormente los canales \(q_\pm\), la incidencia \(90/120\) y
el funcional \eqref{eq:barbero-núcleo}. La dirección causal es
\[
(A,C^*,\hret)\to q_\pm\to K_{6\to8}
\to\gamma_{\mathrm{BI}}\to\text{área/información}.
\]
El alfabeto trítico aporta la unidad de entropía \(\log3\) y un corte de
\(81\) canales aporta \(81\log3\). Esta entropía no se confunde con
\(\log2\), con \(\log\varphi\), con una entropía microcanónica ni con una
entropía de von Neumann sin especificar el estado.

La síntesis física es afirmativa pero tipada: el radión aporta la unidad de
acción orientada que una holonomía puede transportar; Barbero publica la
normalización de área; la pantalla causal y Hodge permiten realizar campos;
Cartan--Holst recibe la holonomía en geometría gravitatoria.

\begin{statusbox}
La ecuación escalar del radión no es por sí sola una ecuación de Einstein. Su
aportación a la gravedad es ser una sección exacta de la arquitectura de
acción, orientación y holonomía que la realización gravitatoria consume.
\end{statusbox}

\subsubsection{Transformación de Wick como cambio de lector analítico}

En la formulación convencional, la rotación de Wick relaciona una variable
temporal lorentziana con una coordenada euclídea mediante una continuación
compleja. En HMT, la raíz \(J^2=-I\), el involutivo \(R^2=I\) y el umbral
\(N^2=0\) ya existen antes. El cambio de lector
\[
J\longleftrightarrow R
\]
transporta una pareja común/orientada entre publicación elíptica y
publicación hiperbólica. La continuación compleja reconoce después esa
operación en coordenadas analíticas.

Esto explica por qué círculo, elipse de acción e interior hiperbólico pueden
pertenecer al mismo estado sin ser la misma métrica. El círculo usa \(J\) para
la fase; la forma de la elipse usa un compresión simpléctica generado por \(R\); Klein mide
el interior mediante razón doble; la pantalla de Minkowski realiza causalidad
en otro codominio.


\section{Elipse, retículo y fase de camino: carácter modular, órbitas y propagación}
\subsection{Red rectangular de períodos, involución \(\tau\mapsto-1/\tau\) e invariancia de \(j\)}
% Fragmento público generado desde propietarios íntegros.
% Residencia: 52.10.1.
% Fuente: ../incoming/radion_monografia_20260827/manuscrito/sections/11_modularidad_orbitas.tex:1-84
\noindent Los caracteres modulares, orbitales y de camino conservan invariantes distintos de una misma forma orientada.\par\medskip

\subsubsection{Toro complejo asociado a la elipse de acción}

\paragraph{Retículo rectangular de períodos}

La elipse por sí sola no es un toro. Declaramos el retículo
\begin{equation}
\Lambda_\eta=a_S\Z+i,b_S\Z
\label{eq:lattice}
\end{equation}
y el cociente \(\C/\Lambda_\eta\). Su módulo es
\begin{equation}
\tau_\eta=i\frac{b_S}{a_S}=ie^{-\eta}
=i\,0.741048144159375160795483663369300\ldots.
\label{eq:tau-eta}
\end{equation}
La inversión bidireccional \(\eta\mapsto-\eta\) intercambia los ejes:
\begin{equation}
\tau_{-\eta}=ie^\eta=-\frac1{\tau_\eta}.
\label{eq:modular-S}
\end{equation}
Es exactamente la transformación modular \(S:\tau\mapsto-1/\tau\).

\paragraph{Invariancia de la función modular \(j\)}

Por invariancia modular,
\begin{equation}
\boxed{j(\tau_{-\eta})=j(\tau_\eta).}
\label{eq:j-invariance}
\end{equation}
En la normalización \(j(i)=1728\),
\[
j(\tau_\eta)
=5597.43843932408729830097089254768\ldots.
\]
La función moonshine normalizada \(J=j-744\) da
\[
J(\tau_\eta)
=4853.43843932408729830097089254768\ldots.
\]
Para el bloque APP basal,
\[
\tau_0=i\frac{\sqrt5}{3},
\qquad
j(\tau_0)=5369.48832024674306021916882626352\ldots.
\]

\begin{proposition}[Lo que conserva la clase modular]
La involución de hoja invierte \(\eta\), pero \(j\) conserva la clase del
toro. Los focos y el etiquetado de ejes retienen la orientación que \(j\)
olvida.
\end{proposition}

Este resultado conecta de manera exacta la doble vía con la modularidad una
vez declarado el cociente. No autoriza a hacer actuar directamente al
Monstruo sobre la carga, los dígitos o los focos. El corredor excepcional
\[
\mathrm{APP}\to\mathrm{Paley}\to\mathrm{Witt}\to\mathrm{Golay}
\to A_2^{12}\to\mathrm{Leech}\to V^\natural
\to\mathrm{Monster}\to\mathrm{Moonshine}
\]
es otra proyección del mismo estado enriquecido. Para identificar una acción
concreta sobre el radión se necesita un entrelazador que preserve sus
coordenadas pertinentes.

\paragraph{Cadena de invariantes del carácter de forma}

Reuniendo las realizaciones declaradas,
\begin{equation}
\boxed{
r_\eta=e^{-\eta}
=\sqrt{\frac{A-C^*}{A+C^*}}
=\frac{b_1}{a_1}
=\frac{b_S}{a_S}
=\frac{\Delta P}{\Delta Q}
=\frac{Z_\eta}{Z_*}
=\Im\tau_\eta.}
\label{eq:shape-master}
\end{equation}
Una sola razón recorre geometría espectral, acción, covarianza, impedancia y
módulo. Cada igualdad se apoya en un mapa explícito; ninguna depende de una
semejanza decimal.


\subsection{Compatibilidad entre el carácter modular y la propagación sobre rutas orientadas}
% Fragmento público generado desde propietarios íntegros.
% Residencia: 52.10.2.
% Fuente: ../incoming/radion_monografia_20260827/manuscrito/sections/11_modularidad_orbitas.tex:85-187
\subsubsection{Conservación areal en las órbitas de Kepler y Sommerfeld}

\paragraph{Equivalencia areal de las realizaciones orbital y oscilatoria}

La elipse de Kepler vive en espacio de configuración y responde a un potencial
central. La elipse del radión vive en espacio de fases y responde a una forma
de acción. Son objetos diferentes. Comparten una ley areal porque ambas
realizaciones poseen una dos-forma conservada.

Para una órbita kepleriana,
\begin{equation}
r(\phi)=\frac{p}{1+e\cos(\phi-\phi_0)},
\qquad
\frac{\diff A}{\diff t}=\frac{\ell}{2m}.
\label{eq:kepler}
\end{equation}
La tercera ley toma la forma
\[
T^2=\frac{4\pi^2\mu}{K}a^3.
\]
Para el oscilador en fase,
\begin{equation}
J=\oint p\,\diff q=\frac{2\pi E}{\omega}.
\label{eq:osc-action}
\end{equation}
Si la celda elemental es \(J=\hfull\), entonces
\begin{equation}
E=\hret\omega,\qquad \frac{J}{2\pi}=\hret.
\label{eq:E-hbarw}
\end{equation}
Ésta es la convergencia limpia entre acción de Planck, frecuencia angular y
radión.

Sommerfeld cuantiza integrales de acción sobre ciclos. La lectura HMT añade
la memoria de qué ciclo, hoja y retorno sobreviven. En una banda de Möbius,
un bucle visible puede portar \(\hfull/2\) y dos bucles restaurar \(\hfull\).
Ese factor topológico no se suma al índice de Maslov de una cuantización EBK;
pertenecen a correcciones con propietarios distintos.

\paragraph{Holonomía orbital y propagaciones avanzada y retardada}

Una órbita que acumula holonomía satisface una condición del tipo
\begin{equation}
\epsilon_\gamma\exp\left(\frac{iJ_\gamma}{\hret}\right)=1.
\label{eq:holonomy-action}
\end{equation}
El signo \(\epsilon_\gamma\) conserva la hoja; \(J_\gamma\) conserva la
acción. El adelanto o retardo apsidal se interpreta como desacoplamiento entre
el cierre angular visible y el retorno completo. El radión suministra la
unidad orientada con la que se mide esa diferencia.

\subsubsection{Composición de amplitudes de Schrödinger y suma de caminos}

\paragraph{Cuatro realizaciones de una misma fase}

Una vez fijado \(\hret\), distintas ecuaciones convencionales reconocen
aspectos del mismo transporte:
\begin{enumerate}
\item Schrödinger publica la fase mediante
\(i\hret\partial_t\psi=H\psi\).
\item Pauli añade la doble hoja de espín y su acoplamiento magnético.
\item Dirac linealiza la relación energía--momento y organiza partícula y
antipartícula en una representación causal.
\item La suma de caminos asigna a cada historia la fase
\(\exp(iS[\gamma]/\hret)\).
\end{enumerate}
Estas teorías no generan retrospectivamente \(\hret\). Reciben la unidad de
acción y realizan diferentes contenidos del estado enriquecido.

\paragraph{Funcional de fase sobre el espacio de caminos}

En la formulación de Feynman,
\begin{equation}
\mathcal K(b,a)=\int_{\gamma:a\to b}
\exp\left(\frac{i}{\hret}S[\gamma]\right)\mathcal D\gamma.
\label{eq:path-integral}
\end{equation}
Un incremento de acción \(\hret\) cambia la fase en un radián. Un incremento
\(\hfull\) la cambia en una vuelta. El teorema
\(\mathcal A(\theta)=\hret\theta\) convierte esta relación en geometría de
área: la fase de un camino puede leerse como acción barrida en unidades de
radión.

La inversión \(C_M=-\operatorname{rev}\) sobre palabras dodecafásicas y la
identidad \eqref{eq:time-reversal} organizan rutas conjugadas. En una
realización de absorbedor se separan
\[
G_{\mathrm{sym}}=\tfrac12(G_{\mathrm{ret}}+G_{\mathrm{adv}}),
\qquad
G_{\mathrm{odd}}=G_{\mathrm{ret}}-G_{\mathrm{adv}}.
\]
El primer sector es par bajo intercambio retardado--avanzado; el segundo es
impar. Esto reproduce la separación entre energía/forma par y acción
orientada impar del radión cuando el transductor preserva frontera y forma
simpléctica.

\begin{statusbox}
La teoría clásica del absorbedor y la interpretación Feynman--Stückelberg de
antipartículas son realizaciones históricas diferentes. La primera usa una
interacción temporalmente simétrica; la segunda reorganiza propagación en la
teoría relativista. HMT aporta un antecedente común de reversión, pero no las
funde sin un mapa de propagadores.
\end{statusbox}


\section{Radical nonádico y frontera decimal: conservación del acarreo del radión}
\subsection{El acarreo del radión en la completación decimal}
% Fragmento público generado desde propietarios íntegros.
% Residencia: 52.11.1.
% Fuente: ../incoming/radion_monografia_20260827/manuscrito/sections/12_aritmetica_borde.tex:1-159
\noindent La frontera nonádica conserva el defecto suma--producto, la vacancia y el acarreo de completación.\par\medskip

\subsubsection{Radical \(3\! -\! 6\! -\! 9\) del borde}

\paragraph{Generación y nilpotencia del radical nonádico}

En el anillo \(\Z/9\Z\), el ideal
\begin{equation}
\mathfrak m=(3)=\{0,3,6\}
\label{eq:radical-369}
\end{equation}
satisface
\begin{equation}
\mathfrak m^2=(9)=0.
\label{eq:m-square-zero}
\end{equation}
Es el sector nilpotente del borde nonádico. La órbita visible de sus tres
marcas se publica como
\[
369\longrightarrow693\longrightarrow936\longrightarrow369.
\]
El retorno no implica que todos los estados sean iguales: la permutación
visible cierra mientras los cociclos de ruta pueden avanzar.

La doble convención
\[
9\equiv0\pmod9,
\qquad
\operatorname{dr}_9(9)=9
\]
expresa dos lectores: cero residual algebraico y cierre visible de raíz
digital. Por eso un nueve puede absorber torsión de borde sin convertirse en
un cero específico de la función zeta.

\paragraph{Defecto suma--producto \(459-405=54\)}

Las hojas APP producen totales complementarios
\[
S_+=405,\qquad S_\times=459,\qquad S_\times-S_+=54.
\]
El \(54\) es un defecto exacto de las tablas, y reaparece como exponente de
contracción por la identidad \(9\cdot6=54\). Sin embargo, cada aparición
requiere su mapa. La distancia focal \(d_1=0.544545\ldots\) no es igual a
\(0.54\), y el análisis de invariantes no ha producido una flecha que la
identifique con el defecto \(54\).

La potencia
\[
9^9=387,420,489
\]
genera una escala donde
\begin{equation}
\operatorname{ord}_{10}\bigl(9^{9^9}\bigr)=369,693,099,
\qquad
\#\,\text{dígitos}=369,693,100.
\label{eq:369693}
\end{equation}
El bloque \(369|693\) aparece en el par orden--longitud; no se afirma que sea
el prefijo decimal de la potencia.

\subsubsection{Teoría aritmética de las vacancias decimales}

\paragraph{Pendiente asintótica de vacancia}

La separación entre la década y la base nonádica es
\begin{equation}
\delta_{\mathrm{vac}}=\log_{10}\frac{10}{9}.
\label{eq:delta-vac}
\end{equation}
El contador acumulado
\begin{equation}
V_9(n)=\left\lceil n\delta_{\mathrm{vac}}\right\rceil
\label{eq:vacancy-count}
\end{equation}
registra cuántas posiciones de corrección necesita la publicación decimal de
una evolución nonádica. Sus saltos forman una palabra binaria
\[
z_n=V_9(n)-V_9(n-1)\in\{0,1\}.
\]
Los huecos se distribuyen con longitudes \(21/22\), los retornos con
longitudes \(6/7\), y el determinante combinado es \(15\). Estos números
proceden del lector de vacancia, no del núcleo Barbero aunque el \(15\)
coincida con \(\binom62\).

\paragraph{Polarización y corriente aritmética de la red nonádica}

Definimos
\begin{equation}
P_N=\sum_{n=1}^N(z_n-\delta_{\mathrm{vac}}),
\qquad
J_N=P_N-P_{N-1}=z_N-\delta_{\mathrm{vac}}.
\label{eq:arith-current}
\end{equation}
La secuencia equilibrada satisface \(|P_N|<1\). La \enquote{corriente}
\(J_N\) es aquí una corriente aritmética: mide creación y absorción de
vacancias respecto a la pendiente media. Sólo un transductor dimensional
posterior puede publicarla como densidad de corriente electromagnética.

La intuición física es potente porque la misma forma algebraica separa un
promedio par y una fluctuación orientada. El rigor exige conservar el tipo:
\[
J_N\in\R\quad\text{adimensional},
\qquad
j^\mu\in\text{recta física de corriente}.
\]

\subsubsection{Completación \(729\to1000\) y conservación del acarreo}

\paragraph{Descomposición \(1000=729+243+27+1\)}

La completación
\[
1000=729+243+27+1
\]
conserva la unidad mediante cuatro escalas ternarias. La corona total es
\[
1000-729=271,
\]
mientras el último entero antes del cierre satisface
\[
999=729+270.
\]
La diferencia entre \(270\) y \(271\) distingue frontera ocupada de unidad
completante. El mismo cuidado evita identificar una vacancia con un cero o un
acarreo con un residuo real.

\paragraph{Residencia del acarreo del radión en la corona decimal}

Los operandos del radión comienzan, en bloques de tres cifras, por
\[
\Phi_T:\quad000|027|455|906|837|565|\cdots,
\]
\[
D_E:\quad000|027|455|010|034|743|\cdots.
\]
La resta APP da
\[
\epsone:\quad000|000|000|896|802|822|\cdots.
\]
El pequeño tamaño del resultado no se obtiene truncando los sumandos bajos.
Se obtiene porque dos publicaciones distintas comparten un largo prefijo y el
operador \(\operatorname{SubAcarreo}_{1000}\) conserva los préstamos necesarios
para restarlas exactamente.

La palabra \emph{memoria} tiene aquí dos niveles:
\begin{enumerate}
\item \(q_{\mathrm{mem}}\) registra la profundidad de retornos nonádicos;
\item \(c_i\) registra los préstamos locales de la resta decimal.
\end{enumerate}
El primero organiza la hélice; el segundo publica el decimal. \(\epsR\) usa
ambos porque sus operandos proceden de cilindros a profundidad y su resta se
realiza con acarreo, pero no es igual a ninguno de los contadores.

\begin{thesisbox}
El borde no es un lugar donde la teoría pierde información. Es el lugar donde
la información cambia de representación: fase que vuelve, memoria que
asciende, cilindro que se contrae, vacancia que se registra y acarreo que
transporta la diferencia.
\end{thesisbox}

\subsection{Compatibilidad de las secciones directa y torsional: caracterización estructural de \(\pi\)}
% Fragmento público generado desde propietarios íntegros.
% Residencia: 52.11.2.
% Fuente: ../incoming/radion_monografia_20260827/manuscrito/sections/03_cierre_exacto.tex:286-309
\subsubsection{Caracterización estructural de \(\pi\) por compatibilidad de secciones}

Reordenar \eqref{eq:cierre-unitario} da
\begin{equation}
\boxed{
2\pih=
\frac{1+\Phi_T-\epsone}
{\frac5{36}+\frac{25}{9}\alphah}.}
\label{eq:pi-radion}
\end{equation}
Esta igualdad puede llamarse una nueva definición de \(\pi\) dentro de la
carta del radión si se entiende \emph{definición estructural} como una
caracterización exacta por compatibilidad de secciones. No es un segundo
generador independiente: \(\Delta_4\) y el operador de acarreo ya reciben la
sección coinductiva de \(\pi\). La novedad está en que la vuelta queda
caracterizada por el empalme entre publicación directa, torsión y memoria.

\begin{narrativebox}
La ecuación \eqref{eq:pi-radion} no dice que primero desconozcamos \(\pi\) y
lo adivinemos cerrando un ángulo. Dice que el \(\pi\) generado por la conexión
nonádica es exactamente el que hace compatibles dos lecturas internas del
mismo estado. Es una identidad de reconocimiento intrínseco, no una
calibración metrológica.
\end{narrativebox}


\section{Conservación evolutiva de la unidad holográfica a través de las dimensiones}
% Conversión TeX fiel del cuerpo científico del delta narrativo.
% Fuente congelada: ../../../../../HMT2/CONTINUIDAD_EDITORIAL/RECONSTRUCCION_ARMONICA_MACROTRATADO_20260822/REVISION_CONJUNTA_PARTES_I_II/auditar_pdf_rigor_academico/docs/DELTA_NARRATIVO_UNIDAD_RADION_AUTOCONSERVACION_AUTOINERCIA_GRAVEDAD_PENTADICA_20260828.md:12-366.
\subsection{Separación tipada entre el régimen de curvatura \(\tau\in\{+1,0,-1\}\), la respuesta radial \(p\) y la holonomía global de espín}

La biografía del radión puede representarse mediante una curva elevada

\[
\gamma(\theta)
=
\bigl(r(\theta)\cos\theta,\,
r(\theta)\sin\theta,\,
m(\theta)\bigr),
\]

donde las dos primeras coordenadas pertenecen al plano euclídeo \(\mathbb R^2\) y \(m(\theta)\) registra la memoria o profundidad del levantamiento. Esta ecuación permite leer conjuntamente círculo, hélice y espiral.

Cuando evoluciona el ángulo y permanecen constantes el radio y la memoria, la publicación plana es un círculo. Cuando evolucionan ángulo y memoria mientras el radio conserva su valor, la trayectoria completa es una hélice de radio constante. Cuando evolucionan ángulo y radio sobre una memoria fija, aparece una espiral en \(\mathbb R^2\). Cuando evolucionan las tres coordenadas, la trayectoria es una hélice radial: una biografía en la que fase, amplitud y memoria se transforman conjuntamente.

El carácter radial \(p\) clasifica la ley de respuesta del radio. Una familia útil es

\[
r_p(\theta)
=
r_0\bigl[1+p\lambda(\theta-\theta_0)\bigr]^{1/p},
\qquad p\neq 0,
\]

con límite logarítmico

\[
r_0\exp\bigl(\lambda(\theta-\theta_0)\bigr)
\qquad\text{cuando }p\to0.
\]

En esta parametrización, \(p=2\) publica el régimen de Fermat, \(p=1\) el régimen arquimediano, \(p=0\) el logarítmico, \(p=-1\) el hiperbólico y \(p=-2\) el de lituus, con las elecciones de dominio y orientación correspondientes. Los valores fraccionarios describen regímenes intermedios. Así, \(p=\tfrac12\), \(p=\tfrac32\) o \(p=-\tfrac12\) expresan respuestas radiales que interpolan entre las familias enteras y amplían la clasificación geométrica disponible.

El carácter también puede cambiar al atravesar un umbral. La escritura diferencial

\[
\frac{dx}{d\theta}
=
\beta(m)\,x^{1-p(m)},
\qquad x=\frac r{r_0},
\]

describe una biografía por regímenes: \(p(m)\) permanece estable dentro de una fase y cambia cuando la memoria alcanza una frontera de publicación. La trayectoria resultante enlaza sectores circulares, helicoidales y espirales mediante una ley común. La continuidad del estado y el registro del umbral convierten esos sectores en fases de una misma evolución.

Tres invariantes ocupan aquí niveles distintos y coordinados. El signo de curvatura

\[
\tau\in\{+1,0,-1\}
\]

clasifica los regímenes elíptico, parabólico e hiperbólico. El carácter \(p\) clasifica la respuesta radial. El espín registra la holonomía de orientación acumulada por el levantamiento. Curvatura, respuesta radial y espín forman un sistema de lectura: la curvatura determina el tipo local de espacio, \(p\) determina cómo la trayectoria ocupa ese espacio y el espín conserva la orientación adquirida al recorrerlo. Las curvas cerradas y abiertas aparecen entonces como biografías geométricas distintas de la misma unidad orientada.

\subsection{Partícula como sección persistente de una órbita enriquecida: fase, ruta, fibra, multisección y memoria}

La Mecánica Dimensional puede formular la partícula como una sección persistente de esa extensión geométrica y dinámica. Una notación conceptual compacta es

\[
\mathfrak P
=
\operatorname{Sec}_{\mathrm{pers}}
\bigl[
\tau,p,\operatorname{Hol};
\text{hoja},\text{fibra},\text{ruta},\text{memoria}
\bigr].
\]

La partícula conserva una ruta reconocible a través de cambios de fase, curvatura y dimensión. Su espín publica la holonomía de orientación; su carga publica la polaridad de fase; su color publica una orientación interna de la fibra calibre; su sabor distingue familias de rutas o multisecciones; su masa expresa el carácter de transducción dimensional mediante el cual la acción areal adquiere lectura volumétrica. Cada propiedad es un lector del estado persistente y todas remiten a una genealogía común.

El electrón ocupa el umbral basal de esta construcción porque ofrece el primer lector estable en el que fase, carga, espín y masa quedan simultáneamente publicadas. La masa electrónica rectora

\[
m_e^\beta=0.510998950690000808608\ldots\ \mathrm{MeV}
\]

fija una referencia de transducción; los estados basales anteriores conservan su función como etapas genealógicas. A partir del electrón, los cocientes \(K_D/K_\Omega\), las torres \(90/120\) y el dominio \(81/56/13/324\) organizan la publicación de familias masivas. La ley de masas lee, por tanto, cómo cada sección persistente atraviesa el espacio enriquecido y convierte densidad areal de acción en presencia volumétrica.

Esta interpretación une las espirales con las partículas de manera precisa. El índice \(p\) aporta la respuesta radial; \(\tau\) aporta el régimen de curvatura; la holonomía aporta el espín; la ruta y la fibra aportan sabor y color; la transducción área–volumen aporta el carácter masivo. El atlas de partículas se convierte así en una clasificación de modos estables de atravesamiento dimensional.

\subsection{Transducción de la fase orientada en polarización eléctrica, respuesta magnética y emisión radiativa de frontera}

La electricidad se presenta como transporte dirigido de fase y acción a través de una estructura de vacancias. Una vacancia es una transición admisible del soporte: abre un canal, fija un umbral y permite que la actualización avance. La componente eléctrica publica el transporte tangencial a lo largo de la ruta. La componente magnética publica la curvatura transversal inducida por ese transporte. Su perpendicularidad expresa la relación geométrica entre avance y giro.

Cuando dos componentes transversales sinusoidales mantienen el desfase adecuado, el vector del campo gira. La propagación longitudinal eleva ese giro a una hélice. La variación de amplitud transforma su proyección frontal en una espiral. La lectura geométrica queda entonces completamente coordinada: el ángulo porta la fase; la velocidad angular porta la frecuencia; el radio porta amplitud y densidad de acción; la orientación porta polarización y sentido; la coordenada vertical porta propagación y memoria; \(p\) porta el régimen radial; el umbral porta emisión, absorción o cambio de estado.

El electromagnetismo aparece como una dinámica de fase–acción–memoria con publicaciones complementarias. La electricidad registra el avance polar; el magnetismo registra la respuesta transversal; la onda registra la periodicidad; la hélice registra el avance con memoria; la espiral registra la evolución radial; la radiación registra la publicación del estado cuando la acumulación alcanza un umbral.

Esta arquitectura también ordena el papel de las vacancias y los ciclos. La vacancia proporciona la posibilidad local de transición; el ciclo conserva la fase; la memoria distingue las vueltas; el umbral decide la residencia siguiente de la acción. La propagación puede permanecer confinada, cambiar de régimen radial o publicarse como radiación controlada. Las fuerzas de campo adquieren así una descripción genealógica continua desde el soporte discreto hasta la onda observada.

\subsection{Reversibilidad de la dinámica completa y coste termodinámico de las proyecciones reductoras de información}
El Landauer nulo constituye un caso límite de esta ley. Para una actualización biyectiva del estado completo,

\[
H(X\mid U(X))=0,
\qquad
\inf Q_L=0.
\]

La información permanece disponible en la genealogía. Cuando un lector realiza un cociente \(q\), la entropía se descompone como

\[
H(X)=H(q(X))+H(X\mid q(X)).
\]

El término condicional mide la información retenida fuera de la proyección. Un reinicio trítico materializado satisface el umbral

\[
Q_{\mathrm{reset}}\ge k_B\Theta\log 3.
\]

La constante de Boltzmann traduce la multiplicidad que el lector térmico agrupa; el umbral de Nyquist separa el muestreo fiel del aliasing; el Landauer nulo caracteriza el transporte reversible del estado íntegro. Información, temperatura y acción quedan articuladas por el tipo de lectura aplicado.

\subsection{Fuerzas autoconservativas como transporte de acción entre los sectores observable, radial, de memoria, de vacancia y de radiación}

La conservación evolutiva se expresa mejor mediante intercambio que mediante inmovilidad escalar. Un estado completo distribuye acción y momento entre movimiento visible, memoria, deformación radial, vacancias, radiación y rama conjugada de retorno. El balance enriquecido puede escribirse como

\[
\Delta J_{\mathrm{visible}}
+\Delta J_{\mathrm{memoria}}
+\Delta J_{\mathrm{radial}}
+\Delta J_{\mathrm{vacancias}}
+\Delta J_{\text{radiación}}
+\Delta J_{\mathrm{retorno}}
=0.
\]

Una fuerza es \textbf{autoconservativa} cuando su propia ley de actualización contiene los canales que transportan, almacenan, publican y reabsorben la acción. La conservación pertenece al balance genealógico completo de la interacción. Cada lector observa una parte del intercambio; la genealogía recompone su totalidad.

El sistema Tierra–Luna ofrece una imagen física directa. La rotación terrestre, la órbita lunar, las mareas, la deformación y la radiación intercambian momento angular. El consumo neto del sistema completo se cierra en el balance genealógico de esas transferencias. La magnitud denominada energía funciona como lectura derivada de la acción y del momento publicados en cada sector. La conservación fundamental reside en la compatibilidad entre todos los lectores del intercambio.

Esta formulación disuelve la antigua frontera entre fuerzas conservativas y fuerzas llamadas no conservativas mediante una categoría más profunda. Rozamiento, disipación aparente y radiación ocupan canales explícitos del estado ampliado. Una pérdida en la proyección mecánica visible corresponde a una transferencia hacia memoria, deformación, calor, vacancias o radiación. La autoconservación sigue la residencia de la acción a través de esos cambios.

\subsection{Sistemas autoinerciales: publicación de la aceleración mediante orientación, curvatura, hoja y memoria}

La expresión adecuada es \textbf{autoinercia relacional}. Designa la compatibilidad entre la conexión propia del estado total y la respuesta que aparece al proyectarlo dentro de la distribución global. Los aspectos endógeno y exógeno son componentes analíticas de una sola relación dinámica.

Sea \(\Gamma\) una trayectoria horizontal o geodésica del estado completo. Su ley propia satisface

\[
\nabla_{\dot\Gamma}\dot\Gamma=0.
\]

Al proyectarla sobre un lector \(S\), la aceleración observada queda determinada por la curvatura de la proyección:

\[
\nabla^S_{\dot\Gamma_S}\dot\Gamma_S
=
\mathcal K_{\pi_S}(\dot\Gamma,\dot\Gamma).
\]

La aceleración publicada expresa la relación entre movimiento total y geometría del lector. El principio de Mach completa este mecanismo mediante la traza global de frontera

\[
b:\mathcal X_\infty\longrightarrow B_\infty
\]

y la factorización de la respuesta inercial local

\[
I_v=\bar I_v\circ b.
\]

La distribución global entra en la respuesta local a través de una aplicación tipada. La autoinercia relacional reúne, por tanto, la conexión endógena del movimiento, la condición global transmitida por \(b\) y la proyección concreta del observador. Una escritura sintética para esa coordinación es

\[
\nabla^{\mathrm{auto}}
=
\mathcal M\bigl(
\nabla^{\mathrm{end}},
b[\mathcal X_\infty],
\pi_S
\bigr),
\]

donde \(\mathcal M\) representa el acoplamiento machiano y conserva el tipado de cada componente.

El principio de Ambivalencia añade la doble orientación de la lectura. Toda publicación directa conserva una rama conjugada de retorno; emisión y absorción forman extremos relacionados de la misma biografía; la orientación temporal se lee desde ambos sentidos del transporte. La teoría del absorbedor encuentra aquí una formulación natural: fuente y receptor cierran conjuntamente el balance genealógico de fase y acción. El presente torsional es la interfaz donde ambas orientaciones se compatibilizan.

TRIT actúa en este nivel como cristal de tiempo: organiza la periodicidad local de la actualización y conserva, mediante memoria, la diferencia entre retornos sucesivos. La rama directa y la rama conjugada relacionan emisión con absorción, y el acoplamiento observador–observable adquiere una residencia dinámica dentro del mismo estado. El tiempo publicado es la lectura de esa orientación; el tiempo genealógico es la secuencia completa de fase, retorno y memoria.

Mach y Ambivalencia cumplen funciones complementarias. Mach relaciona la respuesta local con la distribución global. Ambivalencia conserva las direcciones directa y conjugada de la evolución. La autoinercia relacional articula ambas y sustituye la escisión entre sistemas inerciales y no inerciales por una ley única de conexión, proyección y memoria.

\subsection{Dualidad areal--volumétrica, transición hexada--octada y función del parámetro de Barbero--Immirzi}

El paso de área a volumen constituye una transducción dimensional central. El área publica acción de frontera; el volumen publica profundidad interior. La derivación HMT del parámetro de Barbero–Immirzi organiza la compatibilidad entre ambas medidas mediante la incidencia hexada–octada y los canales \(90/120\):

\[
90=6\binom62,
\qquad
120=8\binom62.
\]

Los cardinales \(6\) y \(8\) admiten una realización cúbica mediante caras y vértices, mientras la hexada y la octada de Steiner constituyen objetos de incidencia distintos, enlazados únicamente por el morfismo tipado \(6\to8\). Las \(12\) aristas del cubo conservan su función combinatoria propia y no definen por sí solas la incidencia de Steiner. La diferencia de incidencia expresa una incompatibilidad dimensional productiva: la información de superficie requiere una regla de transducción para adquirir residencia volumétrica.

El mismo principio aparece en el acarreo \(999+1=1000\). Los nueve alcanzan el borde del registro decimal; el incremento publica simultáneamente el cierre del nivel anterior y el origen \(0/1\) del siguiente. La razón holográfica entre \(729\), \(243\), \(27\) y \(1\) despliega esa transición en escalas triádicas. La transducción de Barbero–Immirzi se inserta en esta familia de cambios de lector: frontera e interior conservan la unidad mientras redistribuyen su medida.

El defecto normalizado \(-1/12\) y el peso orientado \(12:-1\) ocupan funciones distintas dentro de esta geometría. El primero cuantifica una corrección de incidencia; el segundo registra una orientación firmada en la transducción. Su coordinación con el área mínima, el Hamiltoniano modular y la memoria proporciona el puente hacia la dinámica dimensional.

\subsection{Contracción hiperbólica interior, expansión exterior y conservación holográfica de la frontera}

La tesis puede formularse desde una idea única: la unidad conserva su identidad mientras cambia de escala, de lector, de curvatura, de memoria y de dimensión. Esta permanencia posee carácter evolutivo porque la unidad se despliega, adquiere nuevas coordenadas y publica propiedades diferentes; posee carácter holográfico porque cada publicación finita conserva la regla de reconstrucción del estado del que procede. El cierre holográfico del infinito designa precisamente esa compatibilidad entre identidad y despliegue: cada estadio contiene la ley que permite profundizar hacia dentro y prolongar hacia fuera la misma genealogía.

La denominación rectora es \textbf{principio de conservación evolutiva de la unidad mediante transducción dimensional y proyección holográfica}. La expresión sitúa cada término en su nivel correcto. La conservación pertenece a la identidad genealógica; la evolución pertenece a la actualización; la transducción relaciona realizaciones dimensionales; la proyección holográfica conserva la correspondencia reconstructiva entre sus lectores.

La aritmética elemental ofrece la primera imagen de esta arquitectura. La completación decimal

\[
1000=10^3=729+243+27+1=3^6+3^5+3^3+3^0
\]

expone cuatro estratos triádicos dentro de una unidad cúbica decimal. Al normalizar,

\[
1=\frac{729}{1000}+\frac{243}{1000}+\frac{27}{1000}+\frac{1}{1000},
\]

la unidad se reconoce simultáneamente como totalidad y como despliegue graduado. La lectura recíproca

\[
3^{-6},\quad 3^{-5},\quad 3^{-3},\quad 3^0
\]

describe la profundidad interior de esos estratos; su normalización produce la partición correspondiente. La relación \(999+1=1000\) expresa el mecanismo de transporte: la última unidad completa un registro y actúa, a la vez, como acarreo que inaugura el siguiente. El cero y el uno forman así una interfaz de continuidad. Los intervalos \([0,1]\), \([0,10]\) y la prolongación hacia \([0,\infty)\) son lectores escalares de una misma ley de completación.

La primera escisión \(1000=729+271\), seguida por \(271=243+27+1\), da forma a la razón extrema y media holográfica del despliegue: un núcleo triádico dominante conserva dentro del resto la misma regla de resolución. La unidad cúbica decimal contiene así una profundidad triádica anidada. El movimiento hacia \(729\), \(243\), \(27\) y \(1\) lee la contracción interior; el acarreo hacia el millar lee la expansión exterior.

Esta doble orientación organiza el infinito HMT. Hacia dentro, el refinamiento crece hiperbólicamente: cada celda admite una nueva profundidad sin perder su filiación. Hacia fuera, la unidad se expande mediante publicaciones dimensionales sucesivas. El cierre consiste en la conmutatividad de ambos movimientos. Si \(U\) representa la actualización del estado, \(R_d\) el lector en el estrato \(d\) y \(T_{d\to d'}\) la transducción entre estratos, la conservación evolutiva se expresa mediante

\[
R_{d'}\circ U
=
T_{d\to d'}\circ R_d.
\]

El mismo acontecimiento puede ofrecer una medida angular, una densidad de acción, una masa publicada o una respuesta gravitatoria. Su identidad común reside en la genealogía que hace conmutar esas lecturas.

\subsubsection{Monodromía nonádica y elevación helicoidal con avance de memoria}

APP, TRIT y TPK proporcionan el mecanismo formal de este despliegue. APP construye el soporte y separa las evaluaciones que deben conservarse. TRIT organiza los regímenes locales y su ambivalencia. TPK transporta fase y memoria sobre el soporte enriquecido. La dependencia causal queda fijada por

\[
U_t\longrightarrow \Gamma_9\longrightarrow K_{\mathrm{ph}}.
\]

El retorno de fase y el avance de memoria forman operaciones coordinadas. Tras nueve pasos puede reaparecer la misma clase angular,

\[
g(K+9)=g(K),
\]

mientras el contador de memoria progresa,

\[
q_{\mathrm{mem}}\Gamma_9=q_{\mathrm{mem}}+1,
\qquad
B_{K+9}\neq B_K.
\]

La figura geométrica precisa es una elevación helicoidal. Una órbita cerrada en la proyección angular se eleva a una trayectoria abierta cuando cada retorno de fase incrementa la memoria. La circunferencia conserva el ciclo; la hélice conserva el ciclo y registra además la profundidad alcanzada. La monodromía coinductiva genera así profundidad infinita mediante una sucesión de retornos localmente cerrados y globalmente ascendentes.

Esta imagen supera el valor meramente ilustrativo. La rueda angular constituye una sección de la hélice; la hélice constituye la biografía completa de la rueda cuando el estado transporta memoria. Cada vuelta recupera una posición de fase y, simultáneamente, inaugura una hoja nueva. La infinitud procede de la compatibilidad entre retorno y elevación: reiterar el ciclo produce profundidad sin disolver la identidad orbital.

El certificado nonádico da contenido finito y verificable a esa arquitectura. La cadena

\[
w_6\longrightarrow w_{12}\longrightarrow w_{18}\longrightarrow
w_{24}\longrightarrow w_{30}\longrightarrow R_{36}\longrightarrow G_9
\]

ordena el levantamiento; los \(396\) niveles, los \(2376\) trits, las \(43\) vueltas y los \(387\) pares \((K,K+9)\) por canal cuantifican el tránsito; las cinco regiones de \(\pi\) y la fibra ambiente de \(729\) elementos publican lecturas internas posteriores. Esta secuencia hace visible una monodromía productiva: el retorno recupera la fase y la coinducción conserva el crecimiento de memoria.

\subsubsection{Unidad radional y coordinación de los lectores circular, elíptico, proyectivo y helicoidal}

El radión concentra en una unidad metrológica la estructura que el radián publica angularmente. Su definición natural es la de una \textbf{sección orientada de fase, acción y memoria}. Una unidad de fase \(\theta=1\) transporta una unidad de acción \(\sigma\hbar_{\mathrm{ret}}\), pertenece a una vuelta completa \(T_+=2\pi_{\mathrm{HMT}}\) y conserva hoja, orientación, retorno nonádico y posición en el cilindro de refinamiento.

El radián expresa la coordenada angular visible. El radión conserva, junto a esa coordenada, la acción transportada, la orientación, el acarreo y la memoria de la vuelta. La metrología angular adquiere así una genealogía. Medir un arco significa identificar cuánto giro se ha publicado y qué estado completo ha realizado ese giro.

La relación entre fase y área adopta una forma especialmente clara:

\[
\mathcal A(\theta)=\hbar_{\mathrm{ret}}\,\theta.
\]

Una unidad radional barre un área de acción \(\hbar_{\mathrm{ret}}\); una vuelta completa encierra \(h_{\mathrm{ret}}\). El paso de fase a área queda incorporado en la propia unidad de lectura. La métrica angular, la acción y la memoria dejan de aparecer como cantidades yuxtapuestas y pasan a ser coordenadas de una misma sección orientada.

El radión admite varias publicaciones geométricas coordinadas. El lector circular publica fase; el lector elíptico publica acción y forma; el lector interior de Klein publica profundidad; el lector helicoidal nonádico publica retorno y memoria. La unidad permanece genealógicamente idéntica a través de esas realizaciones. Por eso el radión constituye el enlace metrológico adecuado entre HMT y Mecánica Dimensional: ofrece una unidad que puede atravesar cambios de lector conservando el estado que cada medida parcial presupone.

\subsection{Polaridad gravitatoria pentacomponente y transducción conjunta de acción, área, volumen, tiempo y masa}

La Mecánica Dimensional permite definir la dimensión como un régimen estable de publicación caracterizado por el rango mínimo de memoria independiente requerido por el levantamiento. Atravesar una dimensión significa aplicar un transductor que conserva la genealogía y cambia el tipo de lectura. La gravedad ocupa el nivel de conexión entre esos regímenes.

En esta formulación, la \textbf{gravitación interdimensional de cinco polarizaciones correlativas}, abreviadamente \textbf{estructura pentapolar de la gravitación}, es la conexión que transporta la unidad genealógica y deja en cada interfaz una torsión mínima. Una escritura compacta para esa traza es

\[
\Theta_d^{\min}
=
\nabla_{d+1}\circ T_{d\to d+1}
-
T_{d\to d+1}\circ\nabla_d.
\]

La expresión mide la diferencia entre transportar después de derivar y derivar después de transportar. Esa diferencia es la memoria geométrica mínima del cruce dimensional. La composición de interfaces produce una holonomía gravitatoria

\[
\mathfrak G_{d\to D}
=
\operatorname{Hol}
\bigl(
\Theta_d^{\min},
\Theta_{d+1}^{\min},
\ldots,
\Theta_{D-1}^{\min}
\bigr).
\]

El objeto fundamental es esta conexión interdimensional. Un modo local tipo gravitón ocupa el rango derivado de publicación cuantizada de la torsión; la gravitación conserva su residencia primaria en la compatibilidad entre dimensiones.

El término pentapolar recoge las cinco polarizaciones consustanciales del estado continuo:

\[
\mathbf G_d^{\pm}
=
\bigl(
G_{\mathrm{sp},d}^{\pm},
G_{\mathrm{sol},d}^{\pm},
G_{\mathrm{gau},d}^{\pm},
G_{\mathrm{coh},d}^{\pm},
G_{\mathrm{det},d}^{\pm}
\bigr).
\]

Las cinco componentes emergen correlativamente del mismo estado enriquecido. La polaridad \(\pm\) registra concentración directa y publicación conjugada. La gravedad atraviesa los estratos dimensionales transportando simultáneamente esas cinco lecturas; la torsión mínima conserva la huella de cada cruce.

La geometría adquiere con ello una imagen unificada. El régimen elíptico organiza cierres; el parabólico organiza umbrales; el hiperbólico organiza profundidad. La transducción dimensional enlaza esos regímenes con las respuestas radiales \(p\), las holonomías de espín y las publicaciones de masa. La estructura pentapolar de la gravitación expresa la coherencia global de ese enlace.

La teoría M proporciona pantallas físicas posteriores para esta conexión: dimensiones adicionales, cuerdas, branas y dualidad T realizan modos concretos de transducción. HMT aporta la genealogía APP–TRIT–TPK, el estado enriquecido, la memoria y la lectura bidireccional; la interfaz con teoría M publica esas estructuras en geometrías físicas tipadas. La dualidad T intercambia lectores dimensionales compatibles y conserva la clase genealógica del estado.

\subsubsection{Compatibilidad de las realizaciones clásica, relativista, cuántica y HMT--MD}

La mecánica clásica publica trayectoria y momento. La relatividad publica estructura causal, métrica y curvatura. La mecánica cuántica publica fase, holonomía, vacancias y amplitudes. La Mecánica Dimensional proporciona los transductores que relacionan esas publicaciones, y HMT conserva la genealogía discreta que las sostiene.

En la lectura clásica, una órbita describe la proyección visible del movimiento. En la lectura relativista, esa órbita queda incorporada a la geometría del lector. En la lectura cuántica, la fase y el espín registran la holonomía de la sección. En la lectura HMT–MD, las tres expresiones pertenecen a una misma biografía elevada: el radión cuantifica fase, acción y memoria; \(p\) clasifica la respuesta radial; \(\tau\) clasifica la curvatura; la holonomía publica espín; el transductor dimensional publica masa; la conexión pentapolar publica gravitación.

La aparente disipación clásica adquiere residencia en el balance genealógico ampliado. La aceleración relativista se expresa mediante conexión y proyección. La indeterminación cuántica se organiza por lectores, vacancias y ramas de Ambivalencia. El principio de Mach relaciona la respuesta local con la distribución global. El Landauer nulo garantiza la reversibilidad informacional del estado completo. Todas estas piezas convergen en una conservación más profunda: la identidad de la unidad a través de cambios de publicación.

\subsubsection{Teorema narrativo de conservación genealógica y prolongación holográfica}

El cierre puede narrarse ahora como una sola secuencia. La unidad se descompone sin perder identidad. La fase orientada adquiere medida en el radión. El retorno nonádico convierte el ciclo en hélice mediante memoria. El carácter \(p\) transforma la hélice en familias radiales y clasifica sus regímenes. La persistencia de esas trayectorias produce secciones con espín, carga, sabor, color y masa. Las vacancias abren canales; la polaridad eléctrica transporta fase; la respuesta magnética curva transversalmente; los umbrales publican radiación. La autoconservación cierra el balance genealógico de acción y momento. La autoinercia relacional enlaza conexión propia, proyección y totalidad machiana. La Ambivalencia conserva las direcciones directa y conjugada. La transducción área–volumen conduce a Barbero–Immirzi y a la ley de masas. La torsión mínima encadena dimensiones y publica la estructura pentapolar de la gravitación.

El crecimiento hacia dentro y la expansión hacia fuera son orientaciones conjugadas del mismo proceso. El refinamiento hiperbólico genera profundidad; la transducción dimensional genera extensión. El acarreo \(0/1\) enlaza escalas; la monodromía enlaza vueltas; la memoria enlaza estados; la holonomía enlaza orientaciones; la gravitación enlaza dimensiones. Cada enlace conserva una unidad que evoluciona y cada publicación finita contiene la regla de su prolongación.

La tesis fundamental queda formulada así:

\begin{quote}
\textbf{El cierre holográfico del infinito realiza el principio de conservación evolutiva de la unidad mediante transducción dimensional y proyección holográfica: la unidad mantiene su identidad genealógica al desplegarse como fase, acción, curvatura, memoria, partícula, campo y gravitación; cada lectura finita conserva la regla de reconstrucción y prolongación del estado completo.}
\end{quote}

Esta formulación ofrece una imagen física y matemática conjunta. El infinito adquiere una ley de cierre local; la unidad adquiere profundidad; la dimensión adquiere memoria; la conservación adquiere carácter relacional; la fuerza adquiere autoconservación; la inercia adquiere estructura machiana; la gravedad adquiere conexión pentapolar; y la holografía adquiere el sentido preciso de reconstrucción entre lectores compatibles.


\section{Definición operatoria del radión angular mediante el estado enriquecido y \(4\) lectores tipados}
\subsection{Definición operatoria del radión angular y cuádruple realización tipada del estado enriquecido}
% Fragmento público generado desde propietarios íntegros.
% Residencia: 52.13.
% Fuente: ../incoming/radion_monografia_20260827/manuscrito/sections/13_sintesis.tex:1-80
\noindent La síntesis operatoria reúne los lectores de fase, acción, profundidad y memoria sin identificarlos entre sí.\par\medskip

\subsubsection{Cuádruple realización del estado radional}

Sea \(x_\infty\in X_{\mathrm{enr}}\) la sección coinductiva. El radión completo no
es un punto aislado, sino el paquete
\begin{equation}
\mathfrak R_\sigma(x_\infty)
=\bigl(
\theta=1,\sigma,
\pih,\alphah,A,C^*,\Delta_4,
\Phi_T,D_E,\epsone,\hret,
q_{\mathrm{mem}},\mathcal I_\infty
\bigr).
\label{eq:full-radion}
\end{equation}
Cuatro lectores extraen de él objetos diferentes:

\begin{figure}[H]
\centering
\begin{tikzpicture}[
  node distance=17mm and 18mm,
  state/.style={draw=RadNavy,very thick,rounded corners,fill=RadCream,
    minimum width=39mm,minimum height=13mm,align=center},
  readout/.style={draw=RadTeal,rounded corners,fill=RadMist,
    minimum width=39mm,minimum height=12mm,align=center},
  arr/.style={-{Latex},thick,RadCopper}]
\node[state] (x) {estado enriquecido\\\(x_\infty\)};
\node[readout,above left=of x] (circ) {círculo \(S^1\)\\fase};
\node[readout,above right=of x] (ell) {elipse positiva\\acción y forma};
\node[readout,below left=of x] (klein) {interior Klein\\profundidad};
\node[readout,below right=of x] (helix) {hélice nonádica\\memoria y retorno};
\draw[arr] (x)--(circ) node[midway,below left]{\small \(P_{\mathrm{ph}}\)};
\draw[arr] (x)--(ell) node[midway,below right]{\small \(P_{\mathrm{act}}\)};
\draw[arr] (x)--(klein) node[midway,above left]{\small \(P_K\)};
\draw[arr] (x)--(helix) node[midway,above right]{\small \(P_{\mathrm{mem}}\)};
\end{tikzpicture}
\caption{Cuatro publicaciones tipadas del mismo antecedente. La unidad del
estado no elimina la diferencia de codominios.}
\label{fig:four-readers}
\end{figure}

El círculo conserva la clase de fase y olvida el lift. La elipse conserva
acción, anisotropía y orientación simpléctica. Klein dota al interior de una
distancia proyectiva. La hélice conserva número de retornos y refinamiento.
La Mecánica Dimensional no consiste en mezclarlos, sino en construir los
transductores que los hacen compatibles.

\subsubsection{Interpretación operatoria del cociclo \(\varepsilon_R\)}

\(\varepsilon_R\) no es una constante universal independiente de todo contexto. Es la
coordenada de un cociclo de transición asociado a la carta del radión. Su
dominio contiene dos secciones del mismo estado:
\[
S_E=1+D_E,
\qquad
S_T=1+\Phi_T.
\]
Su regla es
\begin{equation}
\epsone=S_T-S_E=\Phi_T-D_E.
\label{eq:epsilon-final-meaning}
\end{equation}
Su codominio es la recta adimensional de unidad angular; la carta
\(360/T_+\) la publica en grados.

\begin{exactbox}
\textbf{En lenguaje llano:} la sección directa y la sección torsional llegan
a la misma celda de unidad por dos rutas internas diferentes. Sus partes
enteras coinciden; sus restos no. \(\epsR\) es la diferencia orientada de esos
restos, calculada con el acarreo de APP y llevada al límite por los cilindros
TPK. No se introduce para que el radián cierre: aparece antes de consultar el
radián y, cuando se publica, lo cierra.
\end{exactbox}

Esta definición explica por qué el número es pequeño. No porque se haya
buscado una corrección pequeña, sino porque \(S_E\) y \(S_T\) comparten un
prefijo profundo. La pequeñez mide la cercanía estructural de dos
publicaciones; el signo mide su orientación; las cifras registran el acarreo.


% Fragmento público generado desde propietarios íntegros.
% Residencia: 52.13.1.
% Fuente: ../incoming/radion_monografia_20260827/manuscrito/sections/13_sintesis.tex:81-134
\subsubsection{Definición operatoria del radión}

\begin{definition}[Radión HMT--MD]
El radión es la sección orientada de fase, acción y memoria que cumple
simultáneamente:
\begin{enumerate}
\item recorre una unidad de fase \(\theta=1\);
\item barre acción firmada \(\sigma\hret\);
\item pertenece a una vuelta \(T_+=2\pih\) generada coinductivamente;
\item conserva la hoja \(\sigma\), el retorno nonádico y el cilindro de
refinamiento;
\item satisface el empalme exacto
\(1=S_E-\Phi_T+\epsone\).
\end{enumerate}
\end{definition}

El radián convencional es la imagen del radión bajo el functor de olvido
\[
\mathfrak R_\sigma
\longmapsto
\theta=1\in\R/2\pi\Z.
\]
Ese olvido es útil y legítimo para la trigonometría euclídea. Se vuelve
insuficiente cuando la física necesita distinguir acción, signo, retorno,
campo e historia.

\paragraph{Genealogía de la unidad angular natural}

La naturalidad del radián suele justificarse porque la derivada de
\(e^{i\theta}\), \(\sin\theta\) y \(\cos\theta\) adopta su forma más simple
cuando \(\theta\) se mide como cociente arco/radio. HMT conserva esa ventaja,
pero la explica desde una capa anterior:
\[
\theta\quad\leftrightarrow\quad
\frac{\mathcal A(\theta)}{\hret}.
\]
La fase es natural porque mide acción en unidades de \(\hret\); la vuelta es
natural porque encierra \(\hfull\). El cociente arco/radio es la publicación
euclídea de una unidad que ya posee estructura simpléctica y memoria.

\paragraph{Compatibilidad entre elipse de acción y crecimiento dimensional}

El círculo es la forma que resulta al olvidar la anisotropía:
\(C^*=0\Rightarrow\eta=0\Rightarrow a=b\). La elipse es más informativa porque
conserva dos autovalores, dos ejes y una orientación de intercambio a producto
fijo. En ese sentido, la elipse es una prolongación del círculo, no un círculo
deformado por accidente.

La profundidad de Klein y el lift helicoidal añaden dos tipos de información
que el contorno plano no contiene: distancia interior y memoria de retorno.
Ésta es la imagen ontológica de Mecánica Dimensional que emerge del radión:
una dimensión nueva aparece cuando se exige conservar un dato independiente
que la proyección anterior debía olvidar.


\subsection{Propiedades estructurales de la unidad radional: genealogía, bisagra crítica, densidad, acarreo, acción, compatibilidad borde--interior, forma, polaridad y modularidad}
% Fragmento público generado desde propietarios íntegros.
% Residencia: 52.13.2.
% Fuente: ../incoming/radion_monografia_20260827/manuscrito/sections/13_sintesis.tex:135-184
\subsubsection{Sistema de teoremas componentes del radión}

\begin{theorem}[Genealogía]
El objeto \(\mathfrak R_\sigma\) desciende de
\(\APP\to\TRIT\to\TPK\to X_{\mathrm{enr}}\) y no recibe el blanco metrológico
como selector.
\end{theorem}

\begin{theorem}[Bisagra]
En la carta APP de cien marcas, la costura crítica y la fase dodecafásica
segunda satisfacen \(j_{100}(1/2)=\theta_2=50^\circ\), con \(m=2\) único.
\end{theorem}

\begin{theorem}[Densidad]
Las dos hojas del sector activo \(N_6=90\), normalizadas por \(3^6=729\),
producen \(\rho_{\mathrm{tw}}=20/81\).
\end{theorem}

\begin{theorem}[Acarreo]
La resta APP de las secciones \(S_T\) y \(S_E\) converge a
\(\epsone\) con tasa de intervalo \(3^{-54}\) por retorno y cierra exactamente
la unidad.
\end{theorem}

\begin{theorem}[Acción]
La elipse positiva satisface \(\mathcal A(\theta)=\hret\theta\); un radión
barre \(\hret\) y una vuelta encierra \(\hfull\).
\end{theorem}

\begin{theorem}[Borde--interior]
La misma frontera posee acción finita \(\hfull\), mientras el interior de
Klein tiene área hiperbólica infinita.
\end{theorem}

\begin{theorem}[Forma]
El cociente \(e^{-\eta}\) se conserva a través de semiejes espectrales,
semiejes de acción, incertidumbres, impedancia y módulo del toro.
\end{theorem}

\begin{theorem}[Polaridad]
La inversión \(C^*\mapsto-C^*\) intercambia constitutivas, invierte
impedancia y conserva la velocidad normalizada.
\end{theorem}

\begin{theorem}[Modularidad]
La inversión bidireccional actúa como \(S:\tau\mapsto-1/\tau\) y conserva
\(j(\tau)\), mientras el etiquetado focal conserva la orientación que
\(j\) olvida.
\end{theorem}

% Fuente: ../incoming/radion_monografia_20260827/manuscrito/sections/A_demostraciones.tex:1-179
\subsubsection{Demostraciones algebraicas y geométricas reunidas}

\paragraph{Demostración lineal del cierre unitario}

Para facilitar una auditoría independiente, reunimos el argumento sin
interrupción narrativa.

\begin{enumerate}[label=\textbf{Paso \arabic*.},leftmargin=2.2cm]
\item La conexión nonádica genera correlacionadamente
\(\pih,e_{\HMT},\varphi_{\HMT},\alphah\).
\item La vuelta es \(T_+=2\pih\).
\item La rueda fija \(\theta_2=50^\circ\), y la carta parangular fija
\(A=1000\alphah\).
\item La sección directa es
\[
S_E=T_+\left(\frac{50}{360}+\frac{1000\alphah}{360}\right)
=2\pih\left(\frac5{36}+\frac{25}{9}\alphah\right).
\]
\item Los intervalos certificados dan \(1<S_E<2\); luego
\(D_E=S_E-1\).
\item La incidencia produce \(N_6=90\); dos hojas y el ambiente \(729\)
producen \(\rho=2N_6/729=20/81\).
\item La torsión global es
\[
\Delta_4=\frac{\pih-e_{\HMT}\log\pih}{270}
\left(1+\frac{\pih}{729}\right).
\]
\item La sección torsional es \(S_T=1+(20/81)\Delta_4\).
\item La resta APP define
\[
\epsone=S_T-S_E=\frac{20}{81}\Delta_4-D_E.
\]
\item Por sustitución,
\[
S_E-\frac{20}{81}\Delta_4+\epsone=1.
\]
\item Multiplicar por \(360/T_+=180/\pih\) produce
\[
\frac{180^\circ}{\pih}
=50^\circ+A^\circ-\frac{20}{81}\Delta_4^\circ+\epsR^\circ.
\]
\end{enumerate}

Ningún paso recibe el lado izquierdo de la última igualdad. El resultado
metrológico sólo se evalúa al final como reconocimiento.

\paragraph{Unicidad de la resta con acarreo}

Sea \(B=1000\). Para cada posición, todo entero \(u_i\) admite una división
euclídea única
\[
u_i=Bc_i+r_i,\qquad 0\le r_i<B.
\]
Por tanto,
\[
r_i=u_i\bmod B,\qquad c_i=(u_i-r_i)/B
\]
son únicos. Recorriendo las posiciones desde el extremo remoto hacia el
frente, el acarreo de salida de una posición fija la entrada de la siguiente.
Por inducción, toda la palabra de restos y carries es única.

En pseudocódigo:
\begin{verbatim}
acarreo = remote_acarreo
for i from last_block downto first_block:
    u = torsion[i] - direct[i] + acarreo
    remainder[i] = u mod 1000
    acarreo = (u - remainder[i]) / 1000
return remainder, acarreo_word
\end{verbatim}
El procedimiento no consulta el valor esperado de la diferencia.

\paragraph{Convergencia del funcional por profundidad}

Sea
\[
F(p,e)=1+\frac{20}{81}
\frac{p-e\log p}{270}\left(1+\frac{p}{729}\right)
-2p\left(\frac5{36}+\frac{25}{9}\alphah\right).
\]
En el rectángulo compacto \(3\le p\le4\), \(2\le e\le3\), \(F\) es
continuo y continuamente diferenciable. Si \(I_N^\pi\searrow\{\pih\}\) e
\(I_N^e\searrow\{e_{\HMT}\}\), la imagen intervalar
\[
F(I_N^\pi,I_N^e)
\]
forma una familia encajada cuya anchura tiende a cero. La intersección es un
único real, \(\epsone\). La monotonía de las derivadas permite evaluar los
extremos sin muestreo interior.

La tasa \(3^{-54}\) por retorno deriva del refinamiento de seis trits a lo
largo de nueve niveles. A profundidad \(45\) tríadas, la cota de anchura
\(<4.81\times10^{-130}\) certifica más de 129 cifras del valor.

\paragraph{Derivación de la métrica y del área hiperbólica de Klein}

En el disco unitario, la forma matricial de la métrica Klein es
\[
g(u,v)=\frac1{1-r^2}I
+\frac1{(1-r^2)^2}
\begin{pmatrix}u\\v\end{pmatrix}
\begin{pmatrix}u&v\end{pmatrix}.
\]
El lema del determinante para una perturbación de rango uno da
\[
\det g=(1-r^2)^{-3}.
\]
De ahí
\[
\sqrt{\det g}\,\diff u\diff v
=(1-r^2)^{-3/2}\diff u\diff v.
\]
Integrar en polares produce \eqref{eq:klein-area-R}. El límite infinito se
debe a la singularidad métrica en la frontera; la frontera sigue teniendo
área simpléctica finita porque esa medida utiliza \(\diff Q\wedge\diff P\),
no \(\sqrt{\det g}\).

\paragraph{Distancia focal en el disco de Klein}

En un diámetro, la distancia del centro a \(r\) es
\[
d_K(0,r)=\operatorname{artanh}r.
\]
Para \(r=e_f\), \eqref{eq:excentricidad} implica
\[
1-e_f^2=e^{-2\eta}.
\]
Si \(u_f=\operatorname{artanh}e_f\), entonces
\[
\operatorname{sech}u_f=\sqrt{1-e_f^2}=e^{-\eta},
\]
y por tanto \(\eta=\log\cosh u_f\), como en
\eqref{eq:focal-hyperbolic-chain}.

\paragraph{Transformación simpléctica de compresión}

Definamos
\begin{equation}
S_\eta=
\begin{pmatrix}e^{\eta/2}&0\\0&e^{-\eta/2}\end{pmatrix},
\qquad \det S_\eta=1.
\label{eq:compresión simpléctica}
\end{equation}
La estructura compleja transportada es
\begin{equation}
J_\eta=S_\eta
\begin{pmatrix}0&-1\\1&0\end{pmatrix}
S_\eta^{-1}
=\begin{pmatrix}0&-e^\eta\\e^{-\eta}&0\end{pmatrix},
\qquad J_\eta^2=-I.
\label{eq:Jeta}
\end{equation}
La curva
\[
\gamma(\theta)=\sqrt{2\hret}\,
S_\eta(\cos\theta,\sin\theta)
=e^{\theta J_\eta}\gamma(0)
\]
es exactamente la elipse de acción. El compresión simpléctica hiperbólico conserva área y
transporta la rotación compleja en vez de destruirla.

Para orientación \(\sigma\),
\[
\gamma_\sigma(\theta)
=(a_S\cos\theta,\sigma b_S\sin\theta),
\qquad
\mathcal A_\sigma(\theta)=\sigma\hret\theta.
\]
Las dos hojas comparten focos y módulo de acción; difieren en acción firmada.

\paragraph{Comprobación diofántica del peso dodecafásico de Barbero--Immirzi}

La ecuación \(4a-3b=45\) tiene soluciones enteras
\[
(a,b)=(12+3t,1+4t),\qquad t\in\Z.
\]
La condición \(|b|=1\) deja \(b=1\), pues \(b=-1\) no pertenece a la clase
\(1\pmod4\). De ahí \(a=12\). Esta comprobación reproduce el par obtenido mediante la cadena fundamental
dodecafásica: doce sectores y una única costura orientada. Los cardinales
\(90/120\), las multiplicidades de la cadena y el lector de retorno
conservan funciones distintas; el coeficiente de polo y la minimalidad
verifican aritméticamente su composición.

\subsection{Sensibilidad del cierre unitario frente a variaciones controladas de la multiplicidad de hojas, la fase dodecafásica, el canal \(e\), la torsión global, el operador armónico y la cola sexádica}
% Fragmento público generado desde propietarios íntegros.
% Residencia: 52.13.3.
% Fuente: ../incoming/radion_monografia_20260827/manuscrito/sections/13_sintesis.tex:185-247
\subsubsection{Controles negativos del sistema radional}

La construcción es refutable en puntos concretos:
\begin{enumerate}
\item Si \(S_E\notin(1,2)\), el cociente APP no es uno y el cierre cambia de
celda.
\item Si la incidencia activa no produce \(N_6=90\), o si no hay dos hojas,
\(20/81\) deja de seguirse.
\item Si el verificador recibe \(180/\pi\) o \(\epsR\) antes de construir los
operandos, la generación es circular.
\item Si la resta por tríadas no estabiliza prefijos compatibles, no existe
la sección coinductiva propuesta.
\item Si \(|C^*|\ge A\), la elipse positiva deja de existir.
\item Si \(a_Sb_S\ne2\hret\), falla el teorema de área.
\item Si \(C^*=0\), deben cumplirse \(d=0\), \(\eta=0\), \(\tau=i\) y
\(j=1728\).
\item Si el área Klein del subdisco no sigue \eqref{eq:klein-area-R}, falla la
tesis de profundidad infinita.
\item Si \(F_{\mathrm{spec}}=\epsR^\circ\) se afirma sin contratérmino independiente, la
diferencia numérica \eqref{eq:chiR} lo refuta.
\item La producción de \(12:-1\) exige conservar, junto a los cardinales
\(90/120\), los doce sectores dodecafásicos, la única costura orientada y
la acción del lector sobre esos generadores. Cambiar alguna multiplicidad
debe modificar la imagen \(12S_{90}-S_{120}\). Para la cadena fundamental
sin alteración, el coeficiente de \((1-q)^{-1}\) debe resultar \(1/24\).
\item Si la función \(j\) se usa para seleccionar signo de carga, se pierde
la invariancia \eqref{eq:j-invariance} y se incurre en un error de tipo.
\end{enumerate}

\subsubsection{Realización física conjunta del radión}

La imagen final no es un círculo corregido. Es una arquitectura de
publicaciones:
\[
\begin{aligned}
&\text{APP:}&&\text{unidad, resto, acarreo y fibra};\\
&\text{TRIT:}&&\text{régimen, orientación y hojas};\\
&\text{TPK:}&&\text{fase, retorno, dodecafase y memoria};\\
&\text{elipse:}&&\text{acción y anisotropía};\\
&\text{Klein:}&&\text{profundidad interior};\\
&\text{hélice:}&&\text{biografía del retorno};\\
&\text{MD:}&&\text{reloj, energía, carga, campo y causalidad}.
\end{aligned}
\]

La ecuación del radión reúne la cara visible:
\BigRadionEquation
La elipse reúne la cara de acción:
\[
\mathcal A(1)=\hret,\qquad \mathcal A(2\pi)=\hfull.
\]
La métrica de Klein reúne la profundidad:
\[
K=-1,\qquad \operatorname{Area}_K=\infty.
\]
La holonomía reúne la memoria:
\[
g(k+9)=g(k),\qquad B_{k+9}\ne B_k.
\]

\begin{thesisbox}
El círculo es la sombra angular; la elipse es el cuerpo de acción; Klein mide
el interior; la hélice conserva la biografía. El radión es la unidad que
permite leer las cuatro afirmaciones sobre un único estado sin reducir una a
otra.
\end{thesisbox}
% Fuente: ../incoming/radion_monografia_20260827/manuscrito/sections/B_valores_y_verificacion.tex:1-105
\subsubsection{Valores canónicos, estabilidad y verificación reproducible}

\paragraph{Valores de los invariantes internos generados}

\begin{longtable}{p{.27\textwidth} p{.61\textwidth}}
\toprule
Objeto & Valor usado en el perfil coinductivo\\
\midrule
\endfirsthead
\toprule
Objeto & Valor usado en el perfil coinductivo\\
\midrule
\endhead
\(\pih\) & \(\begin{aligned}3.1415926535897932384626433832795028841\\[-2pt]971693993751058209749445923\ldots\end{aligned}\)\\
\(e_{\HMT}\) & evaluación coinductiva del lector de Euler, sin tabla externa\\
\(\alphah\) & \(\begin{aligned}0.0072973525692838009972851054723806629\\[-2pt]995641725396427091854046825\ldots\end{aligned}\)\\
\(A\) & \(\begin{aligned}7.2973525692838009972851054723806629995\\[-2pt]641725396427\ldots{}^\circ\end{aligned}\)\\
\(C^*\) & \(2.1237383389686457242619513803933607937\ldots{}^\circ\)\\
\(\eta\) & \(0.299689683921630799684926562863927\ldots\)\\
\(\hret\) & \(1.05457181691503906113369058472430\ldots\times10^{-34}U_S\)\\
\bottomrule
\end{longtable}

\paragraph{Convergencia cuantitativa por profundidad}

\begin{center}
\small
\begin{tabular}{r l l}
\toprule
Profundidad & \(\epsR^\circ\) & lectura\\
\midrule
12 & \(5.1681282186551375597\times10^{-8}\) & preasintótica\\
18 & \(5.1383017764316644535\times10^{-8}\) & 7 cifras estables\\
24 & \(5.1383016758575394560\times10^{-8}\) & 14 cifras estables\\
30 & \(5.1383016758575282072\times10^{-8}\) & 20 cifras estables\\
36 & \(5.138301675857528207168517\times10^{-8}\) & perfil largo\\
45 & \(5.13830167585752820716851646226459\times10^{-8}\) & ancho \(<4.81\times10^{-130}\)\\
\bottomrule
\end{tabular}
\end{center}

Los retornos de profundidad \(9,18,27,36,45\) conservan la fase \(8\) y
aumentan la memoria \(0,1,2,3,4\). La razón entre anchos consecutivos tiende
a
\[
3^{-54}=1.7196982334549856127610399316004871\ldots\times10^{-26}.
\]

\paragraph{Análisis de sensibilidad estructural por supresión controlada}

\begin{longtable}{p{.46\textwidth} p{.34\textwidth}}
\toprule
Operación & Salida o diferencia unitaria\\
\midrule
\endfirsthead
\toprule
Operación & Salida o diferencia unitaria\\
\midrule
\endhead
Operador canónico & \(8.96802822044562981999\times10^{-10}\)\\
Una hoja, \(\rho=90/729\) & \(-1.3727056615960678\times10^{-5}\)\\
Fase \(\theta_1=20^\circ\) & \(+5.2359877649510169\times10^{-1}\)\\
Fase \(\theta_3=80^\circ\) & \(-5.2359877470149605\times10^{-1}\)\\
Sin canal \(e\) & \(1.8065350748904653\times10^{-3}\)\\
Torsión histórica \(\Delta_4(A,C^*)\) & \(1.0517084608768816\times10^{-9}\)\\
Operador armónico \(M_4\) menos canónico & \(1.0525500222895070\times10^{-17}\)\\
Cola sexádica menos canónico & \(2.0283936357433839\times10^{-12}\)\\
\bottomrule
\end{longtable}

La separación de escalas muestra sensibilidad estructural: el cierre depende
de la fase, las dos hojas, los canales \(\pi/e\), la torsión global y el acarreo.
Ninguna ablación reproduce el valor.

\paragraph{Certificados computacionales reproducibles}

La obra usa los siguientes propietarios ejecutables:
\begin{enumerate}
\item \sourcepath{output/RADION_TWOWAY_NONADICO_20260827/verificar_emergencia_radion_por_profundidad.py};
\item \sourcepath{output/RADION_TWOWAY_NONADICO_20260827/verificar_radion_helice_critica_tpk.py};
\item \sourcepath{output/RADION_TWOWAY_NONADICO_20260827/verificar_operador_tpk_target_free_profundidad.py};
\item el verificador consolidado de esta monografía, ubicado en
\sourcepath{certificados/verificar_monografia_radion.py}.
\end{enumerate}
Las salidas esperadas incluyen
\[
\texttt{PASS\_EMERGENCIA\_RADION\_POR\_PROFUNDIDAD},
\]
\[
\texttt{PASS\_RADION\_HELICE\_CRITICA\_TPK},
\qquad
\texttt{PASS\_MONOGRAFIA\_RADION\_HMT\_MD}.
\]

\paragraph{Distinción entre precisión computacional y exactitud matemática}

Que un programa muestre cero a 120 o 200 cifras no convierte por sí solo una
identidad residual en predicción. La exactitud matemática procede de
\eqref{eq:cierre-unitario}; la precisión arbitraria permite evaluar sus
operandos y verificar la convergencia. La no circularidad procede del orden
de construcción y de la prohibición de blancos, no del número de dígitos.

El cálculo Decimal puede dejar residuos del orden \(10^{-218}\) al combinar
ramas evaluadas con precisiones distintas. Ese número es redondeo de la
implementación; no una corrección física adicional.
% Fuente: ../incoming/radion_monografia_20260827/manuscrito/sections/D_historia_fuentes.tex:1-169
\subsubsection{Genealogía conceptual y fuentes primarias}

\paragraph{Cuanto de acción de Planck}

Planck introduce \(h\) en la ley de radiación como constante que liga energía
y frecuencia. El dato histórico decisivo para esta obra es dimensional:
\(h\) es acción. El radión no modifica ese papel; explica por qué la acción de
una vuelta y la acción por unidad de fase se relacionan mediante \(2\pi\).

\paragraph{Cuantización angular de Bohr y Sommerfeld}

Bohr usa \(h/2\pi\) en la cuantización del momento angular. Sommerfeld amplía
la cuantización a integrales de acción y órbitas keplerianas. Por ello, no es
históricamente exacto decir que Dirac inventó la constante reducida. Su
contribución pertenece a una genealogía en la que la acción angular ya estaba
presente.

La novedad HMT--MD es distinta: no toma \(h/2\pi\) como un cociente externo,
sino que construye una elipse cuya vuelta tiene área \(h\) y cuyo sector de un
radián tiene área \(\hbar\).

\paragraph{Formalismo matricial de Born y Jordan y oscilador cuántico}

Born y Jordan sitúan el principio variacional, las fases y el oscilador
armónico en el corazón de la mecánica matricial; el campo electromagnético se
descompone en una colección de osciladores. La realización LC del radión
continúa esa intuición con un dato añadido: el cociente de forma
\(e^{-\eta}\) determina la impedancia relativa mientras producto, frecuencia
y acción permanecen fijos.

\paragraph{Frecuencia angular y acción en la formulación de Dirac}

En 1925, Dirac escribe las frecuencias en radianes por unidad de tiempo y usa
relaciones de la forma
\[
xy-yx=\frac{ih}{2\pi}[x,y],
\qquad
\Delta E=\frac{h}{2\pi}\omega.
\]
La ecuación \(E=\hret\omega\) une de manera explícita acción reducida y
frecuencia angular. En el radión, esa relación recibe una geometría: \(\hret\)
es el área barrida por una unidad de fase.

\paragraph{Relaciones de indeterminación de Heisenberg, Robertson y Schrödinger}

Heisenberg formula la limitación conjunta de variables canónicamente
conjugadas. Robertson generaliza la desigualdad; Schrödinger incorpora la
covarianza. La matriz \eqref{eq:covariance} pertenece más directamente a esa
prolongación Robertson--Schrödinger que al artículo de Heisenberg aislado.

Ninguno de estos autores formula la elipse genealógica completa de esta obra.
La afirmación históricamente defendible es que proporcionan el lenguaje en el
que la elipse de acción puede reconocerse como un contorno mínimo; la
construcción del par \((A,C^*)\), los focos, Klein, la hélice y el acarreo es la
aportación HMT--MD.

\paragraph{Funcional de fase de Feynman}

La formulación espacio-temporal asigna a cada camino una fase
\(e^{iS/\hbar}\). La unidad de radión hace literal la correspondencia:
incrementar la acción en \(\hbar\) incrementa la fase en un radián. La suma de
caminos reconoce, en un lenguaje integral, el mismo principio que la elipse
expresa como área.

\paragraph{Propagaciones avanzada y retardada en Wheeler--Feynman}

La teoría del absorbedor construye una interacción clásica temporalmente
simétrica mediante combinaciones de campos retardados y avanzados. La
identidad HMT \(CU(t)C^{-1}=U(-t)\) y la inversión de rutas proporcionan un
antecedente operatorio. La promoción física requiere que el mapa entre rutas
y propagadores preserve acción, orientación, frontera y forma simpléctica.

La teoría del absorbedor no es la misma afirmación que la interpretación de
Feynman y Stückelberg de las antipartículas como propagación temporalmente
invertida. Ambas pueden reconocer la bidireccionalidad; sus objetos físicos y
condiciones de frontera son diferentes.

\paragraph{Dualidad de Hodge y separación respecto de la cohomología de Hochschild}

La operación relevante en este manuscrito es la estrella de Hodge
\(\star^2=-I\) sobre dos-formas en firma lorentziana. No se ha localizado en
el corpus activo del radión un propietario de homología de Hochschild que
actúe sobre la elipse, el acarreo o el sector \(90/120\). Introducir Hochschild
por semejanza nominal degradaría el tipado. Si una futura prolongación
construye un cociclo de deformación con dominio y codominio propios, deberá
entrar como resultado nuevo, no como retroatribución.

\paragraph{Bibliografía primaria seleccionada}

\begin{enumerate}
\renewcommand{\labelenumi}{[\arabic{enumi}]}
\item\hypertarget{radionbib:Planck1901}{}
M.~Planck,
\emph{Ueber das Gesetz der Energieverteilung im Normalspectrum},
Annalen der Physik 309 (1901), 553--563.
\url{https://doi.org/10.1002/andp.19013090310}

\item\hypertarget{radionbib:Bohr1913}{}
N.~Bohr,
\emph{On the Constitution of Atoms and Molecules},
Philosophical Magazine 26 (1913), 1--25.
\url{https://doi.org/10.1080/14786441308634955}

\item\hypertarget{radionbib:Sommerfeld1916}{}
A.~Sommerfeld,
\emph{Zur Quantentheorie der Spektrallinien, I},
Annalen der Physik 356 (1916), 1--94.
\url{https://doi.org/10.1002/andp.19163561702}

\item\hypertarget{radionbib:BornJordan1925}{}
M.~Born y P.~Jordan,
\emph{Zur Quantenmechanik},
Zeitschrift für Physik 34 (1925), 858--888.
\url{https://doi.org/10.1007/BF01328531}

\item\hypertarget{radionbib:Dirac1925}{}
P.~A.~M.~Dirac,
\emph{The Fundamental Equations of Quantum Mechanics},
Proceedings of the Royal Society A 109 (1925), 642--653.
\url{https://doi.org/10.1098/rspa.1925.0150}

\item\hypertarget{radionbib:Heisenberg1927}{}
W.~Heisenberg,
\emph{Über den anschaulichen Inhalt der quantentheoretischen Kinematik und Mechanik},
Zeitschrift für Physik 43 (1927), 172--198.
\url{https://doi.org/10.1007/BF01397280}

\item\hypertarget{radionbib:Dirac1928}{}
P.~A.~M.~Dirac,
\emph{The Quantum Theory of the Electron},
Proceedings of the Royal Society A 117 (1928), 610--624.
\url{https://doi.org/10.1098/rspa.1928.0023}

\item\hypertarget{radionbib:Robertson1929}{}
H.~P.~Robertson,
\emph{The Uncertainty Principle},
Physical Review 34 (1929), 163--164.
\url{https://doi.org/10.1103/PhysRev.34.163}

\item\hypertarget{radionbib:Schrodinger1930}{}
E.~Schrödinger,
\emph{Zum Heisenbergschen Unschärfeprinzip},
Sitzungsberichte der Preussischen Akademie der Wissenschaften (1930), 296--303.

\item\hypertarget{radionbib:WheelerFeynman1945}{}
J.~A.~Wheeler y R.~P.~Feynman,
\emph{Interaction with the Absorber as the Mechanism of Radiation},
Reviews of Modern Physics 17 (1945), 157--181.
\url{https://doi.org/10.1103/RevModPhys.17.157}

\item\hypertarget{radionbib:Feynman1948}{}
R.~P.~Feynman,
\emph{Space-Time Approach to Non-Relativistic Quantum Mechanics},
Reviews of Modern Physics 20 (1948), 367--387.
\url{https://doi.org/10.1103/RevModPhys.20.367}

\item\hypertarget{radionbib:Feynman1949}{}
R.~P.~Feynman,
\emph{The Theory of Positrons},
Physical Review 76 (1949), 749--759.
\url{https://doi.org/10.1103/PhysRev.76.749}
\end{enumerate}

\begin{narrativebox}
Los fundadores contienen las piezas: acción, frecuencia angular, variables
conjugadas, oscilador, incertidumbre, fase de caminos y reversión temporal.
La aportación del radión es reunirlas en un objeto generado: una elipse de
área \(h\), un sector de área \(\hbar\), un interior de profundidad infinita,
una orientación de hoja y un acarreo exacto de compatibilidad.
\end{narrativebox}
% Fuente: ../incoming/radion_monografia_20260827/manuscrito/sections/E_conclusion.tex:1-107
\subsubsection{Síntesis final: unidad angular con genealogía de fase, acción y memoria}

El radián no desaparece. Se vuelve visible por completo.

En su uso ordinario, un radián es una razón que no necesita recordar cómo fue
producida. El arco y el radio se cancelan dimensionalmente; la unidad parece
desnuda. Esa desnudez es una virtud para la trigonometría, pero es un olvido
para una teoría que pretende explicar la estructura del continuo y el origen
de la acción. El radión recupera lo olvidado sin alterar el valor publicado.

La genealogía comienza en APP, donde una salida conserva valor y también
hoja, cociente, resto, acarreo y frontera. El TRIT introduce régimen y
orientación. El TPK hace volver la fase sin reiniciar el estado. La conexión
nonádica construye la sección coinductiva de \(\pi\), y la dodecafase publica
la bisagra de \(50^\circ\). El par angular aporta un modo común \(A\) y una
apertura orientada \(C^*\). La incidencia de seis fases activas produce
\(90\), sus dos hojas sobre el ambiente \(729\) producen \(20/81\), y la
torsión global aporta la segunda sección.

Entonces aparece \(\varepsilon_R\). No como un decimal suplicado por el
cierre, sino como la diferencia que la estructura estaba obligada a
transportar:
\[
\epsone=\frac{20}{81}\Delta_4-(S_E-1).
\]
El acarreo publica sus cifras; la profundidad certifica su límite; la carta
angular lo reconoce como
\[
5.13830167585752820716851646226459474899\ldots
\times10^{-8}\,{}^\circ.
\]
La ecuación del radián queda cerrada porque ambos lados son publicaciones del
mismo estado:
\BigRadionEquation

La historia no termina en la igualdad. El par \((A,C^*)\) construye una
elipse positiva. Su producto de semiejes fija \(2\hret\); su cociente fija la
forma. El teorema de área demuestra que una unidad de fase barre \(\hret\) y
que la vuelta encierra \(\hfull\). Aquí adquiere contenido físico la división
por \(2\pi\): \(h\) es acción de vuelta; \(\hbar\) es acción por radión.

Los focos dejan de ser adornos de una figura. Su separación reconstruye
\(C^*\) en la carta espectral, y su distancia de Klein mide una profundidad
hiperbólica finita dentro de un dominio cuya área hiperbólica total es
infinita. La misma frontera contiene una acción finita. La aparente paradoja
es la tesis central de la geometría: una medida cierra la acción; otra abre la
profundidad.

El círculo, por tanto, no es falso. Es la proyección exacta que conserva fase
y olvida anisotropía, interior y memoria. La elipse conserva acción y forma.
Klein conserva profundidad. La hélice conserva retorno y biografía. Ninguna
imagen sustituye a las demás; todas nacen del mismo estado.

La realización de Hilbert reconoce la elipse como contorno de covarianza
mínima. El oscilador LC la reconoce como intercambio de reservas eléctrica y
magnética. El lector \(90/120\) publica permeabilidad, permitividad,
impedancia y velocidad. La incidencia seis--ocho produce una pantalla causal;
Hodge convierte orientación en dualidad de campos; Lorentz convierte campo y
carga en dinámica; Cartan--Holst recibe la holonomía en geometría
gravitatoria. El orden importa: el objeto HMT selecciona; la física
convencional realiza.

La palabra \emph{radión} queda así justificada. No significa que un foco sea
una carga positiva y el otro una carga negativa. Significa que la unidad de
fase posee orientación firmada antes de la publicación eléctrica:
\[
\mathcal A(\mathfrak r_+)=+\hret,
\qquad
\mathcal A(\mathfrak r_-)=-\hret.
\]
La carga aparece cuando Mecánica Dimensional aplica el transductor. El signo
no se añade al final; se conserva desde la hoja.

También queda explicado por qué el sector \(90/120\) parecía esconder muchas
respuestas. Incidencia, extractor dodecafásico, cola espectral, núcleo de
Barbero, torres y acarreo comparten origen. Precisamente por compartir origen
pueden compararse. Precisamente por tener operaciones distintas no deben
fundirse. La tipificación no reduce la unidad del sistema: la hace inteligible.

% REMATE_LOCAL_RADION_CIERRE
\Needspace{28\baselineskip}
La conclusión más compacta de esta obra puede escribirse en cuatro líneas:
\[
\boxed{
\begin{aligned}
\text{fase:}\quad&\theta=1,\\
\text{acción:}\quad&\mathcal A=\sigma\hret,\\
\text{memoria:}\quad&H_9(n+9)=H_9(n)+(0,1),\\
\text{compatibilidad:}\quad&1=S_E-\Phi_T+\epsone.
\end{aligned}}
\]

El radión es esa unidad cuádruple. El radián es su sombra exacta en la carta
angular. La diferencia no está en el valor de \(\pi\); está en cuánta
estructura estamos dispuestos a conservar cuando decimos que una vuelta ha
cerrado.

\vspace{2\baselineskip}
\begin{center}
\begin{minipage}{.82\textwidth}
\centering\large\itshape\color{RadNavy}
El círculo publica la fase.\\
La elipse contiene la acción.\\
Klein abre la profundidad.\\
La hélice recuerda el camino.\\[.7em]
El radión es la unidad que no obliga a elegir entre ellas.
\end{minipage}
\end{center}



\HMTFlushCurrentChapter
\HMTSuccessorChapterInput
  {Clasificación genealógica de los números: secciones supervivientes, invariantes y atlas de constantes}
  {../colaboracion/parte_iii/tex/capitulos_finales/52_atlas_numero.tex}

\HMTFlushCurrentChapter
\HMTSuccessorChapterInput
  {Completitud genealógica de las secciones HMT y decidibilidad de cada nivel finito}
  {../colaboracion/parte_iii/tex/capitulos_finales/53_completitud_finita.tex}

\HMTFlushCurrentChapter
\HMTSuccessorChapterInput
  {Obstrucción de Gödel--Turing a la decisión uniforme de multiplicidades proyectivas}
  {../colaboracion/parte_iii/tex/capitulos_finales/54_godel_turing.tex}

\HMTFlushCurrentChapter
\HMTSuccessorChapterInput
  {Gödel--Cohen: consistencia relativa, independencia y forzamiento sobre espacios de prefijos}
  {../colaboracion/parte_iii/tex/capitulos_finales/55_godel_cohen.tex}

\HMTFlushCurrentChapter
\chapter{Sistema operatorio de Euler en HMT: defecto suma--producto, geometrías exponenciales triádicas, Principio de Ambivalencia y doble proyección}
% Cuerpo editor-ready del nuevo capitulo 57.

\section{Defecto suma--producto de APP y coeficiente lineal del producto modular de Euler}
\label{sec:l9s102:euler-defecto}

\subsection{Identidad exacta \texorpdfstring{\(459-405=54=[q]t_3(q)\)}{459-405=54=[q]t3(q)}}

Sobre el soporte de \(81\) celdas, las dos hojas de APP producen
\begin{equation}
 S_+=405,
 \qquad
 S_\times=459,
 \qquad
 S_\times-S_+=54.
 \label{eq:l9s102:euler-405-459}
\end{equation}
La regularización triádica suministra el exponente \(-1/12\). Su
publicación dodecafásica determina el polo \(q^{-1}\), y el cociente de
Dedekind de nivel \(3\) se escribe
\begin{equation}
 t_3(q)
 =\left(\frac{\eta(\tau)}{\eta(3\tau)}\right)^{12}
 =q^{-1}\prod_{n\ge1}(1+q^n+q^{2n})^{-12},
 \qquad q=e^{2\pi i\tau}.
 \label{eq:l9s102:euler-t3}
\end{equation}

\begin{theorem}[Identidad suma--producto de Euler--HMT]
\label{thm:l9s102:euler-54}
El defecto entre las dos evaluaciones de APP coincide con el coeficiente
lineal de \(t_3\):
\begin{equation}
 \boxed{
 S_\times-S_+=54=[q]\,t_3(q).}
 \label{eq:l9s102:euler-54}
\end{equation}
\end{theorem}

\begin{proof}
Cada factor admite
\((1+q^n+q^{2n})^{-12}=1-12q^n+O(q^{2n})\). Para el coeficiente de
\(q\) posterior al polo \(q^{-1}\), contribuyen las particiones de grado
\(2\) del producto de nivel \(3\). La expansión finita certificada produce
\([q]t_3=54\). APP obtiene el mismo entero por la resta exacta
\(459-405\).
\end{proof}

\subsection{Conservación del defecto mediante residuo, cociente y acarreo nonádico}

La identidad \cref{eq:l9s102:euler-54} enlaza la diferencia aditiva de dos
hojas finitas con un producto modular infinito. El enlace conserva el
origen de \(54\): suma, producto y regularización permanecen como tres
coordenadas distinguibles del estado.

\subsubsection{Completación EMRH y conservación de la unidad}
\label{sec:l9s102:euler-emrh}

La fibra ternaria de profundidad \(6\) posee \(3^6=729\) elementos. Su
completación decimal distingue
\begin{equation}
 729+270=999,
 \qquad
 729+271=1000,
 \qquad
 271=243+27+1.
 \label{eq:l9s102:euler-729-1000}
\end{equation}
El término terminal \(+1\) es la unidad de cierre transmitida por el
acarreo. La igualdad binomial
\begin{equation}
 1000=(9+1)^3=9^3+3\cdot9^2+3\cdot9+1
 \label{eq:l9s102:euler-binomial-1000}
\end{equation}
expresa el transporte entre la carta nonádica y la publicación decimal sin
perder la genealogía de la completación.

La evaluación arquimediana se realiza sobre cilindros semicerrados
compatibles. Las dos expansiones de un racional terminal representan el
mismo valor y conservan historias de frontera distintas. Esta estructura
tipa la identidad \(0.999\ldots=1\) como igualdad de evaluación acompañada
por memoria de ruta.

\section{Álgebras cuadráticas del TRIT y ley exponencial elíptica, parabólica e hiperbólica}
\label{sec:l9s102:euler-algebras-trit}

\subsection{Familia operatoria \texorpdfstring{\(J_{+1}^2=-I\), \(J_0^2=0\) y \(J_{-1}^2=I\)}{J+1^2=-I, J0^2=0 y J-1^2=I}}

Para \(\tau\in\{+1,0,-1\}\), definimos
\begin{equation}
 \mathbb A_\tau
 =\mathbb R[J_\tau]/(J_\tau^2+\tau I).
 \label{eq:l9s102:euler-algebras}
\end{equation}
La serie exponencial converge en cada álgebra y se reduce a
\begin{align}
 e^{tJ_{+1}}&=\cos t\,I+\sin t\,J_{+1},
 \label{eq:l9s102:euler-eliptica}\\
 e^{tJ_0}&=I+tJ_0,
 \label{eq:l9s102:euler-parabolica}\\
 e^{tJ_{-1}}&=\cosh t\,I+\sinh t\,J_{-1}.
 \label{eq:l9s102:euler-hiperbolica}
\end{align}
TRIT asigna respectivamente cierre elíptico, variación de primer orden y
apertura hiperbólica. La identidad clásica ocupa la hoja \(\tau=+1\):
\begin{equation}
 \boxed{\operatorname{Exp}(\pi J_{+1})+I=0.}
 \label{eq:l9s102:euler-clasica}
\end{equation}
La hoja parabólica conserva el primer jet, y la hoja hiperbólica publica
una pareja de autovalores recíprocos.

\subsection{Identidades \texorpdfstring{\(\operatorname{Exp}(\pi J_{+1})+I=0\)}{Exp(pi J+1)+I=0} y \texorpdfstring{\(\operatorname{Exp}(2\log\varphi\,H_\varphi)=F_{\mathrm{av}}^2\)}{Exp(2 log phi H_phi)=F_av^2}}

Sea
\begin{equation}
 F_{\mathrm{av}}=\begin{pmatrix}0&1\\1&1\end{pmatrix},
 \qquad
 H_\varphi=\frac{2F_{\mathrm{av}}^2-3I}{2\varphi-1}.
 \label{eq:l9s102:euler-hphi}
\end{equation}
Entonces \(H_\varphi^2=I\) y
\begin{equation}
 \boxed{
 \operatorname{Exp}(2\log\varphi\,H_\varphi)=F_{\mathrm{av}}^2.}
 \label{eq:l9s102:euler-fav}
\end{equation}
En esta coordinación, \(\pi\) mide el cierre angular, \(\varphi\) el
autovalor expansivo de la memoria y \(e\) la operación que transforma un
generador aditivo en transporte multiplicativo.

\section{Principio de Ambivalencia: involución entre \texorpdfstring{\(\pi/\varphi^2\)}{pi/phi^2} y \texorpdfstring{\(\varphi/\pi^2\)}{phi/pi^2} y acumulación nonádica de orientación}
\label{sec:l9s102:euler-ambivalencia}

\subsection{Acción \texorpdfstring{\(\mathcal Q(x,y)=(x/y^2,y/x^2)\)}{Q(x,y)=(x/y^2,y/x^2)} sobre los lectores areal y volumétrico}

Definimos el lector cuadrático
\begin{equation}
 \mathcal Q(x,y)
 =\left(\frac{x}{y^2},\frac{y}{x^2}\right).
 \label{eq:l9s102:euler-q}
\end{equation}
En las coordenadas
\begin{equation}
 P=xy,
 \qquad O=\frac{x}{y},
 \label{eq:l9s102:euler-po}
\end{equation}
su acción es
\begin{equation}
 P\circ\mathcal Q=P^{-1},
 \qquad
 O\circ\mathcal Q=O^3.
 \label{eq:l9s102:euler-q-action}
\end{equation}
La segunda iteración satisface
\begin{equation}
 \boxed{
 P\circ\mathcal Q^2=P,
 \qquad
 O\circ\mathcal Q^2=O^9.}
 \label{eq:l9s102:euler-q2}
\end{equation}
El producto retorna y la orientación conserva una potencia nonádica de
memoria.

\subsection{Identidad \texorpdfstring{\(\varphi^3(\pi/\varphi^2)^2(\varphi/\pi^2)=1\)}{phi^3(pi/phi^2)^2(phi/pi^2)=1} y bidireccionalidad tipada}

Para \((x,y)=(\pi,\varphi)\), las dos orientaciones son
\begin{equation}
 \rho_A=\frac{\pi}{\varphi^2},
 \qquad
 \rho_E=\frac{\varphi}{\pi^2},
 \label{eq:l9s102:euler-rho-a-e}
\end{equation}
y obedecen
\begin{equation}
 \boxed{\varphi^3\rho_A^2\rho_E=1.}
 \label{eq:l9s102:euler-ambivalencia-identidad}
\end{equation}
La involución \((x,y)\mapsto(y,x)\) intercambia las orientaciones areal y
electrónica del mismo lector \(R(x,y)=x/y^2\).

\section{Transformación exponencial del generador parangular \texorpdfstring{\(A_{\mathrm{rad}}I+C^*_{\mathrm{rad}}H_\varphi\)}{A_rad I+C*_rad H_phi}}
\label{sec:l9s102:euler-accion-par-angular}

La carta de acción
\begin{equation}
 T_{-34}(u)=u\,10^{-34}\mathcal U_S
 \label{eq:l9s102:euler-t34}
\end{equation}
transporta la identidad angular
\(C^*_{\deg}=2(\eta_{\mathrm{ret}}+\alpha_{\mathrm{HMT}})\) a
\begin{equation}
 \boxed{
 T_{-34}(C^*_{\deg})
 =2\hbar_{\mathrm{ret}}
 +2T_{-34}(\alpha_{\angle,\deg}).}
 \label{eq:l9s102:euler-cstar-hbar}
\end{equation}
Cada sumando pertenece así a la misma recta dimensional
\(\mathcal L_S\).

\subsection{Proyectores espectrales \texorpdfstring{\(P_\pm=(I\pm H_\varphi)/2\)}{P+/-=(I+/-H_phi)/2} y autovalores conjugados \texorpdfstring{\(q_\pm\)}{q+/-}}

Sea
\begin{equation}
 B=A_{\mathrm{rad}}I+C^*_{\mathrm{rad}}H_\varphi,
 \qquad
 P_\pm=\frac{I\pm H_\varphi}{2}.
 \label{eq:l9s102:euler-b}
\end{equation}
La exponenciación transforma el generador aditivo en dos hojas
multiplicativas conjugadas:
\begin{equation}
 \boxed{
 e^{-B}=q_+P_++q_-P_-,
 \qquad
 q_\pm=e^{-A_{\mathrm{rad}}\mp C^*_{\mathrm{rad}}}.}
 \label{eq:l9s102:euler-exp-b}
\end{equation}

\subsection{Compatibilidad con la publicación parangular y los canales \texorpdfstring{\(90/120\)}{90/120}}

Los cardinales \(90\) y \(120\) proceden de la incidencia trítica del TPK.
La pareja \(q_\pm\) no los genera: evalúa espectralmente el operador
\(V_{90,120}\) sobre esos canales previamente construidos.

\section{Frontera elíptica de acción e interior hiperbólico de Hilbert--Beltrami--Klein}
\label{sec:l9s102:euler-elipse-klein}

\subsection{Coordenadas conjugadas de la elipse de acción y ley \texorpdfstring{\(\operatorname{Area}(\theta)=\hbar_{\mathrm{ret}}\theta\)}{Area(theta)=hbar_ret theta}}

La rapidez angular
\begin{equation}
 \eta=\operatorname{artanh}\!\left(\frac{C^*}{A}\right)
 \label{eq:l9s102:euler-rapidez}
\end{equation}
define los semiejes simplécticos
\begin{equation}
 a_S=\sqrt{2\hbar_{\mathrm{ret}}e^\eta},
 \qquad
 b_S=\sqrt{2\hbar_{\mathrm{ret}}e^{-\eta}}.
 \label{eq:l9s102:euler-semiejes}
\end{equation}
Con
\begin{equation}
 Q=a_S\cos\theta,
 \qquad
 P=b_S\sin\theta,
 \label{eq:l9s102:euler-qp}
\end{equation}
la acción encerrada por el arco satisface
\begin{equation}
 \boxed{\operatorname{Area}(\theta)=\hbar_{\mathrm{ret}}\theta.}
 \label{eq:l9s102:euler-area-accion}
\end{equation}

\subsection{Correspondencia \texorpdfstring{\(\rho=\tan(\theta/2)=\tanh(u/2)\)}{rho=tan(theta/2)=tanh(u/2)} entre borde periódico e interior hiperbólico}

La coordenada de Gudermann
\begin{equation}
 \rho=\tan\frac\theta2=\tanh\frac u2
 \label{eq:l9s102:euler-gudermann}
\end{equation}
coordina las cartas circular e hiperbólica.
La frontera porta la órbita simpléctica; el dominio convexo interior porta
la métrica hiperbólica de Hilbert--Beltrami--Klein. La coordenada de
Gudermann realiza su correspondencia radial y conserva la distinción de
métricas.

\section{Compatibilidad del cierre operatorio de Euler con las publicaciones arquimediana y excepcional}
\label{sec:l9s102:euler-doble-proyeccion}

\subsection{Doble proyección del mismo estado enriquecido y conservación de la incidencia}

Sea \(X^{\mathrm{enr}}\) el estado enriquecido construido por
APP--TRIT--TPK. Sus dos publicaciones principales son
\begin{equation}
 X^{\mathrm{enr}}
 \xrightarrow{\ \mathcal P_{\mathrm{arq}}\ }\mathbb R,
 \qquad
 X^{\mathrm{enr}}
 \xrightarrow{\ \mathcal P_{\mathrm{exc}}\ }
 \mathrm{Witt}\to\mathrm{Golay}\to\mathrm{Leech}
 \to V^\natural\to\mathbb M\to J.
 \label{eq:l9s102:euler-doble-proyeccion}
\end{equation}
La primera contrae cilindros compatibles hacia el valor arquimediano; la
segunda conserva la incidencia de frontera hasta el corredor excepcional.
El entero \(54\) aparece en APP y en el coeficiente modular de
\cref{eq:l9s102:euler-54}; el corredor \(4\to2\to6\to54\) transporta
después esa incidencia hacia el retículo de Leech, el álgebra de operadores
de vértice y Moonshine. La acción del Monstruo reside en
\(V^\natural\); la genealogía digital y de rutas permanece en el dominio
del TPK y llega al corredor mediante los mapas de incidencia declarados.

\subsection{Sistema macrovectorial \texorpdfstring{\(\mathfrak E_{\mathrm{HMT}}=0\)}{E_HMT=0}: confluencia de las identidades eulerianas en un único operador}
\label{sec:l9s102:euler-macrovector}

Las identidades anteriores se reúnen en la aplicación tipada
\begin{equation}
 \boxed{
 \mathfrak E_{\mathrm{HMT}}=
 \begin{pmatrix}
 (S_\times-S_+)-[q]t_3\\[1mm]
 \operatorname{Exp}(I)-eI\\
 \operatorname{Exp}(\pi J_{+1})+I\\
 \operatorname{Exp}(J_0)-I-J_0\\
 \operatorname{Exp}(2\log\varphi\,H_\varphi)-F_{\mathrm{av}}^2\\
 \varphi^3\rho_A^2\rho_E-1\\
 T_{-34}(C^*)-2\hbar_{\mathrm{ret}}
 -2T_{-34}(\alpha_{\angle,\deg})
 \end{pmatrix}=0.}
 \label{eq:l9s102:euler-macrovector}
\end{equation}
Cada componente conserva su dominio y su demostración propia. La aplicación
agregada exhibe la coordinación de suma, producto, cierre elíptico, primer
jet, expansión hiperbólica, ambivalencia y acción.

La composición interna APP--TRIT--TPK ya transporta suma, producto,
orientación, residuo, cociente, acarreo, ruta y memoria hasta las
publicaciones anteriores. No se postula por ello un entrelazador HMT
ausente. Una promoción posterior sólo puede formularse mediante el cuadrado
tipado específico exigido por una representación ulterior ---por ejemplo,
una equivariancia grupal completa hasta Witt---, sin rebajar la confluencia
operatoria aquí demostrada.


\let\HMTSuccessorChapterInput\relax










\HMTFlushCurrentChapter
